
\documentclass[UTF8,10.5pt,openany]{ctexbook}
\usepackage[a4paper,left=2.1cm,right=2.1cm,top=2.5cm,bottom=2.2cm,headheight=23pt,headsep=9pt]{geometry}
\usepackage{fontspec}
\usepackage{amsmath}
\usepackage{amssymb}
\usepackage{xeCJK}
\setCJKmainfont{Noto Serif CJK SC}[BoldFont={Noto Serif CJK SC},BoldFeatures={FakeBold=2.0}]
\setCJKsansfont{Noto Sans CJK SC}
\setCJKmonofont{Noto Sans Mono CJK SC}
\setmainfont{Noto Serif}
\setsansfont{Noto Sans}
\setmonofont{Noto Sans Mono}[Scale=0.86]
\setCJKfamilyfont{zhkai}{FandolKai-Regular.otf}
\newcommand{\kai}{\CJKfamily{zhkai}}
\usepackage{xcolor}
\definecolor{ink}{RGB}{17,17,17}
\definecolor{grayt}{RGB}{110,110,110}
\definecolor{linet}{RGB}{150,150,150}
\definecolor{accent}{RGB}{20,60,120}
\usepackage{titlesec}
\usepackage{titletoc}
\usepackage{fancyhdr}
\usepackage{enumitem}
\usepackage{longtable}
\usepackage{booktabs}
\usepackage{array}
\usepackage{needspace}
\usepackage{multicol}
\usepackage[most]{tcolorbox}
\usepackage[bookmarks=true,bookmarksnumbered=true,bookmarksopen=true,bookmarksopenlevel=1,hidelinks,unicode]{hyperref}
\linespread{1.30}
\setlength{\parindent}{0pt}
\setlength{\parskip}{2pt}
\pagestyle{fancy}
\fancyhf{}
\renewcommand{\headrulewidth}{0.4pt}
\fancyhead[L]{\footnotesize\sffamily\color{grayt}AI Infra 面试笔记}
\fancyhead[R]{\footnotesize\sffamily\color{grayt}\leftmark}
\fancyfoot[C]{\footnotesize\color{grayt}\thepage}
\renewcommand{\chaptermark}[1]{\markboth{#1}{}}
\ctexset{chapter/name={第,篇},chapter/number=\arabic{chapter}}
\titleformat{\chapter}[display]
  {\normalfont\centering}
  {\footnotesize\sffamily\color{grayt}\CTEXthechapter}
  {4pt}
  {\sffamily\bfseries\fontsize{18}{23}\selectfont}
  [\vspace{3pt}{\color{linet}\rule{0.42\textwidth}{0.5pt}}]
\titlespacing*{\chapter}{0pt}{-8pt}{14pt}
\titlecontents{section}[2.2em]{\addvspace{1pt}\sffamily\footnotesize}{}{}{\hfill\contentspage}
\titlecontents{chapter}[0pt]{\addvspace{5pt}\sffamily}
  {\thecontentslabel\quad}{}{\hfill\contentspage}
\newcommand{\lec}[2]{\needspace{10\baselineskip}\vspace{10pt}%
  {\sffamily\bfseries\fontsize{12.5}{16}\selectfont\color{ink}#1\quad#2}\par
  \vspace{-3pt}{\color{linet}\rule{\linewidth}{0.35pt}}\par\vspace{4pt}}
\newcommand{\qlabel}[1]{{\sffamily\footnotesize\bfseries\color{accent}#1}\;}
\newcommand{\askline}[1]{\noindent{\qlabel{问}}%
  \begingroup\sloppy{\fontsize{10}{14.6}\selectfont #1\par}\endgroup\vspace{1.5pt}}
\newcommand{\numline}[2]{\vspace{2pt}\noindent{\qlabel{#1}}%
  \begingroup\sloppy\raggedright{\sffamily\footnotesize\color{grayt}#2\par}\endgroup}
\newcommand{\fol}[2]{\vspace{3pt}\noindent{\qlabel{追问}}%
  \begingroup\sloppy{\kai\fontsize{10}{14.6}\selectfont #1\par}\endgroup\vspace{1pt}%
  \noindent\hspace*{1.15em}%
  \begingroup\sloppy{\fontsize{10}{14.6}\selectfont #2\par}\endgroup}

\begin{document}

\thispagestyle{empty}
\begin{center}
\vspace*{3.2cm}
{\sffamily\fontsize{10}{14}\selectfont\color{grayt}面试问答版 \quad 153 讲 · 20 篇}\par
\vspace{14pt}
{\sffamily\bfseries\fontsize{30}{40}\selectfont AI Infra}\par
\vspace{6pt}
{\sffamily\bfseries\fontsize{22}{30}\selectfont 面试笔记}\par
\vspace{18pt}
{\color{linet}\rule{0.5\textwidth}{0.6pt}}\par
\vspace{16pt}
{\kai\fontsize{12}{19}\selectfont
每一讲都是一段可以开口讲的回答：\\[3pt]
面试官会怎么问、先说什么、凭什么让人信、被追问怎么接\par}
\vfill
{\sffamily\footnotesize\color{grayt}
内容整理自孙晓辉《AI Infra Handbook》V1.0（153 讲 / 1342 页）\\[2pt]
只保留面试要答的部分，推导与代码去原书看\par}
\vspace{1.2cm}
\end{center}
\clearpage


\chapter*{怎么用}
\addcontentsline{toc}{chapter}{怎么用}
\markboth{怎么用}{}

这本笔记按「面试现场」组织。每一讲给三样东西：\textbf{面试官可能怎么问}、\textbf{你的答题主线}、
\textbf{被追问时怎么接}。答题主线是可以直接说出口的完整句子，不是提纲。

\vspace{4pt}
{\sffamily\bfseries 三轮复习法}\par\vspace{2pt}
\begin{itemize}[leftmargin=1.6em,itemsep=2pt,topsep=2pt]
\item \textbf{第一轮（两天）}：把每讲的「问」和「答」各读一遍，先记住 153 个问题的骨架。
\item \textbf{第二轮（三天）}：只看「问」，自己开口答，答完对照主线补漏；重点记「数字」那行。
\item \textbf{第三轮（一小时）}：翻附录的问题清单，看到问题能立刻说出结论 + 一个数字即可。
\end{itemize}

\vspace{4pt}
{\sffamily\bfseries 答题的通用结构}\par\vspace{2pt}
面试里任何一讲的答案都可以套这个顺序：\textbf{结论 $\rightarrow$ 机制（数据怎么流）$\rightarrow$ 为什么更好（对比 + 量级）
$\rightarrow$ 代价与边界 $\rightarrow$ 工程取舍}。笔记里每条 answer 基本就是这个顺序，背的时候按层次记，别背字。

\vspace{4pt}
{\sffamily\bfseries 说明}\par\vspace{2pt}
讲号与原书一致（第 1\textasciitilde 153 讲），篇的划分照原书（20 篇）。数字、公式、器件型号均来自原书；
原书标「选读」的几讲也收录了。若与原书有出入，以原书为准。

\vspace{6pt}
{\footnotesize\color{grayt}共 153 讲，每讲 2\textasciitilde 3 个问法、4\textasciitilde 6 条答题主线、2\textasciitilde 3 个追问。}
\clearpage

\tableofcontents
\clearpage
\chapter{数学基础}
{\footnotesize\sffamily\color{grayt}第 1--11 讲}\par\vspace{-4pt}
\phantomsection\addcontentsline{toc}{section}{1\quad 向量}
\lec{1}{向量}
\askline{大模型怎么把文字变成数字的？}
\askline{向量到底要占多少显存，怎么估？}
\askline{为什么注意力里到处是线性组合？}
\vspace{3pt}
\noindent{\qlabel{答}}\begingroup\sloppy\begin{enumerate}[leftmargin=1.6em,itemsep=2.5pt,topsep=0pt,parsep=0pt,label={\sffamily\footnotesize\color{accent}\arabic*.}]
\item {\fontsize{10}{14.8}\selectfont 向量就是一串按固定顺序排好的数，把事物变成向量最大的好处，是它们之间的关系（像不像、差多少）变成了可以算出来的数字，所以大模型做的第一件事就是给词表里每个词准备一个向量。\par}
\item {\fontsize{10}{14.8}\selectfont 向量有长度，等于各分量平方和再开方，苹果 (7,4) 是 $\surd$65$\approx$8.06；基本运算只有加法和数乘，加法是对应分量相加、维度必须相同，数乘是每个分量乘 k、方向不变长度变 |k| 倍。\par}
\item {\fontsize{10}{14.8}\selectfont 它占多少内存也从这个 d 算起：每个分量是一个浮点数，BF16 占 2 字节，4096 维向量就是 4096$\times$2=8192 字节即 8 KiB；1000 个 token 约 7.8 MiB，12.8 万词的词表约 0.98 GiB。\par}
\item {\fontsize{10}{14.8}\selectfont 维度 d 是整个课程里最基本的数，显存要多少、计算量有多大，几乎都要从「每个向量有 d 个数」开始算；而且维度再高规则一字不改，3 维能画、4 维以上画不出来但加法、数乘、长度公式完全一样。\par}
\item {\fontsize{10}{14.8}\selectfont 代价是单个分量通常没人能读懂：水果的每个分量都有名字，大模型学出来的向量单独看某个分量说明不了什么，有意义的是整个向量和其他向量之间的关系，所以真正要算的是点积和距离。\par}
\item {\fontsize{10}{14.8}\selectfont 反复用到的其实只有线性组合 k$_1$a+k$_2$b，第 26 讲的注意力输出就是把若干向量按注意力权重做线性组合，第 29 讲的残差连接是把输入和输出直接相加，所以基本运算只讲加法和数乘这两种。\par}
\end{enumerate}\endgroup
\numline{公式}{$\lVert v\rVert=\sqrt{v_1^2+v_2^2+\cdots+v_n^2}$；$k\cdot v=(kv_1,kv_2,\ldots,kv_n)$}
\numline{数字}{4096 维；BF16 每分量 2 字节；单向量 8192 字节 = 8 KiB；1000 token $\approx$ 7.8 MiB；词表 128256 词 $\approx$ 0.98 GiB；苹果 (7,4) 长度 $\surd$65 $\approx$ 8.06}
\fol{那维度再高算法会变吗？}{不会。加法、数乘、长度平方和开方这些规则一个字都不用改，只是分量变多；4 维以上画不出来，但 4096 维和 2 维跑的是同一套代码。}
\fol{一个向量才 8 KiB，为什么要算它？}{单看很小，但模型有几十层、要处理上千 token、还要乘 batch；1000 个 token 就是 7.8 MiB，一张词表就 0.98 GiB，量级立刻变成显存问题。}
\fol{你怎么验证这些估算？}{用 NumPy 三行就够：d*2 得 8192 字节，1000*d*2/2**20 得 7.8125 MiB，128256*d*2/2**30 得 0.9785 GiB，和手算一致就说明估法没错。}
\numline{别答错}{别把向量的「维」和数组的 ndim 混为一谈：4096 维向量只是形状 (4096,) 的一维数组。}
\phantomsection\addcontentsline{toc}{section}{2\quad 点积与余弦相似度}
\lec{2}{点积与余弦相似度}
\askline{判断两个向量像不像，为什么不用距离？}
\askline{点积为什么能反映两个向量的方向？}
\askline{注意力里为什么用点积打分？}
\vspace{3pt}
\noindent{\qlabel{答}}\begingroup\sloppy\begin{enumerate}[leftmargin=1.6em,itemsep=2.5pt,topsep=0pt,parsep=0pt,label={\sffamily\footnotesize\color{accent}\arabic*.}]
\item {\fontsize{10}{14.8}\selectfont 用距离量会把「长度不同」误当成「内容不同」：GPU 短文 (4,1) 到长文 (12,3) 的距离是 8.25，到美食文章 (1,5) 只有 5，按距离反而更像美食文，和常识正好相反。\par}
\item {\fontsize{10}{14.8}\selectfont 点积的定义是对应分量相乘再相加，结果是一个数而不是向量：(4,1)$\cdot$(12,3)=51，(4,1)$\cdot$(1,5)=9，这次的排序才符合常识，但只看这个数还说明不了太多，因为长文的数字本身就大。\par}
\item {\fontsize{10}{14.8}\selectfont 用余弦定理把 |a-b|$^2$ 展开，就能推得 a$\cdot$b=|a||b|cos$\theta$，这个式子在 n 维一字不改；长度总是正的，所以点积的符号就等于 cos$\theta$ 的符号，正同向、零垂直、负反向。\par}
\item {\fontsize{10}{14.8}\selectfont 把两个长度除掉就得到只含方向的余弦相似度，取值范围 -1 到 1：短文与长文正好是 1，短文与美食文约 0.43；如果把向量预先归一化成长度 1，点积就直接等于余弦，连除法都省了。\par}
\item {\fontsize{10}{14.8}\selectfont 大模型偏爱点积是因为它便宜又适合成批算：d 维点积只要约 2d 次浮点运算，4096 维约 8192 次，没有开方也没有除法，而矩阵乘法本质上就是一大张点积表，能直接交给 GPU 的矩阵乘硬件。\par}
\item {\fontsize{10}{14.8}\selectfont 代价是计算量随长度平方增长：注意力里每个词都要和每个词做一次点积，s=1000 时按 32 层、32 头、每头 128 维粗算约 2.6$\times$10$^1$$^1$ FLOPs，s 变 10 倍点积个数变 100 倍，这是长文本又慢又贵的第一个伏笔。\par}
\end{enumerate}\endgroup
\numline{公式}{$a\cdot b=\sum_{i}a_i b_i=\lVert a\rVert\,\lVert b\rVert\cos\theta$；$\cos\theta=\frac{a\cdot b}{\lVert a\rVert\,\lVert b\rVert}$}
\numline{数字}{距离 8.25 与 5；点积 51 与 9；余弦 1.0 与 0.43；4096 维点积约 8192 FLOPs；1000$^2$$\times$32$\times$2$\times$128$\times$32 $\approx$ 2.6$\times$10$^1$$^1$ FLOPs；TinyLM「我」与「推理」余弦 0.59、夹角约 54$^{\circ}$}
\fol{余弦为 1 就说明两个向量一样吗？}{不是。余弦只看方向，(4,1) 和 (12,3) 余弦是 1 但两个向量并不相等；长度本身也有意义时（比如注意力打分）就不能只看余弦。}
\fol{余弦相似度为 0 是相反吗？}{不是，0 表示垂直，可以理解成「没关系」；方向相反是 -1。点积为 0 和夹角 90$^{\circ}$ 说的是同一件事。}
\fol{为什么说长文本推理又慢又贵？}{打分个数与 s$^2$ 成正比：s 变 10 倍点数变 100 倍，1000 token 单层打分约 8.2$\times$10⁹ FLOPs，乘 32 层到 2.6$\times$10$^1$$^1$，瓶颈先卡在注意力这一步。}
\numline{别答错}{别把「对应分量相乘」当成点积：a*b 得到的还是向量，再加起来才是点积 a@b。}
\phantomsection\addcontentsline{toc}{section}{3\quad 矩阵乘法}
\lec{3}{矩阵乘法}
\askline{矩阵乘法在模型里到底在算什么？}
\askline{为什么权重矩阵形状必须对得上？}
\askline{逐 token 生成为什么 GPU 跑不满？}
\vspace{3pt}
\noindent{\qlabel{答}}\begingroup\sloppy\begin{enumerate}[leftmargin=1.6em,itemsep=2.5pt,topsep=0pt,parsep=0pt,label={\sffamily\footnotesize\color{accent}\arabic*.}]
\item {\fontsize{10}{14.8}\selectfont 矩阵乘法 C=XW 的每一格都是 X 的一行和 W 的一列的点积，整张 C 就是一张点积表：4 个词对 4 个探测方向共 4$\times$4=16 个点积，按词成行、方向成列排好就是结果。\par}
\item {\fontsize{10}{14.8}\selectfont 形状规则是 [m,k]@[k,n]$\rightarrow$[m,n]，左边矩阵的列数必须等于右边矩阵的行数，因为点积要求两个向量同维；中间的 k 相消、两边的 m 和 n 留下，大模型里就是 [s,d]@[d,d\_out]$\rightarrow$[s,d\_out]。\par}
\item {\fontsize{10}{14.8}\selectfont 换一个角度一行一行看，W 是一台变换机器，X 的每一行被送进去变成 C 的对应一行，所有行用同一台机器而且互不干扰；Q/K/V 投影、FFN 的升降维、最后的输出层全是这件事。\par}
\item {\fontsize{10}{14.8}\selectfont 计算量约 2mnk 次浮点运算，因为 C 有 m$\times$n 格、每格是一个 k 维点积要做 k 次乘加；TinyLM 那个 4$\times$4 乘 4$\times$4 就是 2$\times$4$\times$4$\times$4=128 FLOPs，后面估计算量都用这个单位。\par}
\item {\fontsize{10}{14.8}\selectfont 4096$\times$4096 的 BF16 权重占 32 MiB，不管一次处理几个 token 都要整块读一遍：s=1 时只做 3.4$\times$10⁷ FLOPs、每字节算 1 次；s=1000 时每字节算 1000 次，比值正好约等于 s。\par}
\item {\fontsize{10}{14.8}\selectfont 算术强度低就意味着 GPU 大部分时间在等数据：逐 token 生成恰恰是 s=1 的情况，所以比读提示慢得多；把权重换成 FP8 少搬一半数据，这也是低精度能提速的原因之一。\par}
\end{enumerate}\endgroup
\numline{公式}{$C_{ij}=\sum_{t}A_{it}B_{tj}$；一次矩阵乘法约 $2mnk$ FLOPs；$\text{算术强度}=\frac{2skn}{kn\cdot 2}=s$}
\numline{数字}{16 个点积；16$\times$4=64 次乘加 = 128 FLOPs；4096$^2$ 权重 BF16 占 32 MiB；s=1 时 3.4$\times$10⁷ FLOPs、s=1000 时 3.4$\times$10$^1$⁰ FLOPs；算术强度 = s}
\fol{XW 和 WX 的结果一样吗？}{一般不等。XW 是 X 的行和 W 的列做点积，WX 是 W 的行和 X 的列，参与的向量根本不同；形状非方阵时连结果形状都变，[2,3]@[3,2] 得 2$\times$2，反过来得 3$\times$3。}
\fol{NumPy 里 * 和 @ 差在哪？}{* 是逐元素相乘，两个 4$\times$4 出来还是 4$\times$4，每格只做一次乘法；@ 或 matmul 才是矩阵乘法，每一格是一整个点积。}
\fol{算术强度那个 s 是怎么来的？}{FLOPs 是 2skn，读的字节数是 k$\cdot$n 乘每数 2 字节，相除只剩 s；所以它衡量的是读一遍权重能做多少计算，s 越大 GPU 越忙。}
\numline{别答错}{「行乘列」不是「列乘行」：C\_ij 用的是左边第 i 行和右边第 j 列，写成列乘行说明该转置了。}
\phantomsection\addcontentsline{toc}{section}{4\quad 线性变换}
\lec{4}{线性变换}
\askline{矩阵乘法几何上到底在做什么？}
\askline{行列式有什么几何意义？}
\askline{为什么神经网络一定要激活函数？}
\vspace{3pt}
\noindent{\qlabel{答}}\begingroup\sloppy\begin{enumerate}[leftmargin=1.6em,itemsep=2.5pt,topsep=0pt,parsep=0pt,label={\sffamily\footnotesize\color{accent}\arabic*.}]
\item {\fontsize{10}{14.8}\selectfont 矩阵就是线性变换，W 的第 i 行就是第 i 个坐标方向被送去的终点：e1=(1,0) 乘 W 得到的正是 W 第 1 行，e2 得到第 2 行，所以知道这两行整个空间怎么变就完全确定了。\par}
\item {\fontsize{10}{14.8}\selectfont 因为先组合再变换等于先变换再组合，任意 (x$_1$,x$_2$) 的像就是 x$_1$ 乘 W 第 1 行加 x$_2$ 乘 W 第 2 行，也就是按原来的分量对 W 的各行做线性组合，两种看法算的是同一件事。\par}
\item {\fontsize{10}{14.8}\selectfont 这类变换保持原点不动、直线还是直线、平行和等距都保持，拉伸、旋转、剪切、镜像、压扁都是它的特例；反过来说每个线性变换都能写成一个矩阵，所以「线性变换」和「矩阵」是同一件事的两种说法。\par}
\item {\fontsize{10}{14.8}\selectfont 2$\times$2 的行列式 det=ad-bc 就是单位正方形变换后的面积，也就是整个平面面积的放大倍数：默认的 W 是 2，旋转和剪切是 1，压扁是 0；镜像得到 -1，表示面积大小没变但平面被翻了过来。\par}
\item {\fontsize{10}{14.8}\selectfont det=0 说明空间被压成一条线，原来不同的点会落到同一个位置，(1,0) 和 (-1,1) 都被送到 (1,0.5)，信息丢了、变换不可逆；第 40 讲的 MLA 就是有控制地把向量压到低维再还原。\par}
\item {\fontsize{10}{14.8}\selectfont 两次线性变换可以合成一次，(x$\cdot$A)$\cdot$B = x$\cdot$(A$\cdot$B)，但顺序不能交换；也正因为叠一百层线性层仍然只是一次线性变换、掰不弯一条直线，网络中间必须插入激活函数。\par}
\end{enumerate}\endgroup
\numline{公式}{$x\cdot W=x_1\cdot(W\text{ 第 }1\text{ 行})+x_2\cdot(W\text{ 第 }2\text{ 行})$；$(x\cdot A)\cdot B=x\cdot(A\cdot B)$；$\det W=ad-bc$}
\numline{数字}{det W = 2$\times$1 - 1$\times$0 = 2；旋转、剪切 det=1；压扁 det=0；镜像 det=-1；4096$\times$4096 约 1680 万个数，BF16 占 32 MiB；旋转 30$^{\circ}$ 用 [[0.866, 0.5], [-0.5, 0.866]]}
\fol{det 为 0 在实际模型里意味着什么？}{意味着这一层把信息压没了、不可逆，自由度直接丢失；MLA 那种技术是故意压到还能还原的低维，和 det=0 的彻底丢失不是一回事。}
\fol{两个变换的顺序可以交换吗？}{一般不能。先拉伸再旋转，(1,0) 到 (0,2)；先旋转再拉伸，到 (0,1)，两个结果对应的矩阵 S$\cdot$R 和 R$\cdot$S 也不相等。}
\fol{为什么代码里是行向量乘矩阵？}{课本把向量写成竖的一列、变换写成 W$\cdot$x，这里跟代码习惯写成一横行 x$\cdot$W，两者差一个转置，所以去向在 W 的行，旋转矩阵看起来也「反了」。}
\numline{别答错}{平移不是线性变换：原点会动，一个矩阵乘法表示不了，所以一层网络要写成 x$\cdot$W + b。}
\phantomsection\addcontentsline{toc}{section}{5\quad 张量与形状}
\lec{5}{张量与形状}
\askline{代码里 [b, s, d] 具体指什么？}
\askline{拆头为什么必须用 transpose？}
\askline{广播是怎么对齐形状的？}
\vspace{3pt}
\noindent{\qlabel{答}}\begingroup\sloppy\begin{enumerate}[leftmargin=1.6em,itemsep=2.5pt,topsep=0pt,parsep=0pt,label={\sffamily\footnotesize\color{accent}\arabic*.}]
\item {\fontsize{10}{14.8}\selectfont 张量就是任意维度的数的阵列，标量、向量、矩阵是 0、1、2 维的特例；[b,s,d] 表示 b 个序列、每个 s 个 token、每个 token d 维，x[i,j] 取出来的是第 i 句第 j 个 token 的整个 d 维向量。\par}
\item {\fontsize{10}{14.8}\selectfont 矩阵乘法只看最后两维，前面的维度都当作批一一对应地算：[…,m,k]@[…,k,n]$\rightarrow$[…,m,n]；所以 [2,4,4]@[4,4] 还是 [2,4,4]，相当于两句话各自做一次第 3 讲的矩阵乘法。\par}
\item {\fontsize{10}{14.8}\selectfont 拆多头先 reshape 把 d 拆成 h$\times$d\_h，再 transpose 换到 token 前面得到 [b,h,s,d\_h]；q@kᵀ 用最后两维 [s,d\_h]@[d\_h,s] 算出每个头 s$\times$s 的打分 [b,h,s,s]。\par}
\item {\fontsize{10}{14.8}\selectfont 逐元素运算形状不变，归约则是沿哪一维哪一维就消失：对 [2,4,4] 沿最后一维求和得到 [2,4]，沿 token 维求平均也得到 [2,4]；softmax 和归一化都必须说清沿哪一维做。\par}
\item {\fontsize{10}{14.8}\selectfont 广播从最后一维开始对齐，两维相等或其中一个是 1 就能算，缺的前面补 1：所以 [2,4,4]+[4] 就是把同一个 bias 加到每句话的每个 token 上，省去显式复制也不额外占显存。\par}
\item {\fontsize{10}{14.8}\selectfont 张量大小是各维乘积乘字节数：[8,2048,4096] 是 6700 万个数、BF16 下 134 MB；注意力打分 [8,32,2048,2048] 一层就 2.1 GB，含 s$\times$s，这正是 FlashAttention 要解决的问题。\par}
\end{enumerate}\endgroup
\numline{公式}{$[\ldots,m,k]\mathbin{@}[\ldots,k,n]\to[\ldots,m,n]$；$\text{字节数}=\prod_{i}\text{shape}_i\times\text{每数字节}$}
\numline{数字}{[8,2048,4096] = 6700 万个数 = 134 MB；打分 [8,32,2048,2048] = 2.1 GB；TinyLM 从 [2,4] 到 [2,2,4,4]；拆头 [2,4,4]$\rightarrow$[2,4,2,2]$\rightarrow$[2,2,4,2]}
\fol{拆头时只 reshape 够吗？}{不能。[s,h,d\_h] 直接 reshape 成 [h,s,d\_h] 只是把同一串数重新切分，数据会错位；要让每个头的 token 排在一起必须用 transpose。}
\fol{广播是从左还是从右对齐？}{从最后一维往前对齐。所以 [4,4]+[4] 是给每一行加同一个向量；要给每一列加，就得把形状写成 [4,1]。}
\fol{怎么快速判断一个算子开销大不大？}{先看形状里有没有 s$\times$s 这种平方项，再算总个数乘字节数；134 MB 对 2.1 GB，一眼就能看出瓶颈在注意力打分这一层。}
\numline{别答错}{reshape 只换看法不动数据，transpose 才交换维度；两者互换会把数据错位。}
\phantomsection\addcontentsline{toc}{section}{6\quad 指数、对数与softmax}
\lec{6}{指数、对数与softmax}
\askline{softmax 为什么要先减去最大值？}
\askline{温度这个参数改的到底是什么？}
\askline{注意力里的 softmax 和输出层有何不同？}
\vspace{3pt}
\noindent{\qlabel{答}}\begingroup\sloppy\begin{enumerate}[leftmargin=1.6em,itemsep=2.5pt,topsep=0pt,parsep=0pt,label={\sffamily\footnotesize\color{accent}\arabic*.}]
\item {\fontsize{10}{14.8}\selectfont softmax 把任意一组有正有负的打分变成概率：先对每个分数取指数，让它们全变成正数并且保持大小顺序，再除以总和，让它们加起来等于 1，这就是名字里「soft 的 max」的意思。\par}
\item {\fontsize{10}{14.8}\selectfont 用 TinyLM 算一遍，读到「我」后打分是 1.64、0.83、-0.39、-0.17、-0.46、0.37，取指数得 5.155、2.293…总和 11.048，除完是 0.4666、0.2076、0.0613…加起来正好是 1。\par}
\item {\fontsize{10}{14.8}\selectfont 指数会放大差距：1.64 和 0.83 只差 0.81，概率却差了一倍多；对数反过来把乘法变加法，1000 个 0.01 相乘真值 10$^-$$^2$⁰⁰⁰ 会下溢成 0，取对数相加是 -4605.2，完全存得下。\par}
\item {\fontsize{10}{14.8}\selectfont 温度是在取指数之前先把分数除以 T：T=1 是原始 softmax，T=0.5 时最高项从 0.467 升到 0.749 分布变尖，T=2 时降到 0.302 分布变平，T 趋近 0 就退化成只取最高分。\par}
\item {\fontsize{10}{14.8}\selectfont 所有分数同加一个数结果不变，因为分子分母都多乘了 e\textasciicircum{}c 被约掉；所以实现时先减最大值，直接算 e\textasciicircum{}1000 会溢出成 inf，减完再算 softmax([1000,999]) 正常得到 [0.7311, 0.2689]。\par}
\item {\fontsize{10}{14.8}\selectfont 它在大模型里出现两次：输出层对整个词表（Llama-3 有 128256 个词）做一次，注意力则对 [b,h,s,s] 沿最后一维做；后者要在看不全一整行时算，FlashAttention 靠的正是减去最大值这个性质的推广。\par}
\end{enumerate}\endgroup
\numline{公式}{$p_i=\frac{e^{z_i/T}}{\sum_{j}e^{z_j/T}}$，实现时先减去 $\max(z)$}
\numline{数字}{e\textasciicircum{}z 总和 11.048；概率 0.4666/0.2076/0.0613；T=0.5 $\rightarrow$ 0.749，T=2 $\rightarrow$ 0.302；熵 1.48，ln6=1.79；softmax([1000,999]) = [0.7311, 0.2689]；Llama-3 词表 128256}
\fol{减去最大值会改变结果吗？}{不会。每个 z 都加常数 c，分子分母同时多乘 e\textasciicircum{}c，约掉后概率完全不变；代码里把打分整体加 100，概率一模一样。}
\fol{为什么说乘 2 等于温度减半？}{因为 (1,2,3) 和 (2,4,6) 概率不同，分数整体乘 2 等价于除以 T=0.5；而整体加常数不改变结果，所以 softmax 只关心差值。}
\fol{训练为什么用 log-softmax？}{概率都小于 1，连乘会下溢，取对数变成加法就不会溢出；框架把 softmax 和 log 合成一步算，又快又稳，交叉熵损失直接用它。}
\numline{别答错}{别把 softmax 说成取最大值：它给每一项都分概率，只有温度趋近 0 才退化成 argmax。}
\phantomsection\addcontentsline{toc}{section}{7\quad 概率与采样}
\lec{7}{概率与采样}
\askline{按概率采样在代码里怎么做？}
\askline{温度、Top-k、Top-p 是什么关系？}
\askline{为什么同一个问题每次答得不一样？}
\vspace{3pt}
\noindent{\qlabel{答}}\begingroup\sloppy\begin{enumerate}[leftmargin=1.6em,itemsep=2.5pt,topsep=0pt,parsep=0pt,label={\sffamily\footnotesize\color{accent}\arabic*.}]
\item {\fontsize{10}{14.8}\selectfont 模型对下一个词的输出就是一个离散分布，每个概率在 0 到 1 之间而且总和为 1；采样就是按这个分布随机选一个词，概率 0.47 的词大约一半的时候被选中，0.06 的词偶尔也会被选中。\par}
\item {\fontsize{10}{14.8}\selectfont 做法是先算累积概率，0.4666、0.6742、0.7355、0.8119、0.869、1，把 [0,1) 切成长度正好等于各个概率的几段，再抽一个 [0,1) 上均匀分布的随机数 u，看它落在哪一段。\par}
\item {\fontsize{10}{14.8}\selectfont 为什么这样就对了：u 在 [0,1) 上均匀分布，它落在某一段的可能性正好等于这一段的长度，也就是那个词的概率；u=0.3 落在第一段选「我」，u=0.7 选「学」，u=0.95 选 <eos>。\par}
\item {\fontsize{10}{14.8}\selectfont 抽得越多频率越接近概率，这就是大数定律：抽 10 次「我」是 0.3、AI 一次没抽到，100 次到 0.42，10000 次是 0.468，和 0.4666 只差千分之几，所以一次生成的结果代表不了模型的全部倾向。\par}
\item {\fontsize{10}{14.8}\selectfont 随机数其实是伪随机数，是用确定公式算出来的一串数，起点就是随机种子；种子相同序列就相同，np.random.default\_rng(0) 能让结果完全复现，换种子或换 GPU 的并行顺序结果才会变。\par}
\item {\fontsize{10}{14.8}\selectfont 贪心永远取概率最大的词，确定、稳定但容易单调重复；采样多样却可能抽到低概率的词，折中是先用温度调尖调平、只在最大的几个词里抽（Top-k、Top-p）；熵衡量分布有多分散，这个分布是 1.479、6 词均匀时最大 1.792。\par}
\end{enumerate}\endgroup
\numline{公式}{$\text{cdf}_i=\sum_{j\le i}p_j$，取 $u\sim U[0,1)$ 落段即取样；$H=-\sum_{i}p_i\ln p_i$}
\numline{数字}{累积 0.4666/0.6742/0.7355/0.8119/0.869/1；抽 10 次 0.3、100 次 0.42、10000 次 0.468；熵 1.479，ln6 = 1.792；温度 0.5 / 1 / 2}
\fol{概率高就一定会被选中吗？}{不一定。0.47 只表示大约一半的时候被选中，另一个半选的是别的词；单次生成只是一次抽样，不代表模型的全部倾向。}
\fol{评测时怎么处理这种随机性？}{同一道题生成多次再统计，抽 100 次和 10000 次的频率差别很大；固定种子能复现，但报告结果时要一并说明温度和 Top-p 这些采样参数。}
\fol{贪心是不是一定更好？}{不是。贪心确定、稳定，但同一个问题永远同一个答案、还容易陷进重复；要多样性就得采样，代价是偶尔抽到不合适的低概率词。}
\numline{别答错}{频率不等于概率：一次生成只是一次抽样，评测要同题多次生成再统计。}
\phantomsection\addcontentsline{toc}{section}{8\quad 导数与链式法则}
\lec{8}{导数与链式法则}
\askline{梯度下降里的梯度是怎么算出来的？}
\askline{链式法则在反向传播里怎么用？}
\askline{梯度消失到底是怎么发生的？}
\vspace{3pt}
\noindent{\qlabel{答}}\begingroup\sloppy\begin{enumerate}[leftmargin=1.6em,itemsep=2.5pt,topsep=0pt,parsep=0pt,label={\sffamily\footnotesize\color{accent}\arabic*.}]
\item {\fontsize{10}{14.8}\selectfont 导数是函数在一点的变化率，几何上就是切线的斜率，也可以理解成一个放大倍数：x$^2$ 在 3 处的导数是 6，意味着 x 增加一点点 Δ，f 大约增加 6Δ。\par}
\item {\fontsize{10}{14.8}\selectfont 算的办法是取很小的 h 用割线逼近：(f(3+h)-f(3))/h 在 h=0.1、0.01、0.001 时依次是 6.1、6.01、6.001，越来越接近求导公式给出的 2x=6，梯度下降每一步问的就是这个变化率。\par}
\item {\fontsize{10}{14.8}\selectfont 多变量时偏导数只让一个变量动：L=(wx+b-y)$^2$ 在 w=b=1 处 $\partial$L/$\partial$w=-16、$\partial$L/$\partial$b=-8，w 加 0.01 时损失从 16 变成 15.8404 正好约等于 -16$\times$0.01；所有偏导排成梯度，指向损失增长最快的方向。\par}
\item {\fontsize{10}{14.8}\selectfont 一层套一层时用链式法则，一串影响接起来，总的放大倍数就是每一段局部导数相乘：x$\rightarrow$u=3x+1$\rightarrow$v=u$^2$$\rightarrow$y=ln v 在 x=0.5 处是 0.16$\times$5$\times$3=2.4，数值验证也是 2.4。\par}
\item {\fontsize{10}{14.8}\selectfont 一个变量通过多条路径影响结果时，路径上的局部导数相乘、路径之间相加；第 18 讲的反向传播就是把这个规则系统地用在整张计算图上，从损失出发沿边往回乘、在分叉处把各条路径加起来。\par}
\item {\fontsize{10}{14.8}\selectfont 代价是链式法则本质是连乘：几十层里如果每层局部导数都小于 1（sigmoid 的导数最大只有 0.25），乘起来就趋近 0，前面的层几乎学不到东西；都大于 1 则爆炸，这正是激活函数、残差连接和归一化要解决的问题。\par}
\end{enumerate}\endgroup
\numline{公式}{$f'(x)=\lim_{h\to 0}\frac{f(x+h)-f(x)}{h}$；$\frac{dy}{dx}=\frac{dy}{dv}\cdot\frac{dv}{du}\cdot\frac{du}{dx}$}
\numline{数字}{割线 6.1/6.01/6.001 $\rightarrow$ 导数 6；$\partial$L/$\partial$w = -16、$\partial$L/$\partial$b = -8，w 加 0.01 损失变化 -0.1596；0.16$\times$5$\times$3 = 2.4；σ′(0) = 0.25；Llama-3-8B 80 亿参数、80 亿维梯度}
\fol{梯度下降为什么每次只走一小步？}{因为导数只是局部的，离这一点远了就不准，步子太大梯度指的方向就错了，所以要用学习率把步长压小。}
\fol{为什么 sigmoid 容易梯度消失？}{它的导数 σ(1-σ) 最大只有 0.25，几十层连乘就趋近 0；导数为 1 的 ReLU、残差连接和归一化都是为了缓解这个问题。}
\fol{推理阶段需要算导数吗？}{不需要。导数只在训练时用，推理只做前向计算，不用存中间结果去算梯度，所以推理比训练省得多。}
\numline{别答错}{同一条路径上的局部导数相乘、多条路径之间相加，这两个规则别混。}
\phantomsection\addcontentsline{toc}{section}{9\quad 三角函数与旋转}
\lec{9}{三角函数与旋转}
\askline{RoPE 为什么用旋转来编码位置？}
\askline{两个向量各转不同角度，点积还稳定吗？}
\askline{位置编码为什么要用一组不同转速？}
\vspace{3pt}
\noindent{\qlabel{答}}\begingroup\sloppy\begin{enumerate}[leftmargin=1.6em,itemsep=2.5pt,topsep=0pt,parsep=0pt,label={\sffamily\footnotesize\color{accent}\arabic*.}]
\item {\fontsize{10}{14.8}\selectfont 单位圆上角度 $\theta$ 的点是 (cos $\theta$, sin $\theta$)，sin、cos 周期 2$\pi$，cos$^2$$\theta$ + sin$^2$$\theta$ = 1；把 (a, b) 逆时针转 φ 得到 (a cos φ - b sin φ, a sin φ + b cos φ)。\par}
\item {\fontsize{10}{14.8}\selectfont 旋转是线性变换：(1, 0)$\rightarrow$(cos φ, sin φ)、(0, 1)$\rightarrow$(-sin φ, cos φ)，(a, b) 转完就是 a 乘前者加 b 乘后者；(1, 0.5) 转 30$^{\circ}$ 得 (0.616, 0.933)，长度不变。\par}
\item {\fontsize{10}{14.8}\selectfont 旋转保长度、保夹角，所以 q 和 k 一起转 50$^{\circ}$，点积仍是 -0.15；分别转 mφ 和 nφ 时，夹角只改变 (m - n)φ，点积只取决于 m - n，φ=10$^{\circ}$ 时相差 2 步的三种情况都是 0.0643。\par}
\item {\fontsize{10}{14.8}\selectfont 这正是 RoPE 的做法：第 m 个位置的 q 转 mφ、第 n 个位置的 k 转 nφ，注意力打分就只看相对距离 m - n；整句话一起往后挪几个位置，词与词之间的关系不变。\par}
\item {\fontsize{10}{14.8}\selectfont 实现上把 d 维向量两两分组，第 i 组每步转 $\theta$$_i$ = 10000\textasciicircum{}(-2i/d) 弧度；快组对近处敏感但转一圈会绕回，慢组能区分远处，d=128 时最快 6 步一圈、最慢 54,410 步一圈，像秒针分针时针搭配。\par}
\end{enumerate}\endgroup
\numline{公式}{$x'=a\cos\varphi-b\sin\varphi$，$y'=a\sin\varphi+b\cos\varphi$；$\theta_i=10000^{-2i/d}$；转一圈步数 $=\frac{2\pi}{\theta_i}$}
\numline{数字}{特殊角 0$^{\circ}$/30$^{\circ}$/90$^{\circ}$/180$^{\circ}$ $\rightarrow$ (1,0)/(0.866,0.5)/(0,1)/(-1,0)；(1,0.5) 转 30$^{\circ}$ 得 (0.616,0.933)、长度 1.1180 不变；q$\cdot$k=-0.15 同转 50$^{\circ}$ 仍是 -0.15；φ=10$^{\circ}$ 相差 2 步点积 0.0643；d=4 第 0 组 $\theta$=1、6 步一圈；d=128 第 63 组 $\theta$$\approx$0.000115、54,410 步一圈}
\fol{为什么 RoPE 只转 q 和 k？}{因为位置信息只需要进注意力打分。打分由 q$\cdot$k 决定，转 q、k 让点积只依赖相对距离；v 是被加权求和的内容，转它等于给输出也加了一个绝对角度，反而丢信息。}
\fol{为什么不能只留一组转速？}{只留快组，隔 6 个和 12 个位置的角度会绕回同一处，分不清；只留慢组，差一两个位置角度几乎不变，对近处不敏感。快慢搭配才能同时覆盖近处和远处。}
\fol{怎么验证你的 RoPE 实现是对的？}{用相对位置不变性测：q、k 同时平移一段，点积应当基本不变（只差浮点误差）；再对比 Hugging Face 的 rotary embedding，误差应在 1e-5 量级，差大就是写法错了。}
\numline{别答错}{点积不变的前提是一起转同样的角度；各转 mφ、nφ 时只保证依赖 m-n，别说成绝对值也不变。}
\phantomsection\addcontentsline{toc}{section}{10\quad 浮点数与精度}
\lec{10}{浮点数与精度}
\askline{FP16 和 BF16 到底差在哪，训练用哪个？}
\askline{低精度训练为什么会训不动？}
\askline{推理为什么敢用 FP8 甚至 4 位？}
\vspace{3pt}
\noindent{\qlabel{答}}\begingroup\sloppy\begin{enumerate}[leftmargin=1.6em,itemsep=2.5pt,topsep=0pt,parsep=0pt,label={\sffamily\footnotesize\color{accent}\arabic*.}]
\item {\fontsize{10}{14.8}\selectfont 浮点数就是二进制的科学计数法：一个数拆成符号、指数、尾数三部分，值是 (-1)\textasciicircum{}符号 $\times$ (1 + 尾数/2\textasciicircum{}m) $\times$ 2\textasciicircum{}(指数-bias)；尾数位决定精度、指数位决定范围，这两件事要分开记。\par}
\item {\fontsize{10}{14.8}\selectfont 同样 16 位，FP16 是 5 位指数加 10 位尾数，精度高但最大只到 65504；BF16 是 8 位指数加 7 位尾数，精度只有 2 到 3 位十进制，但范围和 FP32 一样到 3.4e38，训练选 BF16 是为了不溢出。\par}
\item {\fontsize{10}{14.8}\selectfont 精度不够有四种典型故障：舍入误差累积、上溢成 $\infty$（70000 在 FP16）、下溢成 0（1e-8 在 FP16）、大数吃小数，0.01 在 FP16 里加 10000 次停在 32.0，FP32 才能得到 100.003。\par}
\item {\fontsize{10}{14.8}\selectfont 所以混合精度要另存一份 FP32 主权重：前向反向用 BF16 算，更新加在 FP32 上再转回 BF16。权重 1.0 每步加 0.0001，BF16 存永远是 1.0，FP32 存才能到 1.1。\par}
\item {\fontsize{10}{14.8}\selectfont 推理没有微小更新的累积，只要前向够准，BF16、FP8 甚至 4 位都还能保持效果；位数减半，显存和读取字节减半、算力翻倍，这是推理优化最重要的手段之一，但 FP8、FP4 必须先做缩放。\par}
\end{enumerate}\endgroup
\numline{公式}{$\text{值}=(-1)^{s}\times\left(1+\frac{\text{尾数}}{2^{m}}\right)\times 2^{e-\text{bias}}$；尾数 $m$ 位时 $1$ 附近的间隔是 $2^{-m}$，相对误差约为它的一半}
\numline{数字}{FP32 8/23 位、FP16 5/10 位、BF16 8/7 位、FP8 E4M3 4/3 位、FP4 E2M1 2/1 位；最大值 3.4e38/65504/3.38e38/448/6；机器精度 1.19e-7/9.77e-4/7.81e-3；0.1 在 FP16 存 0.0999755859、BF16 存 0.1000976562；0.01 加 10000 次 FP16 停在 32.0、FP32 100.003；权重 1.0 每步加 0.0001 加 1000 次，BF16 仍是 1.0、FP32 得 1.1000}
\fol{BF16 精度更低，为什么训练反而用它？}{不是更准，是更不容易炸。BF16 误差大，但指数和 FP32 一样宽，中间结果不会一超 65504 就变 $\infty$；早年 FP16 训练要配 loss scaling 就是这个原因。}
\fol{那推理为什么不需要 FP32 主权重？}{推理只有一次前向，没有「小更新被舍掉」的问题，误差不跨步累积。权重压到 FP8 的误差是固定的，实测困惑度只涨零点几，换来的显存和带宽减半更划算。}
\fol{怎么量一次量化掉了多少精度？}{在同一份校准文本上算量化前后的困惑度，看涨幅；再抽查几层权重的相对误差。7B 模型在 WikiText-2 上 PPL 通常 5 到 6，好的量化只涨零点几，涨一个点以上就要查缩放因子。}
\numline{别答错}{别把 BF16 说成更好的 FP16，它是用精度换范围；也别把精度理解成小数点后几位，它指有效数字。}
\phantomsection\addcontentsline{toc}{section}{11\quad 计算量与 FLOPs}
\lec{11}{计算量与 FLOPs}
\askline{给你一个模型，怎么估它一次要跑多久？}
\askline{训练一个大模型要多少算力、多少卡？}
\askline{算出来 15 微秒，为什么线上要几毫秒？}
\vspace{3pt}
\noindent{\qlabel{答}}\begingroup\sloppy\begin{enumerate}[leftmargin=1.6em,itemsep=2.5pt,topsep=0pt,parsep=0pt,label={\sffamily\footnotesize\color{accent}\arabic*.}]
\item {\fontsize{10}{14.8}\selectfont 先把「算多少次」和「每秒算多少次」分开：计算量单位是 FLOPs，算力单位是 FLOPS，时间约等于 FLOPs ÷ FLOPS；H100 的 BF16 算力约 989 TFLOPS，这就是估算的基准。\par}
\item {\fontsize{10}{14.8}\selectfont 矩阵乘法 [m,k] @ [k,n] 约 2mnk FLOPs，因为每个结果格子做 k 次乘加、一次乘加按 2 次运算算；TinyLM 的 4$\times$4 逐次是 112，按 2mnk 是 128，k 很大时这点差别可以忽略。\par}
\item {\fontsize{10}{14.8}\selectfont 一个 token 过一个权重矩阵约是它参数量的 2 倍，所以整个前向约 2N；Llama-3-8B 去掉嵌入表后 75.0 亿，每 token 约 15.0 GFLOP，上下文 2048 时注意力再加 1.07 GFLOP，占 7\%。\par}
\item {\fontsize{10}{14.8}\selectfont 训练还要加反向，反向约是前向的 2 倍，算对输入的梯度和算对权重的梯度各一次矩阵乘，所以每 token 约 6N。8B 模型 15 万亿 token 是 7.2$\times$10$^2$$^3$ FLOP，1000 张 H100 按 40\% 利用率约 21 天。\par}
\item {\fontsize{10}{14.8}\selectfont 换算成时间：prefill 1000 个 token 是 15.0 TFLOP、满算力 15.2 ms，生成 1 个 token 只要 15.2 $\mu$s；实际是毫秒级，慢几百倍，因为要把 16 GB 权重从显存读一遍，算力只是故事的一半。\par}
\end{enumerate}\endgroup
\numline{公式}{$\text{时间}\approx\frac{\text{FLOPs}}{\text{FLOPS}}$；$\text{MatMul}\approx 2mnk$；推理 $2N$/token，训练 $6N$/token，训练总量 $\approx 6NT$}
\numline{数字}{H100 BF16 约 989 TFLOPS；Llama-3-8B 80.3 亿参数、去嵌入表 75.0 亿、每 token 15.0 GFLOP；上下文 2048 注意力 +1.07 GFLOP（7\%）；prefill 1000 token 15.0 TFLOP$\rightarrow$15.2 ms；生成 1 token 15.2 $\mu$s；6NT=7.2$\times$10$^2$$^3$ FLOP、1000 卡 40\% 利用率约 21 天；DeepSeek-V3 总 6710 亿、每 token 激活约 370 亿}
\fol{MoE 模型也按总参数算吗？}{不能，要按激活参数算。DeepSeek-V3 总参数 6710 亿，但每个 token 只激活约 370 亿，计算量就按 370 亿估；总参数只决定显存要装多少，两个数要分开看。}
\fol{6N 里的反向为什么正好是 2 倍前向？}{对每个权重矩阵，反向要算对输入的梯度、再算对权重的梯度，各相当于一次前向规模的矩阵乘，所以反向约 4N，加上前向 2N 就是 6N。实际还受重算和通信影响。}
\fol{估算和实测差很多，先查什么？}{先看有没有把带宽算进去：单 token decode 是带宽受限，16 GB 除以 3.35 TB/s 约 4.8 ms，和实测同量级；再查 MFU，大矩阵能到 50\% 以上，小矩阵会掉到个位数。}
\numline{别答错}{FLOPs 是运算次数、FLOPS 是每秒次数；按 MAC 计数的规格只有本讲一半，比较前先对齐。}
\chapter{神经网络基础}
{\footnotesize\sffamily\color{grayt}第 12--20 讲}\par\vspace{-4pt}
\phantomsection\addcontentsline{toc}{section}{12\quad 感知机}
\lec{12}{感知机}
\askline{感知机是什么，它为什么学不会异或？}
\askline{权重和分界线在几何上是什么关系？}
\askline{多层感知机怎么就解决异或了？}
\vspace{3pt}
\noindent{\qlabel{答}}\begingroup\sloppy\begin{enumerate}[leftmargin=1.6em,itemsep=2.5pt,topsep=0pt,parsep=0pt,label={\sffamily\footnotesize\color{accent}\arabic*.}]
\item {\fontsize{10}{14.8}\selectfont 感知机是最早的神经网络，只做两件事：给输入打一个分 z = w$\cdot$x + b，再看 z 是否大于 0，大于 0 输出 1、否则输出 0，也就是一次点积加一个偏置，再过阶跃函数。\par}
\item {\fontsize{10}{14.8}\selectfont 几何上它用一条直线把平面切成两半，分界线是 w$\cdot$x + b = 0；w 垂直于这条线并指向输出 1 的一侧，b 不改变方向、只把线平行推来推去，所以调 b 就是调门槛。\par}
\item {\fontsize{10}{14.8}\selectfont 学习规则很朴素：样本一个个看，分对了不动，分错了就让 w $\leftarrow$ w + (y - 输出)$\cdot$x、b $\leftarrow$ b + (y - 输出)；改一次能让这个样本的 z 至少涨 1，所以线性可分时一定在有限步内停下。\par}
\item {\fontsize{10}{14.8}\selectfont 异或为什么不行：要求 b $\leq$ 0、w1 + b > 0、w2 + b > 0、w1 + w2 + b $\leq$ 0；前两式相加得 w1 + w2 + 2b > 0，由 b $\leq$ 0 推出 w1 + w2 + b > 0，和最后一式矛盾。\par}
\item {\fontsize{10}{14.8}\selectfont 两层就够了：h1 算「或」把 (0,0) 隔在左下，h2 算「与非」把 (1,1) 隔在右上，第三个感知机对 (h1, h2) 算「与」；本质上是先用隐藏层换一套坐标，在新坐标里一条直线就能分开。\par}
\item {\fontsize{10}{14.8}\selectfont 还有一点容易被追问：分界线不唯一。w 和 b 同时乘 2，z 处处翻倍但正负不变，分界线一模一样；不同的起点和样本顺序也会学出不同的线，感知机只保证分对，不保证最好。\par}
\end{enumerate}\endgroup
\numline{公式}{$z=w\cdot x+b$，$y=1\ (z>0)$ 否则 $0$；错时 $w\leftarrow w+(y-\hat{y})x$，$b\leftarrow b+(y-\hat{y})$；分界线 $w\cdot x+b=0$，$w$ 垂直指向输出 $1$ 一侧}
\numline{数字}{水果 7 个点、w=(1,-1)、b=-1 全对；代码 2 轮学到 w=(5,-5)、b=0；与第 6 轮 w=(2,1)、b=-2；或第 4 轮；异或跑满 20 轮仍有错、最多分对 3/4；收敛定理由 Novikoff 1962 年证明}
\fol{学习率去哪了，为什么直接加 x？}{感知机规则里的步长被 x 吸收了，等价于学习率 1；改一次 z 会增加 |x|$^2$ + 1，所以能证收敛。实际用的时候会再乘一个学习率，方向不变，只是步子变小。}
\fol{线性可分这个条件实际常见吗？}{几乎不成立。真实数据有噪声、有重叠，硬要线性可分只能过拟合；所以现代做法是换平滑损失加梯度下降，允许分错几个点，用间隔和正则去控制，而不是要求零错误。}
\fol{怎么验证你的实现没写错？}{拿与、或、异或各 4 个点跑：与应在第 6 轮收敛到 w=(2,1)、b=-2，或第 4 轮收敛，异或跑满 20 轮仍有错且每一轮 4 个全错；三个都对上，更新式和停止条件就都对。}
\numline{别答错}{异或学不会是模型只能画一条直线的限制，不是训练不够久；分界线也不唯一。}
\phantomsection\addcontentsline{toc}{section}{13\quad 神经元与线性层}
\lec{13}{神经元与线性层}
\askline{神经元和线性层是什么关系？}
\askline{一个线性层有多少参数、怎么估？}
\askline{PyTorch 里权重形状为什么是反的？}
\vspace{3pt}
\noindent{\qlabel{答}}\begingroup\sloppy\begin{enumerate}[leftmargin=1.6em,itemsep=2.5pt,topsep=0pt,parsep=0pt,label={\sffamily\footnotesize\color{accent}\arabic*.}]
\item {\fontsize{10}{14.8}\selectfont 神经元就是「点积 + 偏置 + 激活」：z = w$\cdot$x + b，y = f(z)；能调的只有权重 w 和偏置 b，输入 x 是外面给的，所以它是网络里最小的可学习单元，训练就是不停调这两样。\par}
\item {\fontsize{10}{14.8}\selectfont 先把感知机的阶跃函数拿掉，输出才对参数的小改动有反应：阶跃下 w1 加 0.1 输出纹丝不动，去掉之后 w1 加 0.1 输出就加 0.7，因为 wi 动 Δ 时输出动 Δ $\times$ xi，这正是梯度的起点。\par}
\item {\fontsize{10}{14.8}\selectfont 把 n 个神经元的权重竖着排成 n 列，它们对同一个输入打分就是一次矩阵乘法 y = x$\cdot$W + b；W 是 din$\times$dout、b 有 dout 个，几何上是先做线性变换再平移，把输入送到新的打分空间。\par}
\item {\fontsize{10}{14.8}\selectfont 参数个数是 din$\times$dout + dout：水果 2 进 3 出是 9 个，TinyLM 的 4$\times$8 是 40 个；维度一大偏置就微不足道，这也是很多大模型干脆不带偏置的原因之一。\par}
\item {\fontsize{10}{14.8}\selectfont 输入不止一个时把样本叠成矩阵 X，一次 X$\cdot$W 就把整批算完，b 靠广播加到每一行；GPU 上一次算一大块比算很多小块快，所以推理总把多个词、多个请求叠起来送进线性层。\par}
\item {\fontsize{10}{14.8}\selectfont Llama-3-8B 每层 7 个线性层都不带偏置，一层合计 218,103,808 个参数，32 层约 70 亿，占 80.3 亿的 86.9\%；算上把向量变成词表打分的输出层就超过 93\%。\par}
\end{enumerate}\endgroup
\numline{公式}{$y=x\cdot W+b$；参数个数 $=d_{\text{in}}d_{\text{out}}+d_{\text{out}}$；$(x,1)\cdot[W;b]=x\cdot W+b$；PyTorch：$y=x\cdot W^{\mathsf{T}}+b$，$W$ 形状 $[d_{\text{out}},d_{\text{in}}]$}
\numline{数字}{水果 2$\times$3+3=9；TinyLM 4$\times$8+8=40；Llama 每层 7 个线性层：q/o 各 4096$\times$4096、k/v 各 4096$\times$1024、gate/up 各 4096$\times$14336、down 14336$\times$4096；每层 218,103,808、32 层 6,979,321,856、占 86.9\%，算上 4096$\times$128256 的输出层超 93\%；归一化层只有 26 万多参数}
\fol{偏置到底能不能省，省了会怎样？}{能省，Llama 就全省了。偏置只占 din$\times$dout 里的 dout 项，比例小到可以忽略；它的作用只是给每个神经元一个固定门槛，靠归一化和后续层基本能补回来，省下的是访存和一次加法。}
\fol{为什么说推理时间主要花在线性层？}{参数九成在线性层，decode 每生成一个 token 要把这些权重从显存读一遍。8B 模型 16 GB 权重、H100 3.35 TB/s，理论下限约 4.8 ms，差距全在搬数据。}
\fol{nn.Linear 的权重形状怎么对？}{先看权重形状：nn.Linear 存 [dout, din]，要用 x @ W.T；q\_proj 写 [4096, 4096] 看不出，但 k\_proj 写 [1024, 4096] 就能确认。}
\numline{别答错}{PyTorch 的 nn.Linear 把权重存成 [dout, din]，读权重文件先看清存法。}
\phantomsection\addcontentsline{toc}{section}{14\quad 激活函数}
\lec{14}{激活函数}
\askline{为什么神经网络一定要有激活函数？}
\askline{ReLU、GELU、SiLU 该怎么选？}
\askline{激活计算量这么小，为什么还要优化？}
\vspace{3pt}
\noindent{\qlabel{答}}\begingroup\sloppy\begin{enumerate}[leftmargin=1.6em,itemsep=2.5pt,topsep=0pt,parsep=0pt,label={\sffamily\footnotesize\color{accent}\arabic*.}]
\item {\fontsize{10}{14.8}\selectfont 因为纯线性层叠多少层都塌缩成一层：(x$\cdot$W1 + b1)$\cdot$W2 + b2 展开就是 x$\cdot$(W1W2) + (b1W2 + b2)，还是「乘一个矩阵再加一个向量」，叠十层和一层表达能力一样，多出来的层只白费参数。\par}
\item {\fontsize{10}{14.8}\selectfont 打破塌缩要靠层间插一个非线性、逐元素的函数：ReLU(z) = max(0, z) 把负数截成 0、正数原样留下，形状和长度都不变，但每个神经元变成跟着输入亮或不亮的开关，层与层就再也合并不了。\par}
\item {\fontsize{10}{14.8}\selectfont 有了 ReLU，两个神经元就能算异或：h1 = ReLU(x1 + x2)、h2 = ReLU(x1 + x2 - 1)、y = h1 - 2h2；只有 (1,1) 时 h2 才亮，减 2 倍拉回 0，且正侧是斜线，能靠梯度自己学。\par}
\item {\fontsize{10}{14.8}\selectfont 它还带来万能近似：一层 N 个 ReLU 神经元最多 N 个折点，折点越多折线越贴着曲线，N=6 拟合 sin(1.5x) 的误差就到 0.0067，不用激活时是 0.48，再多折点就几乎分不出来。\par}
\item {\fontsize{10}{14.8}\selectfont 大模型现在用平滑版：GELU(z) = z$\cdot$Φ(z)、SiLU(z) = z$\cdot$σ(z)，都是 z 乘一个 0 到 1 的软开关；负侧保留一点斜率，缓解了死亡 ReLU，Llama 用的是 SiLU 配 SwiGLU。\par}
\item {\fontsize{10}{14.8}\selectfont 代价是它访存密集：每个数只算几次却要读一次写一次，单独跑时间几乎全花在读写上；所以推理引擎把它融合进前后的算子，在数据还在芯片上时顺手算掉，省掉一次显存往返。\par}
\end{enumerate}\endgroup
\numline{公式}{$\operatorname{ReLU}(z)=\max(0,z)$；$\operatorname{GELU}(z)=z\,\Phi(z)$；$\operatorname{SiLU}(z)=z\,\sigma(z)$，$\sigma(z)=\frac{1}{1+e^{-z}}$；$h_1=\operatorname{ReLU}(x_1+x_2)$、$h_2=\operatorname{ReLU}(x_1+x_2-1)$、$y=h_1-2h_2$}
\numline{数字}{N=6 拟合 sin(1.5x) 误差 0.0067，不加激活 0.48；z=-1 处 ReLU 0、GELU -0.159、SiLU -0.269；最低点 GELU z$\approx$-0.75/f=-0.17、SiLU z$\approx$-1.28/f=-0.28；Llama-3-8B 每个词每层 14336 次 SiLU，对比同层 3$\times$4096$\times$14336$\approx$1.76 亿次乘加；Sigmoid 斜率最大 0.25}
\fol{ReLU 在 0 处不可导，反向怎么办？}{框架里约定那一点的斜率取 0，实际中 z 恰好等于 0 的概率极低，不影响训练。真正要小心的是整片输出都为负的神经元，梯度恒为 0，权重再也改不动，也就是死亡 ReLU。}
\fol{GELU 精确版和近似版能混用吗？}{不能混。两者数值差很小，但训练用哪个推理就得用哪个，否则输出有细微偏差。看配置里的 gelu、gelu\_new、gelu\_pytorch\_tanh 就知道是哪种，换实现时先对齐。}
\fol{怎么判断激活函数是不是瓶颈？}{看 profiler 里这个 kernel 的时间和它搬的字节：FLOPs 占比不到万分之一却占几个百分点的时间，就是访存受限；这时把激活和前后的矩阵乘、归一化融合成一个 kernel，收益最直接。}
\numline{别答错}{softmax 不是这里说的激活函数：它要跨整个向量归一化，每个输出依赖所有输入，不是逐元素。}
\phantomsection\addcontentsline{toc}{section}{15\quad 多层感知机}
\lec{15}{多层感知机}
\askline{为什么深网络比宽网络更省参数？}
\askline{训练好的 MLP 每层到底在干什么？}
\askline{大模型里的 FFN 和 MLP 是一回事吗？}
\vspace{3pt}
\noindent{\qlabel{答}}\begingroup\sloppy\begin{enumerate}[leftmargin=1.6em,itemsep=2.5pt,topsep=0pt,parsep=0pt,label={\sffamily\footnotesize\color{accent}\arabic*.}]
\item {\fontsize{10}{14.8}\selectfont MLP 就是把「线性层 + 激活函数」一层层串起来，最后接一个不加激活的线性层输出：h1 = f(x$\cdot$W1 + b1)、h2 = f(h1$\cdot$W2 + b2)、y = h2$\cdot$W3 + b3；数据只往前流，所以也叫前馈网络。\par}
\item {\fontsize{10}{14.8}\selectfont 拆开一个训练好的 1$\rightarrow$8$\rightarrow$8$\rightarrow$1 网络看：第 1 层每个神经元只把输入折一次，折点在 x = -bj/wj，8 个折点全挤在 -0.74 到 1.41 之间；第 2 层把这些折线加权组合，在穿过 0 的地方再折一次。\par}
\item {\fontsize{10}{14.8}\selectfont 关键就在「再折一次」：新拐角不受第 1 层折点位置限制，可以落在任何地方，所以 8 个折点最后折出 17 段折线，单靠第 1 层最多 9 段，多出来的拐角全是第 2 层给的。\par}
\item {\fontsize{10}{14.8}\selectfont 锯齿网络把这个推到极致：每层只用 2 个神经元，t(u) = 2$\cdot$ReLU(u) - 4$\cdot$ReLU(u - 0.5)，k 层就折出 2\textasciicircum{}k 段，从 2 到 64 段只用 12 个神经元；单层要做到 64 段至少要 63 个。\par}
\item {\fontsize{10}{14.8}\selectfont 这就是深度让段数指数增长、宽度只能线性增长：后面的层可以反复复用前面层算好的东西，像先认笔画、再认字、再认词；但深网络以前训不好，要等 ReLU、归一化和残差连接出现。\par}
\item {\fontsize{10}{14.8}\selectfont 大模型的 FFN 就是这种 MLP：Llama-3-8B 是 4096$\rightarrow$14336$\rightarrow$4096，隐藏层是输入的 3.5 倍，用 SiLU 加 SwiGLU 门控所以有三个权重矩阵，占每层线性层的 81\%。\par}
\end{enumerate}\endgroup
\numline{公式}{$h_1=f(x\cdot W_1+b_1)$；$h_2=f(h_1\cdot W_2+b_2)$；$y=h_2\cdot W_3+b_3$；$t(u)=2\operatorname{ReLU}(u)-4\operatorname{ReLU}(u-0.5)$，$k$ 层折出 $2^k$ 段}
\numline{数字}{1$\rightarrow$8$\rightarrow$8$\rightarrow$1 参数 97（16+72+9）；输出 17 段、MSE 0.00017；x=1 时 y=0.998 vs sin(1.5)=0.997；第 1 层 8 个折点 -0.74\textasciitilde{}1.41；第 2 层只有 6 个神经元有输出；锯齿 k=1..6 段数 2/4/8/16/32/64、神经元 2\textasciitilde{}12 个；FFN 3$\times$4096$\times$14336=176,160,768，占 218,103,808 的 81\%}
\fol{关掉第 2 层的神经元会怎样？}{第 2 层每个神经元负责曲线的一段：关掉 x=1 时输出 1.25 的那个，中间波峰塌掉，MSE 从 0.00017 涨到 0.25；关掉只在 x<-1 亮的那个，左边波谷就没了。}
\fol{为什么不是把一层做得很宽？}{宽度只线性增加折点，深度指数增加。同样 64 段，深 6 层每层 2 个只要 12 个神经元，单层要 63 个；而且深层能复用中间表示，参数效率高，所以大模型叠 32 层而不是把一层做到几万维。}
\fol{「32 层」到底数的是什么？}{数的是 Transformer 块，每个块里注意力加 FFN 又有七八个线性层。所以「层」在论文里不统一，1$\rightarrow$8$\rightarrow$8$\rightarrow$1 可以数成 2 层、3 层或 4 层，读资料先确认它数的是权重矩阵还是块。}
\numline{别答错}{「几层」数法不统一：可以数隐藏层、数权重矩阵，大模型的 32 层数的是 Transformer 块。}
\phantomsection\addcontentsline{toc}{section}{16\quad 损失函数}
\lec{16}{损失函数}
\askline{分类为什么用交叉熵，不用平方误差？}
\askline{困惑度到底是什么，怎么算？}
\askline{训练日志里 loss 从 11 开始正常吗？}
\vspace{3pt}
\noindent{\qlabel{答}}\begingroup\sloppy\begin{enumerate}[leftmargin=1.6em,itemsep=2.5pt,topsep=0pt,parsep=0pt,label={\sffamily\footnotesize\color{accent}\arabic*.}]
\item {\fontsize{10}{14.8}\selectfont 损失函数是把「错了多少」量成一个数，训练就是让它变小。拟合一个数用均方误差：先把误差平方再求平均，误差不会正负抵消，大错罚得重，而且处处光滑，适合做梯度下降。\par}
\item {\fontsize{10}{14.8}\selectfont 输出概率时用交叉熵：损失 = -ln p，p 是模型给正确答案的概率；可以理解成看到答案时有多吃惊。p=0.99 时 0.010，p=0.01 时 4.605，p 每缩小十倍，损失加 ln 10 $\approx$ 2.3。\par}
\item {\fontsize{10}{14.8}\selectfont 不用平方误差的原因是修正信号：交叉熵配 Sigmoid，对 logit 的斜率恰好是 p - 1，p=0.01 时是 -0.99，错得越离谱推得越用力；平方误差多乘了一个 p(1-p)，同样情况只剩 -0.0196，是前者的五十分之一。\par}
\item {\fontsize{10}{14.8}\selectfont 语言模型是从整个词表里挑一个词：最后一层给 V 个 logit，softmax 成概率，损失仍是 -ln(正确词的概率)；二分类只是 V=2 的特例，修正信号 p - y 也是同一个式子的推广。\par}
\item {\fontsize{10}{14.8}\selectfont 一句话的损失是每个位置交叉熵的平均，取指数就是困惑度：TinyLM 四个位置 1.572、2.323、1.601、2.421，平均 1.979，PPL 7.24，比均匀乱猜的 ln 6 = 1.792、PPL 6 还差，因为权重还是随机的。\par}
\item {\fontsize{10}{14.8}\selectfont 到了大模型，预训练目标就这一个：让下一个词的平均交叉熵尽量小。Llama-3 词表 128256，ln V $\approx$ 11.76，所以日志里 loss 从 11 左右往下掉；远大于这个数通常说明初始化或代码有问题。\par}
\end{enumerate}\endgroup
\numline{公式}{$\text{MSE}=\frac{1}{n}\sum_{i}(\hat{y}_i-y_i)^2$；$\text{CE}=-\ln p$；$\text{PPL}=e^{\text{平均交叉熵}}$；softmax：$p_i=\frac{e^{z_i}}{\sum_{j}e^{z_j}}$；修正信号 $\frac{\partial L}{\partial z_i}=p_i-y_i$}
\numline{数字}{MSE 例 0.125；p=0.99/0.9/0.5/0.1/0.01 $\rightarrow$ CE 0.010/0.105/0.693/2.303/4.605；修正信号 p-1 = -0.010/-0.100/-0.500/-0.900/-0.990，平方误差只有 -0.0002/-0.018/-0.25/-0.162/-0.0196；TinyLM 4 个位置 1.572/2.323/1.601/2.421 $\rightarrow$ 平均 1.979、PPL 7.24；均匀猜 ln 6=1.792、PPL 6；ln 128256$\approx$11.76；8192$\times$128256 的 logits 按 4 字节约 3.91 GiB；量化后 7B 模型 WikiText-2 上 PPL 约 5\textasciitilde{}6}
\fol{困惑度为什么要先平均再取指数？}{因为交叉熵是对数尺度，先平均才是平均每步的损失，取指数才还原成「等效候选词数」。先把每个位置的 PPL 算出来再平均，等于对几何平均做算术平均，得到的是另一个更大的数。}
\fol{准确率也能当指标，为什么不用它训练？}{准确率只看最大概率那个词对不对，正确答案从 0.2 涨到 0.4 只要还不是最大就不变，没有梯度；交叉熵对「分到更多概率」有反应，能给出连续的方向，所以训练用它、评估才看准确率。}
\fol{为什么 logits 要分块算？}{8192$\times$128256 的 logit 矩阵用 4 字节存要 3.91 GiB，还只是一个中间结果；所以框架把算 logit、softmax、取正确词概率、求对数融合成一个算子并分块计算。}
\numline{别答错}{困惑度要先对交叉熵求平均再取指数；训练用 ln（奈特），信息论用 log$_2$（比特），差一个常数倍。}
\phantomsection\addcontentsline{toc}{section}{17\quad 梯度下降}
\lec{17}{梯度下降}
\askline{梯度下降到底在做什么？一句话讲清}
\askline{学习率调大了会怎样，上限怎么定？}
\askline{为什么梯度指向上坡，更新却要减号？}
\vspace{3pt}
\noindent{\qlabel{答}}\begingroup\sloppy\begin{enumerate}[leftmargin=1.6em,itemsep=2.5pt,topsep=0pt,parsep=0pt,label={\sffamily\footnotesize\color{accent}\arabic*.}]
\item {\fontsize{10}{14.8}\selectfont 训练就是把损失从预测的函数换成参数的函数，梯度下降则是对每个参数求偏导、排成梯度$\nabla$L，再沿它的反方向走η倍，一步只改一点点，反复迭代就走到损失最小的那组参数。\par}
\item {\fontsize{10}{14.8}\selectfont 单参数时斜率就是导数，步长等于学习率乘斜率，越靠近谷底斜率越小、步子自动变小，所以不会冲过头；那四个点从w=0出发斜率是-11.2，走一步w到0.56，损失从9.01降到3.84。\par}
\item {\fontsize{10}{14.8}\selectfont 学习率决定能不能收敛：对这条抛物线每步到谷底的距离乘以(1-7η)，η小于1/7时单侧下降，1/7到2/7之间来回震荡但收敛，超过2/7约0.286就发散；实测0.05和0.2都在20步内到1.6，0.3跑到-9.164。\par}
\item {\fontsize{10}{14.8}\selectfont 参数多了以后每个参数按自己的偏导数走，排成的梯度指向损失上升最快的方向、且垂直等高线；从(-0.5,-2)出发梯度是(-24,-12.5)，η=0.1第一步就到(1.9,-0.75)，损失44.625降到2.12。\par}
\item {\fontsize{10}{14.8}\selectfont 但路线不是直奔谷底：这个碗是扁的，先垂直冲向谷底线，再沿谷慢慢挪，η超过0.12就在陡方向震荡、超过约0.24发散，比单参数的0.286还小，因为两个参数一起动时最陡方向更弯。\par}
\item {\fontsize{10}{14.8}\selectfont 到80亿参数上公式一个字不用改，$\theta$和$\nabla$L都是80亿维向量；代价是训练要为每个参数多存一份同样大的梯度，BF16下权重16 GB、梯度16 GB，这还没算优化器状态和激活值。\par}
\end{enumerate}\endgroup
\numline{公式}{$\theta\leftarrow\theta-\eta\nabla L$；单参数收敛条件 $\lvert 1-7\eta\rvert<1$，即 $\eta<\frac{2}{7}\approx 0.286$}
\numline{数字}{w=1.6、b=1.1、L=0.05；斜率-11.2$\rightarrow$-7.28$\rightarrow$-4.732$\rightarrow$-3.076；损失9.01$\rightarrow$3.84$\rightarrow$1.65$\rightarrow$0.73；η=0.3走20步到-9.164；单参数上限2/7$\approx$0.286，双参数0.12震荡、0.24发散；8B权重与梯度各16 GB}
\fol{学习率是不是越小越好？}{不是。太小不仅慢，在真实网络里还可能卡在平坦区域出不来；要在不发散的前提下尽量大，这正是调学习率最难的地方。这里η=0.05要20步，0.2更快但会震荡。}
\fol{扁碗里梯度下降该怎么办？}{两条路：一是把数据摆正，输入居中、尺度统一，就是归一化要干的事；二是换优化器，Momentum压震荡、Adam按梯度平方给每个参数定步长，那个例子里SGD要228步、Momentum只要39步。}
\fol{你怎么判断学习率调对了？}{看损失曲线：前期降得太慢说明偏小，在谷底附近来回跳或出现loss spike说明偏大，直接发散就是超过上限；再配合梯度范数，正常训练它应该平稳而不是忽大忽小。}
\numline{别答错}{梯度指向上坡，更新必须减去它，写成加号会越训越差；求导对象是参数，不是输入。}
\phantomsection\addcontentsline{toc}{section}{18\quad 反向传播}
\lec{18}{反向传播}
\askline{80 亿个参数怎么一次算出梯度？}
\askline{反向传播为什么只多花两倍代价？}
\askline{训练为什么要存那么多激活值？}
\vspace{3pt}
\noindent{\qlabel{答}}\begingroup\sloppy\begin{enumerate}[leftmargin=1.6em,itemsep=2.5pt,topsep=0pt,parsep=0pt,label={\sffamily\footnotesize\color{accent}\arabic*.}]
\item {\fontsize{10}{14.8}\selectfont 反向传播是在计算图上做一次前向加一次反向：前向从左往右算出每个节点的值并留下中间结果，反向从损失出发往回走，每个节点做一次「上游梯度 $\times$ 局部导数」，把梯度传给自己的输入。\par}
\item {\fontsize{10}{14.8}\selectfont 省下来的是重复乘法。逐个参数从头乘整条路径，越靠前路径越长而且大量乘法重复；反向换了顺序后每个节点只算一次，从L走回ŷ那一段被所有下游参数共用。\par}
\item {\fontsize{10}{14.8}\selectfont 一个节点被用几次，梯度就要加几次。PyTorch的backward()是累加而不是覆盖，所以每个训练步更新前必须optimizer.zero\_grad()清零，否则这一步的梯度会和上一步叠在一起。\par}
\item {\fontsize{10}{14.8}\selectfont 一整层还是矩阵乘法。线性层Y=XW的反向是两次矩阵乘法，$\partial$L/$\partial$W=Xᵀ$\cdot$dY、$\partial$L/$\partial$X=dY$\cdot$Wᵀ，梯度的形状永远和对应变量相同；前向一次乘、反向两次乘，决定了训练的计算量。\par}
\item {\fontsize{10}{14.8}\selectfont 代价是算力和显存。每token约6N次浮点运算，前向2N、反向4N，Llama-3-8B每token合计4.82e10，是推理的3倍；前向的激活值要留到反向用完，32层FFN那个14336维向量在batch 8192下就要约7 GiB。\par}
\item {\fontsize{10}{14.8}\selectfont 省显存的办法是激活重计算：前向只存每层输入，其余中间结果扔掉，反向到这一层再从输入重算一遍，用大约多一次前向的计算换大量显存；推理不需要反向，所以显存和算力都远小于训练。\par}
\end{enumerate}\endgroup
\numline{公式}{$\frac{\partial L}{\partial W}=X^{\mathsf{T}}dY$；$\frac{\partial L}{\partial X}=dY\,W^{\mathsf{T}}$；训练 $\text{FLOPs}\approx 6N=2N+4N$}
\numline{数字}{小网络L=0.04、梯度(-0.9,-0.6,-0.4,-0.4)；Llama-3-8B每token前向1.61e10、反向3.21e10、合计4.82e10；FFN 14336$\times$2字节$\times$8192$\times$32$\approx$7.0 GiB}
\fol{数值梯度能直接拿来训练吗？}{不能。每核对一个参数就要多做一次完整前向，4个参数要5次前向，80亿个参数要80亿零1次；反向传播只要一次前向加一次反向，数值梯度只用来查手写反向有没有错。}
\fol{死亡 ReLU 在反向里怎么出现？}{把b1改成-2让z1为负，前向ReLU输出0，反向时梯度被截成0，w1、b1拿不到任何修正信号，w2的梯度也因为乘的h是0而变成0，只剩b2还能动。}
\fol{反向是把网络倒过来跑吗？}{不是。前向算的是值，反向算的是「这个值变一点、损失变多少」，是从损失出发求对每个中间值的导数，不是把数据从输出传回输入。}
\numline{别答错}{反向传播只负责算梯度，更新参数是优化器的事；把两件事说成一件事就会被抓。}
\phantomsection\addcontentsline{toc}{section}{19\quad 优化器}
\lec{19}{优化器}
\askline{SGD、Momentum、Adam 差在哪？}
\askline{为什么大模型训练都用 AdamW？}
\askline{Adam 的显存开销到底有多大？}
\vspace{3pt}
\noindent{\qlabel{答}}\begingroup\sloppy\begin{enumerate}[leftmargin=1.6em,itemsep=2.5pt,topsep=0pt,parsep=0pt,label={\sffamily\footnotesize\color{accent}\arabic*.}]
\item {\fontsize{10}{14.8}\selectfont 优化器决定拿到梯度后怎么改参数。同一个问题只把x整体加2，最陡和最平方向的弯曲就从相差14倍变成166倍，学习率上限从0.238降到0.069，窄谷里SGD被陡方向卡住、平方向又走得极慢。\par}
\item {\fontsize{10}{14.8}\selectfont SGD每步只用一小批数据估梯度，$\theta$$\leftarrow$$\theta$-η$\cdot$g；噪声平均下来方向是对的，但η=0.03走300步损失还停在0.1267，开到0.065也要228步，而原来的数据用0.1只要39步。\par}
\item {\fontsize{10}{14.8}\selectfont Momentum把过去的梯度累成速度v$\leftarrow$$\beta$v+g，$\beta$=0.9时一致方向的速度最终到g/(1-$\beta$)=10g，相当于把学习率放大10倍；来回震荡的方向正负抵消被压下去，同样的η=0.03只用39步就到(1.6,-2.1)。\par}
\item {\fontsize{10}{14.8}\selectfont Adam再给每个参数单独定步长：m是梯度的滑动平均、v是梯度平方的滑动平均，更新时用m̂/$\surd$v̂定标，梯度大就缩、梯度小就放，每个参数都走大约η那么远，第一步两个参数都正好走0.3。\par}
\item {\fontsize{10}{14.8}\selectfont 代价是状态：混合精度加Adam每个参数要存BF16权重2字节加BF16梯度2字节加FP32主权重4字节加m、v各4字节共16字节，8B约128.5 GB，而推理只要BF16权重16.1 GB，80 GB卡放不下。\par}
\item {\fontsize{10}{14.8}\selectfont 工程上用AdamW，把权重衰减和梯度更新分开做更稳定；还要注意$\beta$2常从0.999调到0.95左右，检查点除权重还得存m和v，否则恢复训练时Adam从头积累，前几步会走得很乱。\par}
\end{enumerate}\endgroup
\numline{公式}{Momentum：$v\leftarrow\beta v+g$，$\theta\leftarrow\theta-\eta v$；Adam：$\theta\leftarrow\theta-\eta\frac{\hat{m}}{\sqrt{\hat{v}}+\varepsilon}$，$\hat{m}=\frac{m}{1-0.9^{t}}$、$\hat{v}=\frac{v}{1-0.999^{t}}$}
\numline{数字}{弯曲14$\rightarrow$166倍；η上限0.238$\rightarrow$0.069；SGD 0.03走300步L=0.1267；Momentum 39步、Adam 93步；第一步梯度(-28,-6.5)$\rightarrow$SGD (0.84,0.195)、Adam (0.3,0.3)；16字节/参数、128.5 GB vs 推理16.1 GB}
\fol{SGD 和 Adam 的 η 能比吗？}{不能。SGD的步长是η乘梯度，梯度大步子就大；Adam的步长约等于η，和梯度大小关系不大；大模型用Adam时学习率通常在1e-4量级，比玩具例子小得多。}
\fol{Adam 在噪声下为什么更稳？}{它按梯度的典型大小给步长定标，每步只走η左右，单样本噪声再大也不会把参数甩出去；实测200到300步的平均损失约0.1，三种学习率下都差不多。}
\fol{Momentum 什么时候会变差？}{梯度噪声大时。它把带噪声的梯度累成速度相当于放大10倍，单样本梯度又比整批平均大得多，η=0.03下会在山谷两壁越摆越大，降到0.005才重新稳定。}
\numline{别答错}{不同优化器的学习率不能直接比大小；Adam 步长约等于 η，SGD 是 η$\times$梯度。}
\phantomsection\addcontentsline{toc}{section}{20\quad 训练循环}
\lec{20}{训练循环}
\askline{一个训练 step 到底做了哪几件事？}
\askline{batch size 调大调小分别会怎样？}
\askline{为什么要 warmup 加余弦衰减？}
\vspace{3pt}
\noindent{\qlabel{答}}\begingroup\sloppy\begin{enumerate}[leftmargin=1.6em,itemsep=2.5pt,topsep=0pt,parsep=0pt,label={\sffamily\footnotesize\color{accent}\arabic*.}]
\item {\fontsize{10}{14.8}\selectfont 训练的骨架是两层循环：外层每个epoch先把数据打乱，内层每个step取一个batch做五件事，取批、前向、损失、反向、更新；PyTorch里更新前还要optimizer.zero\_grad()，因为梯度是累加的。\par}
\item {\fontsize{10}{14.8}\selectfont batch是一个step用到的样本数B，step是处理一个batch并更新一次参数，epoch是把整份数据过一遍；64个样本、B=8时每个epoch有8个step，训练400个epoch一共更新3200次。\par}
\item {\fontsize{10}{14.8}\selectfont batch两端都不好：B=64每个epoch只更新1次、400步后损失0.0102，B=1每epoch更新64次前期最快却停不下来，第400个epoch反而回到0.0288；B=8兼顾两边，最后是0.0067。\par}
\item {\fontsize{10}{14.8}\selectfont 损失有地板：数据带标准差0.1的噪声，噪声自身的MSE是0.0101，完美学会sin(1.5x)也降不下去；B=8的0.0067比地板还低，说明在记这64个点的噪声，必须留验证集来看。\par}
\item {\fontsize{10}{14.8}\selectfont 学习率太大会在谷底来回跨，所以用预热加余弦衰减：前20个epoch线性升到0.01，再在400个epoch内余弦降到0；第1个epoch是0.0005、第20个0.00994、第201个正好一半，最后损失稳在0.0084。\par}
\item {\fontsize{10}{14.8}\selectfont 到大模型里循环一字不改，只是按token计batch、一步几百万token，一张卡放不下就靠梯度累积先不清零、累加几次再更新；总计算约6ND，8B乘15T约7.23e23 FLOPs，按40\%利用率约50.7万GPU小时。\par}
\end{enumerate}\endgroup
\numline{公式}{$\text{lr}=\text{base}\cdot\min\left(1,\frac{\text{epoch}+1}{\text{warmup}}\right)\cdot 0.5\cdot\left(1+\cos\frac{\pi\cdot\text{epoch}}{\text{total}}\right)$；总计算 $\approx 6ND$}
\numline{数字}{64样本、97参数、400 epoch、3200步；噪声地板0.0101；B=64 0.0102、B=8 0.0067、B=1 0.0288；调度值0.0005/0.00499/0.00994/0.00854/0.005/0.00146；6ND=7.23e23、5.07e5 GPU小时}
\fol{B=8 的损失比噪声还低，是好事吗？}{不是，是过拟合。0.0067低于噪声地板0.0101，说明网络除了sin(1.5x)还记住了这64个点各自的噪声；判断要看没参与训练的验证集。}
\fol{batch 变大，学习率要不要改？}{要一起调。batch大梯度噪声小，一般可以也用更大的学习率；只改一个不改另一个，结论可能完全相反，这是调参里最常见的坑。}
\fol{损失曲线突然上跳怎么办？}{这叫loss spike，通常是学习率太大或这批数据有问题。工程上是回退到上一个检查点，调低学习率或跳过坏数据重跑，训练时盯得最紧的就是这条曲线。}
\numline{别答错}{横轴是 step 还是 epoch 会得出相反结论；大模型日志的横轴通常是训过的 token 数。}
\chapter{从文本到向量}
{\footnotesize\sffamily\color{grayt}第 21--24 讲}\par\vspace{-4pt}
\phantomsection\addcontentsline{toc}{section}{21\quad 分词与BPE}
\lec{21}{分词与BPE}
\askline{文字是怎么变成模型能算的数字的？}
\askline{BPE 具体怎么学出词表和规则？}
\askline{词表大一点好还是小一点好？}
\vspace{3pt}
\noindent{\qlabel{答}}\begingroup\sloppy\begin{enumerate}[leftmargin=1.6em,itemsep=2.5pt,topsep=0pt,parsep=0pt,label={\sffamily\footnotesize\color{accent}\arabic*.}]
\item {\fontsize{10}{14.8}\selectfont 模型只能处理数字，所以文字先切成token再换成编号。按字切词表小、什么字都能切，但序列长，计算量和显存都随token数涨；按词切含义完整，但词表爆炸、未登录词没法处理，中文还没空格分词。\par}
\item {\fontsize{10}{14.8}\selectfont 实际用的是折中的子词分词：常见的字组合成一个token，不常见的拆开。BPE只有两步反复执行，数语料里所有相邻两个token各出现多少次，把次数最多的那一对合并成新token并记下规则。\par}
\item {\fontsize{10}{14.8}\selectfont 拿那段小语料看：初始11个单字、27个token，4次合并依次是「推理」4次、「AI」3次、「我爱」2次、「学AI」2次，词表变成15个、语料缩短到16个token；次数相同时选语料里最先出现的那一对。\par}
\item {\fontsize{10}{14.8}\selectfont 切新句子先拆单字，再按学到的顺序逐条应用合并规则：「你爱学AI推理」切成你、爱、学AI、推理四个token，编号是1、14、11，其中「你」不在词表里，这正是按字按词都会遇到的未登录问题。\par}
\item {\fontsize{10}{14.8}\selectfont 真实分词器从256个UTF-8字节起步，任何文字最坏退回字节，永不OOV。英文占1个字节、常用汉字占3个字节，常见汉字很快被合并成一个token，这就是GPT、Llama用的字节级BPE。\par}
\item {\fontsize{10}{14.8}\selectfont 词表大小是权衡：Llama-3的词表128256、隐藏4096，嵌入表加输出层约10.5亿参数、占模型13\%；但词表越大同样文字切出的token越少，而上下文长度、KV Cache、API计费全按token算。\par}
\end{enumerate}\endgroup
\numline{公式}{词表 $=256+$ 合并次数；嵌入表与输出层共 $2Vd$ 个参数}
\numline{数字}{27 token$\rightarrow$16 token；11单字加4次合并=15；合并次数4/3/2/2；词表128256、隐藏4096、嵌入加输出1.05B参数占13\%；汉字3字节、英文1字节}
\fol{为什么要从字节开始，不从字符？}{字符表覆盖不了所有文字，没出现过的字就没有编号。字节只有256种取值，任何文字、符号、表情都能切成词表里的token，最坏退回字节，永远不OOV。}
\fol{分词器可以随便换吗？}{不能。分词器和模型是绑定的，换一个分词器同样的编号对应的就是完全不同的token，输出会变成乱码；训练用的是哪个，推理就必须用同一个。}
\fol{BPE 切出来的结果像词吗？}{不一定。BPE不懂语法，只看哪些字常挨在一起，小语料里「我爱」「学AI」都被合并了；语料大以后真正常用的组合优先合并，结果才大体像词。}
\numline{别答错}{token 不等于字或词，token 数要先说清是哪个分词器切的；编号只是位置，没有大小含义。}
\phantomsection\addcontentsline{toc}{section}{22\quad 词嵌入}
\lec{22}{词嵌入}
\askline{为什么不能把 token 编号直接喂给模型？}
\askline{嵌入层到底是什么运算？}
\askline{嵌入里的「含义」是哪来的？}
\vspace{3pt}
\noindent{\qlabel{答}}\begingroup\sloppy\begin{enumerate}[leftmargin=1.6em,itemsep=2.5pt,topsep=0pt,parsep=0pt,label={\sffamily\footnotesize\color{accent}\arabic*.}]
\item {\fontsize{10}{14.8}\selectfont 编号自带大小和距离：爱是1、推理是11、学AI是14，看起来推理离学AI很近、离爱很远，这些关系纯粹是分词编号的顺序碰巧造成的；拿编号直接做乘加，模型就会把虚假关系当真。\par}
\item {\fontsize{10}{14.8}\selectfont 独热向量倒是让所有token一律平等，但平等得过头：任意两个不同独热向量的点积和余弦相似度都是0，苹果和香蕉、苹果和推理一样远；而且Llama-3词表128256维几乎全是0，太长。\par}
\item {\fontsize{10}{14.8}\selectfont 做法是准备一张V行d列的嵌入表E，第i行就是编号i的向量，用编号取出对应那一行；一句话s个token就得到s$\times$d的矩阵X，这就是模型真正的输入，后面所有计算都从这里开始。\par}
\item {\fontsize{10}{14.8}\selectfont 查表本身就是矩阵乘法：onehot(id)$\cdot$E只剩乘1那一项，结果正好是E的第id行；所以嵌入层就是一个输入为独热、没有偏置的线性层，表里每个数都能用反向传播和梯度下降训练。\par}
\item {\fontsize{10}{14.8}\selectfont 反向时梯度只加回用到的那几行，一个batch里没出现的token这一步收不到梯度、原样不动；这也说明嵌入表虽然大，单步更新的行很少，是多卡时按词表切分的依据。\par}
\item {\fontsize{10}{14.8}\selectfont 含义是训练出来的。小语料上训练前和苹果最像的是学0.67、香蕉0.33，毫无道理；训练后变成香蕉0.99、鱼0.97、AI0.93，因为它们在句子里作用相同，都是句末的宾语，前文是吃、学、喜欢。\par}
\end{enumerate}\endgroup
\numline{公式}{$x=E[\text{id}]$；$\text{onehot}(\text{id})\cdot E= E$ 的第 $\text{id}$ 行}
\numline{数字}{TinyLM 6$\times$4；小语料14词$\times$8维；苹果$\rightarrow$训练前学0.67、香蕉0.33、你0.27，训练后香蕉0.99、鱼0.97、AI 0.93；Llama-3-8B嵌入表128256$\times$4096=525M参数、BF16 0.98 GiB、每token 8 KiB}
\fol{同一个词在不同句子里嵌入会变吗？}{查表查出来永远是同一行，苹果在「我吃苹果」和「苹果发布新手机」里一样；区分这两种意思是后面注意力层的工作，它按上下文改造向量。}
\fol{嵌入的某一维代表什么含义？}{单独一维一般没有明确意思。所谓相近指的是整个向量的方向和距离，不是第3维代表动物；含义分散在很多维的组合里，训练只关心整体效果。}
\fol{嵌入表那么大，能省吗？}{可以按词表切到多张卡上做张量并行；有的模型让输出层和嵌入表共用一份参数，叫权重绑定，GPT-2就这么做，Llama-3-8B没绑，两张表合计10.5亿参数。}
\numline{别答错}{嵌入里既没有位置信息也没有上下文；同一个词永远查同一行，别说成「带语义的输入」。}
\phantomsection\addcontentsline{toc}{section}{23\quad 预测下一个token}
\lec{23}{预测下一个token}
\askline{大模型预训练到底在学什么？}
\askline{为什么说这是自监督学习？}
\askline{只猜下一个词怎么学会推理？}
\vspace{3pt}
\noindent{\qlabel{答}}\begingroup\sloppy\begin{enumerate}[leftmargin=1.6em,itemsep=2.5pt,topsep=0pt,parsep=0pt,label={\sffamily\footnotesize\color{accent}\arabic*.}]
\item {\fontsize{10}{14.8}\selectfont 预训练只做一件事：看前面所有token，给出下一个token的概率分布P(下一个|前文)。输出不是一个词，而是词表上每个token各一个概率、加起来等于1，最后一层打分出logit再softmax。\par}
\item {\fontsize{10}{14.8}\selectfont 一段文本有多少token就有多少道题，答案就是文本自身，所以叫自监督学习；「我学AI」补上开头的句号就能拆成4道题，互联网几万亿token就是几万亿道现成带答案的练习题。\par}
\item {\fontsize{10}{14.8}\selectfont 损失是各位置交叉熵的平均-ln p，训练让模型给真实出现的下一个词分配尽量高的概率；小语料上96道题，模型平均交叉熵0.6681，正好等于理论下限，困惑度1.951。\par}
\item {\fontsize{10}{14.8}\selectfont 下限不是0，因为语言本身有不确定性：「我吃」后面是苹果还是香蕉只能五五开，每次都要付-ln0.5约0.69；模型学到的几乎就是语料频率，苹果0.501和香蕉0.499对应语料里的2/4和2/4。\par}
\item {\fontsize{10}{14.8}\selectfont 训练时一次前向就把所有位置的题一起做，靠因果掩码保证预测第i个位置时看不到第i个之后的词；推理时只关心最后一个位置的分布，但前面的token仍要完整算一遍，这就是prefill。\par}
\item {\fontsize{10}{14.8}\selectfont 为什么这目标这么厉害：「法国的首都是」要给巴黎高概率得有事实，「37+25=」得会算术，花瓶打碎了「他很」得懂因果；任何一种能力缺失都会在某类文本上猜不准、损失降不下去，规模足够大就把这些能力逼出来。\par}
\end{enumerate}\endgroup
\numline{公式}{$\text{loss}=-\frac{1}{T}\sum_{t}\ln p(y_t)$；$\text{perplexity}=\exp(\text{loss})$}
\numline{数字}{96道题；交叉熵0.6681、下限0.6681、困惑度1.951；我吃$\rightarrow$苹果0.501/香蕉0.499；小猫吃$\rightarrow$鱼1.0；我学$\rightarrow$AI 0.667；Llama-3在15T token上训练、约7e23 FLOPs}
\fol{训练损失降了，模型就会对话吗？}{不会。只做过预测下一个token的是基座模型，它会接着你的话往下写而不是回答；会对话、听指令是后面监督微调和偏好对齐的结果。}
\fol{「只是统计」等于「只是记忆」吗？}{不等于。小模型确实只记住96道题的频率，但8B参数对上15T token，参数远远不够逐字存数据，它必须找出可复用的规律才能在新文本上猜准。}
\fol{为什么推理只看最后一个位置？}{因为我们要的是整段话之后接什么，输出层只要对最后一个位置做一次4096到128256的打分；但前面每个token仍要完整前向，最后一个位置得看到它们。}
\numline{别答错}{输出是分布不是词，选词规则独立于模型；基座模型只会续写，不会回答问题。}
\phantomsection\addcontentsline{toc}{section}{24\quad 自回归生成}
\lec{24}{自回归生成}
\askline{一段话是怎么一个词一个词出来的？}
\askline{贪心解码和采样分别什么时候用？}
\askline{温度调高调低对输出有什么影响？}
\vspace{3pt}
\noindent{\qlabel{答}}\begingroup\sloppy\begin{enumerate}[leftmargin=1.6em,itemsep=2.5pt,topsep=0pt,parsep=0pt,label={\sffamily\footnotesize\color{accent}\arabic*.}]
\item {\fontsize{10}{14.8}\selectfont 自回归生成是一个循环，每圈三步：把已有token送进模型得到下一个token的分布、按规则挑一个、接到末尾，直到挑出结束标记或达到最大长度；小模型从空白出发4次前向就生成「我吃苹果。」。\par}
\item {\fontsize{10}{14.8}\selectfont 注意第3步的输入里「吃」是第2步自己生成的，每一步的输出都成为下一步输入的一部分，这就是自回归的意思；也正因如此，生成n个token要做n次前向，一次只能出一个，没法并行。\par}
\item {\fontsize{10}{14.8}\selectfont 贪心每步取概率最大的，简单确定：从头开始永远是「我吃苹果。」，从「小狗」开始永远是「小狗吃骨头。」；但苹果0.501、香蕉0.499几乎一样，贪心永远选苹果，合理的话一辈子说不出来，大模型里还容易陷入重复。\par}
\item {\fontsize{10}{14.8}\selectfont 采样按分布概率抽，能说的话就和语料一样丰富；风险是偶尔抽中概率很低又不合理的token，接进去后面就越走越偏，所以采样前先用温度T除logit再softmax，控制这种冒险程度。\par}
\item {\fontsize{10}{14.8}\selectfont 温度就是调分布的尖平：T=0.3时「我」从0.416升到0.782，分布变尖趋近贪心；T=3时降到0.29，冷门词也有机会，会出现「小狗小猫骨头。」这种病句；框架里温度为0一律按贪心处理。\par}
\item {\fontsize{10}{14.8}\selectfont 这就是推理瓶颈的来源：decode每生成一个token都要把全部权重读一遍，16.1 GB在3.35 TB/s下至少4.8 ms，单请求最多约209 token/s，不管算力多强；解法是把多个请求凑成batch，一遍权重服务上百个请求。\par}
\end{enumerate}\endgroup
\numline{公式}{$p_i=\frac{\exp(z_i/T)}{\sum_{j}\exp(z_j/T)}$}
\numline{数字}{贪心4次前向生成「我吃苹果。」；T=0.3我0.782、T=1我0.416、T=3我0.29；不缓存1,499,500个位置 vs 缓存1,999个；16.1 GB÷3.35 TB/s=4.8 ms、209 token/s}
\fol{同样的问题答案为什么不一样？}{采样造成的，不是模型变了。权重在推理时不改，同样输入分布永远一样，不同的是抽中了哪个token，而一个token不同后面整段话就走向不同方向。}
\fol{温度设成 0 是怎么实现的？}{框架里temperature=0都按贪心处理，直接取概率最大的那个，不是真的除以0；所以不考虑浮点误差时，同样的输入会得到同样的输出。}
\fol{没有缓存为什么会多算这么多？}{每一步都把整段话重新算一遍：提示1000个token再生成1000个，依次处理1000到1999个位置，共约150万个；存下KV Cache后每步只算新token，总共1999个。}
\numline{别答错}{温度设 0 按贪心处理，不是除以 0；答案不同是采样造成的，权重并没有变。}
\chapter{Transformer核心}
{\footnotesize\sffamily\color{grayt}第 25--35 讲}\par\vspace{-4pt}
\phantomsection\addcontentsline{toc}{section}{25\quad 位置编码}
\lec{25}{位置编码}
\askline{注意力本身为什么不看顺序？}
\askline{位置编码怎么加进去，为什么是相加？}
\askline{正弦位置编码为什么用不同频率？}
\vspace{3pt}
\noindent{\qlabel{答}}\begingroup\sloppy\begin{enumerate}[leftmargin=1.6em,itemsep=2.5pt,topsep=0pt,parsep=0pt,label={\sffamily\footnotesize\color{accent}\arabic*.}]
\item {\fontsize{10}{14.8}\selectfont 自注意力只按内容比较词与词，公式里没有任何一项用到“第几个位置”，所以把输入打乱，输出只是跟着同样打乱，每个词的数值一个不变。要让模型知道顺序，只能把位置显式塞进输入：每个位置一个和词嵌入同维的向量，直接加到词嵌入上。\par}
\item {\fontsize{10}{14.8}\selectfont 顺序丢掉的后果很直接。用一个复读模型验证：输入 a b c 。 a b c，要它在句号后把第一遍的三个词按位置找回来。有位置嵌入时正确率 100\%，把位置向量拿掉只剩 16.4\%，它知道前面有哪几个词，却不知道谁在前谁在后。\par}
\item {\fontsize{10}{14.8}\selectfont 做法有两条路。一条把位置表 P 当参数和词嵌入一起训练，GPT-2 就是一张 1024 行的可学习表，长度被表的大小锁死；另一条用正弦余弦直接算，P(p,2i)=sin(p/10000\textasciicircum{}(2i/d))，不需要训练，任意位置都能算出向量。\par}
\item {\fontsize{10}{14.8}\selectfont 不同频率是这么用的：分母 10000\textasciicircum{}(2i/d) 随维度组号指数变大，越靠后的维度变化越慢，像快慢不同的钟表指针------只看快针，位置 0 和位置 6 分不清，叠上慢针就唯一了。它还带来一个好性质：位置越近编码越相似。\par}
\item {\fontsize{10}{14.8}\selectfont 代价在长度和一致性上。正弦编码对任意位置都有定义，但模型在没见过的长度上表现好不好是另一回事，可学习表连定义都没有。所以今天主流换成 RoPE：不加到输入上，而是在注意力里旋转 Q、K，Llama-3 还把 rope\_theta 从 10000 调到 500000。\par}
\end{enumerate}\endgroup
\numline{公式}{$x_p=E[\text{token}]+P[p]$；$P(p,2i)=\sin\left(\frac{p}{10000^{2i/d}}\right)$，$P(p,2i+1)=\cos\left(\frac{p}{10000^{2i/d}}\right)$}
\numline{数字}{复读正确率 100\% 对 16.4\%；GPT-2 位置表 1024 行；d=4 时分母 1 与 100；d=64 位置 10 的自点积 32、隔 1 为 30.92、隔 5 为 23.5、隔 30 为 17.07；Llama-3 rope\_theta 10000$\rightarrow$500000}
\fol{那 RoPE 为什么比相加更好？}{相加只在入口注入一次位置，后面每层都可能把它改掉；RoPE 在每层注意力里旋转 Q、K，点积自然只依赖相对距离，长度外推更好。代价是每个头都要现算旋转，多了一点访存。}
\fol{位置编码能外推到训练长度之外吗？}{正弦编码对任意位置都有定义，可学习表连定义都没有；但模型没见过的长度上表现未必好。工程上靠调 rope\_theta、位置插值或 YaRN，4K 训的模型直接上 32K 会明显掉点。}
\fol{推理时位置是怎么传进去的？}{每个新 token 的位置等于它前面已有的 token 数；不同长度的请求拼成一个 batch 时，每个请求各算各的位置，不能混。推理引擎里这是一个独立输入，叫 position ids。}
\numline{别答错}{位置编码是加在词嵌入上、不是拼接，维度不变；注意力公式也没变，变的是输入已带位置。}
\phantomsection\addcontentsline{toc}{section}{26\quad 自注意力}
\lec{26}{自注意力}
\askline{自注意力的公式为什么这么写？}
\askline{为什么要除以根号 d\_k？}
\askline{注意力权重是怎么算出来的？}
\vspace{3pt}
\noindent{\qlabel{答}}\begingroup\sloppy\begin{enumerate}[leftmargin=1.6em,itemsep=2.5pt,topsep=0pt,parsep=0pt,label={\sffamily\footnotesize\color{accent}\arabic*.}]
\item {\fontsize{10}{14.8}\selectfont 自注意力让每个词用自己的一套 q 去和所有词的 k 打分，softmax 成权重后对所有词的 v 加权求和；一句话里所有位置同时做，就是打分 QKᵀ、softmax、加权 AV 三次矩阵运算，没有循环。\par}
\item {\fontsize{10}{14.8}\selectfont 每个词的 q、k、v 都从同一个 x 乘三个矩阵得到：Q=X$\cdot$Wq、K=X$\cdot$Wk、V=X$\cdot$Wv。拆成三份是因为“我在找什么”和“我能提供什么”是两回事：一个代词要找前面的名词，它能提供的只是“这里有个代词”。\par}
\item {\fontsize{10}{14.8}\selectfont 缩放是必需的一步。q、k 各分量方差为 1 时点积方差是 d\_k，实测标准差 d=4 为 2.00、64 为 8.06、128 为 11.23；不缩放时 128 维下一个 q 对 8 个 k 的最大权重平均 0.92，一边倒后梯度几乎为 0。\par}
\item {\fontsize{10}{14.8}\selectfont 算例可以串起来讲：q$\cdot$k 是 0.0986、0.142、-0.0606、-0.1876，除以 $\surd$4 得 0.0493、0.071、-0.0303、-0.0938，softmax 后权重 0.262、0.268、0.242、0.227，加权求和就是输出。\par}
\item {\fontsize{10}{14.8}\selectfont 代价在平方项。打分矩阵是 s$\times$s，Llama-3-8B 一层 32 个头在 8192 序列下，BF16 存要 4 GiB，而这只是中间结果；计算上打分加加权要 1.10 TFLOPs，已经超过 QKV 与输出投影合计的 0.69 TFLOPs。\par}
\item {\fontsize{10}{14.8}\selectfont 学到的是矩阵，不是规则。未训练的 TinyLM 每行权重都在 0.2 到 0.3 之间近乎平均，因为 q 和 k 之间没有学到的对应；训练过的复读模型头 1 在第二遍精确“回头看 3 格”，可见看哪里是训练出来、运行时现算的。\par}
\end{enumerate}\endgroup
\numline{公式}{$\operatorname{Attention}(Q,K,V)=\operatorname{softmax}\left(\frac{QK^{\mathsf{T}}}{\sqrt{d_k}}\right)V$；$q$、$k$ 各分量方差为 $1$ 时 $\operatorname{Var}(q\cdot k)\approx d_k$}
\numline{数字}{q$\cdot$k = 0.0986 / 0.142 / -0.0606 / -0.1876；点积标准差 d=4/64/128 $\rightarrow$ 2.00/8.06/11.23；最大权重不缩放 0.92、缩放后 0.38；Llama-3-8B 投影 0.69 TFLOPs、打分 1.10 TFLOPs、注意力矩阵 4 GiB}
\fol{不除以根号 d\_k 会怎样？}{高维下分数波动大，softmax 接近 one-hot，最大权重 0.92，梯度几乎为 0，训练早期就卡住。除以 $\surd$d\_k 后标准差回到 1 附近，最大权重降到 0.38，分布柔和得多。}
\fol{注意力权重能当解释性看吗？}{只能弱参考。输出是所有 v 的加权和，权重小的项也可能影响很大；而且同一个模型面对不同句子，图案完全不同。它可以用来调试，不能当成模型“在看什么”的结论。}
\fol{长序列为什么贵，贵在哪一步？}{贵在 s$\times$s 的打分和加权。序列翻倍，投影只翻倍，注意力翻四倍；8192 时一层 32 个头的注意力矩阵就有 4 GiB。这正是 FlashAttention 分块计算、不把矩阵落盘的动机。}
\numline{别答错}{「注意」不是从句子里挑一个词，输出是所有 v 的加权平均，每行权重和为 1。}
\phantomsection\addcontentsline{toc}{section}{27\quad 因果掩码}
\lec{27}{因果掩码}
\askline{语言模型为什么必须要因果掩码？}
\askline{不加掩码训练会出什么问题？}
\askline{KV Cache 和因果掩码什么关系？}
\vspace{3pt}
\noindent{\qlabel{答}}\begingroup\sloppy\begin{enumerate}[leftmargin=1.6em,itemsep=2.5pt,topsep=0pt,parsep=0pt,label={\sffamily\footnotesize\color{accent}\arabic*.}]
\item {\fontsize{10}{14.8}\selectfont 因果掩码就是在 softmax 之前，把分数矩阵里 j 大于 i 的位置全部换成 -$\infty$，让第 i 个位置只能看自己和前面的词。实现只有一行，却同时决定了两件事：训练能不能一次前向算完所有位置，推理能不能用 KV Cache。\par}
\item {\fontsize{10}{14.8}\selectfont 不加掩码会教会模型抄答案。训练时整段文本已知，第 i 个位置的输出用来预测第 i+1 个词，而自注意力恰好让它看到第 i+1 个词；模型很快发现不用理解前文、照着抄就行，损失降得很低，生成时答案不存在就直接失效。\par}
\item {\fontsize{10}{14.8}\selectfont 机制上很干净：e\textasciicircum{}-$\infty$ 等于 0，softmax 后被遮位置权重恰好是 0，其余重新归一化仍和为 1，注意力矩阵只剩含对角线的下三角。TinyLM 里第 1 个词权重变成 1，第 2 个词在“我”和“爱”之间分 0.491 和 0.509，最后一行不受影响。\par}
\item {\fontsize{10}{14.8}\selectfont 可以直接测出来。把第 4 个词的向量换成另一个，不加掩码时前 3 个位置的输出最多变 0.3804，说明它们确实看到了第 4 个词；加上掩码后变化是 0，前面的词对后面发生了什么一无所知。\par}
\item {\fontsize{10}{14.8}\selectfont 这也是 KV Cache 成立的前提。第 i 个位置只依赖前 i 个词，往后接新词时前面的输出一个数都不会变，逐层传递下去，每层每个词的 k、v 一旦算出就永远不变，于是可以缓存：decode 时新词只算自己的 q、k、v，用 1$\times$s 的一行去打分，本来就不需要掩码。\par}
\item {\fontsize{10}{14.8}\selectfont 缓存很贵。Llama-3-8B 每 token 每层存 8 个 KV 头的 k 和 v（每头 128 维、BF16），2$\times$32$\times$8$\times$128$\times$2 字节就是 128 KiB；8192 token 要 1 GiB，50 个请求 50 GiB，比 16 GB 权重还大。\par}
\end{enumerate}\endgroup
\numline{公式}{$A=\operatorname{softmax}\left(\frac{QK^{\mathsf{T}}}{\sqrt{d_k}}+M\right)$，$M_{ij}=0\ (j\le i)$、$-\infty\ (j>i)$}
\numline{数字}{改第 4 个词：不加掩码前 3 个输出最多变 0.3804，加掩码变 0；Llama-3-8B 每 token 128 KiB，8192 token 1 GiB，50 请求 50 GiB，权重仅 16 GB}
\fol{掩码能不能加在 softmax 之后？}{不行。先 softmax 再把后面权重改成 0，剩下的权重加起来不到 1，而且 softmax 的分母已经用到了后面词的分数，信息早混进去了。必须在 softmax 之前加 -$\infty$。}
\fol{BERT 这类模型也需要吗？}{不需要。BERT 用双向注意力，每个词看整句，代价是不能像 GPT 那样逐词生成；因果掩码只属于 decoder-only。填充掩码是另一回事，管的是 batch 里补齐的 padding。}
\fol{分块 prefill 的掩码怎么处理？}{按全局位置生成，不是按 chunk 内位置。后面的 chunk 要能看到前面所有 chunk 的 KV，所以掩码是全局下三角；这也是分块后结果能和整段一次算对齐的原因。}
\numline{别答错}{掩码必须在 softmax 之前加，否则剩余权重和不为 1，后面词的信息已经进了分母。}
\phantomsection\addcontentsline{toc}{section}{28\quad 多头注意力}
\lec{28}{多头注意力}
\askline{多头注意力多花了几倍参数？}
\askline{为什么要把注意力拆成多个头？}
\askline{GQA 和多头是什么关系？}
\vspace{3pt}
\noindent{\qlabel{答}}\begingroup\sloppy\begin{enumerate}[leftmargin=1.6em,itemsep=2.5pt,topsep=0pt,parsep=0pt,label={\sffamily\footnotesize\color{accent}\arabic*.}]
\item {\fontsize{10}{14.8}\selectfont 多头注意力把 Q、K、V 按列切成 n\_h 段，每段交给一个头，各自独立做打分、缩放、掩码、softmax 和加权求和，再把各头输出拼接起来乘 W\_o。它不引入新结构，只是把一次“看”拆成几次“从不同角度看”。\par}
\item {\fontsize{10}{14.8}\selectfont 参数并没有变多。Wq、Wk、Wv、Wo 还是 d$\times$d，TinyLM 里各 16 个参数，和一个头时一模一样；变的只是用法：一个头是 4 维整体做一次点积，两个头是前 2 维和后 2 维各做一次，乘法次数相同，多出来的只是两遍 softmax。\par}
\item {\fontsize{10}{14.8}\selectfont 它解决的是单个 softmax 表达不了的问题。一个 softmax 很难同时把很大的权重放在两个毫不相关的词上，两个头就可以各盯一个。复读模型里，头 1 在第二遍“回头看 3 格”负责复读，头 2 第 4 到 6 行几乎全盯着句号，像在标记“复读从这里开始”。\par}
\item {\fontsize{10}{14.8}\selectfont 但头不一定都被用满。关头实验：两个头都在时复读正确率 100\%，关掉头 1 跌到 8.1\%，关掉头 2 仍是 100\%，说明这个任务全靠头 1。大模型里也一样，一部分头冗余，剪掉几乎不掉点，分工是训练自己选的，不是设计出来的。\par}
\item {\fontsize{10}{14.8}\selectfont 头也不是越多越好。每个头只拿 d/n\_h 维，头越多每头越窄，q 和 k 的点积分辨不出足够多的模式。Llama-3-8B 是 32 个查询头、每头 128 维，32$\times$128 正好是隐藏维度 4096。\par}
\item {\fontsize{10}{14.8}\selectfont 工程上有两个直接影响。各头计算完全独立，GPU 上可以一次算完，多卡张量并行最常见的切法就是按头切，8 张卡每卡 4 个头，最后乘 W\_o 时合并；而 KV Cache 按 8 个 KV 头存，靠 GQA 让 4 个查询头共享一组 k、v，缓存降到四分之一。\par}
\end{enumerate}\endgroup
\numline{公式}{输出 $=[\text{head}_1,\ldots,\text{head}_{n}]\cdot W_o$；$\text{head}_h=\operatorname{softmax}\left(\frac{Q_hK_h^{\mathsf{T}}}{\sqrt{d_h}}+M\right)V_h$，$d_h=\frac{d}{n_h}$}
\numline{数字}{TinyLM 每头 2 维、四个矩阵各 16 个参数；关掉头 1 $\rightarrow$ 8.1\%，关掉头 2 $\rightarrow$ 100\%；Llama-3-8B 32 头 $\times$ 128 维 = 4096；GQA 让 KV Cache 降到 1/4}
\fol{头数是不是越多越好？}{不是。每头只拿 d/n\_h 维，头太多每头就太窄，点积分辨力下降；而且 softmax 和掩码的固定开销会摊大。实践上每头 128 维左右是常见甜点，再靠 GQA 去省 KV。}
\fol{W\_o 能不能去掉？}{不能。拼接只是把各头结果并排放在一起，W\_o 才真正把它们混合后写回残差流。去掉它，各头信息互不相干，后面的层用不到跨头的组合，效果明显下降。}
\fol{头的分工是设计出来的吗？}{不是。我们从来没告诉模型头 1 负责复读，是训练自己选出来的；换个随机种子可能变成头 2 负责，或者两头分担。所以别把“第几个头干什么”当固定结论。}
\numline{别答错}{头的维度是切出来的不是复制出来的；W\_o 不是可有可无，它才真正混合各头。}
\phantomsection\addcontentsline{toc}{section}{29\quad 残差连接}
\lec{29}{残差连接}
\askline{深层网络为什么训不动？}
\askline{残差连接为什么能解决梯度消失？}
\askline{只加残差就够了吗，还缺什么？}
\vspace{3pt}
\noindent{\qlabel{答}}\begingroup\sloppy\begin{enumerate}[leftmargin=1.6em,itemsep=2.5pt,topsep=0pt,parsep=0pt,label={\sffamily\footnotesize\color{accent}\arabic*.}]
\item {\fontsize{10}{14.8}\selectfont 残差连接的做法简单得离谱：每层不再只输出 F(x)，而是输出 x + F(x)，输入绕过这一层原封不动加到输出上，这条路叫捷径。正是这个加号让几十层、上百层的网络变得可以训练。\par}
\item {\fontsize{10}{14.8}\selectfont 问题出在连乘。梯度每回传一层要乘一次这一层的局部导数，30 层就是 30 个矩阵连乘；每层都是 0.5 时 0.5$^3$⁰ 等于 9.31$\times$10$^-$$^1$⁰，几乎是 0，这叫梯度消失；每层都是 1.5 时 1.92$\times$10⁵，参数一步被推飞，这叫梯度爆炸。\par}
\item {\fontsize{10}{14.8}\selectfont 加号的作用在导数里看得很清楚：对 x 求导，第一项是 1、第二项是 F′，所以局部导数变成 1+F′。F′ 再小哪怕为 0，这一层也还剩 1，梯度至少能沿捷径原样通过，不会在某层被截断或压缩到 0。\par}
\item {\fontsize{10}{14.8}\selectfont 换个角度更直观。把 30 个 (1+F′) 连乘展开，里面有大量路径：有一项是 1$\times$1$\times$…$\times$1，对应每层都走捷径；有的项只含一个 F′，对应只在某一层走 F。梯度是所有路径之和，其中总有很短的路径，不会因为连乘而消失。\par}
\item {\fontsize{10}{14.8}\selectfont 实测差距很大。30 层、每层 16 维的网络，从最后一层传回长度为 1 的梯度：tanh 下普通网络到第 0 层只剩 0.0506，加残差是 19.57；ReLU 下普通网络直接变成 0，因为一路上总有死亡 ReLU 把梯度截断，加残差则是 3704.18。\par}
\item {\fontsize{10}{14.8}\selectfont 但残差不解决数值膨胀。加了残差的那条曲线不是平的，而是往前越来越大：每层都往信号上加东西，前向越叠越大，反向也跟着越大。所以它必须和归一化、合适的初始化一起用，GPT-2 等模型还会把残差分支的权重初始化得更小。\par}
\end{enumerate}\endgroup
\numline{公式}{$x_{l+1}=x_l+F(x_l)$；$\frac{\partial x_{l+1}}{\partial x_l}=1+F'$}
\numline{数字}{0.5$^3$⁰ = 9.31$\times$10$^-$$^1$⁰，1.5$^3$⁰ = 1.92$\times$10⁵；tanh 第 0 层 0.0506 对 19.57；ReLU 0 对 3704.18；32 层共 64 次残差相加；残差流 4096 维}
\fol{残差用相加还是拼接？}{必须相加。拼接每过一层维度就翻倍，参数和显存都会失控；相加维度不变，同一条 4096 维的残差流才能原样穿过 32 层。这也是子模块输出必须还原形状的原因。}
\fol{残差在推理时开销大吗？}{计算量可以忽略，就是逐元素加法；但它要把两个 [s, 4096] 的矩阵读进来再写出去，是访存密集型。引擎通常把残差相加和后面的归一化融合成一个算子，省一次读写。}
\fol{有了 Pre-Norm 还需要残差吗？}{需要。Pre-Norm 只是把归一化挪到子层入口，让捷径上完全没有归一化、梯度传得更顺，骨架仍是 x 加子层输出。它改的是归一化的位置，不是残差本身。}
\numline{别答错}{残差是相加不是拼接，拼接会让维度每层翻倍；它只管梯度，不管数值膨胀。}
\phantomsection\addcontentsline{toc}{section}{30\quad LayerNorm}
\lec{30}{LayerNorm}
\askline{LayerNorm 具体在算什么？}
\askline{为什么不用 BatchNorm？}
\askline{归一化之后数都在 -1 到 1 吗？}
\vspace{3pt}
\noindent{\qlabel{答}}\begingroup\sloppy\begin{enumerate}[leftmargin=1.6em,itemsep=2.5pt,topsep=0pt,parsep=0pt,label={\sffamily\footnotesize\color{accent}\arabic*.}]
\item {\fontsize{10}{14.8}\selectfont LayerNorm 对每个词的向量单独操作：设这一行有 d 个数，先算均值 $\mu$，每个数减 $\mu$，再算方差 σ$^2$、开根号得标准差，每个数除以标准差，最后乘可训练的 γ、加 $\beta$，把这一行拉回均值 0、方差 1。\par}
\item {\fontsize{10}{14.8}\selectfont 它是被残差逼出来的。残差每层都往残差流里加东西，30 层残差网络第 1、10、20、30 层的输出长度分别是 3.4、42.9、2006.4、82441.6；数值这样失控，激活会饱和、softmax 会一边倒，训练极不稳定。\par}
\item {\fontsize{10}{14.8}\selectfont 加上它之后尺度就固定了。方差为 1 意味着 d 个数平方和为 d，所以归一化后的向量长度总在 $\surd$d 附近：TinyLM 是 2，16 维网络是 4。同一个网络每层输出长度全都稳定在 4.0，不管前面推到多大。\par}
\item {\fontsize{10}{14.8}\selectfont 统计量按词算，这一点是关键。均值和方差在一个词的 d 个维度之间算，每个词各算各的，和 batch 里有什么、有多少完全无关，所以训练和推理的计算一模一样。\par}
\item {\fontsize{10}{14.8}\selectfont BatchNorm 方向正好相反：对同一维、在一个 batch 的所有样本之间算。它在图像上很成功，但不适合语言模型：句子长短不一，样本差异大；更麻烦的是推理时一次可能只生成一个 token，根本没有“一批”可以算统计量。\par}
\item {\fontsize{10}{14.8}\selectfont 它还有两个性质值得说。对整体缩放和平移不敏感：整行乘 3，均值和标准差也乘 3，相除后不变；整行加 5 同理。参数也很少，γ、$\beta$ 各 d 维，初始全 1 和全 0，只占 2d 个，Llama-3-8B 全部 65 个归一化合计 266240 个。\par}
\end{enumerate}\endgroup
\numline{公式}{$\operatorname{LN}(x)=\gamma\odot\frac{x-\mu}{\sqrt{\sigma^2+\varepsilon}}+\beta$，$\mu$ 和 $\sigma^2$ 在每行的 $d$ 个维度上算}
\numline{数字}{AI 行均值 -0.0434、方差 0.4165、标准差 0.6454、单个值最大 1.3123；ε=1e-5；长度 3.4/42.9/2006.4/82441.6 $\rightarrow$ 全 4.0；Llama-3 65 个归一化共 266240 参数}
\fol{RMSNorm 和 LN 差在哪？}{RMSNorm 只除以均方根、不减均值，少一次归约，实测效果几乎不差还更快。Llama-3 每个归一化只有一个 4096 维的 γ、没有 $\beta$，参数减半。}
\fol{Pre-Norm 为什么更好训？}{深层一律 Pre-Norm。Post-Norm 把归一化放在残差之后，梯度要穿过归一化；Pre-Norm 让捷径完全干净，32 层以上明显更好训，代价是出口要补一个最后的归一化。}
\fol{这一层归一化能省掉吗？}{不能。30 层残差网络输出长度能到 8 万多，没有归一化训练直接发散。它能省的是参数（只有 2d 个），不是这个算子；工程上做的是把它和残差、矩阵乘融合。}
\numline{别答错}{归一化后不是所有数都在 -1 到 1，均值 0 方差 1 只是统计量，单个数可以到 1.3123。}
\phantomsection\addcontentsline{toc}{section}{31\quad 前馈网络}
\lec{31}{前馈网络}
\askline{FFN 到底在算什么？}
\askline{为什么要先升维再降维？}
\askline{FFN 和注意力的分工是什么？}
\vspace{3pt}
\noindent{\qlabel{答}}\begingroup\sloppy\begin{enumerate}[leftmargin=1.6em,itemsep=2.5pt,topsep=0pt,parsep=0pt,label={\sffamily\footnotesize\color{accent}\arabic*.}]
\item {\fontsize{10}{14.8}\selectfont FFN 对每个词的向量做三步：乘 W1 从 d 维升到 d\_ff（通常约 4 倍），逐元素过 ReLU 把负数变 0，再乘 W2 降回 d 维，然后把结果加回残差流。降回 d 维是必须的，因为要和输入相加。\par}
\item {\fontsize{10}{14.8}\selectfont 它和注意力是互补关系。注意力让词与词交换信息，依赖 s$\times$s 的打分矩阵；FFN 里词与词完全不交流，整句一起算和一个词一个词算结果完全相同，原论文叫 position-wise，一个常用比喻是注意力开会、FFN 各自回去干活。\par}
\item {\fontsize{10}{14.8}\selectfont 先宽再窄是为了表达力。两个线性层之间没有激活会合并成一个线性层，有了 ReLU 才有非线性；中间层越宽，能拼出的折线越多，能表达的函数越丰富。TinyLM 的 d\_ff=8，8 个神经元里只亮 4 个。\par}
\item {\fontsize{10}{14.8}\selectfont 它还能看成一张键值表。把降维拆开：FFN(x) 等于 $\Sigma$\_j ReLU(x$\cdot$W1 第 j 列) $\times$ W2 第 j 行。W1 的每一列是一个待匹配的“键”，匹配程度决定这个神经元多亮；W2 的对应行是“值”，亮了就按亮度加回输出。\par}
\item {\fontsize{10}{14.8}\selectfont TinyLM 的算例能对上：亮的是第 1、3、5、8 个神经元，亮度 1.598、0.083、1.342、1.528，把 W2 这四行按亮度加起来正好是输出 (2.828, 1.673, 2.837, -0.636)。大模型里大量知识就存在这些权重里。\par}
\item {\fontsize{10}{14.8}\selectfont 它是参数和计算的大头。Llama-3-8B 的 d\_ff 是 14336（3.5 倍），用带门控的 SwiGLU、三个矩阵，每层约 1.762 亿参数占这一层 81\%，每个 token 约 3.52 亿次浮点运算；序列不长时远超注意力，是量化最想压的对象。\par}
\end{enumerate}\endgroup
\numline{公式}{$\operatorname{FFN}(x)=\operatorname{ReLU}(x\cdot W_1)\cdot W_2=\sum_{j}\operatorname{ReLU}(x\cdot W_1[:,j])\,W_2[j,:]$}
\numline{数字}{TinyLM FFN 64 参数 = 注意力 64；Llama-3-8B d\_ff=14336（3.5 倍）、每层 176.2M 参数占 81\%、每 token 352M 次运算；SwiGLU 常用 8/3 倍}
\fol{d\_ff 一定是 4 倍吗？}{不一定。原始 Transformer 是 4 倍；SwiGLU 多了一个矩阵，为保持参数量相近常用 8/3 倍再取整，Llama-3-8B 就是 14336/4096 等于 3.5 倍。}
\fol{FFN 的神经元有明确含义吗？}{部分有。训练好的模型里有些神经元对应可解释的模式，看到某类词或某种语境就亮，写回的值会抬高某些下一个词的概率；但很多神经元同时参与好几种模式，单独看往往说不清。}
\fol{MoE 为什么能省计算？}{它把一个大 FFN 换成很多小专家，每个 token 只激活其中一两个，参数总量可以很大而每 token 计算量不变。代价是路由容易不均、显存要装下全部专家，通信也更复杂。}
\numline{别答错}{d\_ff 不一定 4 倍：SwiGLU 常用 8/3 倍，Llama-3 是 3.5 倍。}
\phantomsection\addcontentsline{toc}{section}{32\quad Transformer Block}
\lec{32}{Transformer Block}
\askline{一个 Block 里有哪些部件？}
\askline{为什么层可以一层层叠起来？}
\askline{现代模型相对原始结构改了什么？}
\vspace{3pt}
\noindent{\qlabel{答}}\begingroup\sloppy\begin{enumerate}[leftmargin=1.6em,itemsep=2.5pt,topsep=0pt,parsep=0pt,label={\sffamily\footnotesize\color{accent}\arabic*.}]
\item {\fontsize{10}{14.8}\selectfont 一个 Transformer Block 就是两个子层：先多头自注意力，让词与词交换信息；再 FFN，让每个词各自加工收集到的信息。每个子层的模式都一样，做一件事、加回残差流、再做归一化。\par}
\item {\fontsize{10}{14.8}\selectfont 形状不变是能叠的前提。一层的输入是 [s, d]，输出还是 [s, d]，注意力和 FFN 都把自己还原回原形状，所以同样的块可以一层接一层叠多少都行。但结构相同不代表权重相同，32 层是 32 组各自训练出来的参数。\par}
\item {\fontsize{10}{14.8}\selectfont 完整数据流是：token 编号查词嵌入表、加位置编码，得到 [s, d] 的初始残差流；过 N 个 Block；最后一层输出乘 lm\_head 得到 [s, V] 的 logits；再过 softmax 得到每个位置下一个词的概率。\par}
\item {\fontsize{10}{14.8}\selectfont 参数可以数清楚。TinyLM 一共 304 个：词嵌入 24、两层各 128、输出层 24。Llama-3-8B 每层 218,112,000，32 层共 6,979,584,000，加嵌入和输出层各 525,336,576，合计 8,030,261,248。\par}
\item {\fontsize{10}{14.8}\selectfont 现代模型改了五处，骨架没动：归一化从残差之后（Post-Norm）挪到子层之前（Pre-Norm）并换成 RMSNorm，位置编码改 RoPE，FFN 换 SwiGLU 三个矩阵，K、V 用 GQA 让多个查询头共享，出口再补一个最后的归一化。\par}
\item {\fontsize{10}{14.8}\selectfont 工程上这些性质很值钱。因为每层结构完全一样，推理引擎只需要为一层写高度优化的算子，然后循环 N 次；权重文件也按层组织，从 layers.0 到 layers.31。层间是串行的，切到多卡时要么按层做流水线并行，要么在层内做张量并行。\par}
\end{enumerate}\endgroup
\numline{公式}{Post-Norm：$x\leftarrow\operatorname{LN}(x+\operatorname{Attn}(x))$，$x\leftarrow\operatorname{LN}(x+\operatorname{FFN}(x))$；Pre-Norm：$x\leftarrow x+\operatorname{Attn}(\operatorname{Norm}(x))$，$x\leftarrow x+\operatorname{FFN}(\operatorname{Norm}(x))$}
\numline{数字}{TinyLM 304（嵌入 24 + 每层 128$\times$2 + 输出 24）；Llama-3-8B 每层 218,112,000、32 层 6,979,584,000、合计 8,030,261,248；70B 是 80 层宽 8192}
\fol{32 层是指 32 个线性层吗？}{不是。32 层是 32 个 Block，每个 Block 里有 7 个线性层（Wq、Wk、Wv、Wo 加 FFN 的三个 SwiGLU 矩阵）和 2 个归一化，按线性层数算会多得多。}
\fol{为什么层间只能串行？}{因为第 l+1 层要读第 l 层写好的残差流，依赖是天然的。能并行的是层内的头和 FFN 的神经元，所以单个请求的延迟下限基本就是层数乘以每层耗时。}
\fol{两者差异影响算子实现吗？}{会。Pre-Norm 和 Post-Norm 的融合方式不同，RMSNorm、RoPE、SwiGLU、GQA 各有自己的 kernel 和最优 tile；但每层循环 N 次这个调度骨架是一样的。}
\numline{别答错}{「32 层」指 32 个 Block，结构相同但权重是 32 份，不是一组权重用 32 遍。}
\phantomsection\addcontentsline{toc}{section}{33\quad 输出层}
\lec{33}{输出层}
\askline{输出层到底做了什么？只是打个分吗？}
\askline{lm\_head 有多大，为什么说它贵？}
\askline{推理时为什么要跳过大部分位置？}
\vspace{3pt}
\noindent{\qlabel{答}}\begingroup\sloppy\begin{enumerate}[leftmargin=1.6em,itemsep=2.5pt,topsep=0pt,parsep=0pt,label={\sffamily\footnotesize\color{accent}\arabic*.}]
\item {\fontsize{10}{14.8}\selectfont 输出层只做两件事：一个 d$\times$V 的 lm\_head 把最后一层的 d 维向量变成词表上的 V 个打分，再过 softmax 变成概率，模型从这里给出下一个词。\par}
\item {\fontsize{10}{14.8}\selectfont 每个打分是向量和 lm\_head 某一列的点积，那一列就是一个词的输出向量；TinyLM 里“AI”的输出向量和“推理”那一列点积得 1.184，softmax 后是 0.441。\par}
\item {\fontsize{10}{14.8}\selectfont 训练时每个位置都要算这一行，第 i 行用来预测第 i+1 个词并算交叉熵；推理时只有最后一行有用，所以 prefill 虽然整段提示都过所有层，到输出层只对最后一个位置做矩阵乘。\par}
\item {\fontsize{10}{14.8}\selectfont 它出奇地贵：Llama-3-8B 的 lm\_head 是 4096$\times$128256，5.25 亿参数，每 token 要 1.05 GFLOPs，占 2N 的 6.5\%，比任何单个线性层都大，词表越大这一项越重。\par}
\item {\fontsize{10}{14.8}\selectfont 它的形状正好是嵌入表的转置，可以令 W\_lm = W\_Eᵀ 做权重绑定省一份参数，GPT-2 就这么做；Llama-3-8B 不绑定，多出的 5 亿参数只占 6.5\%，换来更大的表达自由。\par}
\item {\fontsize{10}{14.8}\selectfont 工程上 logits 本身很占地方：一个位置 128256 个数用 FP32 存约 501 KiB，温度缩放和 Top-k、Top-p 采样都在它上面做，引擎会写成专门的 GPU 算子，避免把 logits 搬回 CPU。\par}
\end{enumerate}\endgroup
\numline{公式}{$\text{logits}=x\cdot W_{\text{lm}}$；$p_i=\frac{e^{\text{logit}_i}}{\sum_{j}e^{\text{logit}_j}}$}
\numline{数字}{TinyLM：W\_lm 4$\times$6，“推理”logit 1.184、softmax 后 0.441；Llama-3-8B：4096$\times$128256 = 525M 参数，每 token 1.05 GFLOPs，占 2N 的 6.5\%；一个位置 logits 用 FP32 存约 501 KiB}
\fol{训练时也只用最后一行吗？}{不行。训练要每个位置都算，第 i 行和真实的下一个词算交叉熵；8192 个位置乘 128256 的 logits 本身就是好几 GiB 显存，所以训练比推理更需要在这些地方省。}
\fol{那为什么还保留两份权重？}{两者语义不同：嵌入表是词进入模型时的向量，lm\_head 是词离开时的向量。小模型绑定能省一份参数，大模型里这 5 亿只占 6.5\%，不绑定的表达自由更值钱。}
\fol{TinyLM 的预测为什么没意义？}{它的权重是随机初始化的，从没训练过，0.441 只是数字碰巧这样；这一讲关心的是从向量到概率的每一步计算，训练好的大模型走的是完全相同的路径。}
\numline{别答错}{logits 不是概率，可正可负；Pre-Norm 在 lm\_head 前还有最终 RMSNorm。}
\phantomsection\addcontentsline{toc}{section}{34\quad 前向传播}
\lec{34}{前向传播}
\askline{一句话从进模型到出概率，经过哪些步骤？}
\askline{为什么前缀的结果和整句的那一行相同？}
\askline{前向传播和推理是一回事吗？}
\vspace{3pt}
\noindent{\qlabel{答}}\begingroup\sloppy\begin{enumerate}[leftmargin=1.6em,itemsep=2.5pt,topsep=0pt,parsep=0pt,label={\sffamily\footnotesize\color{accent}\arabic*.}]
\item {\fontsize{10}{14.8}\selectfont 一次前向就是查嵌入表、加位置编码、过 N 个 Transformer Block、乘 lm\_head、softmax；TinyLM 把这条路径拆成 24 站，每一站的每个数都能查。\par}
\item {\fontsize{10}{14.8}\selectfont 每个 Block 里的顺序是：多头注意力、残差相加、归一化，然后 FFN、残差相加、归一化；从第 2 站到第 22 站形状始终是 [4,4]，只有 FFN 中间变宽到 8 维。\par}
\item {\fontsize{10}{14.8}\selectfont 宽度不变的那条路就是残差流，每一站都从它读、往它写；后面几乎所有优化都能放回这张表里看：KV Cache 省第 3 站的重算，FlashAttention 省第 4、5 站的显存写入。\par}
\item {\fontsize{10}{14.8}\selectfont 因果掩码让前缀的预测藏在整句里：只送“我爱”得到的最后一行，和送整句时的第 2 行完全相同，因为第 2 个位置在每一层都只能看到前两个词。\par}
\item {\fontsize{10}{14.8}\selectfont 把前向放进循环就是生成：从“我”开始每次取概率最大的词接回去，TinyLM 三步得到“我推理我推理”，每一步都从头重算整句，这正是 KV Cache 要省掉的部分。\par}
\item {\fontsize{10}{14.8}\selectfont 同一条路径放到 Llama-3-8B 上只是规模不同：32 层、RMSNorm、32 头带 RoPE 和 GQA 的注意力、SwiGLU 的 FFN，最后一个 RMSNorm 加 lm\_head，每 token 约 160 亿 FLOPs。\par}
\end{enumerate}\endgroup
\numline{公式}{$x=E[\text{ids}]+P$；每层 $x=\operatorname{LN}(x+\operatorname{Attn}(x))$，$x=\operatorname{LN}(x+\operatorname{FFN}(x))$；$p=\operatorname{softmax}(x\cdot W_{\text{lm}})$}
\numline{数字}{TinyLM：24 站、d=4、2 层、2 头、FFN 8 维、词表 6，预测“推理”0.441，贪心生成“我推理我推理”；Llama-3-8B：32 层、d=4096、词表 128256、每 token 约 2$\times$80 亿 = 160 亿 FLOPs}
\fol{那为什么说前向不等于推理？}{一次前向是给一串 token 算出每个位置的概率；推理是把它放进循环，prefill 一次、decode 很多次。训练在前向之后还要算损失、做反向，那是另外 4N 的计算。}
\fol{加了 KV Cache 哪些站能跳过？}{第 3 站的 K、V 不用重算，第 4、5 站的注意力矩阵也不必整个写进显存。前者是 KV Cache，后者是 FlashAttention，两者都只是这张 24 站表上的局部改动。}
\fol{你怎么确认每一站算得对？}{用 tinylm.json 里每站的中间结果逐站对拍，浮点用 1e-3 容差；再验证前缀的最后一行等于整句的对应行。这两条过了，整条数据流就基本可信。}
\numline{别答错}{一次前向不等于推理：推理是 prefill 一次加 decode 多次的循环，训练还有反向。}
\phantomsection\addcontentsline{toc}{section}{35\quad 参数量与计算量}
\lec{35}{参数量与计算量}
\askline{一个模型的参数量怎么算出来的？}
\askline{推理一个 token 大概多少 FLOPs？}
\askline{上下文变长，计算量怎么变？}
\vspace{3pt}
\noindent{\qlabel{答}}\begingroup\sloppy\begin{enumerate}[leftmargin=1.6em,itemsep=2.5pt,topsep=0pt,parsep=0pt,label={\sffamily\footnotesize\color{accent}\arabic*.}]
\item {\fontsize{10}{14.8}\selectfont 参数全在线性层、嵌入表和归一化里：每层注意力 2d$^2$+2d$\cdot$kv，SwiGLU 的 FFN 是 3d$\cdot$d\_ff，再加两个 d 的归一化，入口嵌入表和出口 lm\_head 各 V$\cdot$d。\par}
\item {\fontsize{10}{14.8}\selectfont 照结构逐项加，Llama-3-8B 得 8,030,261,248 个，和 Meta 公布的 8.03B 一个不差；其中 FFN 占 70.2\%、注意力 16.7\%、嵌入表和输出层 13.1\%，归一化小到可以忽略。\par}
\item {\fontsize{10}{14.8}\selectfont 计算量的规律是每个参与矩阵乘法的参数对每个 token 做一次乘加，也就是 2 FLOPs，所以前向约 2(N-V$\cdot$d)，嵌入表只是查表不算；8B 是 15.01 GFLOPs。\par}
\item {\fontsize{10}{14.8}\selectfont 还要加上注意力的打分和加权：每个 token 前面有 s 个 token 时，每层多 4$\cdot$s$\cdot$d 次运算；8K 上下文下每 token 19.30 GFLOPs、打分占 22.2\%，13 万上下文占 82.1\%。\par}
\item {\fontsize{10}{14.8}\selectfont 两个能心算的估算：推理每 token $\approx$2N，8B 约 16 GFLOPs，生成 1000 个 token 约 16 TFLOPs；训练每 token $\approx$6N，前向 2N 加反向 4N，8B 约 48 GFLOPs。\par}
\item {\fontsize{10}{14.8}\selectfont 这两个公式忽略了嵌入、打分、归一化和 softmax，只在上下文不太长时准；MoE 的 2N 要用激活参数，DeepSeek-V3 总参 6710 亿，但每 token 只激活约 370 亿。\par}
\end{enumerate}\endgroup
\numline{公式}{$N=Vd+L(2d^2+2d\cdot\text{kv}+3d\,d_{\text{ff}}+2d)+d+Vd$；每 token 前向 $\approx 2(N-Vd)+4Lsd$；训练 $\approx 6N$}
\numline{数字}{Llama-3-8B 8,030,261,248；Llama-3-70B 70,553,706,496；每 token 15.01(s=0)/19.30(s=8192)/83.73(s=131072) GFLOPs；训练 48.18；BF16 权重 14.96 GiB$\approx$16.1 GB；KV Cache 每 token 128 KiB，8192$\times$16 请求 16 GiB}
\fol{这些数字在工程上有什么用？}{直接决定显存、速度和钱。8B 的 BF16 权重 16.1 GB，一张 80 GB 的 H100 放得下；70B 是 141 GB，一张卡放不下，至少两张，或者量化到 INT4 约 35 GB。}
\fol{单请求推理最快能多快？}{上限是显存带宽除以权重字节数：H100 约 3.35 TB/s，8B 每秒最多约 209 个 token，此时算力只用到 3.4 TFLOPS，所以必须靠 batch 摊薄权重读取。}
\fol{FLOPs 和 FLOPS 要分清吗？}{必须分清：FLOPs 是运算次数，FLOPS 是每秒次数。另外按 MACs 统计会差一半，因为一次乘加算两次浮点运算，报数时要说清用的是哪个口径。}
\numline{别答错}{FLOPs 是次数、FLOPS 是速度；按 MACs 统计算会差一半；MoE 的 2N 用激活参数。}
\chapter{现代LLM架构}
{\footnotesize\sffamily\color{grayt}第 36--44 讲}\par\vspace{-4pt}
\phantomsection\addcontentsline{toc}{section}{36\quad Pre-Norm 与 RMSNorm}
\lec{36}{Pre-Norm 与 RMSNorm}
\askline{为什么现代大模型都用 Pre-Norm？}
\askline{RMSNorm 和 LayerNorm 差在哪？}
\askline{归一化算子怎么加速？}
\vspace{3pt}
\noindent{\qlabel{答}}\begingroup\sloppy\begin{enumerate}[leftmargin=1.6em,itemsep=2.5pt,topsep=0pt,parsep=0pt,label={\sffamily\footnotesize\color{accent}\arabic*.}]
\item {\fontsize{10}{14.8}\selectfont Post-Norm 是 x $\leftarrow$ LN(x + F(x))，归一化卡在残差主干道上，梯度每过一层都要被它重新缩放一次；层数一多，这些缩放层层累积，梯度大小就难以控制。\par}
\item {\fontsize{10}{14.8}\selectfont Pre-Norm 改成 x $\leftarrow$ x + F(Norm(x))，只有送进注意力或 FFN 的那一份被归一化，残差流本身不动，成了 x 加、加、再加的纯加法链，梯度可以一路加回第一层。\par}
\item {\fontsize{10}{14.8}\selectfont 实测很直白：96 层网络从最后一层回传梯度，Post-Norm 到第 0 层膨胀到 5508.33，Pre-Norm 只有 21.374；6 层、24 层时两者还差不多，所以是深层才必须换。\par}
\item {\fontsize{10}{14.8}\selectfont 代价是主干道再没被归一化，数值会一层层累积变大，所以 Pre-Norm 模型要在所有层之后、lm\_head 之前补一个最后的归一化，也就是结构图里那个 Final Norm。\par}
\item {\fontsize{10}{14.8}\selectfont RMSNorm 砍掉与均值有关的部分：不减均值、没有 $\beta$，只用 γ 按均方根缩放；同一行数上 LayerNorm 得 -0.5327 起的那组，RMSNorm 得 -0.5986 起的那组，差别来自 -0.0434 的均值。\par}
\item {\fontsize{10}{14.8}\selectfont 工程上有两个硬收益：Fused Add-RMSNorm 把残差相加和归一化放进同一个内核，寄存器里就地算完，省掉一半以上 HBM 往返；FP16 上限 65504，平方和必须在寄存器里 upcast 到 FP32 才不会溢出。\par}
\end{enumerate}\endgroup
\numline{公式}{Post-Norm: $x\leftarrow\operatorname{LN}(x+F(x))$；Pre-Norm: $x\leftarrow x+F(\operatorname{Norm}(x))$；$\operatorname{RMSNorm}(x)=\frac{\gamma\odot x}{\sqrt{\operatorname{mean}(x^2)+\epsilon}}$}
\numline{数字}{96 层回传梯度 Post-Norm 5508.33、Pre-Norm 21.374（24 层 8.383 对 7.135，6 层 1.297 对 2.258）；LayerNorm (-0.5327,-1.3091,0.5295,1.3123)、RMSNorm (-0.5986,-1.3732,0.4612,1.2422)；Llama-3-8B RMSNorm 262,144 参数(1.0 MiB)、LayerNorm 524,288(2.0 MiB)；FP16 上限 65504}
\fol{RMSNorm 不是对平移不敏感吗？}{对，LayerNorm 第一步就减掉均值，所以对整体加常数完全不变；RMSNorm 会跟着变。但论文的实验说明起作用的主要是“除以尺度”，减均值对效果贡献很小。}
\fol{80 层模型到底省了什么？}{每层两个子层，一共 160 处残差相加紧接归一化；朴素实现每次都要写回 HBM 再读回来。融合后相加结果留在寄存器里，访存量减半以上。}
\fol{加载 Gemma 权重要注意什么？}{Gemma 的实现是 x*(1+gamma)，gamma 初始化为全 0，和 Llama 那种 gamma 全 1 的写法不同；按全 1 加载前向会完全错乱，必须先看对方的实现。}
\numline{别答错}{RMSNorm 对整体平移不敏感是错的；Pre-Norm 要在输出层前补 Final Norm。}
\phantomsection\addcontentsline{toc}{section}{37\quad RoPE 旋转位置编码}
\lec{37}{RoPE 旋转位置编码}
\askline{RoPE 凭什么只感知相对位置？}
\askline{base 从 1e4 调到 5e5 有什么用？}
\askline{部署时 RoPE 最容易踩什么坑？}
\vspace{3pt}
\noindent{\qlabel{答}}\begingroup\sloppy\begin{enumerate}[leftmargin=1.6em,itemsep=2.5pt,topsep=0pt,parsep=0pt,label={\sffamily\footnotesize\color{accent}\arabic*.}]
\item {\fontsize{10}{14.8}\selectfont RoPE 不往输入上加位置，而是在注意力里把 q、k 按自己的位置旋转：位置 m 的 q 转 m$\theta$，位置 n 的 k 转 n$\theta$。旋转不改变向量长度，只把两个向量的夹角改了 (m-n)$\theta$。\par}
\item {\fontsize{10}{14.8}\selectfont 所以 q$\cdot$k 只取决于位置差：TinyLM 里 (3,1)、(5,3)、(10,8) 三组点积都是 0.0069，差 3 的那组是 -0.1178；两个词一起往后挪一位，夹角不变，打分也不变。\par}
\item {\fontsize{10}{14.8}\selectfont 高维就是把维度两两配成 d/2 个平面，第 i 对的转速 $\theta$i = base\textasciicircum{}(-2i/d)，和正弦编码是同一组频率；Llama-3 每头 128 维就是 64 对，各转各的，点积是各对之和。\par}
\item {\fontsize{10}{14.8}\selectfont 和加法位置编码比，它只改 q、k，不碰 v 也不碰残差流：位置信息不再混进词向量，也不会经过层层变换被冲淡，而且一个可学参数都没有。\par}
\item {\fontsize{10}{14.8}\selectfont 它还自带远距离衰减：q = k = 全 1 的 128 维向量，距离 0 时点积 128，距离 100 降到 61.1，距离 1000 只剩 20.4，高频那几对已经转乱互相抵消。\par}
\item {\fontsize{10}{14.8}\selectfont base 决定最慢那一对的周期：base = 10000 时约 54410 个位置转一圈，Llama-3 用 500000 把它拉到约 256 万，能区分的距离更远，这是扩长上下文的第一手段。\par}
\end{enumerate}\endgroup
\numline{公式}{旋转 $\varphi$：$(a,b)\to(a\cos\varphi-b\sin\varphi,\ a\sin\varphi+b\cos\varphi)$；$\theta_i=\text{base}^{-2i/d}$}
\numline{数字}{TinyLM 4 维 2 对，$\theta$0=1、$\theta$1=0.01；base=10000 最快一对 6.3 个位置一圈、最慢 54410 个；Llama-3 base=500000 最慢约 2,559,196 个位置；衰减 128.0/124.2/85.6/61.1/20.4/-3.6 对应距离 0/1/10/100/1000/10000}
\fol{为什么不把 v 也一起旋转？}{旋转只是为了让打分带上相对位置，v 是被收集的内容，和位置无关；转了反而把位置信息混进输出，还白算一遍。}
\fol{长上下文扩展是怎么做的？}{基本都在转速上做文章：调大 base（Llama-3 用 500000），或者把低频那几对按比例拉伸，让最慢的指针转得更慢，能区分的距离就更远。}
\fol{部署时最容易踩什么坑？}{配对方式搞反。交替配对和前一半配后一半数学等价，但旋转基准不同；Hugging Face 的 Llama 用半半拆分，接错了模型直接吐乱码。}
\numline{别答错}{只旋转 q、k 不碰 v；配对方式搞反会出乱码；KV Cache 里存的是旋转之后的 k。}
\phantomsection\addcontentsline{toc}{section}{38\quad SwiGLU 与门控 FFN}
\lec{38}{SwiGLU 与门控 FFN}
\askline{Llama 的 FFN 为什么有三个矩阵？}
\askline{多一个矩阵，参数不是变多了吗？}
\askline{SwiGLU 底层怎么算得快？}
\vspace{3pt}
\noindent{\qlabel{答}}\begingroup\sloppy\begin{enumerate}[leftmargin=1.6em,itemsep=2.5pt,topsep=0pt,parsep=0pt,label={\sffamily\footnotesize\color{accent}\arabic*.}]
\item {\fontsize{10}{14.8}\selectfont 门控 FFN 把升维拆成两路：g = x$\cdot$W\_gate 当门、u = x$\cdot$W\_up 当内容，h = SiLU(g)$\odot$u，最后乘 W\_down 降维；Llama 权重里的 gate\_proj、up\_proj、down\_proj 就是这三步。\par}
\item {\fontsize{10}{14.8}\selectfont 它替代的是 ReLU(u) 那一刀：门由输入自己算出来，SiLU 给出连续的开度，可以只放行一部分，还能取到约 -0.28 的小负值，把内容反过来。\par}
\item {\fontsize{10}{14.8}\selectfont TinyLM 的 8 个神经元里，门开到 0.886 的第 3 个几乎全部放行 1.455，得 1.288；内容最大的是第 6 个（2.095），但门只开 0.277，只放出 0.581。\par}
\item {\fontsize{10}{14.8}\selectfont 为什么更好：Shazeer 2020 年的论文在等参数下比较，带门的 GEGLU、SwiGLU 都优于 ReLU、GELU 版本；直觉是“放不放”和“放什么”由两组权重分开学，乘法交互比加法加 ReLU 表达力强。\par}
\item {\fontsize{10}{14.8}\selectfont 参数不会白多：三个矩阵要对齐两个矩阵的 8d$^2$，d\_ff 要取约 8d/3；d=4096 时 10922，参数 134.2M，和 ReLU 的 4d 版本一样，Llama-3-8B 选了 14336（3.5d），每层 176.2M。\par}
\item {\fontsize{10}{14.8}\selectfont 底层不跑两次 GEMM：加载时把 W\_gate 和 W\_up 拼成 [d, 2d\_ff] 一次乘完，2d\_ff 宽的中间结果留在 SRAM 里原地做 SiLU 和逐元素相乘，再进 W\_down，省掉一次输入读取和一次内核发射。\par}
\end{enumerate}\endgroup
\numline{公式}{$g=x\cdot W_{\text{gate}}$，$u=x\cdot W_{\text{up}}$，$h=\operatorname{SiLU}(g)\odot u$，$y=h\cdot W_{\text{down}}$；等参数时 $d_{\text{ff}}\approx\frac{8d}{3}$}
\numline{数字}{TinyLM 8 个神经元，y=(-0.347,1.904,0.364,-1.628)；门偏移 -3/+3 时 |h| 之和 1.271 对 23.745；Llama-3-8B d\_ff=14336（3.5d），每层 FFN 176.2M 参数，权重拼成 4096$\times$28672；等参数 d\_ff$\approx$8d/3=10922 时 134.2M}
\fol{门是 0/1 开关吗？}{不是。SiLU(g) 是连续的，0.2、0.9 都行，还能是小负数（最低约 -0.28），所以既能放行一部分，也能把内容反号，比 ReLU 的全放全关细得多。}
\fol{训练时这三个矩阵有什么代价？}{中间 g、u 都是 14336 维，是全模型最宽的激活，反向要用所以要存下来，激活显存很大；现代训练用选择性重计算，反向时重算 g、u 去换显存。}
\fol{GLU 和 GELU 是一回事吗？}{不是。GELU 是激活函数，GLU 是“一路当门、一路当内容相乘”的结构；门用 GELU 叫 GEGLU，用 Swish（即 SiLU）才叫 SwiGLU。}
\numline{别答错}{门不是 0/1 开关；内容 u 不过激活；GLU 是结构、GELU 是激活函数，别混成一家。}
\phantomsection\addcontentsline{toc}{section}{39\quad MQA 与 GQA}
\lec{39}{MQA 与 GQA}
\askline{为什么只有 8 个 KV 头，不是 32 个？}
\askline{GQA 省的是显存还是计算量？}
\askline{底层真的把 K、V 复制了 4 份吗？}
\vspace{3pt}
\noindent{\qlabel{答}}\begingroup\sloppy\begin{enumerate}[leftmargin=1.6em,itemsep=2.5pt,topsep=0pt,parsep=0pt,label={\sffamily\footnotesize\color{accent}\arabic*.}]
\item {\fontsize{10}{14.8}\selectfont GQA 让多个查询头共用一组 K、V：32 个查询头分成 8 组，每组 4 个头共用一套，所以缓存里只有 8 个 KV 头；查询头还是 32 个，各自的注意力权重也一个没少。\par}
\item {\fontsize{10}{14.8}\selectfont 这笔账很直接：每 token 的 KV Cache = 2$\times$L$\times$n\_kv$\times$d\_head$\times$字节数。Llama-3-8B 用 MHA 是 512 KiB，GQA 的 8 组是 128 KiB，MQA 只剩 16 KiB。\par}
\item {\fontsize{10}{14.8}\selectfont 换算成并发：8192 token $\times$ 16 个请求，MHA 要 64 GiB，是 15 GiB 权重的四倍多；GQA 是 16 GiB，MQA 是 2 GiB，这就是一张 80 GB 卡能装多少请求的差别。\par}
\item {\fontsize{10}{14.8}\selectfont 三者是一族：MHA 是 n\_kv = n\_h，MQA 是 n\_kv = 1，GQA 夹在中间。2019 年 Shazeer 的 MQA 走到极端、效果有损；2023 年的 GQA 取折中，8 组时效果已经很接近 MHA。\par}
\item {\fontsize{10}{14.8}\selectfont 底层绝不在显存里复制：物理上只留 n\_kv 份，Kernel 里第 h 个查询头的线程用 h÷(n\_h/n\_kv) 算出组号，都去读同一个 KV 头，读进 SRAM 后同组 4 个头并发复用。\par}
\item {\fontsize{10}{14.8}\selectfont 省显存只是一半，另一半是速度：MHA 读一个 128 维 k 要用 256 字节，只服务 1 个头，算术强度 1 FLOP/Byte；GQA 同一份数据服务 4 个头，涨到 4 FLOP/Byte，而 decode 正是访存受限。\par}
\end{enumerate}\endgroup
\numline{公式}{$\text{kv\_group}=\frac{h}{n_h/n_{kv}}$；每 token KV Cache $=2\,L\,n_{kv}\,d_{\text{head}}\cdot\text{字节数}$}
\numline{数字}{Llama-3-8B：MHA 每 token 512 KiB、GQA(8) 128 KiB、MQA 16 KiB；8192$\times$16 请求 64/16/2 GiB；W\_k、W\_v 每层 33.6M/8.4M/1.0M 参数；算术强度 1 $\rightarrow$ 4 FLOP/Byte；config 里 num\_key\_value\_heads=8}
\fol{现成模型能直接改成 GQA 吗？}{不能硬改。同组各头的 K、V 直接取平均会明显掉效果，TinyLM 上“AI”的输出整行都变了，所以 GQA 论文要求再用约原训练量 5\% 的数据继续训练，让模型适应新结构。}
\fol{张量并行时会有什么限制？}{KV 头数要能被 TP 度整除。8 个 KV 头切 8 张卡很整齐，切 16 张卡就除不尽，只能在两张卡之间复制一份 KV 头，显存没法线性下降。}
\fol{那计算量真的没变吗？}{几乎没变。打分和加权求和的次数取决于查询头个数，那还是 32 个；GQA 省的只是要存、要读的 K、V，所以省的是显存和带宽，不是 FLOPs。}
\numline{别答错}{说“8 个头”是错的，要说 32 个查询头、8 个 KV 头；省的是显存和带宽，计算量几乎不变。}
\phantomsection\addcontentsline{toc}{section}{40\quad MLA 多头潜在注意力}
\lec{40}{MLA 多头潜在注意力}
\askline{MLA 怎么做到比 GQA 还省？}
\askline{压缩再还原，还原的开销不大吗？}
\askline{为什么还要单独留一段给 RoPE？}
\vspace{3pt}
\noindent{\qlabel{答}}\begingroup\sloppy\begin{enumerate}[leftmargin=1.6em,itemsep=2.5pt,topsep=0pt,parsep=0pt,label={\sffamily\footnotesize\color{accent}\arabic*.}]
\item {\fontsize{10}{14.8}\selectfont MLA 不减少 K、V 的组数，每个头照样有自己的 k 和 v，但它不缓存展开后的 K、V，只缓存一个压缩向量：c = x$\cdot$W\_DKV，用时再 k = c$\cdot$W\_UK、v = c$\cdot$W\_UV 还原。\par}
\item {\fontsize{10}{14.8}\selectfont DeepSeek-V3 的 d=7168、128 个头、每头 k、v 各 128 维，直接缓存是每层每 token 32768 个数；c 只有 512 维，赌的是这些头落在一个低秩子空间里，512 个自由度就够。\par}
\item {\fontsize{10}{14.8}\selectfont 还原的开销靠矩阵吸收消掉：q$\cdot$k = q$\cdot$(c$\cdot$W\_UK) = (q$\cdot$W\_UKᵀ)$\cdot$c，把 W\_UK 提前吸收进 q，直接拿新 token 的 q 和缓存里每个 c 点积，历史 token 的 k 从来不用展开。\par}
\item {\fontsize{10}{14.8}\selectfont Decode 走吸收，prefill 却相反：prefill 是计算受限，引擎在片上把 K、V 临时展开成常规多头，直接调 FlashAttention-2 跑满算力，算完就丢，写回显存的仍只有 c。\par}
\item {\fontsize{10}{14.8}\selectfont RoPE 会打断这条链：k 旋转后中间夹了随位置变化的 R\_n，吸收不了。DeepSeek 的做法是解耦 RoPE：128 维内容部分不旋转、可吸收，另投一小段 64 维 k\_R 做旋转，所有头共享。\par}
\item {\fontsize{10}{14.8}\selectfont 所以每层每 token 只存 512+64 = 576 个数，61 层共 68.6 KiB，是同规模 MHA 的 1/56.89、8 组 GQA 的 1/3.56，只比极限的 MQA 多一倍，效果反而超过同尺寸的 MHA。\par}
\end{enumerate}\endgroup
\numline{公式}{$c=x\cdot W_{\text{DKV}}$，$k=c\cdot W_{\text{UK}}$，$v=c\cdot W_{\text{UV}}$；$q\cdot k=(q\cdot W_{\text{UK}}^{\mathsf{T}})\cdot c+q_R\cdot k_R$}
\numline{数字}{DeepSeek-V3：d=7168、61 层、128 头、每头 k/v 各 128 维；c 512 维、k\_R 64 维，每层每 token 576 个数，61 层 68.6 KiB；MHA 3904 KiB(56.89$\times$)、GQA 244 KiB(3.56$\times$)、MQA 30.5 KiB(0.44$\times$)；q\_lora\_rank 1536}
\fol{吸收会不会改变模型结果？}{不会，这是纯粹的代数等价，只是换了矩阵乘法的结合顺序；可以用代码验证，显式展开和吸收两条路径的打分逐位相同，模型权重一个都没动。}
\fol{那 prefill 为什么不也用吸收？}{吸收后 Q 的维度会变大，还和 RoPE 的位置项拼在一起，套不上高度优化的 FlashAttention-2；prefill 本来就算力受限，展开成多头跑标准算子更划算。}
\fol{FlashMLA 解决的是什么问题？}{传统算子按头划分线程，128 个头会各自从 HBM 读同一份 c，重复 128 次；FlashMLA 把多个查询头打包进一个线程块，同一批 c 只读一次，在片上给所有头复用。}
\numline{别答错}{MLA 不减少头数，减的是缓存；吸收是推理期的代数等价，不改变模型本身。}
\phantomsection\addcontentsline{toc}{section}{41\quad MoE 混合专家}
\lec{41}{MoE 混合专家}
\askline{MoE 为什么能参数多、计算少？}
\askline{一个 token 的专家怎么选出来？}
\askline{专家负载不均怎么办？}
\vspace{3pt}
\noindent{\qlabel{答}}\begingroup\sloppy\begin{enumerate}[leftmargin=1.6em,itemsep=2.5pt,topsep=0pt,parsep=0pt,label={\sffamily\footnotesize\color{accent}\arabic*.}]
\item {\fontsize{10}{14.8}\selectfont MoE 只改 FFN：把一个大 FFN 换成 E 个独立专家，前面加一个 d$\times$E 的路由器，每个 token 只算 Top-K 个，其余完全旁路、不产生浮点运算。总参数定知识容量，激活参数定计算量，两者解耦。\par}
\item {\fontsize{10}{14.8}\selectfont 路由三步：先由 h(x)=x$\cdot$W\_r 得到 E 个分数，取最大的 k 个（Mixtral 8 选 2、DeepSeek-V3 256 选 8），只对这 k 个做局部 Softmax，再加权求和。全局 Softmax 分母上百项容易下溢。\par}
\item {\fontsize{10}{14.8}\selectfont 底层不是逐 token 循环：先把 S$\times$k 个任务按专家编号重排，记下原行号和权重；再用一个 Grouped GEMM kernel 传入各专家权重指针并行算完；最后按索引加权写回 [S, d]，逐 token 循环会被发射延迟卡死。\par}
\item {\fontsize{10}{14.8}\selectfont 新一代加了共享专家：它不参与打分、所有 token 必经，负责吸收连词介词和通用句式；路由专家就不必浪费容量背通用知识，可以做得更小更专。DeepSeek-V3 每层是 1 个共享加 256 个路由选 8 个。\par}
\item {\fontsize{10}{14.8}\selectfont 参数账：Mixtral 8x7B 共 467 亿总参数，单 token 只用 129 亿、占 28\%，计算量等于 130 亿稠密模型，知识容量接近 470 亿。DeepSeek-V3 更极致，6710 亿里只激活 370 亿，占 5.5\%。\par}
\item {\fontsize{10}{14.8}\selectfont 代价在显存和通信：下个 token 用谁事先不知道，所有专家都得常驻显存，DeepSeek-V3 即便 FP8 也要约 671 GB；专家分卡后每层都要 All-to-All，吞吐瓶颈常在网络，所以 MoE 天然吃大 batch。\par}
\end{enumerate}\endgroup
\numline{公式}{$\operatorname{MoE}(x)=\operatorname{FFN}_{\text{shared}}(x)+\sum_{e\in\text{TopK}}w_{e}\operatorname{FFN}_{e}(x)$，$h(x)=x\cdot W_{r}$}
\numline{数字}{DeepSeek-V3 总 6710 亿/激活 370 亿（5.5\%），每层 1 共享 + 256 路由选 8；Mixtral 8x7B 467 亿/129 亿（28\%）；FP8 权重约 671 GB；偏置均衡最忙/平均 1.81$\rightarrow$1.12}
\fol{专家是按学科分工的吗？}{不是。路由是逐 token、逐层独立做的，同一句话前后两个词可能选完全不同的专家，分工更多体现在词法、语法这些抽象特征上，很难用人类学科划分。}
\fol{不加辅助损失怎么均衡？}{给每个专家一个偏置 b\_e，只加在挑前 k 个那步；b\_e 不进 Softmax，权重仍由原始打分定。太忙的下调、太闲的上调，30 步就能把最忙/平均从 1.81 压到 1.12。}
\fol{MoE 吞吐一定比稠密高吗？}{不一定。batch 小时很多专家没活干、权重白占显存，还要付 All-to-All 开销；只有并发足够大、专家都被填满，单 token 计算量只有稠密的几分之一，优势才成立。}
\numline{别答错}{2N 的 N 要用激活参数，路由与 All-to-All 不在其中，故比同激活的稠密模型略慢。}
\phantomsection\addcontentsline{toc}{section}{42\quad 长上下文技术}
\lec{42}{长上下文技术}
\askline{上下文从 8K 扩到 128K，卡在哪？}
\askline{YaRN 是怎么处理位置的？}
\askline{滑动窗口为什么要留开头 token？}
\vspace{3pt}
\noindent{\qlabel{答}}\begingroup\sloppy\begin{enumerate}[leftmargin=1.6em,itemsep=2.5pt,topsep=0pt,parsep=0pt,label={\sffamily\footnotesize\color{accent}\arabic*.}]
\item {\fontsize{10}{14.8}\selectfont 上下文变长同时撞三堵墙：计算上注意力打分随长度涨，128K 时占每 token 计算量 82\%、1M 时 97\%；显存上 KV Cache 线性涨，1M 要 128 GiB；位置上模型只见过训练长度内的 RoPE 转角，超出去就发散。\par}
\item {\fontsize{10}{14.8}\selectfont 位置最隐蔽：直接外推到 32768，最慢那一对维度的转角从 0.0201 弧度变成 0.0805，模型从没见过；位置插值乘 8192/32768=1/4 把它拉回 0.0201，但相邻位置的转角差也缩成 1/4，近处的词就难区分。\par}
\item {\fontsize{10}{14.8}\selectfont YaRN 频段分治：高频（波长小于训练长度）完全不插值，保住局部语法和紧邻词对的分辨力；低频完全插值，把转角锁在训练范围内；中频平滑过渡；再乘温度修正，把被压平的注意力分布重新聚拢。\par}
\item {\fontsize{10}{14.8}\selectfont 显存和计算靠滑动窗口：每 token 只看最近 W 个，KV 上限就是 W。窗口 4096 时 131072 上下文只要 0.5 GiB，完整注意力要 16 GiB；更远的信息靠层间接力，叠 32 层最远碰到 131072 个位置之前。\par}
\item {\fontsize{10}{14.8}\selectfont 但必须钉住开头：softmax 权重和恒为 1，模型预训练时学会了把多余注意力倒到第 0 到 3 个 token 上。纯滑动窗口一旦丢掉它们，权重就泛滥到窗口内的正常词上，生成立刻崩溃复读。\par}
\item {\fontsize{10}{14.8}\selectfont 所以流式推理的标准做法是固定保留加动态滚动：永久锁住最开头 4 个 token 当 Attention Sink，后面接容量 W 的动态窗口，显存上限是 4+W，却能稳定生成几十万 token 不崩。\par}
\end{enumerate}\endgroup
\numline{公式}{位置插值 $\text{pos}\times\frac{L_{\text{train}}}{L_{\text{target}}}$；最慢维转角 $0.0201\to0.0805$（外推）$\to0.0201$（插值）；base 改 $2{,}044{,}497$}
\numline{数字}{8,192$\rightarrow$1.0 GiB/19.3 GFLOPs/22\%；131,072$\rightarrow$16.0 GiB/83.7/82\%；1,048,576$\rightarrow$128.0 GiB/564.8/97\%；窗口 4096 只需 0.5 GiB；32 层$\times$4096 接力 131,072}
\fol{插值和改 base 有什么不同？}{插值把位置整体乘小，所有维度分辨率一起降；改 base 只放大低频维度的波长，最慢那对在 32768 处仍是 0.0201，最快一对相邻差还保持 1.000，高频不受损。}
\fol{滑动窗口会丢远距离信息吗？}{会丢精度。窗口只留最近 W 个，远处只能靠层间接力间接拿到；需要精确引用很远处原文的任务，效果确实不如完整注意力。}
\fol{号称 128K 就真能用满吗？}{两回事。放得下是位置和显存问题，用得好不好要靠大海捞针这类测试验证；位置插值后通常还要用少量长文本继续训练，直接压缩就上线效果会明显掉。}
\numline{别答错}{128K 指 13 万个 token 而非 13 万字，中文一个字可能切出多个 token。}
\phantomsection\addcontentsline{toc}{section}{43\quad 架构演进对照}
\lec{43}{架构演进对照}
\askline{三代模型的结构差在哪？}
\askline{KV Cache 为什么越做越小？}
\askline{给你 config.json 你怎么看？}
\vspace{3pt}
\noindent{\qlabel{答}}\begingroup\sloppy\begin{enumerate}[leftmargin=1.6em,itemsep=2.5pt,topsep=0pt,parsep=0pt,label={\sffamily\footnotesize\color{accent}\arabic*.}]
\item {\fontsize{10}{14.8}\selectfont 五年里骨架几乎没变：都是仅解码器、Pre-Norm、残差、最后归一化加输出层，目标都是预测下一个 token，GPT-2 甚至已经用了 Pre-Norm。变的全是零件，改动有两条清晰的主线。\par}
\item {\fontsize{10}{14.8}\selectfont 主线一是让每层更好用：LayerNorm 换 RMSNorm 更省，绝对位置嵌入换 RoPE 且无 1024 行硬限制，GELU 换门控 SwiGLU 同参数更好，再去掉线性层偏置。这套到 Llama 已成开源标准配置。\par}
\item {\fontsize{10}{14.8}\selectfont 主线二是让推理更省：MHA 到 GQA 到 MLA，每 token 的 KV Cache 从 300 KiB 降到 128 KiB 再到 68.6 KiB。参数涨了 400 多倍，单 token 的缓存反而越来越小。\par}
\item {\fontsize{10}{14.8}\selectfont 参数和计算还脱钩：DeepSeek-V3 用 MoE 做 6710 亿总参数、单 token 只激活约 370 亿，知识容量远大于 Llama-3-8B，计算量只是四倍多；训练就大量用 FP8，推理可直接用 FP8 权重，省一半显存和带宽。\par}
\item {\fontsize{10}{14.8}\selectfont 读 config.json 就能判断推理特点：num\_key\_value\_heads 定 GQA、kv\_lora\_rank 定 MLA、num\_experts 定 MoE、rope\_theta 定上下文。\par}
\item {\fontsize{10}{14.8}\selectfont 但同名模型不同版本结构可能不同：Llama-2-7B 是 MHA、只有 70B 才用 GQA，Llama-3 全系 GQA；8B 和 671B 也不能直接比计算量，MoE 每个 token 只用激活那部分，比推理成本要看激活参数和 KV。\par}
\end{enumerate}\endgroup
\numline{公式}{每 token KV $=2\times\text{层数}\times(\text{kv 头数}\times d_{\text{head}})\times\text{字节数}$；GPT-2 XL：$2\times48\times1600\times2=300$ KiB}
\numline{数字}{15.6 亿/80.3 亿/6710 亿（激活 370 亿）；48/32/61 层；宽度 1600/4096/7168；词表 50257/128256/129280；KV 300/128/68.6 KiB；上下文 1024/8192/128K}
\fol{MLA 凭什么压到 68.6 KiB？}{它把每层 K、V 吸收成 512 维潜向量加 64 维 RoPE 分量，每 token 只缓存这 576 个数，61 层 16 位就是 68.6 KiB；代价是要专门算子做权重吸收。}
\fol{开源模型都长一个样了吗？}{结构上确实越来越像 Llama：RMSNorm、RoPE、SwiGLU、GQA 成了默认组合。差异主要在 KV 用 GQA 还是 MLA、是否上 MoE、位置编码怎么扩、发布精度多少。}
\fol{FP8 对推理到底省在哪？}{省一半显存和带宽。DeepSeek-V3 训练就用 FP8 混合精度，发布后可直接用 FP8 权重跑，6710 亿参数降到约 671 GB，decode 每步读的权重字节数也减半。}
\numline{别答错}{别拿总参数比推理成本：MoE 只看激活参数和 KV Cache；同名模型不同版本结构也可能不同。}
\phantomsection\addcontentsline{toc}{section}{44\quad 混合架构与线性注意力}
\lec{44}{混合架构与线性注意力}
\askline{线性注意力凭什么和 n 无关？}
\askline{状态固定大小会丢什么？}
\askline{为什么工业界用混合架构？}
\vspace{3pt}
\noindent{\qlabel{答}}\begingroup\sloppy\begin{enumerate}[leftmargin=1.6em,itemsep=2.5pt,topsep=0pt,parsep=0pt,label={\sffamily\footnotesize\color{accent}\arabic*.}]
\item {\fontsize{10}{14.8}\selectfont 核心是换乘法顺序：完整注意力先算 n$\times$n 的 QK\textasciicircum{}T，而 K\textasciicircum{}T V 只有 d$\times$d。没有 softmax 时结合律允许 (QK\textasciicircum{}T)V = Q(K\textasciicircum{}T V)，线性注意力就用核函数 φ 把 q、k 映成非负特征，替掉 softmax。\par}
\item {\fontsize{10}{14.8}\selectfont 于是可以写成递推：维护 d$\times$d 的状态 S，每来一个 token 就累加 φ(k\_t)\textasciicircum{}T v\_t，输出时用 φ(q\_t) 读一次，单步代价 d$^2$，与上下文无关。全部历史都被压进这一个固定大小的记忆里。\par}
\item {\fontsize{10}{14.8}\selectfont 生成一个 token，普通注意力每层每头约 n$\cdot$d 且要存 n 份 KV，线性固定 d$^2$。n=131072、d=128 时普通 134.2 MFLOPs、KV 32.0 MiB/层，线性 0.1 MFLOPs、16384 个数，差约一千倍。\par}
\item {\fontsize{10}{14.8}\selectfont 短板是精确召回：状态大小固定，等于把很多张照片叠印在同一张底片上。往 16$\times$16 状态里存 4、16、64 对 (k,v)，取回误差是 0.67、0.95、1.83；它擅长总体把握，不擅长在长文里精确找出一个电话号码。\par}
\item {\fontsize{10}{14.8}\selectfont Mamba 补了两个死穴：一是选择性，Δ\_t、B\_t、C\_t 都由当前输入投影得到，停用词能让 Δ\_t 趋近 0 跳过，遇到关键实体就清空旧记忆写入新特征；二是把并行扫描和状态更新封进单个 kernel，状态只在 SRAM 流转。\par}
\item {\fontsize{10}{14.8}\selectfont 所以现实方案是混合架构：7/8 的层用 Mamba 或线性注意力压缩长文主干，少数层保留完整 softmax 注意力做精确检索。131K 下每层 KV 从 32.0 MiB 降到 4.0 MiB、省 88\%，Jamba 就走这条路。\par}
\end{enumerate}\endgroup
\numline{公式}{$S_{t}=S_{t-1}+\varphi(k_{t})^{\top}v_{t}$，$o_{t}\propto\varphi(q_{t})\cdot S_{t}$；$(QK^{\top})V=Q(K^{\top}V)$}
\numline{数字}{n=131,072、d=128：普通 134.2 MFLOPs、KV 32.0 MiB/层；线性 0.1 MFLOPs、16,384 个数；混合（7:1）16.8 MFLOPs、4.0 MiB、省 88\%；取回误差 0.67/0.95/1.83}
\fol{这算把注意力做成线性了吗？}{不是。它仍在做软寻址，只是用核函数替掉 softmax，让矩阵乘法的结合律生效；代价是打分不再是指数归一化的精确匹配，很远很细的对应关系会变弱。}
\fol{并行扫描解决什么问题？}{解决训练不能串行的问题。每步参数都随输入变，逐 step 循环 GPU 利用率会跌到冰点；它用结合律把动态循环重写成并行前缀和扫描，并把整个流程封进一个 kernel。}
\fol{混合架构里普通层留几层？}{看任务对精确召回的要求。7:1 时 131K 下 KV 还能省 88\%，总代价主要由那几层普通注意力决定；比例再降，长文里大海捞针这类任务会最先掉点。}
\numline{别答错}{线性注意力并没有把注意力线性化，它是用可压缩的固定记忆换掉精确 KV 检索，状态固定不等于无损。}
\chapter{训练全流程}
{\footnotesize\sffamily\color{grayt}第 45--51 讲}\par\vspace{-4pt}
\phantomsection\addcontentsline{toc}{section}{45\quad 预训练与数据}
\lec{45}{预训练与数据}
\askline{大模型预训练数据从哪来？}
\askline{原始网页为什么要洗几遍？}
\askline{重复数据为什么必须去掉？}
\vspace{3pt}
\noindent{\qlabel{答}}\begingroup\sloppy\begin{enumerate}[leftmargin=1.6em,itemsep=2.5pt,topsep=0pt,parsep=0pt,label={\sffamily\footnotesize\color{accent}\arabic*.}]
\item {\fontsize{10}{14.8}\selectfont 来源以网页为主，最常用 Common Crawl，从 2008 年起持续抓取、累积数千亿网页；再加书籍百科论文、公开代码仓库和数学教材。按配比混合，Llama-3 约一半通用知识、四分之一数学推理、17\% 代码、8\% 多语言。\par}
\item {\fontsize{10}{14.8}\selectfont Llama-3 用了 15T 以上 token，按每 token 约 4 字节粗算约 60 TB，一个人每秒读 1000 个 token、不吃不睡要读 476 年。这些靠成千上万个 CPU 核跑几周，是不在 GPU 上的基础设施。\par}
\item {\fontsize{10}{14.8}\selectfont 清洗六步：抽正文扔掉标签脚本导航、语言识别分桶、完全与近似去重、质量过滤（太短、符号多、重复片段多，加一个小分类器判像不像高质量文本）、安全隐私过滤、按配比分词。演示里 9 篇只剩 4 篇。\par}
\item {\fontsize{10}{14.8}\selectfont 近似去重最难：把文本切成 3 字片段，算两篇文档片段集合的 Jaccard 相似度。第 3 篇只把喜欢吃鱼改成爱吃鱼，相似度 0.79 超过 0.7 就被丢掉。几十亿篇不可能两两比，实际先压成 MinHash 指纹再找相似。\par}
\item {\fontsize{10}{14.8}\selectfont 去重之所以关键：重复内容会让模型反复看到同一段文本、倾向死记而不是学可推广的规律，还白烧算力；充分去重后往往用更少 token 就训出更好的模型。质量过滤同理，广告页重复片段能占到 61\%。\par}
\item {\fontsize{10}{14.8}\selectfont 但数据不是越多越好：低质量重复数据会拖累模型，趋势是在质量上多下功夫。而且预训练数据里没有答案，它只教接着往下写，知识还有截止时间，之后的事只能靠检索或联网来补。\par}
\end{enumerate}\endgroup
\numline{公式}{$\operatorname{Jaccard}(A,B)=\frac{|A\cap B|}{|A\cup B|}$，$A$、$B$ 为 3 字片段集合，阈值 $0.7$}
\numline{数字}{15T token $\approx$ 60 TB；每秒 1000 token 要读 476 年；9 篇$\rightarrow$4 篇（3 中文 1 英文）；Jaccard 0.79 > 0.7；广告重复片段占 61\%；Llama-3 配比约 50\%/25\%/17\%/8\%}
\fol{去重能省多少训练算力？}{省的是重复样本上白跑的那部分。同样算力下，充分去重的数据往往用更少 token 就达到同样的损失；重复率高时模型把容量花在死记这些文本上，泛化更差。}
\fol{质量过滤怎么规模化做？}{规则先粗筛长度、符号比例、重复片段比例这些硬指标；再用一个已有的大模型给少量网页打教育价值这类分，训一个小分类器去扫几十亿篇，成本才压得住。}
\fol{中文数据多会影响推理吗？}{会，通过分词器。分词器是在预训练数据上训出来的，中文语料多，中文切出来的 token 就少，同样一段话更短，prefill 计算和 KV Cache 都更省。}
\numline{别答错}{预训练数据里没有答案，它只教接着往下写；会答题是 SFT 的事，且知识有截止时间。}
\phantomsection\addcontentsline{toc}{section}{46\quad Scaling Laws}
\lec{46}{Scaling Laws}
\askline{固定算力，模型该做多大？}
\askline{Chinchilla 的结论是什么？}
\askline{Llama-3-8B 为什么过度训练？}
\vspace{3pt}
\noindent{\qlabel{答}}\begingroup\sloppy\begin{enumerate}[leftmargin=1.6em,itemsep=2.5pt,topsep=0pt,parsep=0pt,label={\sffamily\footnotesize\color{accent}\arabic*.}]
\item {\fontsize{10}{14.8}\selectfont 损失可以预测：对数坐标下损失随参数量和数据量几乎成直线，即幂律，小规模实验就能外推大规模。Chinchilla 拟合式 L = 1.69 + 406.4/N\textasciicircum{}0.34 + 410.7/D\textasciicircum{}0.28，1.69 是语言本身的不确定性。\par}
\item {\fontsize{10}{14.8}\selectfont 训练总计算量约 C $\approx$ 6ND，固定 C 后把 D = C/(6N) 代回去，损失只随 N 变，是一条 U 形曲线：N 太小容量不够，太大又分不到足够数据。C=10$^2$$^2$ 时最优是 5.15B 配 323B token、每参数 62.8 个。\par}
\item {\fontsize{10}{14.8}\selectfont 最低点随算力移动：算力从 10$^2$⁰ 到 10$^2$$^4$，最优模型从 0.64B 涨到 41.35B，数据从 26B 涨到 4030B token。常引用的每参数约 20 个 token 来自论文另一种估计，用这个公式算是几十到上百，方向一致。\par}
\item {\fontsize{10}{14.8}\selectfont Chinchilla 之前主流是把模型做大、数据给少，GPT-3 有 1750 亿参数却只训约 3000 亿 token，这并不划算：一个 700 亿参数、训了 1.4 万亿 token 的模型效果就超过了大得多的模型。\par}
\item {\fontsize{10}{14.8}\selectfont 今天的模型反而故意过度训练：Llama-3-8B 训练算力约 7.23$\times$10$^2$$^3$ FLOPs，每参数约 1868 个 token；同算力下公式给出的最优是约 36B、3.4T token、损失 1.923，比 8B 的 1.948 还低。\par}
\item {\fontsize{10}{14.8}\selectfont 因为 Chinchilla 只算训练的账，没算推理的账。训练一次却要用千百万次，推理每 token 计算量是 2N，36B 的成本是 8B 的 4.4 倍。所以要按训练加推理的总成本选一个负担得起的大小，再用远超比例的数据训到底。\par}
\end{enumerate}\endgroup
\numline{公式}{$L(N,D)=1.69+\frac{406.4}{N^{0.34}}+\frac{410.7}{D^{0.28}}$；$C\approx6ND$}
\numline{数字}{C=10$^2$$^2$$\rightarrow$5.15B/323B/62.8 token 每参数/2.139；C=10$^2$⁰$\rightarrow$0.64B/26B；C=10$^2$$^4$$\rightarrow$41.35B/4030B/97.5/1.911；Llama-3-8B C=7.23$\times$10$^2$$^3$、1868 token/参数、1.948；同算力最优 36B/3.4T/1.923}
\fol{损失更低就代表能力更强吗？}{不一定。公式预测的是平均交叉熵，某些具体能力比如多位数乘法可能在某个规模附近突然出现，和平均损失的平滑下降并不同步，这就是常说的涌现。}
\fol{这套公式能直接套吗？}{常数不能直接搬。E、A、B、$\alpha$、$\beta$ 依赖数据、分词器和模型结构，换一套数据要重新拟合；能推广的是幂律形式和一起放大的结论，做法是在小规模上拟合再外推。}
\fol{推理侧也有 Scaling 吗？}{有。推理时多花计算，比如生成更长的思维链或多次采样投票，效果随 test-time 计算量上升；这把推理成本变成能力的又一个旋钮，也让推理效率更值钱。}
\numline{别答错}{常数不通用，换数据要重拟合；且最优只针对训练算力，算上推理成本后会偏向更小的模型。}
\phantomsection\addcontentsline{toc}{section}{47\quad 监督微调 SFT}
\lec{47}{监督微调 SFT}
\askline{基座模型和指令模型差在哪？}
\askline{SFT 为什么要做损失掩码？}
\askline{对话模板用错了会怎样？}
\vspace{3pt}
\noindent{\qlabel{答}}\begingroup\sloppy\begin{enumerate}[leftmargin=1.6em,itemsep=2.5pt,topsep=0pt,parsep=0pt,label={\sffamily\footnotesize\color{accent}\arabic*.}]
\item {\fontsize{10}{14.8}\selectfont SFT 就是拿一批提示---回答对，继续用预测下一个 token 的方式训练基座模型，方法和预训练几乎一样，只多两个细节：对话模板和只对回答算损失。它教的是格式和风格，知识还是预训练给的。\par}
\item {\fontsize{10}{14.8}\selectfont 数据量几万到几百万条，比预训练少好几个数量级；但质量和多样性远比数量重要，几千条精心挑选的高质量对话就能让一个基座模型学会像助手一样说话。\par}
\item {\fontsize{10}{14.8}\selectfont 模板用特殊 token 标角色和轮次，如 Llama-3 的 start\_header\_id、user、eot\_id，eot\_id 收一轮。推理时必须套同一模板、末尾写 assistant，出 eot\_id 才停；模板错会胡言乱语或停不住。\par}
\item {\fontsize{10}{14.8}\selectfont 损失掩码是关键细节：给每个位置一个 0 或 1 的标记，只有回答部分（含结束标记）参与交叉熵平均，提示不算。否则模型会浪费容量去学用户会怎么提问，甚至学会替用户说话。\par}
\item {\fontsize{10}{14.8}\selectfont ToyLM 实验：一条对话、只训输出层、学习率 0.5、50 步，你学后面就从推理 0.67 变成 AI 1.00；没教过的我学也从 AI 0.67 变成 AI 1.00，因为我和你的嵌入本来就相近，影响会扩散。\par}
\item {\fontsize{10}{14.8}\selectfont 边界也要说清：这种扩散既是举一反三也是风险，数据太窄、训练太猛会灾难性遗忘，所以学习率很小、轮数不多；它还要更新全部参数，显存和预训练同级，这是 LoRA 要解决的问题。\par}
\end{enumerate}\endgroup
\numline{公式}{$L=-\sum_{t\in\text{answer}}\log\pi(y_{t}\mid y_{<t})$，$\text{mask}\in\{0,1\}$}
\numline{数字}{数据几万\textasciitilde{}几百万条，几千条高质量即可；ToyLM 学习率 0.5、50 步；推理 0.67$\rightarrow$AI 1.00；SFT 常只有几亿\textasciitilde{}几十亿 token；全参数微调每参数约 16 字节}
\fol{掩码会让提示不参与计算吗？}{不会。遮的是损失不是注意力，提示 token 照常做前向、照常落在因果掩码内被回答的 token 看到，只是它们自己那个位置的预测不计入损失。}
\fol{SFT 能往模型里灌新知识吗？}{效果通常不好，还容易让它学会自信地编造。SFT 主要教格式和风格，新知识要靠预训练或继续预训练；服务时要补最新事实，走检索更稳。}
\fol{多轮对话为什么提示变长？}{每一轮都要把系统提示和历史消息重新拼进输入。所以多轮请求的 prefill 越来越重，而很多请求共享同一个系统提示，这正是前缀缓存能省下的部分。}
\numline{别答错}{掩码遮的是损失不是注意力；SFT 教的是格式，指望它灌知识往往只得到自信的编造。}
\phantomsection\addcontentsline{toc}{section}{48\quad RLHF 与 DPO}
\lec{48}{RLHF 与 DPO}
\askline{为什么用比较而不用打分？}
\askline{RLHF 为什么要四个模型？}
\askline{DPO 凭什么不用奖励模型？}
\vspace{3pt}
\noindent{\qlabel{答}}\begingroup\sloppy\begin{enumerate}[leftmargin=1.6em,itemsep=2.5pt,topsep=0pt,parsep=0pt,label={\sffamily\footnotesize\color{accent}\arabic*.}]
\item {\fontsize{10}{14.8}\selectfont 偏好数据是三元组：一个提示、两个回答、人的判断哪个更好，回答通常由 SFT 后的模型对同一提示采样两次生成。比较比打分可靠，人很难稳定说出这个回答值 7.3 分，却很容易说这个比那个好。\par}
\item {\fontsize{10}{14.8}\selectfont RLHF 第一步训奖励模型：结构同语言模型，输出层换成单个标量；用 Bradley--Terry，P(A$\succ$B)=σ(r\_A-r\_B)，分差 0/1/3 对应 0.5/0.731/0.953，损失是 -ln σ(r\_好-r\_差)。\par}
\item {\fontsize{10}{14.8}\selectfont 第二步用 PPO：让模型生成回答、奖励模型打分，再调参数让高分回答概率变大；目标里要减 $\beta$ 乘与参考模型的 KL 距离，防止它发现某种怪说法能骗高分就一直说，这就是奖励投机。\par}
\item {\fontsize{10}{14.8}\selectfont RLHF 的代价是复杂：要同时装策略、参考、奖励三个模型，PPO 还要价值模型，四个大模型同时在显存里，训练很不稳定。每步还要生成大量回答，这部分纯推理常占训练时间的大头。\par}
\item {\fontsize{10}{14.8}\selectfont DPO 直接给出损失：隐式奖励 = $\beta$$\cdot$(ln $\pi$(y) - ln $\pi$\_ref(y))，即相对参考模型的对数概率提升，再套 Bradley--Terry 损失 -ln σ(隐式奖励(y\_w) - 隐式奖励(y\_l))。只留策略和参考两个模型。\par}
\item {\fontsize{10}{14.8}\selectfont 这个相对项正好起偏离惩罚的作用，梯度前还乘着 σ(-差值)：好差没分开时权重接近 0.5 推得用力，分得很开时接近 0 几乎不推。ToyLM 上 $\beta$=0.5、30 步把苹果从 0.501 推到 0.981。\par}
\end{enumerate}\endgroup
\numline{公式}{$P(A\succ B)=\sigma(r_{A}-r_{B})$；DPO 损失 $=-\ln\sigma\bigl(\beta[(\ln\pi(y_{w})-\ln\pi_{\text{ref}}(y_{w}))-(\ln\pi(y_{l})-\ln\pi_{\text{ref}}(y_{l}))]\bigr)$}
\numline{数字}{$\beta$=0.5；30 步 0.501$\rightarrow$0.884$\rightarrow$0.981；损失 0.693$\rightarrow$0.309$\rightarrow$0.129；分差 0/1/3$\rightarrow$概率 0.5/0.731/0.953、损失 0.693/0.313/0.049；RLHF 四个模型 vs DPO 两个}
\fol{$\beta$ 调大调小会怎样？}{$\beta$ 越大越容易被推开，好差回答的差距被放大，学得快但更容易偏离参考模型；$\beta$ 小则推得温和保守，常见取值在 0.1 到 0.5，看偏离和效果的平衡。}
\fol{奖励模型被钻空子怎么办？}{一是 KL 惩罚，或者 DPO 里相对参考模型的那一项，限制它跑太远；二是用更多偏好数据、控制长度偏差把奖励模型做稳；三是定期用人工抽检去校准它。}
\fol{线上怎么用奖励模型？}{常见做法是对同一提示采样多个候选回答，用奖励模型打分排序取最高的返回，也就是 best-of-n；代价是推理算力乘以 n，所以一般只在关键场景打开。}
\numline{别答错}{DPO 不是强化学习但目标相同，它用离线偏好数据直接训，训练中不生成新回答，也省掉奖励模型。}
\phantomsection\addcontentsline{toc}{section}{49\quad 推理模型与强化学习}
\lec{49}{推理模型与强化学习}
\askline{推理模型是怎么训出来的？}
\askline{GRPO 的优势怎么算？}
\askline{先想 8000 token，对推理系统意味着什么？}
\vspace{3pt}
\noindent{\qlabel{答}}\begingroup\sloppy\begin{enumerate}[leftmargin=1.6em,itemsep=2.5pt,topsep=0pt,parsep=0pt,label={\sffamily\footnotesize\color{accent}\arabic*.}]
\item {\fontsize{10}{14.8}\selectfont 推理模型靠强化学习学会「先想再答」，奖励不来自奖励模型，而来自可以自动验证的结果：数学题看最终答案和标准答案是否一致，编程题看代码能不能通过测试用例，不管过程怎么写。\par}
\item {\fontsize{10}{14.8}\selectfont 为什么有用，回到自回归生成：每多写一个思考 token，就多一次前向，后面的步骤还能看到前面写下的中间结果，于是思考越长，花在这道题上的计算越多，这就是思维链有效的原因。\par}
\item {\fontsize{10}{14.8}\selectfont 另一条路是多次采样再投票：单次正确率 0.6 时，采 5 次取多数能到 0.771，采 15 次到 0.981；这类在推理时多花计算换正确率的做法统称 test-time scaling。\par}
\item {\fontsize{10}{14.8}\selectfont GRPO 对同一道题采一组回答，优势 =（它的奖励 - 组内平均）/ 组内标准差，为正就推高生成概率；8 个里答对 3 个时均值 0.375、标准差 0.484，答对的优势 +1.291、答错的 -0.775。\par}
\item {\fontsize{10}{14.8}\selectfont 全对或全错时每个回答都等于均值，优势全是 0，这道题对训练没有任何贡献，所以题目难度要合适；和 PPO 比，GRPO 不用单独训一个价值模型来估平均水平，直接用组内均值代替，省下一个大模型的显存和计算。\par}
\item {\fontsize{10}{14.8}\selectfont 代价在推理侧：先想 8000 个 token 再答 500 个，生成的 token 和响应时间都接近直接回答的 17 倍；Llama-3-8B 下 KV Cache 从 62.5 MiB 涨到 1062 MiB，RL 训练的大头也是推理。\par}
\end{enumerate}\endgroup
\numline{公式}{$\text{优势}=\frac{\text{奖励}-\text{组内平均}}{\text{组内标准差}}$；$\text{pass@k}=1-\frac{C(n-c,\,k)}{C(n,\,k)}$}
\numline{数字}{8 采样答对 3：均值 0.375、标准差 0.484，优势 +1.291 / -0.775；多数投票 0.6 $\rightarrow$ 0.771（5 次）$\rightarrow$ 0.981（15 次）；思考 8000 + 回答 500 是 17 倍；KV Cache 62.5 MiB $\rightarrow$ 1062 MiB}
\fol{写出的思考过程可信吗？}{不保证。思考过程只是模型生成的文本，不一定忠实反映它内部怎么得到答案；奖励只看最终结果，中间每一步对不对没有约束。}
\fol{17 倍的 token 成本怎么扛？}{靠长序列优化：KV Cache 管理、chunked prefill、投机解码，外加思考预算------限制最多想多少 token，或在快速回答和深度思考之间切换，毕竟按 token 计费。}
\fol{只有答案可验证的任务能做吗？}{不是。写作、对话这类没有标准答案的任务还要靠第 48 讲的偏好学习，实际的推理模型往往把可验证奖励和偏好学习结合着用。}
\numline{别答错}{「推理模型」的推理是 reasoning 不是 inference；思考过程不保证忠实。}
\phantomsection\addcontentsline{toc}{section}{50\quad LoRA 参数高效微调}
\lec{50}{LoRA 参数高效微调}
\askline{LoRA 为什么能省显存？}
\askline{秩 r 怎么选，参数省多少？}
\askline{训练完合并还是不合并？}
\vspace{3pt}
\noindent{\qlabel{答}}\begingroup\sloppy\begin{enumerate}[leftmargin=1.6em,itemsep=2.5pt,topsep=0pt,parsep=0pt,label={\sffamily\footnotesize\color{accent}\arabic*.}]
\item {\fontsize{10}{14.8}\selectfont 全量微调要为每个参数存梯度和 Adam 状态，混合精度下每参数 16 字节，Llama-3-8B 光这些就 128.5 GB，一张 80 GB 卡放不下；LoRA 冻结原权重，只在旁边训很少的新参数。\par}
\item {\fontsize{10}{14.8}\selectfont 它假设微调需要的修改量 ΔW 是低秩的，写成 A（d\_in$\times$r）乘 B（r$\times$d\_out），r 常用 8、16、64；前向是 y = xW + (x A)B，B 初始化为全 0，所以一开始输出和原模型完全一样，训练平稳出发。\par}
\item {\fontsize{10}{14.8}\selectfont 参数从 d\_in$\times$d\_out 降到 r$\times$(d\_in+d\_out)：4096$\times$4096 的 q\_proj、r = 16 时，W 有 1678 万个参数，A、B 只有 13.1 万个，不到 1\%。\par}
\item {\fontsize{10}{14.8}\selectfont 每层 7 个线性层都加 LoRA，r = 16 时全模型 4194 万个可训练参数、占 0.52\%，r = 8 是 0.26\%，r = 64 是 2.09\%；秩越大能表达的修改越丰富，参数和小数据上的过拟合风险也越大。\par}
\item {\fontsize{10}{14.8}\selectfont 显存账更直观：冻结权重用 BF16 存一份 16.1 GB，不需要梯度和 Adam 状态，只有 4194 万个 LoRA 参数按 16 字节算合 0.67 GB，总共不到 17 GB；QLoRA 把权重压到 4 位后约 4 GB。\par}
\item {\fontsize{10}{14.8}\selectfont 推理有两种用法：合并成 W + A$\cdot$B 后结构和原模型完全一样，零额外开销；不合并就把几十 MB 的适配器单独存，一个基座挂 100 个客户，各自的 (xA)B 小路分别算，S-LoRA、Punica 专门做这种混批算子。\par}
\end{enumerate}\endgroup
\numline{公式}{$\Delta W=A\cdot B$，$A\colon d_{\text{in}}\times r$，$B\colon r\times d_{\text{out}}$；$y=xW+(xA)B$；参数 $r\times(d_{\text{in}}+d_{\text{out}})$}
\numline{数字}{r=8 2100 万 / 0.26\%；r=16 4194 万 / 0.52\%；r=64 1.68 亿 / 2.09\%；q\_proj r=16：1678 万 $\rightarrow$ 13.1 万；显存 128.5 GB $\rightarrow$ 16.1 + 0.67 GB；QLoRA 4.0 + 0.67 GB}
\fol{r 是不是越大越好？}{不是。r 越大能表达的修改越丰富，但参数和显存线性增长，小数据上更容易过拟合；很多任务 r = 8 或 16 就够，2.09\% 的 r = 64 通常没必要。}
\fol{LoRA 能灌进大量新知识吗？}{不太行。它和 SFT 一样擅长调整模型的行为和输出格式；要灌大量新知识，低秩的修改量不够，那得继续预训练或全量微调。}
\fol{不合并的额外开销有多大？}{每个线性层多两次小矩阵乘法，r 很小时算力可忽略，但多一次访存和算子调用，对延迟有一点影响；只服务一个适配器时通常直接合并。}
\numline{别答错}{合并后推理计算量不变，LoRA 省的是训练显存和存储；秩不是越大越好。}
\phantomsection\addcontentsline{toc}{section}{51\quad 模型评测}
\lec{51}{模型评测}
\askline{怎么判断一个模型好不好？}
\askline{基准分数虚高是怎么来的？}
\askline{为什么说评测本身是推理？}
\vspace{3pt}
\noindent{\qlabel{答}}\begingroup\sloppy\begin{enumerate}[leftmargin=1.6em,itemsep=2.5pt,topsep=0pt,parsep=0pt,label={\sffamily\footnotesize\color{accent}\arabic*.}]
\item {\fontsize{10}{14.8}\selectfont 判模型好不好没有单一尺子：困惑度看预测下一个 token 准不准，基准测试拿标准答案算正确率，开放式任务靠模型当裁判或人类对战排名，三类各管一段。\par}
\item {\fontsize{10}{14.8}\selectfont 困惑度就是 e\textasciicircum{}loss，不用人工标注、算得快，适合比较同一模型的不同版本，比如量化前后在 WikiText-2 上的困惑度；但它不告诉你模型会不会解题，不同分词器之间也不能直接比。\par}
\item {\fontsize{10}{14.8}\selectfont 基准按判分方式分三类：选择题如 MMLU 覆盖 57 个学科，把每个选项接在题目后面比概率，不生成 token；答案匹配如 GSM8K 提取最终答案比对；代码测试如 HumanEval 用测试用例算 pass@k。\par}
\item {\fontsize{10}{14.8}\selectfont 最大的坑是数据污染：题库和答案大多公开在网上，预训练数据也是从网上抓的，模型可能背过答案、分数虚高；对策是训练数据里查重剔除、定期换没公开过的新题、用私有测试集。\par}
\item {\fontsize{10}{14.8}\selectfont 开放式回答靠 LLM-as-a-judge 或人类对战排名，后者用 Elo 算分：A 1200 分、B 1000 分时 A 赢的预期概率 0.760，赢了只加约 7.7 分，爆冷输了 A 掉到 1175.7。\par}
\item {\fontsize{10}{14.8}\selectfont 对做推理系统的人，评测就是一次大批量离线推理：选择题只要一次 prefill、四个选项共享前缀还能吃前缀缓存，需要完整解答的数学和代码基准就要大量 decode，推理模型的评测尤其贵。\par}
\end{enumerate}\endgroup
\numline{公式}{$P(\text{A 赢})=\frac{1}{1+10^{(R_{B}-R_{A})/400}}$；$R\leftarrow R+K\times(\text{实际}-\text{预期})$}
\numline{数字}{MMLU 57 个学科；ToyLM 4/5 = 80\%；Elo 1200 vs 1000 $\rightarrow$ 预期 0.760；赢只加 32$\times$0.24 $\approx$ 7.7 分；爆冷后 1175.7 / 1024.3；K = 32}
\fol{困惑度降了就说明变好？}{不一定。它只衡量预测下一个 token，和解题、写代码、回答是否有帮助不直接对应；量化后困惑度涨一点、基准掉几个百分点都很常见，要一起看。}
\fol{Elo 分数能跨榜比较吗？}{不能。它反映的是在同一个竞技场里、和同一批对手的相对强弱，换一批对手分数就变，不同排行榜之间没有可比性。}
\fol{优化后怎么确认没变笨？}{每次量化、换算子、改解码策略都跑回归评测：困惑度涨多少、基准掉几个百分点、输出和原来的数值误差是否可接受，三项一起看。}
\numline{别答错}{同一模型分数差几个百分点，可能只是提示写法、few-shot、答案提取规则不同。}
\chapter{从零实现一个GPT模型}
{\footnotesize\sffamily\color{grayt}第 52--59 讲}\par\vspace{-4pt}
\phantomsection\addcontentsline{toc}{section}{52\quad g0 $\cdot$ 语料与分词器}
\lec{52}{g0 $\cdot$ 语料与分词器}
\askline{为什么要自己训分词器？}
\askline{从字符出发再加字节回退？}
\askline{分词器最容易写错什么？}
\vspace{3pt}
\noindent{\qlabel{答}}\begingroup\sloppy\begin{enumerate}[leftmargin=1.6em,itemsep=2.5pt,topsep=0pt,parsep=0pt,label={\sffamily\footnotesize\color{accent}\arabic*.}]
\item {\fontsize{10}{14.8}\selectfont g0 只做两件事：prepare.py 按顺序删掉 MDX 里的元数据、代码块、公式、组件和 HTML 标签，留下 499,220 个字符的纯文本；tokenizer.py 用 133 行实现一个 BPE，把它变成编号。\par}
\item {\fontsize{10}{14.8}\selectfont 不用现成的 GPT-2 分词器有两个原因：它按英文网页训练，很多中文常用字要两三个字节 token 才拼出来，序列变长；词表 50,257 时嵌入表加输出层就是 2570 万参数，是 4 层 Transformer 的 8 倍。\par}
\item {\fontsize{10}{14.8}\selectfont 我们从字符出发：语料里出现 3 次以上的字符各占一个编号，一个汉字一开始就是一个 token，合并直接用在「模型」「请求」这种词上；不到 3 次的生僻字退回 UTF-8 字节，这就是 byte fallback。\par}
\item {\fontsize{10}{14.8}\selectfont 编号这样排：0--255 是 256 个字节，256 是 <|endoftext|> 分隔符，257 起是字符，再往后是按学到顺序排的合并结果；最终 4096 = 256 字节 + 1 分隔符 + 1356 字符 + 2483 次合并。\par}
\item {\fontsize{10}{14.8}\selectfont 编码前先切块，合并只在块内发生：一串英文字母、单个数字、一串汉字、单个空白、单个符号各成一块；所以数字一位一位切，2048 永远是 2、0、4、8 四个 token，不会因为常见就单独成词。\par}
\item {\fontsize{10}{14.8}\selectfont 编码时必须按学到的先后合并（编号最小的规则优先），不能按块内频次，训练时的频次是对整个语料统计的；同一个块的结果用 cache 记下来，「的」出现几千次也只算一次，41 个字符的例句切成 23 个 token。\par}
\end{enumerate}\endgroup
\numline{公式}{词表大小 $=256\,\text{字节}+1\,\text{分隔符}+\text{字符数}+\text{合并次数}$；切块正则按英文词、单数字、连续汉字、单空白、其他符号分块}
\numline{数字}{499,220 字符（约 1.1 MB）；4096 = 256 字节 + 1 分隔符 + 1356 字符 + 2483 合并；训分词器 22.3 s；41 字符 $\rightarrow$ 23 token；1.51 字符/token；上下文 256 token $\approx$ 390 字符}
\fol{1.51 字符/token 意味着什么？}{它决定模型一次能看多远：验证集 49,922 个字符编成 33,048 个 token，上下文 256 个 token 约等于 390 个字符，差不多是讲义里的两三段。}
\fol{空格为什么单独成 token？}{切块规则把空白单独成块，换来规则简单，代价是序列长一些；GPT-2 把空格粘在后一个单词开头（「 Cache」是一个 token），能省下这些空格 token。}
\fol{怎么验证分词器可逆？}{两个检查：验证集编码再解码和原文一字不差，没见过的字和表情符号退回字节后也能还原；解码要整体拼字节再解，一个个 token 单独解会把生僻字变成乱码。}
\numline{别答错}{训练集和验证集必须在分词之前切开，只用前 90\% 训分词器，否则验证 loss 偏乐观。}
\phantomsection\addcontentsline{toc}{section}{53\quad g1 $\cdot$ 数据加载器与嵌入}
\lec{53}{g1 $\cdot$ 数据加载器与嵌入}
\askline{训练样本是怎么切出来的？}
\askline{步长 stride 怎么选？}
\askline{编号是整数，怎么变成向量？}
\vspace{3pt}
\noindent{\qlabel{答}}\begingroup\sloppy\begin{enumerate}[leftmargin=1.6em,itemsep=2.5pt,topsep=0pt,parsep=0pt,label={\sffamily\footnotesize\color{accent}\arabic*.}]
\item {\fontsize{10}{14.8}\selectfont 语言模型的训练数据不需要人工标注，一段文字本身就是题目和答案：从长序列里取连续的 T+1 个 token，前 T 个是 x，后 T 个是 y，y 就是 x 整体往后错一位，y[t] = x[t+1]。\par}
\item {\fontsize{10}{14.8}\selectfont 一个长度 T 的样本里有 T 道「预测下一个」的题，一次前向就同时给出这 T 个答案，因果掩码保证位置 t 只看得到 x[0..t]；所以语言模型训练效率高：一次前向，T 个训练信号。\par}
\item {\fontsize{10}{14.8}\selectfont 窗口怎么挪由步长决定：stride = T 时首尾相接，T = 256 时训练集切出 1181 个样本；stride 小于 T 会重叠、相邻几乎一样容易被背下来；大于 T 则有 token 永远用不上。\par}
\item {\fontsize{10}{14.8}\selectfont 训练用 stride = T/2 = 128，样本数翻倍，同一段话出现在窗口的不同位置，一轮 147 个 batch；验证集用 stride = T 不重叠，每个 token 只算一次 loss，曲线才不会抖。\par}
\item {\fontsize{10}{14.8}\selectfont 实现上只存窗口起点 range(0, len - T, stride)，取样本时再切：stride = 1 时全切出来要 7700 万个整数，只存起点内存就是原序列那么大；起点上限是 len - T，因为窗口要取 T+1 个 token。\par}
\item {\fontsize{10}{14.8}\selectfont 编号是整数，第一层是两张表：token 嵌入 V$\times$d 取第 i 行，位置嵌入 T\_max$\times$d 取第 t 行，两者相加得 [B, T, d]；同一个「的」在不同位置查出的 token 向量相同，加上位置向量才不同。\par}
\end{enumerate}\endgroup
\numline{公式}{$h[b,t]=E_{\text{tok}}[x[b,t]]+E_{\text{pos}}[t]$；起点 $\in\operatorname{range}(0,\ \text{len}-T,\ \text{stride})$；$y[t]=x[t+1]$}
\numline{数字}{训练集 302,475 / 验证集 33,048 token；stride 8/4/1 $\rightarrow$ 37,809 / 75,617 / 302,467 个样本；T=256、stride=T 时 1181 样本 73 batch，stride=128 时 147 batch；嵌入 4096$\times$256 + 256$\times$256 = 1,114,112 参数}
\fol{位置嵌入为什么用可学的表？}{这是 GPT-2 的做法：位置向量也是参数，从随机数开始学，比正弦余弦公式简单；代价是只能处理不超过 T\_max 的序列，第 T\_max+1 个位置在表里没有对应的行。}
\fol{验证集为什么不能打乱、重叠？}{不影响对错，但验证 loss 会因为每次抽到的样本不同而抖动，也会把同一段话重复算好几遍；所以验证集用 stride = T 不重叠。}
\fol{嵌入层参数占比多少？}{4096$\times$256 的 token 嵌入加 256$\times$256 的位置嵌入共 1,114,112 个参数，token 嵌入占 94\%；这也是词表定 4096 的原因，词表一大就压过 Transformer。}
\numline{别答错}{y 是 x 往后错一位，写反了配上因果掩码模型什么都看不到，loss 降不下去。}
\phantomsection\addcontentsline{toc}{section}{54\quad g2 $\cdot$ 因果多头注意力}
\lec{54}{g2 $\cdot$ 因果多头注意力}
\askline{多头注意力怎么写成矩阵乘？}
\askline{为什么不用 for 循环逐头算？}
\askline{因果掩码写错会怎样？}
\vspace{3pt}
\noindent{\qlabel{答}}\begingroup\sloppy\begin{enumerate}[leftmargin=1.6em,itemsep=2.5pt,topsep=0pt,parsep=0pt,label={\sffamily\footnotesize\color{accent}\arabic*.}]
\item {\fontsize{10}{14.8}\selectfont 单个头是 Q = XW\_q、K = XW\_k、V = XW\_v，打分 QKᵀ/$\surd$h\_d，j > t 填 -$\infty$，每行 softmax 后乘 V；多头就是 H 组独立投影，每头 h\_d = d/H，这里 d=256、H=4，每头 64 维。\par}
\item {\fontsize{10}{14.8}\selectfont 直接照描述写就是 for 循环逐头算，能跑但 GPU 最怕一堆小矩阵乘法，算术强度太低；实际写法把循环压平进张量维度：所有头的 W\_q 并排正好是 d$\times$d，一次矩阵乘法算出所有头的 q。\par}
\item {\fontsize{10}{14.8}\selectfont 分三步：输出 [B,T,d] 先 view 成 [B,T,H,h\_d] 再转置成 [B,H,T,h\_d]，B$\cdot$H 就是一批独立小问题；一次批量矩阵乘算完 [B,H,T,T] 的打分；最后拼回 [B,T,d] 过 out\_proj。\par}
\item {\fontsize{10}{14.8}\selectfont 参数量是 4d$^2$ + d = 262,400，和头数无关，H=1 还是 H=8 都是四个 d$\times$d 矩阵，变的只是怎么切；掩码是 buffer 不算参数，但 B=16、T=256 时权重是 419 万个数、FP32 16 MiB。\par}
\item {\fontsize{10}{14.8}\selectfont 和 F.scaled\_dot\_product\_attention 比，最大差 3$\times$10$^-$⁸，只是浮点加法顺序的舍入误差；改第 3 个 token，位置 0、1、2 的输出差为 0，只有位置 3 变，掩码真的挡住了未来。\par}
\item {\fontsize{10}{14.8}\selectfont 代价在显存：打分、softmax、dropout 的结果都要留着反向传播，一层 16 MiB、4 层上百 MiB，而且和 T$^2$ 成正比，T 翻倍就是 4 倍，所以 FlashAttention 要想办法不把 T$\times$T 矩阵写进显存。\par}
\end{enumerate}\endgroup
\numline{公式}{$\text{Attn}=\operatorname{softmax}\left(\frac{QK^{\top}}{\sqrt{h_{d}}}+\text{mask}\right)V$，$j>t$ 处 $\text{mask}=-\infty$；$h_{d}=\frac{d}{H}$}
\numline{数字}{d=256、H=4、h\_d=64；参数 4d$^2$+d = 262,400；与 PyTorch 最大差 3$\times$10$^-$⁸；权重 [16,4,256,256] 419 万个数、FP32 16 MiB；改第 3 个 token 只影响位置 3}
\fol{拆头为什么不能直接 view？}{形状对得上不报错，但内存里是 [B,T,d] 的顺序，直接 view 会把不同位置的数混进同一个头，结果是错的；必须先 view 成 [B,T,H,h\_d] 再转置。}
\fol{掩码写成 tril 会怎样？}{模型只能看未来看不到过去，训练 loss 会降得出奇地快，因为它在偷看答案，生成时却一塌糊涂；改第 3 个 token 只影响位置 3 的检查能抓住这种错。}
\fol{缩放为什么除 $\surd$h\_d 不是 $\surd$d？}{多头时每个头的点积是 h\_d 维的，方差随 h\_d 增长，要除 $\surd$h\_d 才把打分拉回单位量级；除 $\surd$d 会让 softmax 过于平缓，注意力集中不起来。}
\numline{别答错}{缩放要除 $\surd$h\_d 不是 $\surd$d；掩码方向写反等于偷看答案，训练 loss 掉得飞快、生成全崩。}
\phantomsection\addcontentsline{toc}{section}{55\quad g3 $\cdot$ 组装GPT}
\lec{55}{g3 $\cdot$ 组装GPT}
\askline{Block 由哪些零件组成？}
\askline{530 万参数都花在哪了？}
\askline{Pre-Norm 和 Post-Norm 差在哪？}
\vspace{3pt}
\noindent{\qlabel{答}}\begingroup\sloppy\begin{enumerate}[leftmargin=1.6em,itemsep=2.5pt,topsep=0pt,parsep=0pt,label={\sffamily\footnotesize\color{accent}\arabic*.}]
\item {\fontsize{10}{14.8}\selectfont 还差四样：LayerNorm 把每个 token 的 d 个数减均值除标准差再乘 γ 加 $\beta$；GELU 负半轴不是切成 0 而是平滑趋近 0；FFN 两个线性层夹 GELU，d $\rightarrow$ 4d $\rightarrow$ d；残差把子层输出加回输入，给梯度一条直通路。\par}
\item {\fontsize{10}{14.8}\selectfont Block 用 Pre-Norm：x $\leftarrow$ x + Attn(LN(x))，x $\leftarrow$ x + FFN(LN(x))；原始 Transformer 是 LN(x + f(x))。两种都能跑，但装 GPT-2 权重时结构必须一致。\par}
\item {\fontsize{10}{14.8}\selectfont 参数逐项对：每层注意力 4d$^2$ + d = 262,400，FFN 8d$^2$ + 5d = 525,568，两个 LayerNorm 各 2d 共 1,024，一层 788,992，4 层 3,155,968。\par}
\item {\fontsize{10}{14.8}\selectfont 再加上 token 嵌入 1,048,576、位置嵌入 65,536、最后 LayerNorm 512、输出层 1,048,576，合计 5,319,168；4 层 Transformer 只占 59\%，两张词表表占 39\%。\par}
\item {\fontsize{10}{14.8}\selectfont 输出层是 d$\times$V 的线性层、不带偏置，GPT-2 让它和 token 嵌入共用一份权重：换成 GPT-2 small 的尺寸，不共用是 163,037,184 个参数，共用后是 124,439,808，这就是常说的 124M。\par}
\item {\fontsize{10}{14.8}\selectfont 生成是 generate\_greedy：每步把整段序列送进模型，只取最后一个位置的 logits 做 argmax 接在后面，超过 T\_max 就截断；还没训练的模型权重随机，产出的是词表里随机拼起来的乱码。\par}
\end{enumerate}\endgroup
\numline{公式}{$x\leftarrow x+\operatorname{Attn}(\operatorname{LN}(x))$；$x\leftarrow x+\operatorname{FFN}(\operatorname{LN}(x))$；每层 $\approx12d^{2}$；总参数 $\approx L\cdot12d^{2}+2Vd$}
\numline{数字}{5,319,168 参数（FP32 20.3 MiB）；一层 788,992 = 注意力 262,400 + FFN 525,568 + LN 1,024；词表相关的两张表占 39\%；GPT-2 small 163,037,184 $\rightarrow$ 共用输出层 124,439,808}
\fol{小模型为什么要注意词表？}{词表决定嵌入表和输出层：V=50,257、d=256 时两张表 2570 万参数，是 4 层 Transformer 的 8 倍，而其中绝大多数 token 在语料里一次都不出现，白占参数。}
\fol{因果性怎么验证到整个模型？}{LayerNorm、FFN、残差都是逐位置的，只要注意力因果整个模型就因果；但 LayerNorm 写成在 T 维求均值就会把位置混在一起，形状还对，只能靠改第 10 个 token 的检查发现。}
\fol{生成时最容易踩什么坑？}{忘记截断到 T\_max，序列一长位置嵌入查表就报 index out of range；位置编号还得和模型在同一 device 上，否则查表直接报错。}
\numline{别答错}{Pre-Norm 写成 Post-Norm 都能跑，但装 GPT-2 权重时结构必须一致。}
\phantomsection\addcontentsline{toc}{section}{56\quad g4 $\cdot$ 预训练}
\lec{56}{g4 $\cdot$ 预训练}
\askline{预训练到底在优化什么？}
\askline{验证 loss 为什么先降后升？}
\askline{这次训练算力利用率多少？}
\vspace{3pt}
\noindent{\qlabel{答}}\begingroup\sloppy\begin{enumerate}[leftmargin=1.6em,itemsep=2.5pt,topsep=0pt,parsep=0pt,label={\sffamily\footnotesize\color{accent}\arabic*.}]
\item {\fontsize{10}{14.8}\selectfont 预训练就是每个位置上做「下一个 token 是哪个」的分类题：模型输出 4096 个打分，softmax 后正确 token 的概率是 p，loss 就是 -ln p；一个 batch 是 4096 个位置，每个位置一个 loss 取平均。\par}
\item {\fontsize{10}{14.8}\selectfont 起点可以算出来：完全没把握时在 4096 个 token 里平均猜，p = 1/4096，loss = ln 4096 = 8.318；实测第 1 步是 8.45，比它略高是 PyTorch 默认初始化带来的随机起伏，正常。\par}
\item {\fontsize{10}{14.8}\selectfont loss 的指数就是困惑度：8.318 相当于在 4096 个选项里瞎猜，降到 3 就相当于在 20 个选项里挑；所以一开始 loss 远大于 ln V 说明初始化有问题，远小于则多半是数据泄漏、掩码写反了。\par}
\item {\fontsize{10}{14.8}\selectfont 3000 步、每步 4096 个 token 共 12.3 M 个，等于把 0.30 M 的训练集过了 41 遍；训练 loss 从 8.45 降到 0.98，验证 loss 第 900 步最低 4.843，之后回升到 5.857。\par}
\item {\fontsize{10}{14.8}\selectfont 这就是过拟合：Chinchilla 的经验比例是每参数约 20 个 token，530 万参数需要约 106 M token，我们只有 0.30 M、差 352 倍，同一段话看第十几遍时模型从学规律变成背原文，要翻转只能加数据或缩小模型。\par}
\item {\fontsize{10}{14.8}\selectfont 算力账：每 token 约 6N 次浮点运算（前向 2N、反向 4N），一共 3.92e14 FLOPs，跑了 569 s，平均 0.69 TFLOPS，只有 M2 十核 GPU 标称 FP32 3.6 TFLOPS 的 19\%。\par}
\end{enumerate}\endgroup
\numline{公式}{$\text{loss}=-\ln p$；困惑度 $=e^{\text{loss}}$；训练算力 $\approx6N\times\text{token 数}$}
\numline{数字}{ln 4096 = 8.318；训练 8.45 $\rightarrow$ 0.98，验证最低 4.843（第 900 步）$\rightarrow$ 5.857；过训练集 41 遍；差 352 倍；6N$\times$token = 3.92e14 FLOPs，569 s、0.69 TFLOPS、利用率 19\%}
\fol{验证 loss 回升该改什么？}{先排除 bug：验证集有没有混进训练、估 loss 时有没有 eval()。确认是过拟合，就存验证 loss 最低那版，或者加数据、缩小模型、加大 dropout。}
\fol{为什么 MFU 只有 19\%？}{6N 没算注意力打分（T=256 时约一成），每 100 步要估 40 个 batch，矩阵只有 256$\times$1024 喂不饱计算单元，还有 Python 调度开销；大模型靠大矩阵能到 40\%--60\%。}
\fol{怎么最快确认训练代码是对的？}{拿一个 batch 反复训 50 步，看 loss 能不能降到接近 0；正常模型一定能背下一个 batch，降不下去就说明梯度没清零、学习率没生效或目标错位。}
\numline{别答错}{验证 loss 回升不是 bug 而是过拟合；dropout 只能推迟，不能凭空变出数据。}
\phantomsection\addcontentsline{toc}{section}{57\quad g5 $\cdot$ 解码与存档}
\lec{57}{g5 $\cdot$ 解码与存档}
\askline{训练时怎么只留最好的那一份？}
\askline{温度、Top-k、Top-p 有什么区别？}
\askline{存档里为什么一定要存优化器状态？}
\vspace{3pt}
\noindent{\qlabel{答}}\begingroup\sloppy\begin{enumerate}[leftmargin=1.6em,itemsep=2.5pt,topsep=0pt,parsep=0pt,label={\sffamily\footnotesize\color{accent}\arabic*.}]
\item {\fontsize{10}{14.8}\selectfont 存档里放三样东西：模型配置、state\_dict、优化器状态，再加上步数、已看 token 数和日志。带配置的存档才是自描述的，读档时不用先知道 d、H、L 就能建出形状对得上的模型。\par}
\item {\fontsize{10}{14.8}\selectfont AdamW 给每个参数记两个数：梯度的滑动平均和梯度平方的滑动平均，530 万参数就是 1060 万个额外的数。所以优化器状态是参数的两倍大，存档从 20.6 MiB 涨到 61.2 MiB，约 3 倍。\par}
\item {\fontsize{10}{14.8}\selectfont 训练每 100 步估一次验证 loss，创新低就存一份 gpt\_best.pt。结束时手里有两份：最后一步的 gpt.pt 和验证最好的 gpt\_best.pt，这就是 early stopping 的思路，不必真停训练，只留最好那份。\par}
\item {\fontsize{10}{14.8}\selectfont dropout 从 0.1 加到 0.2，末步验证 loss 从 5.86 降到 5.25，过拟合确实慢了，但最低点还是第 900 步的 4.876。数据只有计算最优的 1/350，正则化只决定背得多快，不决定能学到多少。\par}
\item {\fontsize{10}{14.8}\selectfont 解码侧有三个旋钮：打分除以温度 T 放大或抹平差距，T$\rightarrow$0 退化成贪心；Top-k 留分数最高的 k 个；Top-p 从高到低累加、凑够 p 就截断。去掉的置 -$\infty$，剩下的重新归一化再抽样。\par}
\item {\fontsize{10}{14.8}\selectfont 在 4.876 那份存档上，下一个 token 的分布很平：最高概率只有 0.032，前 12 个加起来才 0.201，T=1 要 573 个 token 才凑够 90\%。所以贪心扎进分隔符循环，加了随机性句子才像讲义。\par}
\end{enumerate}\endgroup
\numline{公式}{$l\leftarrow l/T$；Top-$k$ 取前 $k$ 大；Top-$p$ 按“前面的累计概率”截断；其余置 $-\infty$ 后 softmax 采样}
\numline{数字}{最低验证 loss 4.876（第 900 步）、末步 5.25；存档 20.6 $\rightarrow$ 61.2 MiB；最高概率 0.032，T=1 凑 90\% 要 573 个 token}
\fol{续训时丢了优化器状态会怎样？}{差别很小。从第 900 步接着训 140 步，新建优化器前 20 步平均只高 0.03，第 61 步后已在噪声内。Adam 的偏差修正前几步就把滑动平均放大回来，丢的只是梯度方向记忆。}
\fol{加大 dropout 为什么没用？}{它只拖慢背的速度。第 900 步时模型已把 30 万 token 的训练集看了 12 遍，能学的差不多学完了。正则化不会带来新信息，想让验证 loss 更低只能加数据。}
\fol{Top-p 最容易写错在哪？}{边界要按“前面的累计概率”判断。用“包含自身”会删掉刚好凑够 p 的那个 token；分布极尖、第一个就超过 p 时全删光，softmax 得到 NaN。温度 0 还要单独走贪心。}
\numline{别答错}{Top-p 要按“前面的累计概率”截断；用“包含自身”会删掉凑够 p 的 token，极尖时全删。}
\phantomsection\addcontentsline{toc}{section}{58\quad g6 $\cdot$ 装进 GPT-2 的权重}
\lec{58}{g6 $\cdot$ 装进 GPT-2 的权重}
\askline{别人训练好的权重，怎么装进自己的模型？}
\askline{装完怎么确认真的装对了？}
\askline{哪些权重需要转置、哪些要切开？}
\vspace{3pt}
\noindent{\qlabel{答}}\begingroup\sloppy\begin{enumerate}[leftmargin=1.6em,itemsep=2.5pt,topsep=0pt,parsep=0pt,label={\sffamily\footnotesize\color{accent}\arabic*.}]
\item {\fontsize{10}{14.8}\selectfont 装别人权重的难点不在复制，在对应。GPT-2 small 一共 149 个张量、1.24 亿参数，逐个按名字映射到我们写的 GPT 类里，模型一行不改；vLLM 的 load\_weights 做的也是这件事。\par}
\item {\fontsize{10}{14.8}\selectfont 第一处差在存法：GPT-2 的线性层叫 Conv1D，权重存 [输入, 输出]，nn.Linear 存 [输出, 输入]，互为转置。所以 attn.c\_proj 和 mlp 的两个矩阵都要先转置再复制。\par}
\item {\fontsize{10}{14.8}\selectfont 第二处差在合并：q、k、v 拼在一个 [768, 2304] 的 c\_attn 里，前 768 列是 q，接着 k，最后 v。要按列也就是输出维切成三份再转置，按行切 split(dim=0) 也能出三个方阵，但值是错的。\par}
\item {\fontsize{10}{14.8}\selectfont 还有共用：输出层和 token 嵌入共用一份权重，文件里只有 wte，要复制一份进 out\_head。所以装完是 1.63 亿参数，比 1.24 亿多出的 3860 万就是这份复制出来的输出层。\par}
\item {\fontsize{10}{14.8}\selectfont 验证不用下载 550 MB 真权重。用 GPT-2 small 的配置随机初始化一个模型，名字、形状、存法都一样，装进来和 transformers 比 logits，最大差 5$\times$10$^-$⁶，每个位置的 argmax 全一致。\par}
\item {\fontsize{10}{14.8}\selectfont 误差不是 0，因为两边数学一样但运算顺序不同：Conv1D 用 addmm 一次算 x$\cdot$W + b，我们走 nn.Linear，softmax 和 LayerNorm 也不同。浮点加法不满足结合律，12 层累积到 10$^-$⁶。\par}
\end{enumerate}\endgroup
\numline{公式}{$c_{\text{attn}}[768,2304]\to$ 按列切 $q/k/v$ 各 $[768,768]$；$W_{\text{ours}}=W_{\text{conv1d}}^{\top}$；$\text{out\_head}\leftarrow w_{\text{te}}$}
\numline{数字}{149 个张量 / 1.24 亿参数；装后 1.63 亿（+3860 万）；最大误差 5$\times$10$^-$⁶；忘转置方阵误差 3.29、argmax 一致率 0\%}
\fol{方阵忘了转置会怎样？}{不报错。768$\times$768 的 out\_proj 转置前后形状一样，形状检查拦不住，程序正常跑完，但 logits 最大误差变成 3.29，128 个位置一个都没选对。所以必须和参考实现逐位比对。}
\fol{除了权重还有什么容易漏？}{分词器。嵌入表第 i 行只对它自己分词器的第 i 个 token 有意义，用我们的词表喂 GPT-2 权重，模型看到的是无关编号。另外 ε 要用 10$^-$⁵、GELU 用 tanh 近似。}
\fol{怎么区分浮点噪声和真错了？}{两个指标一起看：最大误差和 argmax 一致率。浮点顺序差异在 10$^-$⁶ 量级、argmax 全一致；误差到 10$^-$$^3$ 以上、argmax 开始不一致，就是映射错了。}
\numline{别答错}{方阵转置错了形状照样对、不报错、结果全错；c\_attn 必须按列切，按行切也是三个方阵但值全错。}
\phantomsection\addcontentsline{toc}{section}{59\quad g7 $\cdot$ 指令微调}
\lec{59}{g7 $\cdot$ 指令微调}
\askline{SFT 和预训练的训练循环差在哪？}
\askline{为什么只在回答上算 loss？}
\askline{小模型做 SFT 能学会什么？}
\vspace{3pt}
\noindent{\qlabel{答}}\begingroup\sloppy\begin{enumerate}[leftmargin=1.6em,itemsep=2.5pt,topsep=0pt,parsep=0pt,label={\sffamily\footnotesize\color{accent}\arabic*.}]
\item {\fontsize{10}{14.8}\selectfont SFT 用的还是预测下一个 token 的循环，只改两件事：数据和哪些位置算 loss。样本套成“\#\#\# 指令：”加“\#\#\# 回答：”的模板，y 里指令和补齐位置设成 -100，cross\_entropy 会跳过。\par}
\item {\fontsize{10}{14.8}\selectfont 指令是用户给的，模型不用学会写指令，只关心看到指令后能不能写出回答。但回答后面的结束符要算 loss，模型得学会“说完了就停”，不然会一路编到长度上限。\par}
\item {\fontsize{10}{14.8}\selectfont 数据从课程目录两栏出：技术点问“解决什么问题”、问题问“哪个技术点回答了它”。112 讲出 224 道，每 10 讲留 1 讲做测试，200 条训练、24 条测试，同讲两题不跨集合。\par}
\item {\fontsize{10}{14.8}\selectfont 从第 900 步的 gpt\_best.pt 起微调 20 轮，学习率 5$\times$10$^-$$^4$、批大小 8，loss 从 5.227 降到 0.112。这个 loss 比预训练低得多，因为只算回答那几个 token，200 条又被看了 20 遍。\par}
\item {\fontsize{10}{14.8}\selectfont 测试集 0/24 一道没对，但格式全对：问“解决什么问题”就答“怎么……”式的答案，问“哪个技术点”就答一个技术名词，说完都停下。SFT 教的是答什么类型、答多长、何时停。\par}
\item {\fontsize{10}{14.8}\selectfont 它教不会模型不知道的东西。测试那 12 讲的对应关系没见过，530 万参数在 30 万 token 上训练，也没能力从正文归纳出“这讲在解决什么”，只能从背下的 200 个答案里挑。\par}
\end{enumerate}\endgroup
\numline{公式}{$y[i]=-100$（指令与补齐位置）；$L=\operatorname{CE}(\text{logits}[\text{回答位置}],\ y[\text{回答位置}])$}
\numline{数字}{112 讲 $\rightarrow$ 224 题（200 训练 / 24 测试）；一条样本 34 token（28 指令 + 6 回答）；loss 5.227 $\rightarrow$ 0.112；训练集 198/200 vs 测试集 0/24}
\fol{屏蔽位置差一位会怎样？}{不报错，只是 loss 算在错的位置上。一条 5 个 token、指令占 3 个的样本，y 应该是 [-100, -100, 4, 5]；补齐那条是 [-100, 3, -100, -100]。}
\fol{想让它测试集答对能怎么做？}{两个方向：让测试题在训练集里有近亲，同一讲换个问法；或换更强的基座，比如 GPT-2，题目换成英文。更省的是 LoRA，冻结原参数只训旁边的小矩阵，还能少背题。}
\fol{大模型的 SFT 为什么能举一反三？}{预训练已经给了它事实和推理方式，SFT 只是教它用问答形式说出来。LIMA 里约一千条精选指令就能让强基座学会对话，要学的主要是格式；基座不知道的，数据再多也只能背下来。}
\numline{别答错}{学习率别照搬预训练：小数据上用 10$^-$$^3$ 能训到训练集全对，但参数离预训练更远，续写能力退化更多。}
\chapter{硬件与性能基础}
{\footnotesize\sffamily\color{grayt}第 60--67 讲}\par\vspace{-4pt}
\phantomsection\addcontentsline{toc}{section}{60\quad GPU 架构}
\lec{60}{GPU 架构}
\askline{GPU 为什么比 CPU 适合矩阵乘法？}
\askline{989 TFLOPS 是怎么算出来的？}
\askline{Tensor Core 比 CUDA 核快多少？}
\vspace{3pt}
\noindent{\qlabel{答}}\begingroup\sloppy\begin{enumerate}[leftmargin=1.6em,itemsep=2.5pt,topsep=0pt,parsep=0pt,label={\sffamily\footnotesize\color{accent}\arabic*.}]
\item {\fontsize{10}{14.8}\selectfont CPU 是几个到几十个强核，靠分支预测、乱序执行和大缓存把一个逻辑复杂的任务做快；GPU 是成千上万个简单核，任务能拆成大量互不依赖、做同样运算的小份就能一起算。矩阵乘法正是这种。\par}
\item {\fontsize{10}{14.8}\selectfont H100 由 132 个 SM 组成，每个 SM 有 4 个 Tensor Core、128 个 FP32 CUDA 核、256 KB 寄存器、256 KB 的 L1 和共享内存，还分 4 个分区各一个调度器，每周期发射一个 warp。\par}
\item {\fontsize{10}{14.8}\selectfont 峰值算力一层层乘出来：132 个 SM $\times$ 4 个 Tensor Core $\times$ 512 次 BF16 乘加 $\times$ 2（一次乘加算两次运算）$\times$ 1.83 GHz $\approx$ 989 TFLOPS。同样套 CUDA 核是 62 TFLOPS，只有 6\%。\par}
\item {\fontsize{10}{14.8}\selectfont 和 CPU 比：64 核、每核 2 个 AVX-512 单元的 CPU 按 2.5 GHz 约 10 TFLOPS，H100 是它的近 100 倍；4096$^3$ 的矩阵乘法按峰值，Tensor Core 0.14 ms、CPU 13 ms。\par}
\item {\fontsize{10}{14.8}\selectfont GPU 的另一个特点是在飞的线程数：每个 SM 最多驻留 2048 个线程，全卡约 27 万个，64 核 CPU 开超线程只有 128 个。单线程访存要等几百周期，靠换 warp 掩盖延迟。\par}
\item {\fontsize{10}{14.8}\selectfont 峰值只是上限，要真达到还要两个条件：数据供得上，Tensor Core 一个周期吃的数得及时送到；任务足够多，132 个 SM 都得有活干。实测到峰值的 60\% 到 80\% 已经很好。\par}
\end{enumerate}\endgroup
\numline{公式}{$P=\text{SM}\times\frac{\text{TC}}{\text{SM}}\times\frac{\text{FMA}}{\text{TC}}\times2\times f=132\times4\times512\times2\times1.83\ \text{GHz}\approx989\ \text{TFLOPS}$}
\numline{数字}{H100：132 SM、每 SM 4 个 Tensor Core、每周期 512 次 BF16 乘加、1.83 GHz、16,896 个 CUDA 核；64 核 CPU 约 10 TFLOPS}
\fol{CUDA 核数能代表算力吗？}{不能。16,896 是 FP32 标量单元的数量，大模型的矩阵乘法主要走 Tensor Core，要看 Tensor Core 的规格。比峰值还必须对齐精度和稠密/稀疏，稀疏标称常是稠密两倍。}
\fol{实测为什么达不到峰值？}{两个条件缺一不可：数据供不上，计算单元就在干等；SM 没活干，算力就闲置。小 batch 的 decode 注意力只有几个工作单元，大部分 SM 闲着；一般能到 60\%\textasciitilde{}80\% 就不错。}
\fol{别家加速器结构也一样吗？}{思路一样：AMD 的 MI300、华为昇腾、Google 的 TPU 都是大量并行单元加专门的矩阵乘法硬件、片上存储和高带宽显存，只是名字不同，叫计算单元、矩阵单元或脉动阵列。}
\numline{别答错}{“CUDA 核数”不代表算力，大模型走 Tensor Core；比峰值要对齐精度和稠密/稀疏。}
\phantomsection\addcontentsline{toc}{section}{61\quad 存储层级}
\lec{61}{存储层级}
\askline{H100 的存储分几层，各多快？}
\askline{从显存读一字节能做多少次运算？}
\askline{SRAM 那么快，为什么不多做点？}
\vspace{3pt}
\noindent{\qlabel{答}}\begingroup\sloppy\begin{enumerate}[leftmargin=1.6em,itemsep=2.5pt,topsep=0pt,parsep=0pt,label={\sffamily\footnotesize\color{accent}\arabic*.}]
\item {\fontsize{10}{14.8}\selectfont 从近到远：寄存器每 SM 256 KB，延迟约 1 周期；L1 和共享内存合计约 31 TB/s；L2 约 50 MB、约 5.5 TB/s；HBM 80 GB、3.35 TB/s；经 PCIe 到 CPU 内存只有 64 GB/s。\par}
\item {\fontsize{10}{14.8}\selectfont 把算力和带宽相除：989 TFLOPS ÷ 3.35 TB/s $\approx$ 295 次运算每字节。显存每读 1 字节的时间，Tensor Core 能做约 295 次运算；同样的比例从 L2 是 180 次，从共享内存是 32 次。\par}
\item {\fontsize{10}{14.8}\selectfont 矩阵乘法：n $\times$ n 的 BF16 乘法做 2n$^3$ 次运算，至少要读写 A、B、C 共 6n$^2$ 字节，每字节约 n/3 次。n = 256 只有 85 次，喂不饱；n = 4096 有 1365 次，喂得饱。\par}
\item {\fontsize{10}{14.8}\selectfont 逐元素运算最吃亏：每个 BF16 数读一次写一次共 4 字节、只做 1 次运算，每字节 0.25 次，比 295 小一千多倍。读写 1 GB 要 0.60 ms，运算本身只要 0.5 $\mu$s，时间全花在搬数据上。\par}
\item {\fontsize{10}{14.8}\selectfont 片上存储不做大是因为贵：一个 SRAM 单元要 6 个晶体管，面积是 DRAM 的好几倍。H100 芯片才 800 多平方毫米，几十 MB 已占相当部分，80 GB 都做 SRAM 得几千颗芯片。\par}
\end{enumerate}\endgroup
\numline{公式}{$n\times n$ BF16 矩阵乘法：$2n^{3}$ 次运算 $\div6n^{2}$ 字节 $=\frac{n}{3}$ 次运算每字节；喂饱比例 $=\frac{\text{峰值算力}}{\text{带宽}}$}
\numline{数字}{寄存器与共享内存全卡各约 35 MB；共享内存约 31 TB/s、L2 5.5 TB/s、HBM 3.35 TB/s、PCIe 64 GB/s；每读 1 字节能做 295 / 180 / 32 次}
\fol{decode 为什么卡在带宽上？}{decode 每一步要把约 16 GB 的权重从 HBM 完整读一遍，却对每个字节只做约 1 次运算，远低于 295。所以速度基本等于权重字节数除以带宽，提高算力没用，只能减字节或加 batch。}
\fol{带宽和延迟哪个更重要？}{大模型运算基本看带宽。GPU 靠大量线程同时访问把延迟掩盖掉了，一个 warp 等数据时调度器马上换另一个。延迟只在小消息通信和依赖链长时才关键，所以第 64 讲要单独讲 $\alpha$。}
\fol{怎么把片上存储用起来？}{算子融合、FlashAttention、分块矩阵乘法做的是同一件事：把数据搬进共享内存或寄存器后反复用，少碰 HBM。A、B 的每个数本来要用 n 次，每次都重读，读写量大得多。}
\numline{别答错}{显存指 GPU 上的 HBM，内存指 CPU 那侧的 DRAM，中间要过 PCIe、慢 50 倍。}
\phantomsection\addcontentsline{toc}{section}{62\quad Roofline 模型}
\lec{62}{Roofline 模型}
\askline{Roofline 怎么判断卡在哪？}
\askline{算子的算术强度怎么算？}
\askline{decode 和 prefill 落在哪一块？}
\vspace{3pt}
\noindent{\qlabel{答}}\begingroup\sloppy\begin{enumerate}[leftmargin=1.6em,itemsep=2.5pt,topsep=0pt,parsep=0pt,label={\sffamily\footnotesize\color{accent}\arabic*.}]
\item {\fontsize{10}{14.8}\selectfont 算术强度 I 等于运算次数除以显存读写的字节数，单位 FLOPs/字节。能达到的速度是 min(P, I $\times$ B)，两条线的交点叫屋脊点 I* = P ÷ B，H100 上是 989 TFLOPS ÷ 3.35 TB/s $\approx$ 295。\par}
\item {\fontsize{10}{14.8}\selectfont 强度低于屋脊点落在斜线上，卡带宽、算力用不满；高于屋脊点落在水平线上，卡算力、带宽用不满。估时间用 max(W ÷ P, Q ÷ B)，第 68 讲以后所有性能估算都用这个式子。\par}
\item {\fontsize{10}{14.8}\selectfont 几个典型的点：逐元素运算强度 0.25，最多 0.8 TFLOPS，只有峰值的 0.1\%；1 个 token 过 4096 $\times$ 14336 的线性层，整个权重读一遍只做一次乘加，强度约 1，最多 3.3 TFLOPS。\par}
\item {\fontsize{10}{14.8}\selectfont 线性层的算术强度约等于一起处理的 token 数：32 个 token 强度 32、106 TFLOPS；256 个强度 237、峰值 80\%；4096 个强度 1792，卡在算力上，这是 prefill。\par}
\item {\fontsize{10}{14.8}\selectfont decode 注意力每读一个历史 token 的 k、v 只和 1 个查询做运算，强度约 1；GQA 让 4 个查询共用一份 KV 也才 4，仍深在带宽区。批处理把点往右移，这就是它提高吞吐的原因。\par}
\item {\fontsize{10}{14.8}\selectfont 用法分两种：在带宽区优化计算没用，要减少读写（融合、量化、复用）或提高强度（批处理、GQA）；在算力区才该用更低精度和更好的 Tensor Core 利用率。\par}
\end{enumerate}\endgroup
\numline{公式}{速度 $=\min(P,\ I\cdot B)$；$I=W/Q$；屋脊点 $I^{*}=P/B$；时间 $=\max(W/P,\ Q/B)$}
\numline{数字}{屋脊点 H100 295、A100 153、B200 281（FP8 翻倍到 590）；256 token 强度 237 $\rightarrow$ 794 TFLOPS（80\%）；decode 注意力强度 1\textasciitilde{}4}
\fol{利用率不到 1\% 是代码差吗？}{不是，是运算性质决定的。判断实现好坏要看它离自己那条屋顶有多远：带宽区的算子如果只到带宽屋顶的一半，才说明数据搬运的方式还有问题。}
\fol{屋脊点和精度有关吗？}{有。FP8 算力翻倍，屋脊点也翻倍，H100 上约 590，同样的运算更容易落进带宽区。新一代卡算力涨得比带宽快，所以访存优化越来越重要。}
\fol{Roofline 为什么总偏乐观？}{它是上限模型：只算 HBM 一层存储，还假设读写和计算完美重叠。实际 kernel 有启动开销、同步和负载不均，达不到屋顶；逐元素运算走 CUDA 核，屋顶本身也更低。}
\numline{别答错}{算术强度只算显存读写，同一数在片上用多次只算一次；分块和融合会提高实际强度，别用理论值否定实现。}
\phantomsection\addcontentsline{toc}{section}{63\quad 低精度计算}
\lec{63}{低精度计算}
\askline{为什么精度减半算力就翻倍？}
\askline{FP8 省的是带宽还是算力？}
\askline{低精度的代价怎么兜住？}
\vspace{3pt}
\noindent{\qlabel{答}}\begingroup\sloppy\begin{enumerate}[leftmargin=1.6em,itemsep=2.5pt,topsep=0pt,parsep=0pt,label={\sffamily\footnotesize\color{accent}\arabic*.}]
\item {\fontsize{10}{14.8}\selectfont 规律很整齐：H100 的 TF32 是 495 TFLOPS、BF16 是 989、FP8 是 1979，B200 再加 FP4 到 9000。精度每减半峰值算力翻一倍，所以看规格表要先问是哪种精度、稠密还是稀疏。\par}
\item {\fontsize{10}{14.8}\selectfont 第一个原因是乘法器小。乘法器核心是把两个尾数相乘，m 位乘 m 位面积大致和 m$^2$ 成正比：BF16 尾数 8 位、FP8 只有 4 位，面积约四分之一，同样芯片能放更多乘法器。\par}
\item {\fontsize{10}{14.8}\selectfont 实际只翻倍而不是四倍，因为加法器、累加器和数据通路不按平方缩小。第二个原因是数据少，每个数从 2 字节变 1 字节，同样带宽送进两倍的数；第三个是能耗更低。\par}
\item {\fontsize{10}{14.8}\selectfont 对 decode 省的是带宽。70B 的 BF16 权重 141 GB、读一遍 5.3 ms，FP8 只有 71 GB、2.6 ms，几乎翻倍。这个收益不需要硬件支持 FP8 计算，读进来转回 BF16 算就行。\par}
\item {\fontsize{10}{14.8}\selectfont 对 prefill 省的是算力，必须让 Tensor Core 直接做 FP8 乘法，H100 起才有这个支持。同样 70B、8192 个 token，8 张 H100 上 BF16 约 146 ms、FP8 约 73 ms。\par}
\item {\fontsize{10}{14.8}\selectfont 代价是精度：E4M3 每个 2 的幂之间只有 8 个格点，FP4 只有 8 个非负值，直接转会明显掉效果，要靠 GPTQ、AWQ、SmoothQuant 和分组缩放来保；累加器仍要用 FP16 或 FP32。\par}
\end{enumerate}\endgroup
\numline{公式}{乘法器面积 $\propto(\text{尾数位数})^{2}$（相对 BF16：FP8 $0.25$、FP4 $0.06$）；decode 时间 $\approx\frac{\text{权重字节数}}{\text{带宽}}$}
\numline{数字}{H100：TF32 495 / BF16 989 / FP8 1979 TFLOPS；B200：BF16 2250 / FP8 4500 / FP4 9000；70B 权重 BF16 141 GB $\rightarrow$ FP8 71 GB}
\fol{存成 FP8 就等于用 FP8 算吗？}{不等于。权重存 FP8、读进来转 BF16 再算，只省带宽不省算力，decode 正好吃这个收益；要省算力必须让 Tensor Core 直接做 FP8 乘法，需要 H100 起的硬件。}
\fol{FP16 和 BF16 该选哪个？}{算力一样，但 FP16 尾数 11 位、范围小，BF16 尾数 8 位、范围和 FP32 一样大。训练大模型时 BF16 不容易溢出，所以成了主力，推理也大多用它或更低精度。}
\fol{规格表的算力怎么比才公平？}{对齐同一种精度、同样稠密。“B200 的 FP4 9000 TFLOPS”和“H100 的 BF16 989”不能直接比；稀疏算力常是稠密两倍，本站估算一律用稠密。}
\numline{别答错}{别把“存成低精度”和“用低精度算”混为一谈：前者只省带宽，后者才省算力，收益来源完全不同。}
\phantomsection\addcontentsline{toc}{section}{64\quad 互联}
\lec{64}{互联}
\askline{机内和机间带宽差多少？}
\askline{一条路径的传输时间怎么估？}
\askline{PCIe、NVLink、IB 各用在哪？}
\vspace{3pt}
\noindent{\qlabel{答}}\begingroup\sloppy\begin{enumerate}[leftmargin=1.6em,itemsep=2.5pt,topsep=0pt,parsep=0pt,label={\sffamily\footnotesize\color{accent}\arabic*.}]
\item {\fontsize{10}{14.8}\selectfont 分四层：PCIe 5.0 x16 每方向 64 GB/s，连 CPU 和网卡，比 HBM 慢 50 倍；NVLink 卡间专用，H100 每卡单向 450 GB/s；机间 IB 或 RoCE 每卡 400 Gb/s，约 50 GB/s。\par}
\item {\fontsize{10}{14.8}\selectfont 单位要看清。H100 标称的“NVLink 900 GB/s”是双向合计，估一次传输要用单向的 450；网卡的 400 Gb/s 是比特，除以 8 才是 50 GB/s；PCIe 和 NVLink 的规格一般直接给字节。\par}
\item {\fontsize{10}{14.8}\selectfont 路径带宽取决于最慢的一段：同机两卡走 NVLink，瓶颈 450 GB/s，传 1 GB 约 2 ms；只有 PCIe 是 64 GB/s、约 16 ms；跨机经 PCIe、网卡、交换机，瓶颈 50 GB/s、约 20 ms。\par}
\item {\fontsize{10}{14.8}\selectfont 整机看：NVLink 每卡 450 $\times$ 8 = 3600 GB/s，网卡 8 $\times$ 50 = 400 GB/s，机内比机间快 9 倍。所以通信频繁的并行要放机内：张量并行几乎总在一台机器里，跨机主要是流水并行、PD 的 KV 传输和专家并行。\par}
\item {\fontsize{10}{14.8}\selectfont 时间用 $\alpha$ + m/$\beta$ 估：小消息几乎全是固定开销 $\alpha$，NVLink 约 3 $\mu$s、PCIe 约 5、IB 约 6、RoCE 约 8。传 64 KB 时各链路都在 3 到 9 $\mu$s，看不出带宽差别；传 1 GB 才由带宽决定。\par}
\item {\fontsize{10}{14.8}\selectfont 集群用胖树拓扑，每层总带宽保持不变；再按轨道优化，每台机器第 i 张网卡接同一组交换机，正好匹配“不同机器同号 GPU 互相通信”的分层 All-Reduce。\par}
\end{enumerate}\endgroup
\numline{公式}{$t=\alpha+\frac{m}{\beta}$；$\alpha$ 取 NVLink $3\ \mu\text{s}$、PCIe $5\ \mu\text{s}$、IB $6\ \mu\text{s}$、RoCE $8\ \mu\text{s}$；瓶颈带宽 $=\min(\text{路径各段})$}
\numline{数字}{PCIe 5.0 x16 64 GB/s；NVLink 单向 A100 300 / H100 450 / B200 900 GB/s；IB NDR 400 Gb/s = 50 GB/s；机内 3600 vs 机间 400 GB/s，差 9 倍}
\fol{机内机间差 9 倍影响什么？}{直接影响并行能不能跑得动。张量并行每层都要 All-Reduce，只有 450 GB/s 的 NVLink 撑得住；跨机的流水并行、PD 分离传的是 KV，量小得多，50 GB/s 的网卡才够用。}
\fol{IB 和 RoCE 怎么选？}{都提供 RDMA，区别在底层网络和生态：IB 专用、延迟低丢包少但贵；RoCE 跑以太网、便宜、云厂商用得多，但要配无损网络，调好了接近 IB。对引擎通常只是换通信库配置。}
\fol{实际带宽怎么测出来？}{用 nccl-tests 的 all\_reduce 从小消息扫到大消息，拟合 $\alpha$ + m/$\beta$：斜率给带宽、截距给延迟，再和标称的 450 GB/s 比。NVLink 上大消息到 80\% 以上算正常。}
\numline{别答错}{NVLink 900 GB/s 是双向，估传输要用单向 450；400 Gb/s 是比特不是字节。}
\phantomsection\addcontentsline{toc}{section}{65\quad 超节点}
\lec{65}{超节点}
\askline{超节点和普通多机集群差在哪？}
\askline{同样 64 张卡，超节点能快多少？}
\askline{为什么 All-to-All 受益最大？}
\vspace{3pt}
\noindent{\qlabel{答}}\begingroup\sloppy\begin{enumerate}[leftmargin=1.6em,itemsep=2.5pt,topsep=0pt,parsep=0pt,label={\sffamily\footnotesize\color{accent}\arabic*.}]
\item {\fontsize{10}{14.8}\selectfont 集群有两层连接：scale-up 用 NVLink 直连，带宽高、延迟低、能直接读写对方显存；scale-out 用 InfiniBand 或 RoCE，规模不限但每卡带宽低一个数量级，还是消息语义。超节点就是把 scale-up 域从 8 卡扩到几十上百卡。\par}
\item {\fontsize{10}{14.8}\selectfont 差距有多大：H100 机内每卡单向约 450 GB/s、B200 约 900 GB/s，出了机器每卡只剩约 50 GB/s，差 9 倍，语义也从直接访存变成打包成消息，RDMA 只是让这个过程更高效。\par}
\item {\fontsize{10}{14.8}\selectfont 域做大受三个物理约束：铜缆只能传一两米、交换芯片的端口数与每端口带宽、几十上百千瓦的功耗与散热；所以 GB200 NVL72 把 72 张 Blackwell 塞进一个机柜全互连，每卡双向 1.8 TB/s，CloudMatrix 384 用 384 张昇腾 910C。\par}
\item {\fontsize{10}{14.8}\selectfont 同样 64 张卡、每卡 1 GB 的 All-Reduce：A 方案 8 台 8 卡机分层做，机内 Reduce-Scatter 缩成 1/8，机间只传 1/8，再机内 All-Gather，8.39 ms；B 方案一个 64 卡域直接做 2.26 ms，只快 3.7 倍。\par}
\item {\fontsize{10}{14.8}\selectfont All-to-All 占不到这个便宜：每张卡要给其余 63 张卡各发一份不同的数据，其中 56 份必须挤过 50 GB/s 的网卡，17.79 ms 对 1.12 ms，差 16 倍；对分带宽也从 1.6 拉到 28.8 TB/s。\par}
\item {\fontsize{10}{14.8}\selectfont 代价也要说清楚：成本、功耗、故障影响范围都变大，软件还要适配；而且 scale-out 不会消失，超节点之间、超节点和存储之间仍要靠以太网或 InfiniBand，两层要一起规划。\par}
\end{enumerate}\endgroup
\numline{公式}{分层 All-Reduce：机内 Reduce-Scatter 缩成 $1/8$ $\to$ 机间对 $1/8$ 做 All-Reduce $\to$ 机内 All-Gather}
\numline{数字}{机内 450 GB/s；机间 50 GB/s（差 9 倍）；NVL72 72 卡、每卡双向 1.8 TB/s；CloudMatrix 384；64 卡 All-Reduce 1 GB 8.39 $\rightarrow$ 2.26 ms（3.7 倍）；All-to-All 17.79 $\rightarrow$ 1.12 ms（16 倍）；对分带宽 1.6 对 28.8 TB/s（18 倍）}
\fol{超节点能让单卡变快吗？}{不能。它改变的是卡与卡之间的连接方式和带宽，单卡自己就能做完的工作一点也不会变快，收益全部来自跨卡通信，尤其是没法压缩的 All-to-All。}
\fol{怎么判断我的负载受不受益？}{看通信里走机间网络的比例。能分层压缩的 All-Reduce 只快 3.7 倍，每份数据都不同的 All-to-All 能快 16 倍，拿通信分析里的流量占比一算就有答案。}
\fol{为什么张量并行必须放在域内？}{张量并行每层都要 All-Reduce 同步，机间 50 GB/s 会把每一层都拖慢；域扩到 64 卡以后 TP 才能超过 8，更大的稠密模型也能只用 TP 而不用流水线并行。}
\numline{别答错}{超节点不是更快的单卡，也不等于一个机柜；判断标准是连接方式和带宽，不是物理边界。}
\phantomsection\addcontentsline{toc}{section}{66\quad 性能分析}
\lec{66}{性能分析}
\askline{线上推理慢，你怎么定位瓶颈？}
\askline{Profiler 时间线该怎么读？}
\askline{GPU 利用率 91\% 算好吗？}
\vspace{3pt}
\noindent{\qlabel{答}}\begingroup\sloppy\begin{enumerate}[leftmargin=1.6em,itemsep=2.5pt,topsep=0pt,parsep=0pt,label={\sffamily\footnotesize\color{accent}\arabic*.}]
\item {\fontsize{10}{14.8}\selectfont 性能分析就是别猜，把真实运行记录下来看时间花在哪。工具分三层：PyTorch Profiler 先粗看算子和对应的 kernel；Nsight Systems 看系统级时间线、谁在等谁；Nsight Compute 深挖单个 kernel 的带宽、算力和停顿原因。\par}
\item {\fontsize{10}{14.8}\selectfont 第一步看 GPU 有没有空等。8B、TP=4、decode 一步的模拟时间线里，一步 3.09 ms、GPU 忙 91\%、空等 0.27 ms，空白出现在一串两三微秒的逐元素 kernel 之间，是 CPU 发射跟不上，换 CUDA Graph 后降到 2.84 ms。\par}
\item {\fontsize{10}{14.8}\selectfont 第二步按类别分账：矩阵乘 1.69 ms 占 55\%，逐元素 0.49 ms、通信 0.35 ms、注意力 0.30 ms。逐元素计算量微不足道却占 16\%，一层 5 个、32 层 160 个小 kernel，开销主要在发射和读写显存，这就指向算子融合。\par}
\item {\fontsize{10}{14.8}\selectfont 第三步按名字排序找大头：gate\_up 投影 0.80 ms、down 投影 0.45 ms、All-Reduce 0.35 ms、注意力 0.30 ms。优化从这些开始，把一个只占 1\% 的 kernel 优化到零，整步也只快 1\%。\par}
\item {\fontsize{10}{14.8}\selectfont 第四步和屋顶比：最大的 gate\_up 一次 24.9 $\mu$s、读 58.7 MB 权重，实际带宽 2.36 TB/s，只有峰值 3.35 的 70\%。差的部分可能是 kernel 选错、批太小并行度不够，也可能是这台机器的实际带宽本来就不到标称。\par}
\item {\fontsize{10}{14.8}\selectfont 测量本身有坑：kernel 异步发射，计时前后要 torch.cuda.synchronize() 或用 CUDA 事件；第一次跑要编译、分配显存、autotune，得先预热；Profiler 有开销会放大空白；要在接近线上的条件下测，再用端到端计时复核。\par}
\end{enumerate}\endgroup
\numline{公式}{还能优化多少 $=\frac{\text{屋顶}}{\text{实测}}$（$2.36/3.35\approx70\%$）；只占 $1\%$ 的 kernel 优化到零也只快 $1\%$}
\numline{数字}{一步 3.09 ms、384 个 kernel、GPU 忙 91\%、空等 0.27 ms；CUDA Graph 2.84 ms；矩阵乘 1.69 / 逐元素 0.49 / 通信 0.35 / 注意力 0.30 ms；gate\_up 0.80 ms、24.9 $\mu$s、58.7 MB、2.36 TB/s；屋顶 3.35 TB/s}
\fol{GPU 利用率 99\% 就到位了吗？}{不够。nvidia-smi 的利用率只说明这段时间有 kernel 在跑，哪怕只用了一小部分带宽；要看实际带宽和屋顶的距离，gate\_up 只有 2.36/3.35 就是典型。}
\fol{为什么先修空等再优化 kernel？}{空等是 GPU 干等着，怎么调 kernel 都是白搭。0.27 ms 大约占整步的 9\%，换 CUDA Graph 一步就省 0.25 ms，比磨一个 kernel 便宜得多。}
\fol{多机的时候时间线怎么看？}{把每张卡的时间线放在一起看，找出最慢的那张。一张卡在等 All-Reduce，往往不是通信本身慢，而是别的卡还没算完、负载不均，或者网络拥塞。}
\numline{别答错}{nvidia-smi 的高利用率只说明有 kernel 在跑，不代表用上了带宽和算力。}
\phantomsection\addcontentsline{toc}{section}{67\quad 理论 vs 实测}
\lec{67}{理论 vs 实测}
\askline{规格表的峰值实际能跑到几成？}
\askline{怎么标定一台机器的真实上限？}
\askline{手算和实测差多少算正常？}
\vspace{3pt}
\noindent{\qlabel{答}}\begingroup\sloppy\begin{enumerate}[leftmargin=1.6em,itemsep=2.5pt,topsep=0pt,parsep=0pt,label={\sffamily\footnotesize\color{accent}\arabic*.}]
\item {\fontsize{10}{14.8}\selectfont 规格表上的峰值算力和带宽是理想条件下的上限，实际能到多少必须测。办法是用两个标定实验量出这台机器真实的屋顶：读写一大块内存标带宽，大矩阵乘标算力，之后所有估算都乘上这两个系数。\par}
\item {\fontsize{10}{14.8}\selectfont 带宽标定就是读一个 4 GB 的张量，它几乎不做计算，时间完全由带宽决定：DGX Spark 理论 14.65 ms、H100 1.19 ms。实测带宽通常只有标称的 70\% 到 90\%，这个比例要一路带进 decode、逐元素这些带宽区的估算。\par}
\item {\fontsize{10}{14.8}\selectfont 算力标定用 4096$^3$ 的 BF16 矩阵乘，算术强度约 1365，远在屋脊点右边，时间由算力决定：DGX Spark 理论 1.10 ms、H100 0.139 ms。调 cuBLAS 一般能到峰值的 60\% 到 80\%，这就是算力区的修正系数。\par}
\item {\fontsize{10}{14.8}\selectfont 两台机器对比：DGX Spark 是 GB10、128 GB 统一内存、BF16 约 125 TFLOPS，屋脊点 458；H100 是 989 TFLOPS、3.35 TB/s，屋脊点 295。带宽差 12 倍而算力只差 8 倍，所以更多运算落在带宽区。\par}
\item {\fontsize{10}{14.8}\selectfont 再扫一遍算术强度：线性层 1、8、64 个 token 在 Spark 上都约 0.43 ms，都在带宽区、时间由读权重决定；512 个 token 约 0.50 ms 才开始接近屋脊点。离屋顶特别远的点，就该用 Profiler 深挖。\par}
\item {\fontsize{10}{14.8}\selectfont 用屋顶预测真实模型：decode 批 1 的上限 $\approx$ 带宽除以权重字节数，8B 的 BF16 在 Spark 上约 17 tok/s、H100 约 209 tok/s，4 bit 约 60 tok/s。实测到上限的 70\% 到 80\% 算不错，30\% 就该看时间线。\par}
\end{enumerate}\endgroup
\numline{公式}{$t=\max\left(\frac{\text{FLOPs}}{\text{峰值算力}},\ \frac{\text{字节}}{\text{带宽}}\right)$；decode 速度上限 $\approx\frac{\text{带宽}}{\text{权重字节数}}$}
\numline{数字}{DGX Spark 273 GB/s、BF16 约 125 TFLOPS、屋脊点 458；H100 989 TFLOPS / 3.35 TB/s、屋脊点 295；读 4 GB 14.65 对 1.19 ms；4096$^3$ 1.10 对 0.139 ms；8B BF16 17 对 209 tok/s；4 bit 60 tok/s；70B 4 bit 39.7 GB 约 7 tok/s}
\fol{实测超过理论值说明什么？}{几乎一定是测量有问题。最常见是没同步就把异步发射的时间当成了执行时间，或者数据还在缓存里没真正读显存，这两条先查，再谈优化。}
\fol{统一内存的机器适合推理吗？}{适合开发和大模型实验，不适合高吞吐服务。128 GB 能装下 4 bit 的 70B，但 273 GB/s 只有 H100 的十二分之一，decode 上限自然低一个量级。}
\fol{怎么校准自己集群的效率系数？}{在目标卡上跑同样两个标定：4 GB 求和量带宽、4096$^3$ 矩阵乘量算力，把实测值替换掉标称值填进模型，之后所有预估值都按这个系数打折。}
\numline{别答错}{理论值是上限不是应该值，实测 70\%\textasciitilde{}90\% 属正常；统一内存不是 HBM，容量大但带宽低。}
\chapter{推理基础}
{\footnotesize\sffamily\color{grayt}第 68--73 讲}\par\vspace{-4pt}
\phantomsection\addcontentsline{toc}{section}{68\quad Prefill 与 Decode}
\lec{68}{Prefill 与 Decode}
\askline{prefill 和 decode 瓶颈为什么不同？}
\askline{decode 为什么慢，怎么估算？}
\askline{为什么 decode 一定要做批处理？}
\vspace{3pt}
\noindent{\qlabel{答}}\begingroup\sloppy\begin{enumerate}[leftmargin=1.6em,itemsep=2.5pt,topsep=0pt,parsep=0pt,label={\sffamily\footnotesize\color{accent}\arabic*.}]
\item {\fontsize{10}{14.8}\selectfont 一次前向的耗时约等于 max(计算量/算力, 读写字节数/显存带宽)，谁更慢取决于算术强度，也就是每读一个字节做多少次运算。H100 是 989 TFLOPS 除以 3.35 TB/s 约 295 FLOPs/字节，强度高于它计算受限，低于它带宽受限。\par}
\item {\fontsize{10}{14.8}\selectfont prefill 把整段提示一次送进去，权重只读一遍却被 1000 个 token 共用。8B 的 BF16 处理 1000 个 token，计算约 15.44 ms、读写约 4.85 ms，算术强度约 927，卡在算力；粗略地说，prefill 的强度约等于提示长度。\par}
\item {\fontsize{10}{14.8}\selectfont decode 每步只处理 1 个新 token，计算量约 160 亿次、只要 0.016 ms，可 16 GB 权重一个字节都不能少读，读写 4.83 ms，算术强度只有约 0.93，计算单元 99\% 以上的时间在等数据，算力只用了 0.3\%。\par}
\item {\fontsize{10}{14.8}\selectfont 放到一个请求上看：提示 1000、生成 200，prefill 15.4 ms 加 decode 200 步约 967 ms，合计约 983 ms，98\% 的时间花在 decode 上。所以优化优先级由这个比例决定，不是按阶段个数平均分。\par}
\item {\fontsize{10}{14.8}\selectfont 两个阶段的优化方向完全不同：prefill 减计算，用 FlashAttention、更低精度的计算、跳过重复前缀；decode 减要读的字节或提高每次读取的利用率，用批处理、量化、GQA/MLA、投机解码。\par}
\item {\fontsize{10}{14.8}\selectfont 最核心的一条是 decode 的批处理：既然每步都要把权重读一遍，就让这一遍同时服务 b 个请求，64 个请求一起 decode 算术强度提到约 64，单位时间生成的 token 接近翻 64 倍；这也是 PD 分离的动机，两批卡各按自己的瓶颈优化。\par}
\end{enumerate}\endgroup
\numline{公式}{耗时 $\approx\max\left(\frac{\text{计算量}}{\text{算力}},\ \frac{\text{读写字节}}{\text{带宽}}\right)$；算术强度 $=\frac{\text{FLOPs}}{\text{字节}}$，H100 分界线 $295$}
\numline{数字}{989 TFLOPS；3.35 TB/s；295 FLOPs/字节；prefill 15.44 ms、强度 927；decode 0.016 ms、读写 4.83 ms、强度 0.93、算力利用率 0.3\%；整请求 983 ms、decode 占 98\%}
\fol{提示很短时 prefill 还卡算力吗？}{不卡了。提示短到 16 或 64 个 token，计算量掉下来但读权重的时间不变，prefill 也变成带宽受限，耗时停在约 4.8 ms，和一步 decode 差不多。}
\fol{decode 慢是因为算得多吗？}{不是。decode 每步的计算量只有 prefill 的千分之一，慢是因为每步都要把 16 GB 权重整个搬一遍，搬数据的时间是计算时间的约 300 倍。}
\fol{大 batch 会是计算受限吗？}{理论会，但长上下文下很难。8B、每请求 2048 token 时，并发 256 的算术强度也只有约 60，离 295 还很远，所以长上下文 decode 始终带宽受限。}
\numline{别答错}{decode 慢不是因为算得多，而是每步都要读完 16 GB 权重；受限与否看强度和 295。}
\phantomsection\addcontentsline{toc}{section}{69\quad KV Cache}
\lec{69}{KV Cache}
\askline{推理为什么要做 KV Cache？}
\askline{KV Cache 省多少，代价是什么？}
\askline{为什么只缓存 K、V，不缓存 Q？}
\vspace{3pt}
\noindent{\qlabel{答}}\begingroup\sloppy\begin{enumerate}[leftmargin=1.6em,itemsep=2.5pt,topsep=0pt,parsep=0pt,label={\sffamily\footnotesize\color{accent}\arabic*.}]
\item {\fontsize{10}{14.8}\selectfont KV Cache 缓存的是每一层、每个历史 token 的 k 和 v。因果掩码保证了一个 token 在每一层的输出只取决于它自己和它前面的 token，后面接什么都不影响它，所以这些结果上一步已经算过，这一步再算一遍结果完全一样，纯属浪费。\par}
\item {\fontsize{10}{14.8}\selectfont decode 每一步因此变成：只把新 token 送进模型，在每一层只算它自己的 q、k、v，把 k、v 追加进这层的缓存，再用它的 q 和缓存里全部 k 打分、对全部 v 加权求和。新 token 排在最后，能看的正好是全部缓存，所以这里不需要掩码。\par}
\item {\fontsize{10}{14.8}\selectfont prefill 时整段提示一起算，把每一层所有 token 的 k、v 一次性写进缓存，之后 decode 每步只写一组。位置编码要用这个 token 在整句里的位置，这就是推理引擎里 position ids 的作用：没有了整句，新 token 得自己知道排在第几。\par}
\item {\fontsize{10}{14.8}\selectfont 省下的计算量可以精确算：8B 提示 1000、生成 1000，无缓存每步重算整句，共约 150 万个位置、22.5 PFLOPs；有缓存是 prefill 1000 加 decode 999 共 1999 个位置、约 30 TFLOPs。\par}
\item {\fontsize{10}{14.8}\selectfont 代价是显存：每 token 每层 2$\times$8 个 KV 头$\times$128 维$\times$2 字节，32 层合计 128 KiB，请求结束时缓存 250 MiB，随序列长度和并发线性增长，16 个 8192 token 的请求就要 16 GiB；decode 每步还要把它整个读一遍。\par}
\item {\fontsize{10}{14.8}\selectfont 所以它本质是一笔用显存换计算的交易，工程上四个方向：把每个 token 的缓存做小（GQA、MLA、量化）、让缓存不浪费（PagedAttention）、让相同前缀只存一份（前缀缓存）、显存不够时换出或丢掉重算（抢占）。\par}
\end{enumerate}\endgroup
\numline{公式}{每 token KV 字节 $=2\times\text{层数}\times\text{KV 头数}\times\text{每头维度}\times\text{字节数}$；Llama-3-8B：$2\times32\times8\times128\times2=128$ KiB}
\numline{数字}{TinyLM 处理位置 7 对 28；无缓存 1,499,500 位置 / 22.5 PFLOPs；有缓存 1,999 位置 / 30 TFLOPs，约 750 倍；每 token 128 KiB；单请求 250 MiB；16$\times$8192 即 16 GiB}
\fol{历史 token 的 q 为什么不用存？}{历史 token 不再提问，只有新 token 的 q 需要和缓存里全部 k 打分。存 q 既没有地方会被用到，又要多占一半显存，所以缓存里只放 k 和 v。}
\fol{缓存能跨请求复用吗？}{只有共享前缀能。请求结束缓存就释放，下一个请求即使开头相同默认也要重新 prefill；前缀缓存按块做哈希匹配，命中才能跳过这部分的 prefill。}
\fol{怎么验证缓存没算错？}{拿带缓存的逐步生成和每步整句重算对比每一步的概率分布，TinyLM 上两者完全一致；再核对处理的位置数，有缓存 7 个、无缓存 28 个，账对得上才算对。}
\numline{别答错}{别把它说成缓存注意力权重：缓存的是每层每个 token 的 k、v，请求结束即释放，不缓存。}
\phantomsection\addcontentsline{toc}{section}{70\quad TTFT、TPOT 与吞吐}
\lec{70}{TTFT、TPOT 与吞吐}
\askline{TTFT 和 TPOT 分别由什么决定？}
\askline{并发变高，吞吐和延迟怎么变？}
\askline{推理服务的吞吐上限在哪？}
\vspace{3pt}
\noindent{\qlabel{答}}\begingroup\sloppy\begin{enumerate}[leftmargin=1.6em,itemsep=2.5pt,topsep=0pt,parsep=0pt,label={\sffamily\footnotesize\color{accent}\arabic*.}]
\item {\fontsize{10}{14.8}\selectfont 三个指标各管一件事：TTFT 是从请求发出到第一个 token，由 prefill 加排队决定；TPOT 也叫 ITL，是之后每个 token 的间隔，由 decode 一步决定；吞吐是整个系统每秒生成多少 token，是服务商口径。前两个看单请求，最后一个看系统。\par}
\item {\fontsize{10}{14.8}\selectfont 一个请求的总耗时约等于 TTFT 加 TPOT 乘以生成的 token 数减 1。人的阅读速度大约是每秒几个到十几个 token，TPOT 在几十毫秒以内用户就感觉很流畅。\par}
\item {\fontsize{10}{14.8}\selectfont 只服务一个请求时，8B 在 H100 上 TTFT 约 15.5 ms、TPOT 约 4.85 ms，相当于每秒 206 个 token。体验很好，可 GPU 几乎闲着：decode 的算术强度约 1，算力只用了 0.3\%，一张很贵的卡每秒只产出 206 个 token。\par}
\item {\fontsize{10}{14.8}\selectfont 批处理让 b 个请求共享一次权重读取，但 KV Cache 不能共享，一步读写量是 16 GB 权重加上 b$\times$192 MiB。并发从 1 到 8，TPOT 只从 4.85 涨到 5.27 ms，总吞吐却从 206 涨到 1,517 token/s，几乎是白捡的。\par}
\item {\fontsize{10}{14.8}\selectfont 并发继续加，KV Cache 读取开始占主导：128 个并发时 TPOT 12.49 ms、单请求掉到 80 token/s，总吞吐 10,251 token/s。显存是硬上限，每请求 2048 token 占 256 MiB，80 GB 的卡最多约 238 个请求。\par}
\item {\fontsize{10}{14.8}\selectfont 所以本质是用延迟换吞吐：先定服务等级目标，比如 TTFT 不超过 500 ms、TPOT 不超过 50 ms，再在满足它的前提下把并发推到最高。注意即使并发 256，算术强度也只有约 60，远不到 295，长上下文的 decode 始终带宽受限。\par}
\end{enumerate}\endgroup
\numline{公式}{总耗时 $\approx\text{TTFT}+\text{TPOT}\times(N_{\text{out}}-1)$；总吞吐 $=b/\text{TPOT}$；一步读写 $\approx16\ \text{GB}+b\times\text{每请求 KV}$}
\numline{数字}{TTFT 15.5 ms；TPOT 4.85 / 5.27 / 12.49 ms；吞吐 206 / 1,517 / 10,251 token/s；显存上限约 238 个请求（TPOT 19.1 ms、12,463 token/s）；每请求 KV 256 MiB；并发 256 强度约 60}
\fol{TTFT 里排队占多少？}{系统空闲时几乎全是 prefill，8B 一个 1024 token 的提示约 15.5 ms；繁忙时排队往往是大头，一个很长的 prefill 插进来还会把所有人的 TPOT 顶上去。}
\fol{为什么要看 P99 而不是平均值？}{平均好看、P99 很差，说明总有一部分用户等得很久。线上通常同时报 TTFT 和 TPOT 的 P50、P90、P99，压测也是看分布，不能只看一个均值。}
\fol{吞吐高就能说明服务好吗？}{不是。并发 128 时总吞吐是并发 1 的 50 倍，但单请求从 206 掉到 80 token/s；要按 SLO 一起看，还要对齐口径是每请求速度、系统吞吐还是含输入。}
\numline{别答错}{吞吐高不等于每个用户都快；「每秒 token 数」要看清单请求速度、系统吞吐还是含输入。}
\phantomsection\addcontentsline{toc}{section}{71\quad 显存账本}
\lec{71}{显存账本}
\askline{80 GB 的卡能放多少请求？}
\askline{权重装得下就能部署了吗？}
\askline{KV Cache 到底能占多少显存？}
\vspace{3pt}
\noindent{\qlabel{答}}\begingroup\sloppy\begin{enumerate}[leftmargin=1.6em,itemsep=2.5pt,topsep=0pt,parsep=0pt,label={\sffamily\footnotesize\color{accent}\arabic*.}]
\item {\fontsize{10}{14.8}\selectfont 把一张卡的显存当成一本账，四项：权重、KV Cache、激活、其他开销，再留引擎的预留。80 GB 的卡按 90\% 只有 67.1 GiB 可用，因为 vLLM 这类引擎默认 gpu\_memory\_utilization=0.9，剩下留给意外和碎片。\par}
\item {\fontsize{10}{14.8}\selectfont 权重最好算也最确定，参数量乘字节数：8B 的 BF16 是 80.3 亿$\times$2 字节约 16.06 GB 即 15.0 GiB，FP8 减半、4 bit 再减半；张量并行 n 卡时每张卡放 1/n。\par}
\item {\fontsize{10}{14.8}\selectfont KV Cache 每 token 是 2$\times$层数$\times$KV 头数$\times$每头维度$\times$字节数，8B 的 BF16 是 128 KiB，乘上所有请求的 token 总数就是总量。激活按一次处理的 token 数预留，prefill 8192 token 约 0.91 GiB。\par}
\item {\fontsize{10}{14.8}\selectfont 8B 单卡算下来：67.1 GiB = 权重 15.0 + 激活 0.91 + 其他 1.9 + KV Cache 49.3，49.3 GiB 能放约 40.4 万个 token，每请求 4096 token 就是 99 个并发请求。\par}
\item {\fontsize{10}{14.8}\selectfont 权重和 KV 都换 FP8：权重降到 7.5 GiB、每 token 降到 64 KiB，KV 空间涨到 56.8 GiB，能放约 93 万个 token 即 227 个请求。这就是低精度在推理里最直接的价值：同样的卡服务更多人。\par}
\item {\fontsize{10}{14.8}\selectfont 70B 更能说明问题：BF16 权重 131.4 GiB 单卡放不下；FP8 的 65.7 GiB 勉强塞进 67.1 GiB，却一点 KV 空间都不剩，服务不了任何请求；四卡 BF16 每张卡 32.9 GiB 权重、剩 31.9 GiB，才够放 102 个请求。\par}
\end{enumerate}\endgroup
\numline{公式}{显存 $=\text{权重}+\text{KV Cache}+\text{激活}+\text{其他}$（再留约 $10\%$）；最大并发 $=\frac{\text{剩余 KV 空间}}{\text{每请求 KV}}$}
\numline{数字}{可用 67.1 GiB；8B 权重 15.0 / 激活 0.91 / 其他 1.9 / KV 49.3 GiB；40.4 万 token、99 请求；FP8 227 请求；70B FP8 65.7 GiB 无 KV 空间；四卡 102 请求；80 GB $\approx$ 74.5 GiB}
\fol{引擎启动时怎么定 KV 池？}{先加载权重，再用一批假数据跑一次最大规模的前向测出激活峰值，然后可用显存减权重、减激活、减其他，把剩下的全部申请下来切成固定大小的块，日志会打印能放多少块。}
\fol{为什么 GB 和 GiB 一定要分清？}{标称 80 GB 是 80$\times$10⁹ 字节，约 74.5 GiB，差 7\%；引擎日志按 GiB 算，混着用会把并发高估一成，账本从一开始就错。}
\fol{长上下文的并发为什么掉得厉害？}{每个请求 131072 token 时 8B 的 KV Cache 就要 16 GiB，一张 80 GB 的卡只能放 3 个。并发低吞吐就低，每个 token 的成本也就高得多。}
\numline{别答错}{能装下权重不等于能部署：70B 的 FP8 吃满显存，没地方放 KV Cache。}
\phantomsection\addcontentsline{toc}{section}{72\quad 推理成本模型}
\lec{72}{推理成本模型}
\askline{每百万 token 的成本怎么估？}
\askline{为什么输出比输入贵好几倍？}
\askline{怎么把这个成本降下来？}
\vspace{3pt}
\noindent{\qlabel{答}}\begingroup\sloppy\begin{enumerate}[leftmargin=1.6em,itemsep=2.5pt,topsep=0pt,parsep=0pt,label={\sffamily\footnotesize\color{accent}\arabic*.}]
\item {\fontsize{10}{14.8}\selectfont 成本逻辑只有一条：GPU 按时间收费，一小时的租金除以这一小时能处理多少 token，再乘 10⁶ 就是每百万 token 的成本。prefill 和 decode 瓶颈不同，所以输入和输出一定要分开算，不能用一个平均速度。\par}
\item {\fontsize{10}{14.8}\selectfont 输入 token 走 prefill，计算受限，每个 token 约 2$\times$75 亿次运算，H100 每秒 989 万亿次，理论上每秒能处理约 6.5 万个输入 token，这个数和并发几乎无关，所以输入成本很便宜。\par}
\item {\fontsize{10}{14.8}\selectfont 输出 token 走 decode，带宽受限，吞吐是 b 除以 TPOT，完全由并发决定。还要乘一个效率系数，示例取 60\%，因为真实引擎达不到峰值算力和带宽，这个数要用实测校准。\par}
\item {\fontsize{10}{14.8}\selectfont 按租金 2.5 美元每小时、效率 60\% 算：并发 1 时每百万输出 token 要 5.618 美元，并发 32 是 0.243，并发 128 只要 0.113，同一张卡同一个模型差 50 倍；输入始终约 0.018 美元。并发决定输出成本，这就是大家死磕批处理的原因。\par}
\item {\fontsize{10}{14.8}\selectfont 输出比输入贵是结构性的：并发 128 时输出是输入的 6 倍多，因为 prefill 时同一份权重被很多 token 共用，decode 每步都要把它重读一遍。这也解释了为什么几乎所有大模型 API 的输出定价都比输入贵好几倍。\par}
\item {\fontsize{10}{14.8}\selectfont 换成绝对数：一个请求 1024 进 1024 出、并发 128 约 0.00013 美元，一万次约 1.3 美元；70B 用四张卡、并发 64，每百万输出约 0.936 美元，是 8B 的 8 到 9 倍，和参数量的比例差不多。\par}
\end{enumerate}\endgroup
\numline{公式}{每百万 token 成本 $=\frac{\text{每小时租金}}{\text{每小时 token 数}}\times10^{6}$；输出吞吐 $=\frac{b}{\text{TPOT}}\times\text{效率}$}
\numline{数字}{租金 2.5 美元/小时；效率 60\%；输出 5.618 $\rightarrow$ 0.243 $\rightarrow$ 0.113 美元/百万（并发 1/32/128）；输入 0.018；0.00013 美元/请求；70B 四卡输入 0.164、输出 0.936；估算与实测差 2 倍以内}
\fol{量化成 FP8 能省多少？}{权重字节减半、KV Cache 也减半，带宽受限时 TPOT 接近减半，并发上限还提高了，输出成本大约降一半。这是低精度在成本上最直接的收益。}
\fol{这个模型什么时候会算错？}{它假设整张卡只在做 prefill 或只在做 decode，忽略了请求到达不均匀、夜里 GPU 闲着也照计费；所以量级对，和实测通常差在两倍以内。}
\fol{最便宜就是把并发拉满吗？}{单价上是的，但并发推到显存上限时 TPOT 也最长。成本要在满足延迟目标的前提下比，否则省下的钱会被用户感受到的变慢吃掉。}
\numline{别答错}{API 价格不等于成本，含利润与闲时浪费；最便宜不是并发拉满，而是在 SLO 内拉满。}
\phantomsection\addcontentsline{toc}{section}{73\quad 解码策略}
\lec{73}{解码策略}
\askline{温度、Top-k、Top-p 到底在调什么？}
\askline{Top-p 比 Top-k 好在哪？}
\askline{线上为什么不能直接纯采样？}
\vspace{3pt}
\noindent{\qlabel{答}}\begingroup\sloppy\begin{enumerate}[leftmargin=1.6em,itemsep=2.5pt,topsep=0pt,parsep=0pt,label={\sffamily\footnotesize\color{accent}\arabic*.}]
\item {\fontsize{10}{14.8}\selectfont 模型每步给出十几万维的分布，解码策略只做一件事：抽样之前把不可能的长尾砍掉，再在剩下的词里按概率抽。砍长尾换可靠性，留几个候选换多样性，这就是全部取舍。\par}
\item {\fontsize{10}{14.8}\selectfont 温度把每个 logit 除以 T 再 softmax，只改分布尖平，不删任何词：「好」的概率 T=0.5 时从 0.321 升到 0.544，T=1.5 时降到 0.235，而「猫」从 0.001 涨到 0.004，长尾始终大于 0。\par}
\item {\fontsize{10}{14.8}\selectfont Top-k 只留概率最高的 k 个，其余置 0 再归一化。k=3 时只剩「好、不错、晴朗」，概率变成 0.472、0.316、0.212。问题在 k 固定：模型确定时留一堆垃圾，不确定时又砍掉好词。\par}
\item {\fontsize{10}{14.8}\selectfont Top-p 从高到低累加，累计第一次达到 p 就停。p=0.8 加到第 5 个「冷」是 0.849，于是保留 5 个；模型确定时一两个词就凑够，不确定时要很多词，能自适应，所以最常用。\par}
\item {\fontsize{10}{14.8}\selectfont Min-p 按最大概率定门槛，min-p=0.1 时门槛是 0.321$\times$0.1=0.0321，留下 7 个词；模型越自信砍得越狠，高温度下比 Top-p 更不容易放进怪词。常见顺序是温度、Top-k、Top-p、Min-p、归一化。\par}
\item {\fontsize{10}{14.8}\selectfont 采样在 GPU 上做，一个 batch 里每个请求的参数还可能都不一样；128256 维排序不便宜，引擎用只找前若干个的专用算子，不把 logits 搬回 CPU。要严格复现就老老实实用贪心并固定配置。\par}
\end{enumerate}\endgroup
\numline{公式}{$p_{i}=\operatorname{softmax}(\text{logit}_{i}/T)$；Top-$p$ 取最小 $n$ 使 $\sum_{i\le n}p_{i}\ge p$；Min-$p$ 保留 $p_{i}\ge\text{min\_p}\cdot\max p$}
\numline{数字}{T=0.5「好」0.544、T=1.5「好」0.235 / 猫 0.004；Top-k=3 $\rightarrow$ 0.472/0.316/0.212；Top-p=0.8 保留 5 个；Min-p=0.1 门槛 0.0321 留 7 个；T=0.7+Top-p=0.9「好」0.458，抽 1000 次 448/274/127/86/65；词表 128256}
\fol{温度调到 2 再配 Top-p 会怎样？}{分布被抹平、长尾概率抬头，Top-p 要累加更多词才凑够 p，保留集合变大，怪词更容易进来。这种高温场景 Min-p 更稳，因为门槛随最大概率一起缩放。}
\fol{温度为 0 是除以 0 吗？}{不是。引擎判断 T=0 就直接取 argmax 走贪心，不做除法，否则 exp(logit/0) 会溢出。所以 T=0 和「很小的 T」行为不同，前者完全确定。}
\fol{怎么验证参数真的生效？}{固定种子抽 1000 次，看命中次数是否贴合归一化概率，448/274/127/86/65 就是 T=0.7+Top-p=0.9 的期望；再和线上日志比，偏移超过 5\% 就查是否被默认值覆盖。}
\numline{别答错}{Top-p 不是「概率大于 p 的词」，是从高到低累加到 p；温度 0 是贪心，不是除 0。}
\chapter{推理引擎与调度}
{\footnotesize\sffamily\color{grayt}第 74--82 讲}\par\vspace{-4pt}
\phantomsection\addcontentsline{toc}{section}{74\quad Continuous Batching}
\lec{74}{Continuous Batching}
\askline{线上 batching 和训练时差在哪？}
\askline{连续批处理为什么能提吞吐？}
\askline{它会让单个请求变慢吗？}
\vspace{3pt}
\noindent{\qlabel{答}}\begingroup\sloppy\begin{enumerate}[leftmargin=1.6em,itemsep=2.5pt,topsep=0pt,parsep=0pt,label={\sffamily\footnotesize\color{accent}\arabic*.}]
\item {\fontsize{10}{14.8}\selectfont 真实服务里请求陆续到达、每个请求要生成多少 token 事先不知道。静态批处理整批一起进、一起出，短请求结束后槽位只能空等最长的那个，新请求哪怕有空位也要排队。\par}
\item {\fontsize{10}{14.8}\selectfont 连续批处理把调度粒度从「一批请求」缩到「一步迭代」：每步结束的请求立刻移出，空出的位置立刻从等待队列补进来，批的组成一直在流动，所以也叫迭代级调度。\par}
\item {\fontsize{10}{14.8}\selectfont 10 个请求、4 个槽位、共 45 个 token：静态要 20 步、平均延迟 7.9 步、利用率 56\%；连续只要 13 步、5.0 步、87\%。理论上限是 12 步，连续已经很接近。\par}
\item {\fontsize{10}{14.8}\selectfont 实现上不做 padding：把这一步所有请求的 token 首尾相接拼成一张 [总 token 数, d] 矩阵，线性层本来就逐 token 独立，一次大矩阵乘就够；只有注意力按每个请求的起止位置分开算。\par}
\item {\fontsize{10}{14.8}\selectfont 代价在新请求的 prefill：它加入的那一步同时有 decode 和几千 token 的 prefill，计算量涨几百倍，所有正在 decode 的请求都被拖一步，TPOT 从 6.05 ms 跳到 138.9 ms。\par}
\item {\fontsize{10}{14.8}\selectfont 它不改变任何请求的结果，只决定谁在同一步一起算。批的上限由显存和延迟目标决定，vLLM 里是 max\_num\_seqs 和 max\_num\_batched\_tokens，想法来自 2022 年的 Orca。\par}
\end{enumerate}\endgroup
\numline{公式}{槽位利用率 $=\frac{\text{busy}}{\text{slots}\times\text{steps}}$；步数下限 $=\left\lceil\frac{\text{总 token 数}}{\text{槽位数}}\right\rceil$}
\numline{数字}{10 请求 / 4 槽位 / 45 token；静态 20 步、7.9 步、56\%；连续 13 步、5.0 步、87\%；下限 12 步；Orca 2022；TPOT 6.05 $\rightarrow$ 138.9 ms}
\fol{动态批处理是一回事吗？}{不是。dynamic batching 是等一小段时间凑够一批再发，粒度还是请求级；连续批处理每步都重新组批，粒度是迭代级，利用率能差一倍以上。}
\fol{batch 越大越好吗？}{不是。批越大权重读取摊得越薄、吞吐越高，但每个请求的 TPOT 也跟着变长，同时能跑几个还受 KV Cache 显存限制，上限要按延迟目标和显存一起定。}
\fol{怎么确认它真的生效了？}{看运行队列是否长期贴着 max\_num\_seqs、GPU blocks 占用是否稳定在高位；再看 TPOT 分布，如果利用率只有 50\% 左右，说明调度还在等整批结束。}
\numline{别答错}{连续批处理不改变结果，只改「谁在同一步算」；它不是等一小段凑批的 dynamic batching。}
\phantomsection\addcontentsline{toc}{section}{75\quad PagedAttention}
\lec{75}{PagedAttention}
\askline{KV Cache 的显存浪费在哪？}
\askline{分页之后并发为什么翻三倍？}
\askline{块表怎么定位一个 token？}
\vspace{3pt}
\noindent{\qlabel{答}}\begingroup\sloppy\begin{enumerate}[leftmargin=1.6em,itemsep=2.5pt,topsep=0pt,parsep=0pt,label={\sffamily\footnotesize\color{accent}\arabic*.}]
\item {\fontsize{10}{14.8}\selectfont 传统做法按最大上下文 4096 给每个请求预留一整块连续显存，浪费有三种：预留了却没写满的内部碎片、还没生成就被占住的后半段、请求进出释放后凑不出连续空间的外部碎片。\par}
\item {\fontsize{10}{14.8}\selectfont 64 个请求的快照实际只用到 86151 个 token：按最大长度预留占 262144 个，利用率 32.9\%；就算按最终长度预留（实际做不到）也只有 85.5\%；分页每块 16 个 token，能到 99.5\%。\par}
\item {\fontsize{10}{14.8}\selectfont 分页把整个 KV 池切成固定大小的块按需分配，块表记录第几个逻辑块落在哪个物理块：第 i 个 token 在 $\lfloor$i/16$\rfloor$ 块的第 i mod 16 个位置，浪费只剩每个请求最后一块的空位。\par}
\item {\fontsize{10}{14.8}\selectfont 块内 16 个 token 是连续的，GPU 按 32 或 128 字节缓存行合并访存，每跨 16 个 token 才查一次块表。Llama-3-8B 每 token 的 KV 是 128 KiB，一块正好 2 MiB，是个干净的对齐单位。\par}
\item {\fontsize{10}{14.8}\selectfont 省下的显存直接变成并发：49.3 GiB 的池子按 4096 预留只能放 98 个请求，分页时平均每请求占约 1352 个 token，能放约 299 个，并发三倍，这就是 vLLM 吞吐提升的主要来源。\par}
\item {\fontsize{10}{14.8}\selectfont 代价是注意力算子要一边查块表一边去不同地方读：块越小查表开销越大，块越大最后一块浪费越多、共享粒度也越粗，16 是经验折中；FlashInfer 这类算子库已原生支持块表。\par}
\end{enumerate}\endgroup
\numline{公式}{逻辑块 $=\left\lfloor i/16\right\rfloor$，块内偏移 $=i\bmod 16$；每 token KV $=2\times\text{层数}\times\text{KV 头数}\times\text{头维}\times2\ \text{字节}=128$ KiB}
\numline{数字}{利用率 32.9\% / 85.5\% / 99.5\%；占用 262144 / 100758 / 86560 个 token；池 403865 个 token；98 $\rightarrow$ 299 个请求；每 token 128 KiB、每块 2 MiB；共享 4800 $\rightarrow$ 1840 省 62\%}
\fol{分页能减小每 token 的 KV 吗？}{不能。Llama-3-8B 每 token 还是 128 KiB，分页减的是「预留了没用上」的浪费；要减小单 token 的大小得靠 GQA、MLA 和 KV Cache 量化。}
\fol{共享的块怎么防止互相污染？}{每个块记引用计数，被多个请求引用时只读；谁要写最后那个没装满的块，就先 copy-on-write 复制一份给自己再写，引用计数归 0 才真正还回空闲列表。}
\fol{块大小为什么选 16？}{折中：块小查表跳转多、算子开销大；块大最后一块浪费多、前缀共享粒度粗。2 MiB 一块又刚好对齐显存映射，实测 16 在开销和浪费之间最好用。}
\numline{别答错}{分页不减每 token 的 KV 大小，只减预留浪费；块表映射的是同一块显存池。}
\phantomsection\addcontentsline{toc}{section}{76\quad Prefix Cache}
\lec{76}{Prefix Cache}
\askline{多轮对话为什么要重算前缀？}
\askline{指纹为什么必须链式哈希？}
\askline{整段命中能跳过 prefill 吗？}
\vspace{3pt}
\noindent{\qlabel{答}}\begingroup\sloppy\begin{enumerate}[leftmargin=1.6em,itemsep=2.5pt,topsep=0pt,parsep=0pt,label={\sffamily\footnotesize\color{accent}\arabic*.}]
\item {\fontsize{10}{14.8}\selectfont 因为因果掩码，每个 token 的 k、v 只取决于它自己和它前面的 token。两个请求只要前 n 个 token 完全相同，这 n 个 token 的 KV 就完全相同，不管后面接的是什么内容。\par}
\item {\fontsize{10}{14.8}\selectfont 判断开头相同用块级链式哈希：hash\_k = H(hash\_\{k-1\}, 本块 16 个 token, extra)。这样每个块的指纹就是从第 0 个 token 到它末尾整段历史的摘要，前面任何一处不同都会让后面全变。\par}
\item {\fontsize{10}{14.8}\selectfont 必须链式，因为块内 token 的 KV 已经融合了前面全部上下文。两个请求第 1 块内容一样、第 0 块不同，KV 数值差得很远，只按当前块算哈希就会误命中，输出会彻底跑偏。\par}
\item {\fontsize{10}{14.8}\selectfont 查表从头逐块进行，遇到第一个未命中就停，后面即使碰巧存在也不能用，因果链已经断了。只给装满的块登记哈希，未装满的末尾块还会被写入，共享出去会和其他请求冲突。\par}
\item {\fontsize{10}{14.8}\selectfont 命中的块直接放进新请求的块表、引用计数加 1，只对未命中部分做 prefill。引用计数归 0 的块不释放、数据也不清，进空闲 LRU 队列保留哈希，显存耗尽才从 LRU 头部淘汰叶子块。\par}
\item {\fontsize{10}{14.8}\selectfont 收益很直接：系统提示 2000 token 加问题 100 token，prefill 从 31.9 ms 降到 4.8 ms，TTFT 约六分之一；所以很多 API 对命中缓存的输入 token 打折。\par}
\end{enumerate}\endgroup
\numline{公式}{$\text{hash}_{0}=H(\text{None},\ \text{tokens}[0{:}16],\ \text{extra})$；$\text{hash}_{k}=H(\text{hash}_{k-1},\ \text{tokens}[16k{:}16k+16],\ \text{extra})$}
\numline{数字}{块 16 个 token；4 请求 154 token 命中 52 个、命中率 34\%；prefill 2100 token 31.9 ms $\rightarrow$ 100 token 4.8 ms；Agent 命中率上限 98\%、实践 90\textasciitilde{}95\%}
\fol{整段命中就完全不用算了吗？}{不行。KV Cache 里只有 K 和 V，没有最后一层的 hidden state 和 logits，提示的最后一个 token 必须单步前向一次，才能采样出第一个回答 token。}
\fol{淘汰为什么只从叶子开始？}{中间块一旦被淘汰，因果链就断了，它后面的块即使还在显存里也再不可能被命中；所以按 LRU 从链条末端剪枝，公共的前缀根节点永远不动。}
\fol{多机部署命中率为什么崩？}{哈希表只在单实例内。10 台机器随机路由同一段系统提示，命中率会掉到 10\% 左右；要改成缓存感知路由，把同前缀的请求粘到同一台机器上。}
\numline{别答错}{只能复用从头连续的前缀，中间相同片段不行；复用的是 KV 不是答案，最后一个 token 仍要前向。}
\phantomsection\addcontentsline{toc}{section}{77\quad Chunked Prefill}
\lec{77}{Chunked Prefill}
\askline{长提示为什么把大家卡住？}
\askline{预算 512 是怎么来的？}
\askline{切块会让首字更慢吗？}
\vspace{3pt}
\noindent{\qlabel{答}}\begingroup\sloppy\begin{enumerate}[leftmargin=1.6em,itemsep=2.5pt,topsep=0pt,parsep=0pt,label={\sffamily\footnotesize\color{accent}\arabic*.}]
\item {\fontsize{10}{14.8}\selectfont 32 个请求一起 decode，每步约 6.05 ms，时间主要花在读一遍权重和 KV 上。这时来一份 8000 token 的长文档，整段 prefill 塞进一步，这一步要算 8032 个 token，变成计算受限的 138.9 ms。\par}
\item {\fontsize{10}{14.8}\selectfont 这一百多毫秒里那 32 个请求各只出了 1 个 token，TPOT 从 6 ms 跳到 139 ms。服务通常要求 P99 不超过 50 ms，偶尔来一次长提示就能让所有人的 P99 超标。\par}
\item {\fontsize{10}{14.8}\selectfont 切法是给每步设一个 token 预算：先给每个 decode 请求留 1 个，剩下给 prefill。预算 512 时 32 个 decode 用掉 32 个，剩 480 个切给 prefill，8000 个 token 分 17 步做完。\par}
\item {\fontsize{10}{14.8}\selectfont prefill 和 decode 本来就是同一件事，只是每次处理的新 token 数不同；每块的 KV 写进缓存后，下一块的注意力能看到前面所有块的 KV，最后一块做完的结果和一次做完完全相同。\par}
\item {\fontsize{10}{14.8}\selectfont 预算越小 decode 越稳、TTFT 越差：2048 时 37.4 与 140.4 ms，512 时 9.7 与 146.9 ms，256 因为每步都要重读一遍权重涨到 223.5 ms。512 以上每步基本计算受限，切块不增加总算力。\par}
\item {\fontsize{10}{14.8}\selectfont 还有白捡的算力：decode 一步本来就带宽受限，再塞 363 个 prefill token 耗时仍是 6.06 ms，prefill 几乎搭顺风车。vLLM V1 默认开启，预算参数是 max\_num\_batched\_tokens。\par}
\end{enumerate}\endgroup
\numline{公式}{一步耗时 $\approx\max\left(\frac{\text{计算量}}{\text{算力}},\ \frac{\text{读写字节}}{\text{带宽}}\right)$；每步 prefill token $=\text{预算}-\text{decode 请求数}$}
\numline{数字}{平时 6.05 ms；不切块 138.9 ms；2048$\rightarrow$140.4/37.4；1024$\rightarrow$147.5/19.5；512$\rightarrow$146.9/9.7（17 步）；256$\rightarrow$223.5/6.4（36 步）；白捡 363 个 token；P99 50 ms}
\fol{切块会改变输出结果吗？}{不会。每块的注意力都读前面所有块的 KV，得到的 KV Cache 和第一个 token 与一次做完一致；代价是后面的块要重复读前面的 KV，切太细反而更慢。}
\fol{每步预算到底该怎么定？}{没有。chunked prefill 保护的是别人的 TPOT，长提示的 TTFT 基本不变，预算太小还会变长；要缩短它得靠前缀缓存、更多算力或 PD 分离。}
\fol{每步预算到底该怎么定？}{让一步刚好落在「开始计算受限」的位置。H100 上的 8B 模型约几百到两千 token；在线服务取小保 TPOT，离线批量取大保吞吐，再按 P99 和吞吐曲线微调。}
\numline{别答错}{切块不改变结果，也不缩短长提示自己的 TTFT；预算太小 TTFT 反而变长。}
\phantomsection\addcontentsline{toc}{section}{78\quad 调度与抢占}
\lec{78}{调度与抢占}
\askline{KV Cache 池子满了怎么办？}
\askline{为什么抢占最后进来的？}
\askline{被抢占后是重算还是换出？}
\vspace{3pt}
\noindent{\qlabel{答}}\begingroup\sloppy\begin{enumerate}[leftmargin=1.6em,itemsep=2.5pt,topsep=0pt,parsep=0pt,label={\sffamily\footnotesize\color{accent}\arabic*.}]
\item {\fontsize{10}{14.8}\selectfont 调度器手里有等待队列和运行队列：等待队列放还没开始以及被抢占后等着重来的请求，运行队列放正在 prefill 或 decode 的。每一步都要决定放谁进来、块不够时让谁让位。\par}
\item {\fontsize{10}{14.8}\selectfont 保守做法按「提示长度 + 最大生成长度」把块一次预留好，进来就不会缺块，但大部分预留长期空着，并发很低；引擎都选按需分配、每写满一块再要一块，再用抢占兜底。\par}
\item {\fontsize{10}{14.8}\selectfont 6 个请求、16 个块的例子里：按需分配时第 5 步池子用完，抢占最后加入的 F，第 6 步又抢占 E，13 步完成、平均延迟 8.7 步、抢占 2 次、重算 16 个 token。按最终长度预留是 17 步、10.8 步、零抢占。\par}
\item {\fontsize{10}{14.8}\selectfont 抢最后加入的，一是符合先来先服务的公平性，二是它算得最少、扔掉的 KV 最少重算最省，三是它回到等待队列最前面，一有空位就最先回来，不会被后来的请求饿死。\par}
\item {\fontsize{10}{14.8}\selectfont 发生抢占的那一步不再放新请求，否则刚赶走一个又被占满，请求在进进出出之间抖动；有的实现还在池子里留水位线，空闲块低于阈值就提前停止放行。\par}
\item {\fontsize{10}{14.8}\selectfont 被抢占的 KV 有重算和换出两种：2000 token 的请求 KV 是 262 MB，重算约 30.3 ms、换出换回约 10.5 ms，但引擎仍偏向重算，因为不占 PCIe，还能吃到前缀缓存和分块 prefill。\par}
\end{enumerate}\endgroup
\numline{公式}{需要的块数 $=\left\lceil\frac{\text{提示长度}+\text{已生成 token 数}}{\text{块大小}}\right\rceil$}
\numline{数字}{16 块 $\times$ 4 token；按需 13 步 / 8.7 步 / 抢占 2 次 / 重算 16 token；预留 17 步 / 10.8 步；262 MB、30.3 ms vs 10.5 ms；PCIe 5.0 约 50 GB/s}
\fol{抢占会让请求失败吗？}{不会。它只是暂停，已生成的 token 都保留，回来重新 prefill 后接着生成，结果和没被抢占一致；用户最多感觉到一次停顿。}
\fol{抢占和算力不足是一回事吗？}{不是。算力不够表现为每步变慢、TPOT 拉长；显存不够才触发抢占。抢占如果持续发生，说明负载超过这张卡，该做的是限流或加卡。}
\fol{默认 FCFS 不总是最好？}{队头阻塞：长请求排在前面但池子放不下，后面一串短请求即使放得下也只能干等。可以允许短请求越过，但要防止长请求被无限越过。}
\numline{别答错}{抢占不是取消，结果不变；它解决显存不够而非算力不够，也不是丢弃已生成 token。}
\phantomsection\addcontentsline{toc}{section}{79\quad 多 LoRA 服务}
\lec{79}{多 LoRA 服务}
\askline{100 个客户的模型怎么共用一张卡？}
\askline{LoRA 推理几乎不要钱吗？}
\askline{SGMV 到底省了什么？}
\vspace{3pt}
\noindent{\qlabel{答}}\begingroup\sloppy\begin{enumerate}[leftmargin=1.6em,itemsep=2.5pt,topsep=0pt,parsep=0pt,label={\sffamily\footnotesize\color{accent}\arabic*.}]
\item {\fontsize{10}{14.8}\selectfont 不能合并。合并后的 W′ 是某个客户专属的，换个客户就要换一份 16 GB 权重；所以多 LoRA 服务不合并：大路所有请求一起乘同一个 W，小路每个请求用自己的 A、B。\par}
\item {\fontsize{10}{14.8}\selectfont r=16 时一个适配器 41.9 M 参数、BF16 占 83.9 MB，只有 8.03 G 参数基座的 0.52\%；每个 token 只过自己那一个适配器，计算量也只多 0.52\%，看起来几乎免费。\par}
\item {\fontsize{10}{14.8}\selectfont 但 decode 是带宽受限的，每步要把用到的权重从显存完整读一遍：基座读一遍，batch 里出现几个不同适配器就要多读几份，代价和不同适配器的个数成正比，不是和请求数成正比。\par}
\item {\fontsize{10}{14.8}\selectfont 不同适配器个数 D 的期望是 $\Sigma$[1-(1-p\_i)\textasciicircum{}B]。100 个客户完全均匀、batch 64 时平均 47.4 个、每步多读 3.98 GB；Zipf 的 s=1 时降到 31.4 个、2.64 GB，流量越集中越省。\par}
\item {\fontsize{10}{14.8}\selectfont 逐个适配器算的话，32 层 $\times$ 7 个线性层 $\times$ 2 次矩阵乘 $\times$ D 个适配器，D=30 时每步一万三千多次 kernel。Punica、S-LoRA 的 SGMV 一次 kernel 算完整个 batch，每层固定 2 次。\par}
\item {\fontsize{10}{14.8}\selectfont 显存分两级：GPU 里留 max\_loras 个槽位，CPU 里留 max\_cpu\_loras 个按 LRU 淘汰。84 MB 的适配器经 PCIe 拷进显存约 1.3 ms，约为一步 decode 的四分之一，所以要尽量别每步都在换。\par}
\end{enumerate}\endgroup
\numline{公式}{$y=xW+\frac{\alpha}{r}\cdot(xA)B$；$D=\sum_{i=1}^{N}\left[1-(1-p_{i})^{B}\right]$}
\numline{数字}{r=8 21.0M/41.9MB；r=16 41.9M/83.9MB/0.52\%；r=64 167.8M/335.5MB；s=0 $\rightarrow$ 47.4 个/3.98GB，s=1 $\rightarrow$ 31.4 个/2.64GB；一步 4.79$\rightarrow$5.58 ms；拷 84MB $\approx$ 1.3 ms；kernel 13,440 $\rightarrow$ 448}
\fol{前缀缓存能跨适配器共享吗？}{不能。LoRA 不改 KV 的维度，但不同适配器算出的 K、V 不同，哪怕提示一字不差，跨适配器命中也会出错，所以哈希里要混入适配器标识。}
\fol{max\_loras 设大点不行吗？}{设大上限就高，一步里可能混进更多适配器，那一步被拖慢。vLLM 宁可让用冷门适配器的请求多等几步，也不让一步里适配器太多，是吞吐和公平的取舍。}
\fol{prefill 也是这笔账吗？}{不是。prefill 每个请求有几百上千 token，小路的矩阵乘不再小，计算量按 r 与 d 的比例增加但仍可忽略；读权重不免费主要是 decode 的问题。}
\numline{别答错}{LoRA 计算几乎免费但读权重不免费，代价随不同适配器数涨；SGMV 只减调用不减字节。}
\phantomsection\addcontentsline{toc}{section}{80\quad 约束解码}
\lec{80}{约束解码}
\askline{怎么保证一定输出合法 JSON？}
\askline{每步都检查一遍不慢吗？}
\askline{约束解码能保证内容对吗？}
\vspace{3pt}
\noindent{\qlabel{答}}\begingroup\sloppy\begin{enumerate}[leftmargin=1.6em,itemsep=2.5pt,topsep=0pt,parsep=0pt,label={\sffamily\footnotesize\color{accent}\arabic*.}]
\item {\fontsize{10}{14.8}\selectfont 在选词之前动手脚：把格式写成状态机，每一步检查词表里每个 token，只要把它接上去就不可能再合法，就把它的 logit 设成 -$\infty$，剩下的重新归一化再照常采样。\par}
\item {\fontsize{10}{14.8}\selectfont 判断合法性要把 token 的字符逐个送进状态机，因为 token 往往不是一个字符，一个 token 可能一次跨过好几个字符甚至跨过段边界，只看首字符就会判错。\par}
\item {\fontsize{10}{14.8}\selectfont 结束标记只在格式完整时才允许输出，所以整个输出一定合法。那个手写模型本来会先说「好的，JSON：」、把年龄写成「十二岁」，加约束后第 1 步屏蔽到只剩 \{（原概率 19\%），第 9 步只剩 12（13\%）。\par}
\item {\fontsize{10}{14.8}\selectfont 真实词表有 128256 个 token，每步逐字符检查要几十毫秒，比 GPU 算一步还慢。做法是为每个状态提前算好掩码，Outlines 就是 JSON Schema 转正则、转状态机、再建 token 索引。\par}
\item {\fontsize{10}{14.8}\selectfont 嵌套结构有限状态机描述不了，要用上下文无关文法，检查器是带一个栈的下推自动机。XGrammar、llguidance 把大部分与栈无关的合法性提前算好，每步压到几十微秒。\par}
\item {\fontsize{10}{14.8}\selectfont 代价是屏蔽改变了原来的分布，可能把模型逼到低概率的 token 上。它保证格式不保证内容，生成长度被截断也会破坏格式，所以提示里仍然要把格式和字段含义写清楚。\par}
\end{enumerate}\endgroup
\numline{公式}{若 $\text{tokens}+t$ 送入状态机后无合法后继，则 $\text{logit}_{t}=-\infty$；仅当状态到达终止态才允许 $\langle\text{eos}\rangle$}
\numline{数字}{词表 128256；例子 11 步；第 1 步 \{ 原概率 19\%、第 9 步 12 原概率 13\%；逐 token 检查几十 ms $\rightarrow$ 优化后几十 $\mu$s；JSON 拆成 5 段}
\fol{JSON 模式等于约束解码吗？}{不等。有的 API 的 JSON 模式只保证输出能被解析，字段名和类型都不保证；按 JSON Schema 约束才保证结构完全一致，背后通常也是约束解码。}
\fol{强制先输出 JSON 有什么副作用？}{模型没机会先想。一开口就得是左花括号，思维链被挤掉了。常见做法是在 Schema 里把推理字段放在最前面，或先自由思考再只约束答案部分。}
\fol{怎么验证约束真的生效？}{跑一万条请求全量解析，失败率应为 0；再抽样看被屏蔽 token 的概率是不是真为 0，顺便测每步状态机的耗时，几十微秒以上就说明掩码没预计算。}
\numline{别答错}{只保证格式不保证内容；屏蔽改变分布；长度被截断仍会破坏格式。}
\phantomsection\addcontentsline{toc}{section}{81\quad 请求生命周期}
\lec{81}{请求生命周期}
\askline{一个请求从发出去到出第一个字，经过哪些模块？}
\askline{推理引擎为什么要拆成几个进程？}
\askline{TTFT 的时间都花在哪里？}
\vspace{3pt}
\noindent{\qlabel{答}}\begingroup\sloppy\begin{enumerate}[leftmargin=1.6em,itemsep=2.5pt,topsep=0pt,parsep=0pt,label={\sffamily\footnotesize\color{accent}\arabic*.}]
\item {\fontsize{10}{14.8}\selectfont 请求依次过前端、引擎核心、GPU 工作进程三段：API 服务器校验参数、分配编号、建一个输出通道，然后套聊天模板、分词成 token id，从这一步起引擎里流动的就只有数字了。\par}
\item {\fontsize{10}{14.8}\selectfont 三个进程之间用 ZeroMQ 通消息，因为分工性质不同：HTTP 收发、分词、反分词是 CPU 上琐碎的字符串处理，量大又时快时慢；调度循环一旦被打断 GPU 就要空等，而 Python 一个进程只有一个线程跑 Python 代码。\par}
\item {\fontsize{10}{14.8}\selectfont 引擎核心每转一步就是一轮调度：先给已经在 decode 的请求各排 1 个 token、缺块就向 KV 块管理器要，再用剩下的 token 预算切块排 prefill，最后还有预算和空闲块才从等待队列放新请求。\par}
\item {\fontsize{10}{14.8}\selectfont 放新请求时 KV 块管理器先查前缀缓存，命中的整块直接放进块表、引用计数加一，没命中的按需分配；调度器最后输出的是一张清单，写明这一步哪些请求参加、各处理几个 token、块表长什么样。\par}
\item {\fontsize{10}{14.8}\selectfont 输入准备把清单变成张量：所有请求这一步的 token 首尾拼成一条长序列，附上每个 token 的位置编号、每个请求的块表和新 token 该写进哪个块的哪个位置，GPU 前向只取各请求最后一位算 logits。\par}
\item {\fontsize{10}{14.8}\selectfont 采样按每个请求自己的参数走：有格式约束的先加掩码，再除温度、做 Top-p 截断，最后抽出一个 token；反分词必须增量，因为一个汉字可能被切成好几个 token，得等后面的 token 到了才能输出已确定的文字。\par}
\end{enumerate}\endgroup
\numline{公式}{$\text{TTFT}=\text{排队时间}+\text{prefill 时间}+\text{首个 token 的采样返回}$；总延迟 $=\text{TTFT}+(\text{输出 token 数}-1)\times\text{TPOT}$}
\numline{数字}{3 个进程、9 个模块；示例提示 12 个 token、每步预算 10，于是 prefill 切成 10 + 2 两块；R2 命中 2 块、省 8 个 token，块表 [0, 1, 4]；32 层 Transformer}
\fol{停止词为什么在反分词时才判断？}{停止词是文字不是 token，只有转回文字才能比对。命中时引擎已经多算了一两个 token，只是没返回给用户；要在 token 层拦就得另维护一份映射，不划算。}
\fol{CPU 上的开销真会影响吞吐吗？}{会。decode 一步 GPU 只跑几毫秒，CPU 上调度、输入准备、采样、反分词串行一两毫秒，GPU 就有相当部分时间空等；要重叠第 n 步和 n+1 步，用 CUDA Graph 压启动开销。}
\fol{不流式返回会省时间吗？}{不会。stream 只决定文字是边生成边送还是攒齐一次送，token 照样一个一个生成；只是不流式时 TTFT 对用户没有意义了。}
\numline{别答错}{停止词是文字不是 token，只能在反分词时判断，命中时引擎已多算一两个 token，只是没返回。}
\phantomsection\addcontentsline{toc}{section}{82\quad 推理压测实践}
\lec{82}{推理压测实践}
\askline{一张卡吞吐一万 token/s，这话可信吗？}
\askline{压测该怎么加压、该报哪些数？}
\askline{goodput 是什么，为什么比吞吐重要？}
\vspace{3pt}
\noindent{\qlabel{答}}\begingroup\sloppy\begin{enumerate}[leftmargin=1.6em,itemsep=2.5pt,topsep=0pt,parsep=0pt,label={\sffamily\footnotesize\color{accent}\arabic*.}]
\item {\fontsize{10}{14.8}\selectfont 单独一个吞吐数几乎没有信息量：得说清是什么模型、什么卡、提示和输出多长、同时多少个请求，那时用户的 TTFT 和 TPOT 又是多少；否则把并发堆高，吞吐总能很好看。\par}
\item {\fontsize{10}{14.8}\selectfont 可信的压测给出的是一条曲线：按泊松过程逐级提高到达率，每一级跑足够多的请求，记录输出吞吐和 TTFT、TPOT 的中位数与 P99，找出在满足延迟目标之前能承受多大压力。\par}
\item {\fontsize{10}{14.8}\selectfont 加压分两种：开环按到达率发，间隔服从指数分布，最接近线上真实用户；闭环固定 c 个用户、收到回答就发下一个，服务越慢发得越慢，永远不会压垮服务，也就把排队问题掩盖了。\par}
\item {\fontsize{10}{14.8}\selectfont 吞吐随到达率先线性增长，到某个点饱和；再加压只让请求排队，TTFT 跳一两个量级。Llama-3-8B 加 H100 在每秒 40 个时 8316 token/s，60 个时吞吐只多 6\%，P99 却从 197 ms 到 4176 ms。\par}
\item {\fontsize{10}{14.8}\selectfont goodput 是满足 SLO 的最大吞吐，说法来自 2024 年的 DistServe；比配置、比引擎、比硬件都该比 goodput，因为一张卡能承担的合格流量就是它，成本也就按它算。\par}
\item {\fontsize{10}{14.8}\selectfont 坑主要在数据集和样本量：长度分布不对结果差几倍，输出被结束标记截短就不具可比性，公共前缀让前缀缓存白送 prefill，没预热时前几个请求特别慢，600 个请求的 P99 只由最慢六七个决定。\par}
\end{enumerate}\endgroup
\numline{公式}{$\text{goodput}=\max\{\text{吞吐}\mid\text{TTFT P99}\le\text{SLO}\ \text{且}\ \text{TPOT P99}\le\text{SLO}\}$；过载区延迟随压测时长增长，不可比}
\numline{数字}{Llama-3-8B + H100、KV 池约 40 万 token、分块 prefill 每步 2048：5 个/s 时 1349 token/s；40 个/s 时 8316 token/s、TTFT 中位 59 ms / P99 197 ms、TPOT P99 25.2 ms；45 个/s 时 8570 token/s、P99 1306 ms；60 个/s 时 8780 token/s、TTFT P99 4176 ms}
\fol{为什么延迟一定要看 P99？}{用户的体验由最慢那几次决定。平均 TTFT 50 ms、P99 却是 3 秒，一百次请求里就有一次让人觉得卡死；中位数和 P99 一起报，才看得出尾部有多重。}
\fol{饱和点之后的数字还能比吗？}{不能。过载时队列一直变长，TTFT 取决于压测跑了多久，跑 10 秒和跑 10 分钟差很多；要比的是饱和点在哪、饱和之前的延迟，以及 goodput。}
\fol{怎么让结果可以复现？}{把条件写全：模型、精度、硬件、引擎及版本、批大小上限、token 预算、是否开前缀缓存、数据集、加压方式，再逐级给出吞吐和 P50、P99；漏一项就没法比较。}
\numline{别答错}{TPOT 先按请求平均、ITL 混所有间隔；偶发长间隔只在 ITL 的 P99 里显形。}
\chapter{算子与Kernel}
{\footnotesize\sffamily\color{grayt}第 83--89 讲}\par\vspace{-4pt}
\phantomsection\addcontentsline{toc}{section}{83\quad 算子融合}
\lec{83}{算子融合}
\askline{几个逐元素算子，为什么合成一个就快了？}
\askline{融合到底能省多少，怎么算出来？}
\askline{kernel 快了 2.6 倍，请求能快多少？}
\vspace{3pt}
\noindent{\qlabel{答}}\begingroup\sloppy\begin{enumerate}[leftmargin=1.6em,itemsep=2.5pt,topsep=0pt,parsep=0pt,label={\sffamily\footnotesize\color{accent}\arabic*.}]
\item {\fontsize{10}{14.8}\selectfont 逐元素运算每个元素只做几次乘加，算术强度不到 1，而 H100 的拐点是 295 FLOPs/字节，所以耗时几乎就是读写的字节数除以带宽，算得再快也要等数据；让它变快的唯一办法是少读写几遍显存。\par}
\item {\fontsize{10}{14.8}\selectfont 两个 kernel 之间只能经显存交换数据。不融合时 SwiGLU 拆成 sigmoid、乘 g、乘 u 三个 kernel，整个 117 MB 的张量读 5 遍写 3 遍共约 940 MB，按 3.35 TB/s 要 280 $\mu$s。\par}
\item {\fontsize{10}{14.8}\selectfont 融合成一个 kernel 后，每个线程读进自己负责的那几个 g 和 u，在寄存器里算完 sigmoid、SiLU 和乘积只写 y，中间量 t、s 从来不落地，读写剩 3 遍约 105 $\mu$s，是原来的八分之三。\par}
\item {\fontsize{10}{14.8}\selectfont 残差相加加 RMSNorm 也明显：逐个算子要 7 个 kernel、读写 10 遍约 100 $\mu$s，融合后 4 遍约 40 $\mu$s；放进 Llama-3-8B 的一层，从 481 $\mu$s 降到 185 $\mu$s，同层矩阵乘要约 1812 $\mu$s。\par}
\item {\fontsize{10}{14.8}\selectfont 常见的融合有四类：逐元素链（silu\_and\_mul、fused\_add\_rms\_norm）、矩阵乘的尾部融合、把 QKV 和 gate/up 投影拼成一次矩阵乘、整个注意力融成一个 kernel。\par}
\item {\fontsize{10}{14.8}\selectfont 边界在需要整行或整个张量结果的地方：softmax 的分母、整行共用的量化 scale、下一层的矩阵乘都得断开；而且省的是读写量，峰值显存不一定跟着降多少。\par}
\end{enumerate}\endgroup
\numline{公式}{$\text{耗时}\approx\text{读写遍数}\times\frac{\text{张量字节}}{\text{带宽}}$；整体加速 $=\frac{1}{(1-f)+f/s}$（Amdahl）}
\numline{数字}{SwiGLU 117 MB：8 遍 940 MB / 280 $\mu$s $\rightarrow$ 3 遍 105 $\mu$s；RMSNorm 34 MB：10 遍 100 $\mu$s $\rightarrow$ 4 遍 40 $\mu$s；一层 481 $\rightarrow$ 185 $\mu$s，同层矩阵乘 1812 $\mu$s；Qwen3-8B 请求约 817 ms，激活 13.6 $\rightarrow$ 12.2 ms}
\fol{为什么不能一直融合下去？}{后一个运算要前一个在所有位置上的结果时就得等它全部结束，比如全局最大值、整行共用的 scale、softmax 的分母；片上寄存器和共享内存也放不下太多，所以融合有明确的上限。}
\fol{读写少了，显存峰值也会降吗？}{不一定。三步分开时峰值 72 MiB，三步全融合是 60 MiB：读写少了 96 MiB，峰值只少 12 MiB，因为同一块显存可以反复读写。}
\fol{kernel 快 2.6 倍，请求呢？}{只快约 0.16\%。SwiGLU 激活从 13.6 ms 降到 12.2 ms，而请求约 817 ms；热点只占 20\% 时 Amdahl 上限也只有 1.25 倍，先看热点占比再决定做不做。}
\numline{别答错}{融合不减少计算量，减少的是数据搬运；kernel 的加速比不等于请求的加速比，热点占比决定上限。}
\phantomsection\addcontentsline{toc}{section}{84\quad FlashAttention}
\lec{84}{FlashAttention}
\askline{标准注意力慢在哪，为什么必须融合？}
\askline{softmax 要先看完整行，怎么分块算？}
\askline{块该切多大，怎么权衡？}
\vspace{3pt}
\noindent{\qlabel{答}}\begingroup\sloppy\begin{enumerate}[leftmargin=1.6em,itemsep=2.5pt,topsep=0pt,parsep=0pt,label={\sffamily\footnotesize\color{accent}\arabic*.}]
\item {\fontsize{10}{14.8}\selectfont 先要算出 s $\times$ s 的打分矩阵 S：s = 8192 时一个头 128 MiB；缩放、掩码、softmax、乘 V 让它读写 8 遍，一个头约 1082 MB、323 $\mu$s，而真正的乘加只要 35 $\mu$s，九成时间在等数据。\par}
\item {\fontsize{10}{14.8}\selectfont 难点在 softmax 要先知道整行最大值和总和。它用在线 softmax：对一块查询维护当前最大值 m、指数和 l、加权和 acc，读到更大的打分就把旧值乘 e\textasciicircum{}(m\_旧 - m\_新) 换算到新基准，读完输出 acc/l。\par}
\item {\fontsize{10}{14.8}\selectfont 这个结果是精确的，不是近似。8 个打分分两块走一遍：第二块出现 3.4 时旧的 l、acc 乘 0.2725 缩小，得到 l = 1.8411、acc = 3.5166，输出 1.9100，和一次算完整行完全相同。\par}
\item {\fontsize{10}{14.8}\selectfont 整个注意力塞进一个 kernel 后 S 只在片上出现，读写只剩 Q、K、V、O 和每块查询扫一遍的 K、V：约 277 MB、83 $\mu$s，是标准实现的四分之一；K、V 还能命中 50 MB 的 L2，注意力从卡带宽变回卡算力。\par}
\item {\fontsize{10}{14.8}\selectfont 另一个收益是显存：S 不再落地，注意力额外占用从 s$^2$ 级降到 s 级，每个查询只多存 m 和 l 两个数；s = 32768 时一个头的 S 就有 2 GiB，标准实现根本放不下几个请求，长上下文才因此可行。\par}
\item {\fontsize{10}{14.8}\selectfont 块大小由片上容量决定，也在读写和更新次数之间折中：99 KB 共享内存下 K/V 块取 64 行能放 82 行 Q、读写 404 MiB、更新 12,800 次；取 1 行只读 260 MiB，更新却多 41 倍，实战取 64 或 128。\par}
\end{enumerate}\endgroup
\numline{公式}{$l\leftarrow l\cdot e^{m-m_{\text{new}}}+\sum e^{x-m_{\text{new}}}$，$\text{acc}\leftarrow\text{acc}\cdot e^{m-m_{\text{new}}}+\sum e^{x-m_{\text{new}}}\cdot v$，输出 $=\text{acc}/l$}
\numline{数字}{s = 8192、d = 128：S 128 MiB，标准 1082 MB / 323 $\mu$s $\rightarrow$ FlashAttention 277 MB / 83 $\mu$s，计算 35 $\mu$s；实测 RTX PRO 6000 快约 33 倍、新增显存 848 $\rightarrow$ 22 MiB；A100 上达到峰值 50\%\textasciitilde{}73\%}
\fol{训练时反向传播怎么办？}{标准实现要把 P 存下来，又是 s$^2$ 级显存。FlashAttention 不存 P，只存每个查询的 m 和 l，反向时按块把 S 和 P 重算一遍：多做计算，省下大量读写和显存，整体反而更快。}
\fol{块越小不是读得越少吗？}{读得少但更新次数暴涨。b = 1 时读 260 MiB、更新 524,288 次，是 b = 64 的 41 倍，矩阵乘的一个维度只有 1，Tensor Core 用不上，实战取 64 或 128。}
\fol{它到底减少计算量了吗？}{没有。乘加次数和标准实现一样，训练时因为重算还多了一些；复杂度仍是 s$^2$，它只是让这份 s$^2$ 的计算跑在接近算力上限的地方，而不是被带宽拖住。}
\numline{别答错}{它既不近似也没减少计算量，复杂度仍是 s$^2$；省的只是显存读写和 s$^2$ 级的显存占用。}
\phantomsection\addcontentsline{toc}{section}{85\quad Decode 注意力 Kernel}
\lec{85}{Decode 注意力 Kernel}
\askline{decode 的注意力和 prefill 差在哪？}
\askline{KV 按块散落，kernel 怎么读？}
\askline{批只有 1 时 GPU 怎么才不闲着？}
\vspace{3pt}
\noindent{\qlabel{答}}\begingroup\sloppy\begin{enumerate}[leftmargin=1.6em,itemsep=2.5pt,topsep=0pt,parsep=0pt,label={\sffamily\footnotesize\color{accent}\arabic*.}]
\item {\fontsize{10}{14.8}\selectfont decode 每步每个请求只有 1 个查询，却要和上下文里所有历史 token 的 k、v 算一遍；每读一个 k 只做一次点积，算术强度是个位数，远低于 H100 的 295 FLOPs/字节，所以完全受带宽限制。\par}
\item {\fontsize{10}{14.8}\selectfont KV 按每块 B 个 token 切开、散落在池子各处，只有块表知道位置：第 i 个 token 在物理块 table[i ÷ B] 的第 i mod B 个位置，块内是连续的，所以只要 B 不太小，显存带宽仍然能用得很好。\par}
\item {\fontsize{10}{14.8}\selectfont 最自然的切法是每个请求的每个 KV 头交给一个工作单元，批 64 时就是 512 个单元，132 个 SM 都有活干；可批为 1、上下文又长时只有 8 个单元，其余 124 个 SM 全在等。\par}
\item {\fontsize{10}{14.8}\selectfont 解法是沿上下文再切一刀：每 512 个 token 一段，32K 切成 64 段，8 个 KV 头就是 512 个单元；每段各算一组 (m, l, acc)，最后用在线 softmax 的公式换算再相加，结果和不切完全相同。\par}
\item {\fontsize{10}{14.8}\selectfont 这叫 Flash-Decoding，也叫 split-K，vLLM 的 PagedAttention V2 kernel 按 512 切段；批 1、32K 从 661 $\mu$s 降到 40 $\mu$s，批 64、4K 时切不切都约 321 $\mu$s。\par}
\item {\fontsize{10}{14.8}\selectfont 切段有代价：每段多写一组 (m, l, acc)、多一个归约 kernel，还会改浮点结果------段数和合并次序让同一前缀在不同 batch 里算出不同概率，要 batch invariance 就得固定段数和次序。\par}
\end{enumerate}\endgroup
\numline{公式}{$\text{物理块}=\text{table}[\,i\div B\,]$，块内位置 $=i\bmod B$；合并 $l=\sum_s l_s\cdot e^{m_s-M}$，$\text{acc}=\sum_s \text{acc}_s\cdot e^{m_s-M}$}
\numline{数字}{B = 16、每 token 每层 KV 4096 字节；批 1 / 32K：8 $\rightarrow$ 512 个单元，661 $\rightarrow$ 40 $\mu$s；批 1 / 4K：83 $\rightarrow$ 10 $\mu$s；批 64 / 4K：321 / 321 $\mu$s；RMSNorm 切 8 段：24 $\rightarrow$ 32.1 MiB、1 $\rightarrow$ 3 个 kernel、误差 111 $\rightarrow$ 9 ULP}
\fol{切段合并会不会损失精度？}{不会，合并用的是在线 softmax 的精确公式，和一次算完数学等价。但段数变了结果会有几个 ULP 的差别，因为浮点加法不满足结合律：8 个局部和的 40,320 种合并次序能得到 4 个不同结果。}
\fol{PagedAttention 是什么？}{PagedAttention 首先是一种显存管理方式，讲 KV 怎么按块存；这一讲的 kernel 讲分页之后的 KV 怎么读。两者合起来，才是 vLLM 论文里那个 PagedAttention。}
\fol{decode 注意力还能怎么加速？}{把 KV Cache 存成 FP8，要读的字节减半，注意力快将近一倍，读进片上再转回高精度算；MLA 把所有查询头压到一份 KV 上，算术强度高得多，要用 FlashMLA 这类专门 kernel。}
\numline{别答错}{Flash-Decoding 只补查询太少的并行度；段数随 batch 变，结果还会不可复现。}
\phantomsection\addcontentsline{toc}{section}{86\quad CUDA 编程模型}
\lec{86}{CUDA 编程模型}
\askline{CUDA 里线程、块、网格是怎么组织的？}
\askline{矩阵乘法为什么非要走共享内存？}
\askline{占用率是不是越高越快？}
\vspace{3pt}
\noindent{\qlabel{答}}\begingroup\sloppy\begin{enumerate}[leftmargin=1.6em,itemsep=2.5pt,topsep=0pt,parsep=0pt,label={\sffamily\footnotesize\color{accent}\arabic*.}]
\item {\fontsize{10}{14.8}\selectfont CUDA 只写“一个线程做什么”：i = blockIdx.x $\times$ blockDim.x + threadIdx.x 决定它负责哪个元素，越界的线程跳过；启动时用 <<<块数, 每块线程数>>> 指定开多少线程，所有线程执行同一段代码。\par}
\item {\fontsize{10}{14.8}\selectfont 块这一层对应硬件：一个块整个跑在一个 SM 上，块内线程能用共享内存交换数据、用 \_\_syncthreads() 同步；块与块之间不能直接合作、顺序也没保证，所以上一讲的切段要先写回显存再起一个 kernel 合并。\par}
\item {\fontsize{10}{14.8}\selectfont 块里的线程每 32 个编成一个 warp，同时执行同一条指令，于是有两条铁律：同一个 warp 里别让一半走 if 一半走 else，相邻线程要读相邻地址好让硬件合并访存。\par}
\item {\fontsize{10}{14.8}\selectfont 存储分三级：寄存器每线程私有、最快，共享内存每块私有、只有几十到两百多 KB，全局内存就是显存、大而慢；写快 kernel 的核心问题永远是让数据尽量少从显存读、尽量多在寄存器和共享内存里复用。\par}
\item {\fontsize{10}{14.8}\selectfont 矩阵乘法直接写法是每个线程读 A 的一行和 B 的一列，共要读 2K$^3$ 个元素，K = 4096 时约 274.9 GB、82 ms，而计算只要 0.14 ms；分块复用后降到 1/T，T = 64 时只剩 4.3 GB。\par}
\item {\fontsize{10}{14.8}\selectfont 共享内存分 32 个 bank，行跨度 32 时读一列全落在同一个 bank，要排 32 轮；每行末尾补一个字让跨度变成 33 就一轮完成，只多占 3.1\% 空间，也可以改成第 r 行第 c 列放 c $\oplus$ r 列的 swizzle。\par}
\end{enumerate}\endgroup
\numline{公式}{$i=\text{blockIdx}.x\times\text{blockDim}.x+\text{threadIdx}.x$；显存读取量 $\approx2K^{3}/T$ 个元素；$T=t_{\text{搬}}+(n-1)\cdot\max(t_{\text{搬}},\ t_{\text{算}})+t_{\text{算}}$}
\numline{数字}{H100 132 SM；SM 资源 2048 线程槽 / 65,536 寄存器 / 228 KB 共享内存；n = 1000、每块 256 $\rightarrow$ 4 个块 1024 个线程、24 个越界；K = 4096 读 274.9 GB / 82 ms $\rightarrow$ T = 64 时 4.3 GB / 1.3 ms；128$\times$128 块驻留 2 块、占用率 25\%}
\fol{占用率到底有什么用？}{用来提供在途请求隐藏显存延迟。H100 带宽乘 600 ns 约 2.01 MB，平摊到 132 个 SM 每个要 15.2 KB；16 个 warp 给出 32 KiB，是需求的 2.15 倍就够。}
\fol{块做大了不是读得更少吗？}{读得少但住得少。Qwen3-8B 投影 128$\times$128 只要 1560 MiB，64$\times$64 要 3096 MiB；可大块要 80 KiB 局部存储、驻留块变少，只拿到四成带宽时时间反成 1.26 倍。}
\fol{双缓冲为什么只快了 13\%？}{因为搬比算慢 4.6 倍：每块搬 1.29 $\mu$s、算 0.28 $\mu$s，重叠之后总时间仍由搬移决定，4 块从 6.28 $\mu$s 只降到 5.44 $\mu$s，计算单元只有两成时间在干活，再加第三个缓冲槽也没用。}
\numline{别答错}{CUDA 块是线程组，PagedAttention 块是 KV 段；占用率高不等于快，够用就行。}
\phantomsection\addcontentsline{toc}{section}{87\quad 矩阵乘法 Kernel 与 Tensor Core}
\lec{87}{矩阵乘法 Kernel 与 Tensor Core}
\askline{一个矩阵乘法怎么才能跑到峰值？}
\askline{1536$^2$ 只多 19\% 计算，为什么慢一倍？}
\askline{decode 的矩阵乘为什么只用到 1\% 算力？}
\vspace{3pt}
\noindent{\qlabel{答}}\begingroup\sloppy\begin{enumerate}[leftmargin=1.6em,itemsep=2.5pt,topsep=0pt,parsep=0pt,label={\sffamily\footnotesize\color{accent}\arabic*.}]
\item {\fontsize{10}{14.8}\selectfont 矩阵乘 kernel 有三层分块：线程块负责 C 的 128$\times$128，沿 K 把 A 的 128$\times$32 和 B 的 32$\times$128 搬进共享内存；warp 在寄存器里累加 64$\times$64 的小块；最后交给 Tensor Core 指令。\par}
\item {\fontsize{10}{14.8}\selectfont 指令形状决定下限：A100 的 mma 是一个 warp 一次算 16$\times$8；H100 的 wgmma 是四个 warp 组成 warpgroup、一次算 64$\times$N，A、B 直接从共享内存异步读，配 TMA 搬运，搬和算完全重叠。\par}
\item {\fontsize{10}{14.8}\selectfont 输出的块按波分给 132 个 SM，两种浪费：补零只在边角，M = 1000、块 128 时浪费 2.3\%；波次浪费大得多，块数刚过 132 的整数倍就要多占一波，最后一波大部分 SM 空等，时间直接跳一级。\par}
\item {\fontsize{10}{14.8}\selectfont decode 反过来：16 个请求乘 4096$\times$1024 的权重，算术强度约等于 M = 16，远在拐点 295 左边，理想时间就是读一遍 8.5 MB 权重约 2.6 $\mu$s；可 64$\times$64 的块只有 16 个，大多数 SM 闲着。\par}
\item {\fontsize{10}{14.8}\selectfont 解法和 Flash-Decoding 同源：沿 K 切成 8 段各算部分和再归约，块数变成 128、几乎占满所有 SM，时间从 10.5 $\mu$s 降到 2.9 $\mu$s，接近只读一遍权重；代价是部分和的读写和一次归约。\par}
\item {\fontsize{10}{14.8}\selectfont 库的对策是四件事：cuBLAS 按 M、N、K 用启发式或测好的表选配置，CUTLASS 把配置做成模板参数，Stream-K 把 K 段均分给 132 个 SM 消掉最后一波空闲，推理引擎把 decode 的 batch 补齐到固定档位。\par}
\end{enumerate}\endgroup
\numline{公式}{$\text{波次利用率}=\frac{\text{块数}}{\text{波数}\times\text{SM}}$；算术强度 $\approx\frac{2MNK}{2KN}\approx M$；split-K 的块数 $=\lceil M/M_t\rceil\cdot\lceil N/N_t\rceil\cdot s$}
\numline{数字}{H100 132 SM、BF16 峰值 989 TFLOPS、拐点 295 FLOP/字节：1408$^2$ $\rightarrow$ 121 块 17.9 $\mu$s / 907 TFLOPS，1536$^2$ $\rightarrow$ 144 块 35.8 $\mu$s / 539 TFLOPS；16$\times$1024$\times$4096 $\rightarrow$ 16 块 10.5 $\mu$s / 13 TFLOPS（峰值 1.3\%），split-K = 8 $\rightarrow$ 128 块 2.9 $\mu$s}
\fol{上 Tensor Core 就到峰值了？}{不等于。块形状、波次、流水、数据搬移任何一项没做好，Tensor Core 都在等；同一个 kernel，矩阵从 1408 长到 1536 时间翻倍，再长到 2048 又和 1536 一样。}
\fol{decode 的瓶颈到底是算还是读？}{是读。M = 16 时 wgmma 一次输出 64 行、四分之三的乘加白算，但这不要紧：它本来就受带宽限制、Tensor Core 在等数据；要优化的是让足够多 SM 一起读，以及少读。}
\fol{模型形状上能提前做什么？}{把隐藏维度、FFN 宽度、词表大小取成 64 或 128 的倍数，避免边角补零；Llama 的词表 128,256 = 128 $\times$ 1002 就是一例。}
\numline{别答错}{上 Tensor Core 不等于到峰值；性能随形状呈阶梯，1536 比 1408 慢一倍。}
\phantomsection\addcontentsline{toc}{section}{88\quad Triton 入门}
\lec{88}{Triton 入门}
\askline{Triton 和写 CUDA 有什么不一样？}
\askline{同一个 kernel，块大小怎么挑？}
\askline{推理每一步 token 数都在变，怎么办？}
\vspace{3pt}
\noindent{\qlabel{答}}\begingroup\sloppy\begin{enumerate}[leftmargin=1.6em,itemsep=2.5pt,topsep=0pt,parsep=0pt,label={\sffamily\footnotesize\color{accent}\arabic*.}]
\item {\fontsize{10}{14.8}\selectfont CUDA 的视角是一个线程：我是第几个线程、算哪个元素；Triton 把视角抬到一个程序：我是第几个程序、处理哪一块数据，块里的线程怎么分工、warp 怎么切、要不要共享内存，都交给编译器。\par}
\item {\fontsize{10}{14.8}\selectfont 写 kernel 就三步：program\_id 定位自己负责哪一块，arange 生成整块下标，mask 盖住最后一块的越界；然后 load 整块、整块算、store 整块，合并访存和 Tensor Core 都由编译器安排。\par}
\item {\fontsize{10}{14.8}\selectfont 融合的 SwiGLU 十几行就够：读进 g 和 u，y = g * sigmoid(g) * u，中间结果只在片上，只写最终结果，读写从 8 遍降到 3 遍；RMSNorm 则是一个程序处理整行，一行 tl.sum 就做完了归约。\par}
\item {\fontsize{10}{14.8}\selectfont 矩阵乘法几乎是分块那句话的直译：二维网格定位 C 的一块，沿 K 一块一块 load A、B，tl.dot 的结果累加进 FP32 累加器，编译器把它变成 Tensor Core 指令并自动安排共享内存。\par}
\item {\fontsize{10}{14.8}\selectfont 块大小用 @triton.autotune 试，候选不只块形状，还有 num\_warps 和 num\_stages（提前搬几块让搬运和计算重叠）；挑选顺序是先按共享内存、寄存器、对齐排除放不下的，再估算，最后实测。\par}
\item {\fontsize{10}{14.8}\selectfont 单个形状上的加速不等于整体加速：新实现处理 A 只要 5 $\mu$s、处理 B 却要 20 $\mu$s，A 的调用占比得超过 2/3 才比原来的 10 $\mu$s 划算；搜索本身也要回本，省 5 $\mu$s 的搜索跑 10 分钟要 1.2 亿次调用。\par}
\end{enumerate}\endgroup
\numline{公式}{$\text{平均时间}=5p+20(1-p)\ \mu\text{s}$，优于 $10\ \mu\text{s}$ 需 $p>2/3$；总时间 $=\text{准备}+\text{组数}\times\text{每组时间}$}
\numline{数字}{n = 1000、BLOCK = 256 $\rightarrow$ 4 个程序、末块 232 个有效；转置 16$\times$16 的块 7.9 $\mu$s 对 32$\times$32 的 2.9 $\mu$s；分桶 5632 $\rightarrow$ 8192 行、多算 45\%；通用 / 分桶 / 特化每组 16.9 / 12.3 / 6.8 ms，前 65 组通用、第 66 组起分桶、第 91 组起特化；回本要 1.2 亿次调用}
\fol{Triton 的程序是线程吗？}{不是。一个程序处理一整块数据，背后是一个线程块里的几十到几百个线程，num\_warps 决定用几个 warp；它和 NVIDIA 那个同名的 Triton 推理服务器也毫无关系，只是重名。}
\fol{动态形状怎么选策略？}{三种：通用按当前形状算下标和边界，分桶把输入补齐到少数代表形状，特化给每个常见形状单独编译。越利用已知形状越快，但要编译和缓存的程序越多；跑 70 组时反而是分桶最省，1.26 s。}
\fol{什么时候还该手写 CUDA？}{矩阵乘法和注意力这类 kernel 才值得：CUDA 配 CUTLASS 能精细控制 TMA、wgmma，FlashAttention-3 就是这样写的；新模型新结构先用 Triton 快速实现。}
\numline{别答错}{程序不是线程；块形状是编译时常数、一般取 2 的幂；单个形状快不等于整体快，要看调用比例。}
\phantomsection\addcontentsline{toc}{section}{89\quad CUDA Graph}
\lec{89}{CUDA Graph}
\askline{decode 为什么要用 CUDA Graph？}
\askline{什么情况下 CPU 会成为瓶颈？}
\askline{CUDA Graph 有什么代价？}
\vspace{3pt}
\noindent{\qlabel{答}}\begingroup\sloppy\begin{enumerate}[leftmargin=1.6em,itemsep=2.5pt,topsep=0pt,parsep=0pt,label={\sffamily\footnotesize\color{accent}\arabic*.}]
\item {\fontsize{10}{14.8}\selectfont eager 模式下每个 kernel 都要过 Python、框架和驱动，CPU 花几微秒；当 GPU 执行一个 kernel 比 CPU 发射一个还短时，GPU 算完只能等下一个，一步的时间就由 CPU 决定，GPU 算得再快也没用。\par}
\item {\fontsize{10}{14.8}\selectfont 8 卡张量并行的 Llama-3-8B 一步有 545 个 kernel，GPU 真正在算只要 1.93 ms，CPU 逐个发射要 545 $\times$ 8 $\mu$s $\approx$ 4.36 ms，eager 一步 4.40 ms，GPU 只忙 44\%。\par}
\item {\fontsize{10}{14.8}\selectfont CUDA Graph 分捕获和回放两步：捕获时 CPU 照常执行一遍，但 kernel 不真正执行，只记下有哪些 kernel、什么顺序、什么参数、读写哪些显存地址；回放时 CPU 一次调用把整张图交给 GPU，只花几十微秒。\par}
\item {\fontsize{10}{14.8}\selectfont 同一配置改用图回放，一步从 4.40 ms 降到 1.95 ms、GPU 忙 99\%；1B 模型从 2.08 ms 降到 1.11 ms。规律是模型越小、并行的卡越多、批越小，每个 kernel 越短，CPU 越容易成为瓶颈。\par}
\item {\fontsize{10}{14.8}\selectfont 代价是图和形状、地址都固定：一张图对应一个批大小，实际批大小向上补齐（13 就补到 16 的那张图），输入要先拷进捕获时的静态缓冲，不能有 CPU 端的分支判断。\par}
\item {\fontsize{10}{14.8}\selectfont 它也不是万能：捕获几十张图要占显存、启动多花几十秒，prefill 的 token 数每次都变、kernel 又大，一般不用图；图只填 kernel 之间的空隙，不让 kernel 本身变快。\par}
\end{enumerate}\endgroup
\numline{公式}{第 $i$ 个 kernel 开始 $=\max(\text{它被发射完},\ \text{第 }i-1\text{ 个做完})$；图重放的临界输入 $2X\div B=15\ \mu\text{s}\to X\approx12.8\ \text{MiB}$}
\numline{数字}{481 个 kernel/5.18 ms，eager 5.19 ms（单卡 8B）；545 个/1.93 ms，eager 4.40 ms vs 图 1.95 ms（8 卡）；241 个/1.09 ms，eager 2.08 vs 1.11 ms（1B）；回放省 15 $\mu$s；临界输入 12.8 MiB；vLLM v0.6.0 吞吐 2.7 倍}
\fol{GPU 利用率 100\% 就没瓶颈吗？}{不代表。nvidia-smi 的利用率只说明有 kernel 在跑，几个 SM 的小 kernel 也算忙；看有没有空等，要用 profiler 把 CPU 和 GPU 两条时间线画在一起。}
\fol{大模型大批量也需要图吗？}{几乎不需要。单卡 8B 的 481 个 kernel，GPU 要 5.18 ms，CPU 发射 3.85 ms 跑在 GPU 前面，eager 5.19 ms 和图 5.20 ms 一样。}
\fol{怎么验证图真的省了时间？}{用 nsys 看一步里 GPU 的空隙和 kernel 之间的间隔，再对照 enforce\_eager 开关的吞吐。Python 侧直接计时量到的只是发射时间，要先 synchronize。}
\numline{别答错}{别把图当加速：它只填 kernel 间隙；输入超 12.8 MiB 复制开销吃掉收益。}
\chapter{模型压缩与加速}
{\footnotesize\sffamily\color{grayt}第 90--98 讲}\par\vspace{-4pt}
\phantomsection\addcontentsline{toc}{section}{90\quad 量化基础}
\lec{90}{量化基础}
\askline{量化误差怎么算，上限是多少？}
\askline{离群值为什么会毁掉一组数？}
\askline{W4A16 为什么只加速 decode？}
\vspace{3pt}
\noindent{\qlabel{答}}\begingroup\sloppy\begin{enumerate}[leftmargin=1.6em,itemsep=2.5pt,topsep=0pt,parsep=0pt,label={\sffamily\footnotesize\color{accent}\arabic*.}]
\item {\fontsize{10}{14.8}\selectfont 量化就是选一个 scale 把这组实数缩到整数格点上：q = round(x ÷ scale)，存下这些整数和 scale，用时再乘回去 x̂ = q $\times$ scale；因为取整最多差半个格点，误差不超过 scale ÷ 2。\par}
\item {\fontsize{10}{14.8}\selectfont 对称量化让绝对值最大的数对应到最大的整数 2\textasciicircum{}(b-1) - 1，8 比特是 127、4 比特是 7。示例的 8 个权重 INT8 的 scale 是 0.01819，最大误差 0.0077，在 scale ÷ 2 = 0.0091 以内。\par}
\item {\fontsize{10}{14.8}\selectfont 位宽越低误差越陡：同一组数 INT4 平均误差 0.0725、INT3 是 0.1362，INT2 只剩 -1、0、1 三个值，8 个数里 5 个变成 0，平均误差 0.59，基本不能用了。\par}
\item {\fontsize{10}{14.8}\selectfont 非对称量化多存一个零点 z，把 [min, max] 映射到 [0, 2\textasciicircum{}b - 1]，全是正数的分布也能用满格点；同样 8 个数 INT4 的误差从 0.0725 降到 0.064。\par}
\item {\fontsize{10}{14.8}\selectfont 软肋是 scale 由最大的数决定：把 2.31 换成 12，INT4 的 scale 变成 1.714，其余 7 个数的平均误差从 0.07 涨到 0.49；每 4 个一组各算各的 scale 后回到 0.17。\par}
\item {\fontsize{10}{14.8}\selectfont 收益按字节算：Llama-3-8B 的权重 BF16 是 16.06 GB、decode 读一遍 4.79 ms；INT4 每组 128 个配一个 16 位 scale，平均 4.125 比特，4.14 GB、1.24 ms。\par}
\end{enumerate}\endgroup
\numline{公式}{$\text{scale}=\frac{\max|x|}{2^{b-1}-1}$，$q=\operatorname{round}(x\div\text{scale})$，$\hat{x}=q\times\text{scale}$，误差 $\le\text{scale}\div2$；非对称 $q=\operatorname{round}(x\div\text{scale})+z$，$\hat{x}=(q-z)\times\text{scale}$}
\numline{数字}{Llama-3-8B：16.06 GB/4.79 ms（BF16）、8.03 GB/2.40 ms（INT8）、4.14 GB/1.24 ms（INT4）；scale 0.01819、最大误差 0.0077；INT4 0.0725、INT3 0.1362、INT2 0.59；离群值 0.49 $\rightarrow$ 分组 0.17；4.125 比特；70B 140 GB $\rightarrow$ 35 GB}
\fol{存成 INT4 就是用 INT4 算吗？}{不是。W4A16 的权重以 INT4 存在显存里，kernel 读进片上先反量化成 BF16，再和 BF16 的激活相乘；省的是读取的字节，计算还是 BF16，所以 prefill 几乎没有收益。}
\fol{为什么每组 128 个是常见折中？}{scale 的开销是 16 ÷ 组大小。128 时每个权重多 0.125 比特、实际 4.125 比特；组再小精度更好但开销涨，组太大又压不住离群值，128 是常用的折中。}
\fol{量化后怎么判断能不能上线？}{不看单个数的误差，看输出变了多少和任务指标。权重误差大不一定差，有的通道的误差会被放大；要在自己关心的任务上评测，不能只看困惑度。}
\numline{别答错}{存成 INT4 不等于用 INT4 算：W4A16 读进片上反量化成 BF16，省的是字节不是计算。}
\phantomsection\addcontentsline{toc}{section}{91\quad 权重量化：GPTQ 与 AWQ}
\lec{91}{权重量化：GPTQ 与 AWQ}
\askline{RTN 掉点，GPTQ 和 AWQ 差在哪？}
\askline{AWQ 的等价缩放为什么有效？}
\askline{GPTQ 的补偿矩阵从哪来？}
\vspace{3pt}
\noindent{\qlabel{答}}\begingroup\sloppy\begin{enumerate}[leftmargin=1.6em,itemsep=2.5pt,topsep=0pt,parsep=0pt,label={\sffamily\footnotesize\color{accent}\arabic*.}]
\item {\fontsize{10}{14.8}\selectfont 两者都不换存储格式，量化完还是 INT4 加 scale，推理时的计算和 RTN 的模型一样；它们只是更聪明地决定每个权重取整成哪个整数，目标是最小化这一层的输出误差。\par}
\item {\fontsize{10}{14.8}\selectfont 输出误差是 $\Sigma$ⱼ xⱼ(Wⱼ - Ŵⱼ)：第 j 行权重的误差会被第 j 个输入通道的激活放大。大模型的激活里恰恰有少数通道长期比别的大几十倍，这几行的误差就成了输出误差的大头。\par}
\item {\fontsize{10}{14.8}\selectfont 所以两者都要先拿一小批校准数据（几百条普通文本）跑一遍模型、记录每层输入 X，再按 X 量化这一层的权重。人造的 64$\times$64 层、4 个离群通道时，4 比特 RTN 的输出误差是 10.56\%。\par}
\item {\fontsize{10}{14.8}\selectfont AWQ 用 x$\cdot$W = (x/s)$\cdot$(s$\cdot$W) 把重要的行放大后再量化：同一列的 scale 由最大数决定，放大通常不改变它，取整误差的绝对值也基本不变，反量化后却要除以 s，这一行的误差就小了。\par}
\item {\fontsize{10}{14.8}\selectfont sⱼ = (平均|xⱼ|)\textasciicircum{}$\alpha$，在 0 到 1 之间搜 $\alpha$ 选输出误差最小的，例子里 $\alpha$ = 0.2 最好，输出误差从 10.56\% 降到 6.34\%；s 可以离线合并进前一层 RMSNorm，推理时没有额外开销。\par}
\item {\fontsize{10}{14.8}\selectfont GPTQ 换了个角度：逐行量化，把这一行的取整误差按 H = XᵀX 的逆摊给还没量化的行去补偿，并且先量化 H 对角线最大、也就是激活最大的行。例子里它输出误差 5.30\%，但权重平均误差 0.016 反而最大。\par}
\end{enumerate}\endgroup
\numline{公式}{$x\cdot W=(x/s)\cdot(s\cdot W)$；$s_j=(\text{平均}|x_j|)^{\alpha}$；补偿矩阵来自 $H=X^{\top}X$ 的逆；输出误差 $=\frac{\lVert X\cdot W-X\cdot\hat{W}\rVert}{\lVert X\cdot W\rVert}$}
\numline{数字}{校准 256 token；4 比特 RTN 10.56\%、AWQ 6.34\%、GPTQ 5.30\%；8 比特 0.58\%/0.36\%/0.30\%；3 比特 23.99\%/14.71\%/13.07\%；AWQ 离群行 0.0046 $\rightarrow$ 0.0024；GPTQ 权重误差 0.016；175B 量化几个 GPU 小时}
\fol{AWQ 放大重要行没有代价吗？}{有。放大太猛会把这一列的 scale 抬起来，其他行的取整误差跟着变大；$\alpha$ 在 0 到 1 之间搜，找的就是这个平衡点，例子里 0.2 最优。}
\fol{GPTQ 需要重新训练吗？}{不需要，只做前向、记录激活，不做反向传播，几分钟到几小时就能完成。校准数据最好和真实场景相近，用纯英文校准再去跑中文或代码，效果会差一些。}
\fol{4 比特模型线上怎么跑？}{vLLM、SGLang 直接加载 GPTQ/AWQ 的权重，用 Marlin 这类 W4A16 kernel，读进片上先反量化再算 BF16；小批量 decode 能快两三倍，prefill 不变。}
\numline{别答错}{GPTQ 的权重平均误差反而最大，因为它有意把权重调离原值去抵消输出误差------权重误差不是优化目标。}
\phantomsection\addcontentsline{toc}{section}{92\quad 激活量化与 SmoothQuant}
\lec{92}{激活量化与 SmoothQuant}
\askline{为什么激活量化这么难？}
\askline{为什么激活不能逐通道量化？}
\askline{SmoothQuant 怎么平衡两边？}
\vspace{3pt}
\noindent{\qlabel{答}}\begingroup\sloppy\begin{enumerate}[leftmargin=1.6em,itemsep=2.5pt,topsep=0pt,parsep=0pt,label={\sffamily\footnotesize\color{accent}\arabic*.}]
\item {\fontsize{10}{14.8}\selectfont 要让矩阵乘法本身用低精度算，激活也得量化，这就是 W8A8。难点是模型超过六七十亿参数后，每一层的输入激活里会出现少数几个固定的离群通道，值比别的大几十上百倍，例子里普通通道最大约 3.0、离群通道约 178。\par}
\item {\fontsize{10}{14.8}\selectfont 整数矩阵乘要求 scale 在求和维上不变：Y\_ik $\approx$ a$\cdot$b$\cdot$$\Sigma$ⱼ qx$\cdot$qw，a 和 b 必须能提到求和号外面，所以激活只能每 token 一个 scale、权重只能每输出通道一个。\par}
\item {\fontsize{10}{14.8}\selectfont 于是每个 token 的 scale 都被离群通道定死，普通通道的值只落在少数几个格点上：直接 W8A8 的整体激活误差只有 1.52\%，普通通道的激活误差却高达 17.40\%，输出误差 1.70\%。\par}
\item {\fontsize{10}{14.8}\selectfont SmoothQuant 用同一个等价变换 X$\cdot$W = (X/s)$\cdot$(s$\cdot$W)，sⱼ = max|Xⱼ|\textasciicircum{}$\alpha$ ÷ max|Wⱼ|\textasciicircum{}(1-$\alpha$)；离群通道的 s 取得大，激活那一列被压小，对应的权重行被放大，难度就匀过去了。\par}
\item {\fontsize{10}{14.8}\selectfont $\alpha$ 是迁移强度：$\alpha$ = 0 时激活最大值÷中位数 59.8、权重 1.0，输出误差 1.75\%；$\alpha$ = 0.5 时两边都是 7.7，误差最小 0.50\%；$\alpha$ = 1 又倒过来，权重 59.8、误差 1.88\%。\par}
\item {\fontsize{10}{14.8}\selectfont s 在离线时用校准数据算好、合并进前一层（比如 RMSNorm 的权重），推理时没有额外开销，矩阵乘照常走 INT8；今天 W8A8 大多换成 FP8，但离群通道的问题并没有消失。\par}
\end{enumerate}\endgroup
\numline{公式}{$X\cdot W=(X/s)\cdot(s\cdot W)$；$s_j=\frac{\max|X_j|^{\alpha}}{\max|W_j|^{1-\alpha}}$；$\alpha=0.5$ 时两边都是 $\sqrt{\max|X_j|\cdot\max|W_j|}$}
\numline{数字}{普通通道 max $\approx$ 3.0、离群通道 $\approx$ 178；直接 W8A8 整体 1.52\%、普通通道 17.40\%、输出 1.70\%；$\alpha$ = 0/0.5/1 $\rightarrow$ 1.75\%/0.50\%/1.88\%，难度 59.8:1.0 $\rightarrow$ 7.7:7.7 $\rightarrow$ 1.0:59.8；逐通道对照 0.99\%}
\fol{离群通道能不能直接截断？}{不能。它们不是噪声，去掉模型效果会大幅下降，只能安置：平滑把难度匀给权重，或者用正交旋转把能量摊到所有通道。LLM.int8() 则是把它们拆出来用 FP16 单算。}
\fol{$\alpha$ 必须取 0.5 吗？}{不一定，用校准数据扫一遍选输出误差最小的，论文里大部分模型在 0.5 左右。$\alpha$ = 0 激活太难、$\alpha$ = 1 权重太难，两头的误差都更大。}
\fol{激活的 scale 静态还是动态？}{现在的实现多用每个 token 动态现算，更稳，开销能融合进前一个 kernel；静态用校准数据提前算好最快，但遇到和校准数据差别大的输入可能被截断。}
\numline{别答错}{离群通道不是噪声，去掉模型效果大降，只能安置不能截断；逐通道量化激活硬件做不了。}
\phantomsection\addcontentsline{toc}{section}{93\quad FP8 与 FP4}
\lec{93}{FP8 与 FP4}
\askline{FP8 两种格式怎么选？}
\askline{为什么 FP4 必须分组缩放？}
\askline{MXFP4 和 NVFP4 差在哪？}
\vspace{3pt}
\noindent{\qlabel{答}}\begingroup\sloppy\begin{enumerate}[leftmargin=1.6em,itemsep=2.5pt,topsep=0pt,parsep=0pt,label={\sffamily\footnotesize\color{accent}\arabic*.}]
\item {\fontsize{10}{14.8}\selectfont 浮点数由符号、指数、尾数拼成：(-1)\textasciicircum{}S $\times$ (1 + M/2\textasciicircum{}m) $\times$ 2\textasciicircum{}(E-bias)。整数格点等间隔，浮点格点在对数坐标上均匀，0 附近密、越往外越疏，正合神经网络里大多是小数、偶有几个大数的分布。\par}
\item {\fontsize{10}{14.8}\selectfont FP8 有两种：E4M3 是 4 位指数 3 位尾数，最大 448，1 到 2 之间有 8 个格点；E5M2 范围大得多、最大 57344，但每段只有 4 个格点、精度更粗。推理的权重和激活一般用 E4M3。\par}
\item {\fontsize{10}{14.8}\selectfont FP4 常用的 E2M1 只有 0、0.5、1、1.5、2、3、4、6 八个非负值。4096 个重尾数整体用一个 scale 时，INT4 的误差是 51.25\%、FP4 是 34.02\%，都还不能用。\par}
\item {\fontsize{10}{14.8}\selectfont 所以 4 位离不开分组缩放：MXFP4 每 32 个数一组，scale 用 8 位 E8M0、只能是 2 的幂，好处是乘 scale 只改指数、硬件简单，平均 4.25 比特，误差 15.91\%。\par}
\item {\fontsize{10}{14.8}\selectfont NVFP4 每 16 个数一组，scale 用 FP8 E4M3 存、能取非 2 的幂，平均 4.5 比特，误差降到 9.25\%；这些 scale 由 Blackwell 的 Tensor Core 直接乘进矩阵乘，组再细也不拖慢计算。\par}
\item {\fontsize{10}{14.8}\selectfont 速度上 H100 的 FP8 稠密算力约 1979 TFLOPS，是 BF16 989 的两倍；显存上 8B 用 FP8 约 8 GB、NVFP4 约 4.5 GB。硬件不支持就没有加速，A100 上 FP8 权重只能转回 BF16 再算。\par}
\end{enumerate}\endgroup
\numline{公式}{$(-1)^{S}\times(1+M/2^{m})\times2^{E-\text{bias}}$；次正规数 $(-1)^{S}\times(M/2^{m})\times2^{1-\text{bias}}$；平均位宽 $=4+8\div\text{组大小}$}
\numline{数字}{E4M3：127 个非负值、最大 448、最小 2$^-$⁹；E5M2：124 个、最大 57344；E2M1：8 个非负值；整体 scale INT8 3.22\% vs E4M3 2.44\%、INT4 51.25\% vs FP4 34.02\%；分组 INT4 15.09\%、MXFP4 15.91\%、NVFP4 9.25\%；H100 FP8 1979 vs BF16 989 TFLOPS；8B：FP8 约 8 GB、NVFP4 约 4.5 GB}
\fol{E4M3 和 E5M2 能混用吗？}{不能，位含义不同，存储和计算时要明确用哪一种。E4M3 精度好，推理的权重和激活都用它；E5M2 范围大，主要给训练里数值跨度大的梯度。}
\fol{FP8 在 A100 上还有意义吗？}{有，但只省显存和带宽。A100 没有 FP8 计算单元，权重读进来要先转回 BF16 再乘；8B 的权重照样从 16 GB 降到 8 GB，算力不变。}
\fol{分组为什么不能无限小？}{scale 本身也要存、也要读。MXFP4 每 32 个一个 8 位 scale 是 4.25 比特，NVFP4 每 16 个一个是 4.5 比特，组再小平均位宽就往 5 比特以上走。}
\numline{别答错}{FP8 有两种，说「FP8 推理」一般指 E4M3；FP4 离开分组缩放几乎不能用。}
\phantomsection\addcontentsline{toc}{section}{94\quad KV Cache 量化}
\lec{94}{KV Cache 量化}
\askline{KV Cache 量化能换来什么？}
\askline{为什么字节减半 TPOT 不变？}
\askline{K 和 V 为什么要分开量化？}
\vspace{3pt}
\noindent{\qlabel{答}}\begingroup\sloppy\begin{enumerate}[leftmargin=1.6em,itemsep=2.5pt,topsep=0pt,parsep=0pt,label={\sffamily\footnotesize\color{accent}\arabic*.}]
\item {\fontsize{10}{14.8}\selectfont 先算账：80 GB 的卡跑 8B，权重占 16 GB，KV 池留下约 49.3 GiB，每个 token 的 KV 是 128 KiB，4096 token 的请求只能放 98 个。这时 KV 是权重的三倍，比权重还大。\par}
\item {\fontsize{10}{14.8}\selectfont KV 的大小就是 token 数乘每 token 的字节数，池子固定时字节减半能放的 token 就翻倍：FP8 每 token 64 KiB，能放 197 个请求；INT4 每 128 个数一个 scale，33 KiB、382 个。\par}
\item {\fontsize{10}{14.8}\selectfont 满载时 decode 一步主要花在读 KV 上，98 个请求的 KV 就有 49 GiB；字节减半，同样的时间能读两倍请求，所以 TPOT 几乎不变（20.5 对 20.6 ms），吞吐从 4781 翻倍到 9573 token/s。\par}
\item {\fontsize{10}{14.8}\selectfont 精度上 K 和 V 不一样：K 的少数几维（固定通道，一般认为和 RoPE 转得慢的维度有关）在所有 token 上都特别大，V 比较均匀，所以 K 适合按通道分 scale、V 适合按 token 分。\par}
\item {\fontsize{10}{14.8}\selectfont 这个差别在低位宽下很致命：INT4 时 K、V 都按 token 量化，注意力输出误差 64.71\%，K 改成按通道后降到 18.84\%；FP8 是 3.87\%，INT8 是 5.00\%。\par}
\item {\fontsize{10}{14.8}\selectfont 注意力 kernel 把 FP8 或 INT4 的 KV 读进片上先反量化成 BF16 再算，计算方式不变，所以 FP8 的 KV 并不需要 FP8 计算单元；今天 vLLM 的 kv\_cache\_dtype="fp8" 已经是常规选项。\par}
\end{enumerate}\endgroup
\numline{公式}{$\text{可放 token 数}=\text{KV 池}\div\text{每 token 字节}$；满载一步时间 $\approx\text{KV 总字节}\div\text{带宽}$，故字节减半时 TPOT 不变、吞吐翻倍}
\numline{数字}{每 token 128/64/33/17 KiB；池子 403,865/807,731/1,566,508 个 token；98/197/382/742 个请求；TPOT 20.5/20.6 ms；吞吐 4,781/9,573/18,564/36,042 token/s；注意力误差 FP8 3.87\%、INT8 5.00\%、INT4 64.71\% $\rightarrow$ 18.84\%}
\fol{单个请求的 TPOT 会变快吗？}{不会。并发低时 decode 主要在读权重，KV 小一半对 TPOT 几乎没影响；收益全在同样的显存能服务更多请求、支持更长的上下文。}
\fol{INT4 的 KV 能上线吗？}{18.84\% 的输出误差还太大。要更细的分组、最近一小段 token 保持高精度，或者对 K 先做旋转把离群值摊平，还得在长上下文任务上认真评测。}
\fol{KV 量化和权重量化冲突吗？}{不冲突，两者独立。可以只量化 KV、权重继续 BF16，也可以 FP8 权重配 FP8 KV；还能和 GQA、MLA、分页、前缀缓存叠加着用。}
\numline{别答错}{KV 量化不让单个请求变快：并发低时 decode 在读权重，收益在更多请求和更长上下文。}
\phantomsection\addcontentsline{toc}{section}{95\quad KV Cache 压缩与驱逐}
\lec{95}{KV Cache 压缩与驱逐}
\askline{长上下文 KV 太大，除了量化还能怎样？}
\askline{为什么只留最近的 token 不行？}
\askline{H2O 和 SnapKV 怎么挑？}
\vspace{3pt}
\noindent{\qlabel{答}}\begingroup\sloppy\begin{enumerate}[leftmargin=1.6em,itemsep=2.5pt,topsep=0pt,parsep=0pt,label={\sffamily\footnotesize\color{accent}\arabic*.}]
\item {\fontsize{10}{14.8}\selectfont 另一条路是不存所有的 token：一条 128K 上下文的请求 BF16 的 KV 是 16 GiB、FP8 也有 8 GiB，只留 4096 个 token 就降到 0.5 GiB、是原来的 3.1\%，缓存不再随上下文增长。\par}
\item {\fontsize{10}{14.8}\selectfont 注意力高度集中：Qwen3-0.6B 的实验里，尾巴的注意力有 63.5\% 落在上下文前 4 个 token 上，而它们只占 448 个位置的 0.9\%；最近 64 个拿到 15.2\%，其余 380 个位置分剩下的约五分之一。\par}
\item {\fontsize{10}{14.8}\selectfont 那 4 个是汇聚点（attention sink）：softmax 要求每一行的权重加起来是 1，当前 token 没什么可看时就把用不掉的注意力倒给开头几个，它们的内容不重要，位置才重要。\par}
\item {\fontsize{10}{14.8}\selectfont 所以只留最近 B 个是灾难：B = 32 损失多 4.49，B = 256（留了一半多）还多 2.81，就是丢了开头。保留开头 4 个再加最近的，B = 32 只多 0.54、B = 128 只多 0.13。\par}
\item {\fontsize{10}{14.8}\selectfont H2O 按谁收到的注意力多来留，一半预算给最近的、一半给分数最高的，B = 256 只多 0.013；SnapKV 只用最后 32 个查询的注意力打分，还会把选中位置的邻居带上保住完整的词，B = 32 时最省（+0.50）。\par}
\item {\fontsize{10}{14.8}\selectfont 落地有冲突：PagedAttention 按 16 个 token 一块管理、所有头一起分配，按头按 token 驱逐会在块里留下洞，块表和 kernel 都要改，所以主流引擎默认不做；SnapKV 更适合长文档读完再答的场景。\par}
\end{enumerate}\endgroup
\numline{公式}{损失增量换算困惑度：$e^{\Delta}$ 倍，$\Delta=0.539\to1.71\times$、$0.134\to1.14\times$、$0.013\to1.01\times$}
\numline{数字}{128K 请求 16 GiB，只留 4096 个 0.5 GiB（3.1\%）；前 4 个 token 收 63.5\% 注意力、只占 0.9\%；B = 32 四种策略 +4.49/+0.54/+0.57/+0.50；B = 256 +2.81/+0.10/+0.01/+0.03；不驱逐损失 3.316；Δ0.539 $\rightarrow$ 困惑度 1.71$\times$}
\fol{随机丢为什么有时也很准？}{因为随机抽的预算里恰好包含了第 0 个 token：B = 32、64、256 抽中它，损失就小；B = 128 那次没抽中，损失多 4.72。汇聚点比其他位置加起来都重要。}
\fol{驱逐能省 prefill 吗？}{不能。prefill 时所有 token 都要参与注意力，SnapKV 正是用 prefill 的注意力来选；它省的是之后 decode 的显存和读 KV 的时间。}
\fol{为什么线上引擎默认不开？}{和分页 KV 的块管理冲突，还要改块表和 kernel；更常见的是模型自带滑动窗口（Mistral 7B、Gemma 2 隔层的局部注意力），引擎专门为它管理。}
\numline{别答错}{驱逐扔掉的 token 再也找不回来；损失好看不代表长文问答好，丢了第 0 个就多 4.72。}
\phantomsection\addcontentsline{toc}{section}{96\quad 投机解码}
\lec{96}{投机解码}
\askline{投机解码为什么不掉质量？}
\askline{加速比怎么估算才准？}
\askline{什么时候不该用投机解码？}
\vspace{3pt}
\noindent{\qlabel{答}}\begingroup\sloppy\begin{enumerate}[leftmargin=1.6em,itemsep=2.5pt,topsep=0pt,parsep=0pt,label={\sffamily\footnotesize\color{accent}\arabic*.}]
\item {\fontsize{10}{14.8}\selectfont 前提是小批量 decode 卡在读权重：Llama-3-8B 一次前向处理 1 个、5 个、32 个 token 都约 4.83 ms，因为读一遍权重的钱是固定的，多处理几个只是多用一点本来闲着的算力。\par}
\item {\fontsize{10}{14.8}\selectfont 每轮先让小模型自回归地猜 k 个 token，再把这 k 个接在已有文本后面让大模型做一次前向；因果注意力保证这一次前向同时给出 k+1 个位置的预测。\par}
\item {\fontsize{10}{14.8}\selectfont 从前往后比对，接受到第一个不一致的位置为止，那个位置换成大模型自己的预测，后面的草稿全部作废、KV 位置回退；所以每轮至少前进 1 个、最多 k+1 个，输出和大模型单独生成完全一样。\par}
\item {\fontsize{10}{14.8}\selectfont 期望前进 E = (1 - $\alpha$\textasciicircum{}(k+1)) ÷ (1 - $\alpha$)，一轮耗时是 1 + k$\cdot$c 个大模型步，加速比 = E ÷ (1 + k$\cdot$c)；$\alpha$ = 0.8、k = 4、c = 5\% 时 E = 3.362，加速 2.80 倍。\par}
\item {\fontsize{10}{14.8}\selectfont k 不是越大越好：$\alpha$ = 0.6 时最好的 k 是 4（1.92 倍），$\alpha$ = 0.8 时 k = 8 最好（3.09 倍），$\alpha$ = 0.9 时 k = 10 还能到 4.57 倍；猜得越往后越难全对，小模型的成本却是线性涨。\par}
\item {\fontsize{10}{14.8}\selectfont 采样时用拒绝采样：以 min(1, p/q) 接受草稿，拒绝时从 max(0, p-q) 归一化后重抽，输出分布恰好是 p；20 万次实验里接受率 0.699、理论 $\Sigma$ min(p, q) = 0.700，输出分布和 p 一致。\par}
\end{enumerate}\endgroup
\numline{公式}{$E=\frac{1-\alpha^{k+1}}{1-\alpha}$；加速比 $=E\div(1+k\cdot c)$；接受率 $\alpha=\sum\min(p,\ q)$}
\numline{数字}{1/5/32 个 token 前向都约 4.83 ms；$\alpha$ = 0.8、k = 4 $\rightarrow$ E = 3.362、2.80$\times$；$\alpha$ = 0.8 最佳 k = 8 $\rightarrow$ 3.09$\times$；$\alpha$ = 0.6 最佳 k = 4 $\rightarrow$ 1.92$\times$；$\alpha$ = 0.9、k = 10 $\rightarrow$ 4.57$\times$；接受率 0.699 对理论 0.700；模拟 3.361 对公式 3.362}
\fol{大批量下还能加速吗？}{不能。批量越大，大模型的 decode 越接近计算受限，多验证的那几个 token 不再免费，加速比随之下降，大批量时甚至会比不用更慢。}
\fol{草稿必须来自小模型吗？}{不一定。可以从输入里找现成片段当草稿（n-gram、提示查找，几乎零成本），也可以给大模型加预测头（Medusa、EAGLE），或者用训练时就带的多 token 预测（MTP）。}
\fol{怎么验证线上真有加速？}{在真实流量上测端到端吞吐和 TPOT，而不是只看接受率；写代码、改写文本这类可预测的输出 $\alpha$ 高，开放式创作和高温采样 $\alpha$ 低，加速比要按场景量。}
\numline{别答错}{加速比不等于每轮前进的 token 数，要除以 1 + k$\cdot$c；k 不是越大越好，$\alpha$ 低时多猜浪费。}
\phantomsection\addcontentsline{toc}{section}{97\quad EAGLE 与 Medusa}
\lec{97}{EAGLE 与 Medusa}
\askline{Medusa 和草稿模型比好在哪？}
\askline{为什么要一次验证一棵候选树？}
\askline{EAGLE 为什么在特征层猜？}
\vspace{3pt}
\noindent{\qlabel{答}}\begingroup\sloppy\begin{enumerate}[leftmargin=1.6em,itemsep=2.5pt,topsep=0pt,parsep=0pt,label={\sffamily\footnotesize\color{accent}\arabic*.}]
\item {\fontsize{10}{14.8}\selectfont 单独的草稿模型要和大模型用同一个词表、要占显存、要单独部署，很多大模型根本没有合适的小弟，更可惜的是大模型已经算好的内部表示它一点用不上；Medusa 和 EAGLE 就改成在大模型自己身上接轻量头。\par}
\item {\fontsize{10}{14.8}\selectfont Medusa 在最后一层隐藏状态上再接几个一两层的小头，第 i 个头猜第 t+i+1 个 token，原来的输出层照常预测第 t+1 个，一次前向就同时得到后面几个位置的猜测，不需要额外模型也不需要额外前向。\par}
\item {\fontsize{10}{14.8}\selectfont 越靠后的头越不准，因为它要跳过中间几个 token 直接猜，而中间是什么它并不知道；只取第 1 名时四个头的命中率是 0.60、0.45、0.35、0.28，按上一讲的算法一轮平均只前进 1.99 个 token。\par}
\item {\fontsize{10}{14.8}\selectfont 改进是每个头不只取第 1 名而取前 2 名，各个头前 2 名里含正确答案的概率升到 0.78、0.62、0.52、0.44；但前 2 名组合起来不是一条链而是一棵树，深度 3 就有 2+4+8 共 14 个候选。\par}
\item {\fontsize{10}{14.8}\selectfont 树形注意力只换掩码：14 个候选排成一排送进大模型，每个候选只看得见前文和自己的祖先，看不见同层兄弟，所以每个候选得到的预测和它单独送进去时完全一样，14 个候选一次前向验证完。\par}
\item {\fontsize{10}{14.8}\selectfont 验证完从根开始沿大模型认可的方向往下走，能走多深就接受多少个：示意数字下深度 2、3、4 的树平均前进 2.26、2.52、2.63 个，候选越多收益越小而验证成本越大，实际会挑一棵不对称的树。\par}
\end{enumerate}\endgroup
\numline{公式}{一轮期望前进 $=1+\sum_{i=1}^{k}\prod_{j=1}^{i}p_j$（链取 top-1，树取 top-2）；深度 $d$ 的候选数 $=2^{1}+\cdots+2^{d}$}
\numline{数字}{top-1 命中率 0.60/0.45/0.35/0.28 $\rightarrow$ 链 1.99；top-2 0.78/0.62/0.52/0.44；深度 2/3/4 的树 2.26/2.52/2.63，候选 6/14/30；MTP 第二个 token 接受率 85\%\textasciitilde{}90\%、约 1.8 倍}
\fol{树形注意力是新的注意力吗？}{不是。它只是换了一个掩码，让多条候选路径挤在一次前向里互不干扰，计算方式和普通因果注意力一样，所以 vLLM、SGLang 的 kernel 只要改掩码就能支持。}
\fol{这些方法会改变输出吗？}{不会。留下来的 token 都经过大模型验证，采样时再配合拒绝采样规则，结果和大模型单独生成相同；加速来自一次验证多个 token，不是换了模型。}
\fol{batch 大时还能加速吗？}{收益会变差。这些方法依赖大模型一次多处理几个 token 几乎免费，batch 大时这些额外 token 不再免费，引擎会按负载调小候选树甚至直接关掉投机解码。}
\numline{别答错}{树形注意力不是新机制只是换掩码；Medusa 的头不是 MoE 专家，而且越靠后的头越不准。}
\phantomsection\addcontentsline{toc}{section}{98\quad 蒸馏与剪枝}
\lec{98}{蒸馏与剪枝}
\askline{蒸馏为什么比只用标准答案好？}
\askline{剪枝剪完能直接变快吗？}
\askline{先剪再蒸馏这条路线怎么走？}
\vspace{3pt}
\noindent{\qlabel{答}}\begingroup\sloppy\begin{enumerate}[leftmargin=1.6em,itemsep=2.5pt,topsep=0pt,parsep=0pt,label={\sffamily\footnotesize\color{accent}\arabic*.}]
\item {\fontsize{10}{14.8}\selectfont 交叉熵只告诉模型标准答案是猫、其他都错；训练好的老师给出的却是整个分布，猫 0.744、狗 0.166、老虎 0.082，汽车只有千分之几，这个分布里藏着哪些答案也说得通的暗知识，蒸馏就是让学生去拟合它。\par}
\item {\fontsize{10}{14.8}\selectfont 为了让小概率的信息更明显，蒸馏把师生 logits 都除以温度 T 再做 softmax，损失是两个分布的 KL 散度 KL(p老师‖p学生)；T 取 4 时老师变成猫 0.336、狗 0.231、老虎 0.194，相对关系一目了然。\par}
\item {\fontsize{10}{14.8}\selectfont 差别看梯度最清楚：只用标准答案时梯度是 p-y，狗和汽车一起被压低；蒸馏的梯度是 (p学生-p老师)/T，狗反而被提高、只在压汽车，每个样本都带着更丰富的信号，学生用更少的数据就学得更好。\par}
\item {\fontsize{10}{14.8}\selectfont 大模型多用数据蒸馏：让老师生成大量回答、学生直接在上面做 SFT，不用同词表也不用跑老师前向；DeepSeek 就用 R1 生成的约 80 万条数据微调出一系列 Qwen 和 Llama 小模型，让它们也学会长思维链。\par}
\item {\fontsize{10}{14.8}\selectfont 剪枝的核心是判重要性：只按 |w| 剪会误伤激活大的离群通道，Wanda 用 |w|$\cdot$‖xⱼ‖ 把对应输入通道的激活算进去，同样剪掉一半，输出误差从 27.12\% 降到 10.72\%，而且只要一小批校准数据、不用再训练。\par}
\item {\fontsize{10}{14.8}\selectfont 剪了不一定快：零散的 0 矩阵形状没变，普通矩阵乘法照样要读进来乘一遍；只有 2:4 半结构化有稀疏 Tensor Core 支持（误差 13.92\%、理论快一倍）和整块结构化让矩阵真变小（误差 23.43\%）能在硬件上加速。\par}
\end{enumerate}\endgroup
\numline{公式}{$p=\operatorname{softmax}(z/T)$，$L=\text{KL}(p_{\text{老师}}\,\Vert\,p_{\text{学生}})$；重要性 $=|w|\cdot\|x_j\|$}
\numline{数字}{T=1 猫 0.744 / 狗 0.166；T=4 猫 0.336 / 狗 0.231；剪一半的误差 27.12\%（|w|）、10.72\%（Wanda）、13.92\%（2:4）、23.43\%（结构化）；R1 约 80 万条数据}
\fol{稀疏到 90\% 会变快吗？}{要有专门的稀疏格式和 kernel 才有收益。零散的 0 在普通矩阵乘法里照样要算，工程上通常不如直接量化，量化有统一的硬件支持和成熟的工具链。}
\fol{为什么结构化剪枝误差最大？}{因为它一次去掉整行整列甚至整层，动的东西多、损失难补，误差 23.43\% 比非结构化大一倍多；所以它默认要再训练，通常是拿原来的大模型当老师做蒸馏补回来。}
\fol{部署现成模型该选哪个？}{首选量化。非结构化难加速、2:4 的实际加速不到理论两倍且效果损失比同压缩比的量化明显、结构化要重训，剪枝和蒸馏更多是模型厂商做小模型的环节。}
\numline{别答错}{零散的 0 在普通矩阵乘法里照样要算，稀疏不等于快；蒸馏是训练另一个小模型，不是压缩原模型。}
\chapter{分布式与集合通信}
{\footnotesize\sffamily\color{grayt}第 99--109 讲}\par\vspace{-4pt}
\phantomsection\addcontentsline{toc}{section}{99\quad 为什么需要多卡}
\lec{99}{为什么需要多卡}
\askline{一张卡为什么不够用？}
\askline{70B 模型到底要几张 H100？}
\askline{多卡为什么也能变快？}
\vspace{3pt}
\noindent{\qlabel{答}}\begingroup\sloppy\begin{enumerate}[leftmargin=1.6em,itemsep=2.5pt,topsep=0pt,parsep=0pt,label={\sffamily\footnotesize\color{accent}\arabic*.}]
\item {\fontsize{10}{14.8}\selectfont 原因就三个：放不下、算不快、扛不住。权重等于参数量乘字节数，70B 的 BF16 是 141 GB、405B 812 GB，一张 H100 只有 80 GB，光权重就装不下，V3 用原生 FP8 也要 671 GB。\par}
\item {\fontsize{10}{14.8}\selectfont 再加 KV Cache：70B 每 token 320 KiB、405B 504 KiB，32 个 8K 请求分别要 86 GB 和 135 GB；按每卡 90\% 可用算，70B 要 4 张 H100、405B 14 张、V3 19 张。\par}
\item {\fontsize{10}{14.8}\selectfont 算不快是带宽问题：decode 每生成一个 token 要把权重从显存完整读一遍，H100 带宽 3.35 TB/s，70B 读一遍约 42 ms、单请求每秒最多 24 个 token，405B 要 242 ms、每秒 4 个，用户等不起。\par}
\item {\fontsize{10}{14.8}\selectfont 把权重分到 8 张卡上各读八分之一、同时读，70B 一步降到 5.3 ms、405B 降到 30.3 ms，8 张卡的总带宽就是一张卡的 8 倍；prefill 是计算受限的，道理一样，切到 8 张卡上也快得多。\par}
\item {\fontsize{10}{14.8}\selectfont 第三个原因是流量：线上每秒成百上千个请求，超过单卡饱和点延迟就失控，这时要把整个模型复制几份、请求分摊到各组；所以前两个原因对应切开，第三个对应复制，而切开以后卡与卡之间必须交换数据。\par}
\item {\fontsize{10}{14.8}\selectfont 训练比推理更依赖多卡：混合精度加 Adam 每个参数要 16 字节，8B 就是 128 GB、一张卡放不下，70B 是 1130 GB、至少 16 张；实际预训练动辄用上千甚至上万张卡。\par}
\end{enumerate}\endgroup
\numline{公式}{$\text{权重}=\text{参数量}\times\text{字节数}$；decode 单请求 token/s $\approx\beta\div\text{权重字节}$；训练状态 $\approx16$ 字节/参数；卡数 $=\lceil(\text{权重}+\text{KV})\div(80\ \text{GB}\times0.9)\rceil$}
\numline{数字}{H100 80 GB；权重 70B 141 / 405B 812 / V3 1342 GB（原生 FP8 671）；每 token KV 320 / 504 / 68.6 KiB；卡数 4 / 14 / 19；单卡 42 ms、8 卡 5.3 ms；训练 16 字节/参数}
\fol{换更大显存的卡不行吗？}{单卡带宽终究有限。H200 的 141 GB 能装下 70B 的 FP8，但 405B 就算塞进一张卡，decode 一步也要读几百 GB；真正方向是把更多卡连得更紧。}
\fol{切开和复制能一起用吗？}{能，而且最常见：先把大模型切到 8 张卡上组成一组，再把这一组复制成很多份去服务不同请求，切分解决装不下、复制解决扛不住，两种手段叠加。}
\fol{刚装得下的卡数能上生产吗？}{不能。卡数刚好装下权重时留给 KV Cache 的空间很少，并发上不去；实际部署会比最低要求多留几张卡，而且这里只算了权重和 KV，没算激活和框架本身的开销。}
\numline{别答错}{放得下只是最低要求，此时没给 KV Cache 留空间；多卡也不是免费加速，通信可能吃掉收益。}
\phantomsection\addcontentsline{toc}{section}{100\quad 集合通信原语}
\lec{100}{集合通信原语}
\askline{六种集合通信原语怎么区分？}
\askline{All-Reduce 每卡要传多少？}
\askline{All-to-All 为什么像转置？}
\vspace{3pt}
\noindent{\qlabel{答}}\begingroup\sloppy\begin{enumerate}[leftmargin=1.6em,itemsep=2.5pt,topsep=0pt,parsep=0pt,label={\sffamily\footnotesize\color{accent}\arabic*.}]
\item {\fontsize{10}{14.8}\selectfont 切开后卡间交换的形状很固定：Broadcast 一对多复制、Reduce 多对一求和、All-Reduce 求和后人人都有、Reduce-Scatter 求和后每人一段、All-Gather 拼成完整、All-to-All 互相交换。\par}
\item {\fontsize{10}{14.8}\selectfont 用 4 张卡、卡 r 的数据是 [10r, 10r+1, 10r+2, 10r+3] 看最清楚：按位置相加得 [60, 64, 68, 72]，Reduce 只把结果放在卡 0，All-Reduce 让四张卡都拿到完整的和。\par}
\item {\fontsize{10}{14.8}\selectfont Reduce-Scatter 把和切成 4 段，卡 0 拿 [60]、卡 1 拿 [64]；All-Gather 再把每段拼成完整。两步连起来就是 All-Reduce，ZeRO 只用前一半就省一半通信。\par}
\item {\fontsize{10}{14.8}\selectfont All-to-All 不一样：每张卡把数据切成 p 段，第 j 段发给卡 j，卡 0 最后拿到 [0, 10, 20, 30]，等于把一个 4$\times$4 的数据矩阵转置；这里既没有求和也没有复制，MoE 的专家并行靠它。\par}
\item {\fontsize{10}{14.8}\selectfont 通信量有理论下限：设每张卡的数据量是 n、共 p 张卡，All-Reduce 每卡至少收发 2(p-1)/p$\cdot$n，其余几种是 (p-1)/p$\cdot$n；4 张卡时分别是 1.5n 和 0.75n。\par}
\item {\fontsize{10}{14.8}\selectfont p 越大 (p-1)/p 越接近 1，每卡通信量几乎不随卡数增长，前提是实现让所有卡同时收发、把任务分摊开；另外集合通信要所有卡到场，最慢的那张卡决定了整体速度。\par}
\end{enumerate}\endgroup
\numline{公式}{$\text{All-Reduce}=\text{Reduce-Scatter}+\text{All-Gather}$；All-Reduce $2(p-1)/p\cdot n$，其余三种 $(p-1)/p\cdot n$}
\numline{数字}{p=4：[60, 64, 68, 72]；All-Reduce 每卡 1.5n，其余三种 0.75n；张量并行每层 2 次 All-Reduce、MoE 每层 2 次 All-to-All}
\fol{Reduce 结果只在卡 0 吗？}{差在结果放在哪。Reduce 只在一张卡上留下完整的和，其他卡什么都没有；All-Reduce 每张卡都拿到完整的和，所以数据并行同步梯度必须用后者，省掉一次广播。}
\fol{All-Gather 后数据一样吗？}{All-Gather 之后每张卡的数据完全相同，是完整的拼接；All-to-All 之后每张卡的数据各不相同，各自拿到的是数据矩阵里属于自己的那一列。}
\fol{为什么说 1.5n 是下限？}{因为每张卡至少要收别人那 p-1 段、又至少要参与一次求和，加起来就是 2(p-1)/p$\cdot$n；朴素做法让卡 0 收发 2(p-1)$\cdot$n，卡越多越慢，差距就在这里。}
\numline{别答错}{All-Reduce 每卡都有、Reduce 只在一张卡；集合通信最慢的卡决定整体速度。}
\phantomsection\addcontentsline{toc}{section}{101\quad Ring 与 Tree All-Reduce}
\lec{101}{Ring 与 Tree All-Reduce}
\askline{Ring 怎么把带宽做到下限？}
\askline{小消息为什么 Tree 更快？}
\askline{NCCL 怎么选算法？}
\vspace{3pt}
\noindent{\qlabel{答}}\begingroup\sloppy\begin{enumerate}[leftmargin=1.6em,itemsep=2.5pt,topsep=0pt,parsep=0pt,label={\sffamily\footnotesize\color{accent}\arabic*.}]
\item {\fontsize{10}{14.8}\selectfont 朴素 All-Reduce 把负担全压在一张卡上：所有卡把数据发给卡 0 求和、卡 0 再发回去，卡 0 要收发 2(p-1)$\cdot$n，其他卡大部分时间在等，时间随卡数线性增长。\par}
\item {\fontsize{10}{14.8}\selectfont Ring 把 p 张卡排成环、数据切成 p 块，先 p-1 步 Reduce-Scatter：每块沿环走一圈、被每张卡加一次，再 p-1 步 All-Gather 把完整的块传一圈，一共 2(p-1) 步。\par}
\item {\fontsize{10}{14.8}\selectfont 每一步所有 p 条链路同时传 n/p，没有谁闲着，每张卡一共发 2(p-1)/p$\cdot$n，正好达到理论下限；4 张卡时每张卡 1.5n，而朴素做法里卡 0 要收发 6n。\par}
\item {\fontsize{10}{14.8}\selectfont 但 Ring 的延迟项是 2(p-1)$\cdot$$\alpha$，随卡数线性增长：1 GB 数据、512 张卡时朴素要 20 秒而 Ring 只要 45 ms；换成 64 KB 的小消息，Ring 却要 5.1 ms，几乎全是延迟。\par}
\item {\fontsize{10}{14.8}\selectfont Tree 的思路是压缩步数：数据从叶子向根逐层相加、再从根向叶子广播，深度是 log$_2$p，所以只要约 2$\cdot$log$_2$p 步；512 张卡时 Ring 要 1022 步，树只要 18 步。\par}
\item {\fontsize{10}{14.8}\selectfont 树的代价是根附近要收发整份数据，要把数据切块像流水线一样往上传，让带宽项接近 2n/$\beta$；NCCL 用双二叉树，让每张卡在一棵树里当叶子、在另一棵树里当中间节点，负担均衡。\par}
\end{enumerate}\endgroup
\numline{公式}{$T_{\text{ring}}=2(p-1)\cdot(\alpha+n/(p\cdot\beta))$；$T_{\text{tree}}\approx2\lceil\log_2 p\rceil\cdot\alpha+2n/\beta$}
\numline{数字}{1 GB / 512 卡：朴素 20 s、Ring 45 ms、树 40 ms；64 KB / 512 卡：Ring 5.1 ms、树 0.09 ms；4 卡每卡 1.5n、朴素 6n；步数 1022 vs 18}
\fol{Ring 的环要物理连成圈吗？}{不要。环是逻辑上的，只要求相邻两张卡之间能通信；NCCL 会按实际的 NVLink、PCIe、网卡安排环的顺序，尽量让每一次跳转都走最快的路径。}
\fol{为什么带宽最优还会更慢？}{因为时间分延迟和带宽两项。消息小的时候延迟项占主导，Ring 有 2(p-1) 次 $\alpha$，卡一多就吃亏；Tree 只要 2log$_2$p 次 $\alpha$，所以小消息时反而反超。}
\fol{跨机时怎么把慢网络用好？}{用分层通信：先在每台机器内部做 Reduce-Scatter 走 NVLink，再在机器之间对各自的一份做 All-Reduce 走网络，最后机内 All-Gather，慢速网络上只传八分之一的数据。}
\numline{别答错}{带宽最优不等于最快，小消息时 Tree 反超；环只是逻辑上的，不要求物理成环。}
\phantomsection\addcontentsline{toc}{section}{102\quad 通信代价模型}
\lec{102}{通信代价模型}
\askline{一次通信的时间怎么估？}
\askline{多大的消息算小消息？}
\askline{张量并行为什么要留在机内？}
\vspace{3pt}
\noindent{\qlabel{答}}\begingroup\sloppy\begin{enumerate}[leftmargin=1.6em,itemsep=2.5pt,topsep=0pt,parsep=0pt,label={\sffamily\footnotesize\color{accent}\arabic*.}]
\item {\fontsize{10}{14.8}\selectfont 一条消息的耗时是 T = $\alpha$ + m/$\beta$：$\alpha$ 是每条消息的固定开销，包含通信库准备和启动传输、建立连接、过交换机、对方确认，通常是几微秒到十几微秒；$\beta$ 是链路的单向带宽。\par}
\item {\fontsize{10}{14.8}\selectfont 两项相等时得到分界点 m* = $\alpha$$\cdot$$\beta$：NVLink 上 3 $\mu$s $\times$ 450 GB/s 约 1.35 MB，InfiniBand 上 6 $\mu$s $\times$ 50 GB/s 约 300 KB；比它小的消息卡在延迟上，比它大的才接近标称带宽。\par}
\item {\fontsize{10}{14.8}\selectfont 这和第 62 讲的 Roofline 同构：算术强度低于拐点就卡带宽，消息小于 $\alpha$$\cdot$$\beta$ 就卡延迟。NVLink 上发 8 KB 只要 3.0 $\mu$s，99.4\% 是延迟，实际带宽只有 2.7 GB/s，不到标称的百分之一。\par}
\item {\fontsize{10}{14.8}\selectfont 集合通信把每一步加起来：Ring 是 2(p-1)$\cdot$$\alpha$ + 2(p-1)/p$\cdot$m/$\beta$，延迟项被放大 2(p-1) 倍；一步到位则每卡同时读其余 p-1 张卡的数据自己相加，只付一次 $\alpha$ 但要多收 (p-1)$\cdot$m。\par}
\item {\fontsize{10}{14.8}\selectfont 推理的张量并行是大量小消息：70B 八卡 decode、80 层共 160 次 524 KB 的 All-Reduce，Ring 一步 7.05 ms 里 6.72 ms 是延迟，一步到位只要 1.78 ms。\par}
\item {\fontsize{10}{14.8}\selectfont 训练的梯度同步相反：8B 的 16 GB 梯度在 16 个节点上 Ring 约 0.60 s，几乎全是带宽项，这时该降低精度、把通信和反向计算重叠，而不是去优化延迟；模型只是估算，具体数字要用 nccl-tests 测。\par}
\end{enumerate}\endgroup
\numline{公式}{$T=\alpha+m/\beta$；$m^{*}=\alpha\cdot\beta$；$T_{\text{ring}}=2(p-1)\cdot(\alpha+m/(p\beta))$}
\numline{数字}{$\alpha$=3 $\mu$s / $\beta$=450 GB/s（NVLink）、$\alpha$=6 $\mu$s / $\beta$=50 GB/s（IB）；m* 1.35 MB / 300 KB；512 KB 延迟占 72.5\%；160 次 524 KB：Ring 7.05 ms vs 一步到位 1.78 ms；16 GB 梯度 0.60 s}
\fol{标称带宽能直接代进模型吗？}{不能，要先换算成单向的字节每秒。NVLink 900 GB/s 是所有链路双向的总和，单向约 450 GB/s；400 Gb/s 的网卡单位是比特，除以 8 才约 50 GB/s。}
\fol{$\alpha$ 到底是什么，能优化吗？}{它不只是物理延迟，大头是软件和协议：通信库调度、kernel 启动、同步，所以优化通信库能显著降低 $\alpha$；vLLM 就专门给张量并行写了小消息的 All-Reduce。}
\fol{nccl-tests 的带宽怎么读？}{看总线带宽而不是算法带宽。总线带宽等于算法带宽乘 2(p-1)/p，把每张卡实际要传的数据量算进去；16 GB 用 0.60 s 是算法 26.7 GB/s，总线 50 GB/s 正好等于链路 $\beta$。}
\numline{别答错}{标称常是双向总和（NVLink 900 即单向 450）；要用总线带宽才能和链路带宽比。}
\phantomsection\addcontentsline{toc}{section}{103\quad 数据并行与 ZeRO}
\lec{103}{数据并行与 ZeRO}
\askline{数据并行怎么保证各卡不跑偏？}
\askline{ZeRO 三级各切了什么？}
\askline{ZeRO-3 为什么能省到 1/N？}
\vspace{3pt}
\noindent{\qlabel{答}}\begingroup\sloppy\begin{enumerate}[leftmargin=1.6em,itemsep=2.5pt,topsep=0pt,parsep=0pt,label={\sffamily\footnotesize\color{accent}\arabic*.}]
\item {\fontsize{10}{14.8}\selectfont 数据并行是 N 张卡各放一份完整模型、各算 1/N 的数据，梯度用 All-Reduce 求平均；因为损失对所有样本取平均，各份梯度平均正好等于整个大批次的梯度，所以各卡更新后模型始终一致。\par}
\item {\fontsize{10}{14.8}\selectfont 代价是显存一点没省：每参数 16 字节，BF16 权重 2、梯度 2，FP32 主权重 4，Adam 两个状态 8，8B 模型每卡要 128.5 GB，一张 80 GB 的卡根本放不下。\par}
\item {\fontsize{10}{14.8}\selectfont ZeRO 指出这些状态在 N 张卡上完全一样、是 N 倍冗余：ZeRO-1 只切优化器状态，卡 r 只负责第 r 段参数的更新，更新后用 All-Gather 把各段拼回完整模型，8 卡时每卡降到 44.2 GB。\par}
\item {\fontsize{10}{14.8}\selectfont ZeRO-2 再切梯度：既然卡 r 只更新第 r 段，它也只需要第 r 段的梯度，把梯度的 All-Reduce 换成 Reduce-Scatter 就够了，每卡再降到 30.1 GB。\par}
\item {\fontsize{10}{14.8}\selectfont ZeRO-3 连参数也切：平时只存 1/N 的 BF16 权重，前向算到某层时用 All-Gather 临时拼出完整权重、算完就丢，反向再拼一次，每卡 16.1 GB 即 16Ψ/N，PyTorch 的 FSDP 就是它。\par}
\item {\fontsize{10}{14.8}\selectfont 通信代价：ZeRO-1 和 ZeRO-2 都是 Reduce-Scatter 加 All-Gather，约 2$\times$2Ψ，和数据并行一样；ZeRO-3 多一次 All-Gather，约 3$\times$2Ψ，是它的 1.5 倍。\par}
\end{enumerate}\endgroup
\numline{公式}{每卡状态：数据并行 $16\Psi$；ZeRO-1 $4\Psi+12\Psi/N$；ZeRO-2 $2\Psi+14\Psi/N$；ZeRO-3 $16\Psi/N$}
\numline{数字}{8B、8 卡：128.5 $\rightarrow$ 44.2 $\rightarrow$ 30.1 $\rightarrow$ 16.1 GB；每步每卡通信 32 / 32 / 32 / 48 GB；ZeRO-3 是数据并行的 1.5 倍；每参数 16 字节}
\fol{ZeRO 会改变训练结果吗？}{不会。它只改变状态存在哪里、什么时候通信，每一步的更新和普通数据并行完全相同，代码能验证切分更新再 All-Gather 拼起来和不切分的结果一样。}
\fol{ZeRO-3 和张量并行都是切参数？}{切法不同。ZeRO-3 是存的时候切开、算的时候拼回完整矩阵，每张卡仍然做完整的计算只是数据不同；张量并行连计算也切开，每卡只算矩阵乘法的一部分、只交换激活。}
\fol{推理为什么不用 ZeRO-3？}{因为每一步都要 All-Gather 一遍权重。8B 模型每卡每步要收约 14 GB，走 NVLink 也要 30 ms，比读一遍权重的 4.8 ms 慢得多；推理用不搬权重、只搬激活的张量并行。}
\numline{别答错}{ZeRO 不改变计算结果；卡数增加时没切的部分不缩小，ZeRO-1 的权重和梯度始终 4Ψ。}
\phantomsection\addcontentsline{toc}{section}{104\quad 混合精度训练与激活重计算}
\lec{104}{混合精度训练与激活重计算}
\askline{16 字节/参数是怎么来的？}
\askline{BF16 为什么不需要损失缩放？}
\askline{激活重计算到底多花多少？}
\vspace{3pt}
\noindent{\qlabel{答}}\begingroup\sloppy\begin{enumerate}[leftmargin=1.6em,itemsep=2.5pt,topsep=0pt,parsep=0pt,label={\sffamily\footnotesize\color{accent}\arabic*.}]
\item {\fontsize{10}{14.8}\selectfont 前向反向的矩阵乘法都用 BF16，Tensor Core 上 989 对 67 TFLOPS、数据还小一半；但必须留一份 FP32 主参数，因为 BF16 尾数只 7 位、1 附近间隔 2$^-$⁷，每步加 1e-4 加一千步还是 1.0。\par}
\item {\fontsize{10}{14.8}\selectfont 所以每个参数要存 16 字节：BF16 参数 2、BF16 梯度 2、FP32 主参数 4、Adam 两份 FP32 滑动平均 8；更进一步的 FP8 训练只是把矩阵乘法的输入降到 8 位，主参数和优化器状态仍然是高精度。\par}
\item {\fontsize{10}{14.8}\selectfont FP16 不行是因为指数只有 5 位，2$\times$10$^-$⁸ 的梯度直接下溢成 0，要靠损失缩放把 loss 乘 1024 再在 FP32 里除回来，倍数大了还会上溢成无穷；BF16 指数 8 位、范围和 FP32 一样，省掉整套机制。\par}
\item {\fontsize{10}{14.8}\selectfont 另一半显存是激活：每层约 s$\cdot$b$\cdot$h$\cdot$(34 + 5as/h) 字节，前一项和 s 成正比，后一项是 [a, s, s] 的注意力矩阵、和 s$^2$ 成正比；8B、8192、b=1 时 32 层要 380.1 GB，比参数状态还大三倍。\par}
\item {\fontsize{10}{14.8}\selectfont 全部重计算是前向只留每层输入、中间结果全扔，反向算到这层时重跑一遍前向再算梯度，激活降到 3.29 GB，代价是训练从 6N 变 8N、多算 33\%，并不是多一倍。\par}
\item {\fontsize{10}{14.8}\selectfont 选择性重计算只重算 5as/h 那一项就把平方项去掉，8B、8192 时降到 36.5 GB、只多算 8.3\%；FlashAttention 不把 [s, s] 的矩阵写进显存，天然等效于选择性重计算。\par}
\end{enumerate}\endgroup
\numline{公式}{每层激活 $\approx s\cdot b\cdot h\cdot(34+5as/h)$ 字节；训练 FLOPs $6N\to8N$（多 33\%）}
\numline{数字}{H100 BF16 989 vs FP32 67 TFLOPS；16 字节/参数 = 2+2+4+8；8B s=8192：380.1 / 36.5 / 3.29 GB，多算 8.3\% 与 33\%；s=32768：146.0 / 13.15 GB}
\fol{BF16 也能做损失缩放吗？}{不需要。它的范围和 FP32 一样，2$\times$10$^-$⁸ 存得好好的；损失缩放是 FP16 时代为了绕开 6$\times$10$^-$⁸ 的最小正规数，现在还留着主要是兼容旧代码。}
\fol{重计算和 ZeRO 是一回事吗？}{不是。重计算省的是激活，参数、梯度、优化器状态要靠 ZeRO 切；两者通常一起用，8B 就是 ZeRO-3 的 16.1 GB 加 FlashAttention 的 36.5 GB。}
\fol{序列再长该怎么处理？}{s=32768 时即使用 FlashAttention 也要 146 GB，一张卡放不下，要么全部重计算降到 13.15 GB，要么用张量并行、序列并行把激活切到多张卡上，代价是多算 19\%。}
\numline{别答错}{重计算省的是激活不是参数状态；6N 变 8N 只多 33\%，不是翻倍；BF16 不需要损失缩放。}
\phantomsection\addcontentsline{toc}{section}{105\quad 张量并行}
\lec{105}{张量并行}
\askline{张量并行是怎么切矩阵的？}
\askline{为什么 FFN 先按列切、再按行切？}
\askline{每层要通信几次，为什么？}
\vspace{3pt}
\noindent{\qlabel{答}}\begingroup\sloppy\begin{enumerate}[leftmargin=1.6em,itemsep=2.5pt,topsep=0pt,parsep=0pt,label={\sffamily\footnotesize\color{accent}\arabic*.}]
\item {\fontsize{10}{14.8}\selectfont 张量并行把矩阵乘法本身拆到 t 张卡上，每张卡只用自己那一块权重算出一部分结果，卡与卡之间只交换小得多的激活；它最早由 2019 年的 Megatron-LM 系统化提出，是推理里最常用的模型切分方式。\par}
\item {\fontsize{10}{14.8}\selectfont 一个矩阵乘法有两种切法：A 按列切，每张卡拿一部分列、用完整的 X 乘，得到的是完整值的部分列；A 按行切且 X 也按列切，每张卡只得到部分和，必须 All-Reduce 相加才是完整的 Y。\par}
\item {\fontsize{10}{14.8}\selectfont FFN 用先列后行：A 按列切、B 按行切，ReLU 逐元素各卡独立做，中间列正好对上 B 的行，最后只 All-Reduce 一次；反过来 A 按行切，ReLU 非线性，它之前就得先 All-Reduce，FFN 要通信两次。\par}
\item {\fontsize{10}{14.8}\selectfont 注意力按头切：QKV 投影按列（按头）切，每张卡算出自己那几个头的 q、k、v，输出投影按行切，最后一次 All-Reduce；所以每个 Transformer 层前向只需要 2 次 All-Reduce，KV Cache 也随头一起切开。\par}
\item {\fontsize{10}{14.8}\selectfont 收益是每卡只读 1/t 权重：70B、32 请求一步，TP=8 时每卡 17.6 GB、读权重 5.3 ms 加通信 1.78 ms 共 7.0 ms，TP=4 是 35.2 GB、11.6 ms；代价是通信占比从 3\% 涨到 25\%。\par}
\item {\fontsize{10}{14.8}\selectfont 限制是每张卡都要参与每一层，任何一张慢全组都等，所以 TP 通常只做机内 8 卡；GQA 只有 8 个 KV 头，TP=8 正好每卡 1 个，再大就要复制 KV 头，KV Cache 不再随卡数减少。\par}
\end{enumerate}\endgroup
\numline{公式}{按列切 $Y=[X\cdot A_0,\ X\cdot A_1]$；按行切 $Y=X_0\cdot A_0+X_1\cdot A_1$，需 All-Reduce}
\numline{数字}{Llama-3.1-70B：TP=1 每卡 141 GB、42.1 ms；TP=2 70.5 GB、21.7 ms；TP=4 35.2 GB、11.6 ms；TP=8 17.6 GB、7.0 ms；通信占比 3\%$\rightarrow$25\%；2019 年 Megatron-LM}
\fol{不按头切注意力行吗？}{不行。每个头的注意力要看到该头完整的 q、k、v，按 hidden 维切会把一个头拆散，必须额外通信；按头切正好落在本来就独立的头上，各卡之间零通信。}
\fol{TP 跨机还能做吗？}{能但基本不划算。跨机带宽只有机内的九分之一，70B 在 TP=8 时机内通信 1.78 ms，跨机后涨到 16 ms 以上，TPOT 直接到三倍，所以 TP 一般限制在机内 8 卡。}
\fol{怎么验证切法写对了？}{用小规模 NumPy 对比：4 张卡的部分和相加应等于单卡结果；再故意改成 A 按行切、先 ReLU 再相加，结果会不相等，正好证明非线性之前必须先 All-Reduce。}
\numline{别答错}{别把按行切说成「也得到一部分输出」------它得到的是部分和，不做 All-Reduce 结果就是错的。}
\phantomsection\addcontentsline{toc}{section}{106\quad 流水并行}
\lec{106}{流水并行}
\askline{流水并行怎么切，气泡从哪来？}
\askline{气泡占比怎么算、怎么压？}
\askline{推理用流水并行能降延迟吗？}
\vspace{3pt}
\noindent{\qlabel{答}}\begingroup\sloppy\begin{enumerate}[leftmargin=1.6em,itemsep=2.5pt,topsep=0pt,parsep=0pt,label={\sffamily\footnotesize\color{accent}\arabic*.}]
\item {\fontsize{10}{14.8}\selectfont 流水并行不拆矩阵而是按层切：80 层模型放 4 张卡，每卡 20 层，数据从卡 0 流到卡 3；段与段之间只在交界处用点对点 Send/Recv 传一次激活，通信量很小，适合跨机。\par}
\item {\fontsize{10}{14.8}\selectfont 只有一批数据时，任何时刻只有一张卡在干活，4 段利用率只有 25\%，空等的部分叫气泡；把批次切成 m 个微批次依次送进流水线，灌满排空之后所有卡才同时工作。\par}
\item {\fontsize{10}{14.8}\selectfont GPipe 让所有微批次先做完前向、再做完反向、最后统一更新，气泡只出现在灌满和排空：占比是 (p-1)/(m+p-1)，4 段时 m=1 为 75\%、m=4 为 42.9\%、m=16 为 15.8\%，所以要求 m 远大于 p。\par}
\item {\fontsize{10}{14.8}\selectfont 1F1B 先做 p-s-1 个前向热身，之后每做一个前向就紧跟一个反向，微批次的反向一完就释放激活；每卡同时存的激活从 m 份降到最多 p 份，m=16 时从 16 份降到 4 份，总时长和气泡完全一样。\par}
\item {\fontsize{10}{14.8}\selectfont 推理没有反向，连续批处理把正在跑的请求分组依次送进流水线，组数不少于段数就满载、吞吐几乎无损；但单个请求的每个 token 仍要依次走完 p 段，70B 四段一步 42.1 ms，而 TP=4 同时读只要约 10.5 ms。\par}
\item {\fontsize{10}{14.8}\selectfont 实际部署是机内 TP、机间 PP：405B 用两台 8 卡机器，每台 TP=8，两台之间 PP=2；段间每步只传 32 个 token 的隐藏状态约 524 KB，跨机也只要约 16 $\mu$s。\par}
\end{enumerate}\endgroup
\numline{公式}{$\text{气泡占比}=\frac{p-1}{m+p-1}$}
\numline{数字}{p=4：m=1 总时长 12、气泡 75\%；m=4 21、42.9\%；m=8 33、27.3\%；m=16 57、15.8\%；70B 一步：PP 四段 42.1 ms vs TP=4 约 10.5 ms；段间 524 KB、跨机约 16 $\mu$s}
\fol{微批次 m 越大越好吗？}{不是。气泡随 m 减小但边际递减，m 从 4 到 16 只把气泡从 42.9\% 压到 15.8\%；而 GPipe 的激活随 m 线性涨，所以 1F1B 把激活压到 p 份，m 常取 8 到 32。}
\fol{推理为什么还能用 PP？}{因为推理不用等整个批次结束。连续批处理把在跑的请求分组，组数不少于段数流水线就一直是满的，吞吐几乎不损；代价是单请求要走完所有段，延迟不降。}
\fol{段负载不均衡会怎样？}{公式假设每段耗时相等，实际最慢的那段决定整条流水线的速度，气泡比公式更大；所以第一段带嵌入、最后一段带输出层时，要少分几层来配平。}
\numline{别答错}{别说 PP 能降低单请求延迟------它让模型装得下、让吞吐上去，每个 token 仍要走过所有段。}
\phantomsection\addcontentsline{toc}{section}{107\quad 序列与上下文并行}
\lec{107}{序列与上下文并行}
\askline{长序列放不下，怎么沿序列切？}
\askline{Ring Attention 怎么保证结果不变？}
\askline{因果掩码下为什么要之字形切？}
\vspace{3pt}
\noindent{\qlabel{答}}\begingroup\sloppy\begin{enumerate}[leftmargin=1.6em,itemsep=2.5pt,topsep=0pt,parsep=0pt,label={\sffamily\footnotesize\color{accent}\arabic*.}]
\item {\fontsize{10}{14.8}\selectfont 先分清两种：Megatron 2022 的序列并行（SP）只切张量并行里没切的逐 token 运算，上下文并行（CP）才把注意力本身沿序列切开；文档里同一个词常指不同东西，回答前要先问清切的是什么。\par}
\item {\fontsize{10}{14.8}\selectfont SP 让每卡只对 1/t token 做 RMSNorm 和残差，进注意力或 FFN 前用 All-Gather 拼回完整序列，算完把 All-Reduce 换成 Reduce-Scatter；通信总量不变，激活显存降到 1/t，几乎免费。\par}
\item {\fontsize{10}{14.8}\selectfont Ring Attention 把序列切成 P 段，每卡存自己那段的 Q、K、V，K、V 沿环传 P 步；每步用自己 Q 和刚收到的块算一块，用在线 softmax 维护 m、l、acc，最后输出 acc/l，和一次算完相同。\par}
\item {\fontsize{10}{14.8}\selectfont 因果掩码下按顺序切很不均衡：卡 0 的 token 在最前面，大部分块被掩码挡掉，各卡计算量 36、100、164、228，最忙是最闲的 6.3 倍；之字形切成 2P 段、卡 i 拿第 i 段和倒数第 i 段，四卡都是 132。\par}
\item {\fontsize{10}{14.8}\selectfont Ulysses 用两次 All-to-All 在「按序列切」和「按头切」之间切换，切完每卡对部分头做完整注意力、直接用 FlashAttention；但并行度不能超过头数，GQA 的 KV 头更少。\par}
\item {\fontsize{10}{14.8}\selectfont 代价是通信不便宜：100 万 token 的 Llama-3-8B，KV Cache 122 GiB，8 卡每卡 15 GiB，每层每步要传 488 MiB；所以 CP 通常放在 NVLink 上，靠计算和传输重叠藏住通信。\par}
\end{enumerate}\endgroup
\numline{公式}{Ring：$P$ 步轮转 $K$、$V$，每步按在线 softmax 合并 $(m,\ l,\ \text{acc})$}
\numline{数字}{1M token 的 Llama-3-8B KV 122 GiB，8 卡每卡 15 GiB，每层每步传 488 MiB；因果掩码按顺序切 36/100/164/228，之字形后都是 132；Ulysses 并行度受头数限制}
\fol{SP 会多花通信吗？}{不会。All-Reduce 本来就等于 Reduce-Scatter 加 All-Gather，SP 只是把这两半挪到注意力或 FFN 的前后，总字节数不变，省下来的是这部分激活的显存。}
\fol{CP 是近似计算吗？}{不是。Ring Attention 用在线 softmax 精确合并，和 FlashAttention 一样，结果与完整注意力一致，代码里带不带因果掩码都验证过。}
\fol{双向注意力也要之字形吗？}{不用。双向注意力每个查询都看全部键，按顺序切本来就已经均衡；之字形只是为因果掩码下「前轻后重」的分布设计的。}
\numline{别答错}{「序列并行」一词被用得很乱，先问清它切的是逐 token 运算还是注意力本身再回答。}
\phantomsection\addcontentsline{toc}{section}{108\quad 专家并行与 All-to-All}
\lec{108}{专家并行与 All-to-All}
\askline{MoE 的专家怎么分到卡上？}
\askline{一个 MoE 层要几次 All-to-All？}
\askline{专家并行最大的问题是什么？}
\vspace{3pt}
\noindent{\qlabel{答}}\begingroup\sloppy\begin{enumerate}[leftmargin=1.6em,itemsep=2.5pt,topsep=0pt,parsep=0pt,label={\sffamily\footnotesize\color{accent}\arabic*.}]
\item {\fontsize{10}{14.8}\selectfont MoE 每个 token 只走 top-k 个专家，参数却一个不少：DeepSeek-V3 每层 256 个路由专家、每 token 选 8 个，总参数 6710 亿、每 token 只用约 370 亿，一层就有 113 亿参数。\par}
\item {\fontsize{10}{14.8}\selectfont 用张量并行切专家不划算：每个专家本来只有约 4400 万参数，切开后矩阵乘更小、GPU 根本跑不满；专家并行把不同的专家整块放到不同卡上，每张卡放完整的几个专家。\par}
\item {\fontsize{10}{14.8}\selectfont 一个 MoE 层是四步：每张卡对自己的 token 算路由打分选 top-k，dispatch 把 token 副本发到专家所在的卡，各卡用自己的专家算，combine 再做一次 All-to-All 把结果送回原卡、按门控权重加起来。\par}
\item {\fontsize{10}{14.8}\selectfont 每对卡之间传的都是一批不同的 token，这正是 All-to-All；所以它的数据量取决于路由而不是固定值，同一个模型、同样的批大小，换一批输入各卡之间的流量就变。\par}
\item {\fontsize{10}{14.8}\selectfont 负载不均是 EP 特有的难题：4 卡 8 专家的例子里各专家收到 18、17、7、7、4、5、5、1 个 token，专家 0、1 都在卡 0 上，卡 0 要算 35 个副本而平均只有 16 个，第二次 All-to-All 要等它算完。\par}
\item {\fontsize{10}{14.8}\selectfont 通信量很大：DeepSeek-V3 decode 每卡 128 token、d=7168，每层分发 7.3 MB、合并 14.7 MB，50 GB/s 下约 0.44 ms，58 层约 26 ms；对策是限制跨机、算通重叠和 DeepEP。\par}
\end{enumerate}\endgroup
\numline{公式}{每卡每层发出量 $\approx\text{token 数}\times\text{top-}k\times d\times\text{字节数}$（留在本卡的不发）}
\numline{数字}{每层分发 7.3 MB（FP8）、合并 14.7 MB（BF16），50 GB/s 下 0.44 ms、58 层约 26 ms；prefill 用 32 卡 EP32、decode 用 320 卡 EP320；每 token 最多发往 4 个节点；最忙卡 35 个副本 vs 平均 16 个}
\fol{为什么不把专家集中放？}{放不下也算不快：EP=8 时每卡每层 32 个专家，58 层专家权重加非专家约 99 GB，超过 80 GB；专家集中后每卡读的权重变多，专家计算从 0.053 ms 变成 0.42 ms。}
\fol{负载不均该怎么解决？}{训练时用无辅助损失的均衡去调路由偏置；推理时复制热门专家、分散到别的卡上，并按统计动态调整放置，DeepSeek-V3 的 decode 就额外放了 32 个冗余专家。}
\fol{怎么验证 All-to-All 写对了？}{用 4 卡 8 专家的小例子对比不分布式结果，两次 All-to-All 之后应完全一致；同时打印分发矩阵和每个专家的负载，看是否和路由统计对得上。}
\numline{别答错}{EP 只切 MoE 层，注意力和嵌入还要 DP 或 TP，所以它几乎总是和其他并行组合出现。}
\phantomsection\addcontentsline{toc}{section}{109\quad 并行策略与拓扑映射}
\lec{109}{并行策略与拓扑映射}
\askline{五种并行怎么映射到硬件拓扑？}
\askline{同样 TP=8，分组不同差多少？}
\askline{超节点改变了哪些东西？}
\vspace{3pt}
\noindent{\qlabel{答}}\begingroup\sloppy\begin{enumerate}[leftmargin=1.6em,itemsep=2.5pt,topsep=0pt,parsep=0pt,label={\sffamily\footnotesize\color{accent}\arabic*.}]
\item {\fontsize{10}{14.8}\selectfont 互联是分层的：8 卡机内经 NVSwitch 全连，每卡每向约 450 GB/s、延迟几微秒；机间 400 Gb/s InfiniBand 约 50 GB/s，带宽差 9 倍，系统越大层级越多、越往外越慢。\par}
\item {\fontsize{10}{14.8}\selectfont 原则是按通信频率和同步等待程度从里到外排：TP 每层两次 All-Reduce、下一层还要等它，最难重叠，放最里层 NVLink；CP、EP 量大也尽量留在 NVLink 域；PP 段间点对点量小、DP 每步同步一次梯度，放外层跨机。\par}
\item {\fontsize{10}{14.8}\selectfont 两台 8 卡 H100 跑 70B：TP=8 机内两副本 7.0 ms、9083 token/s；TP 组跨机则 21.6 ms、2958 token/s；TP=16 跨机通信 31.95 ms、TPOT 34.6 ms。\par}
\item {\fontsize{10}{14.8}\selectfont TP=8 加 PP=2（PP 跨机）几乎不受影响：TPOT 7.1 ms、9062 token/s，而且每卡只占 8.8 GB 权重，留给 KV Cache 的空间更大。\par}
\item {\fontsize{10}{14.8}\selectfont 延迟和吞吐要分开看：TP=4 机内乘 4 份单请求更慢（TPOT 11.6 ms），但份数多，全集群约 11,071 token/s，反而比 TP=8 高；选哪个取决于服务看延迟还是看吞吐。\par}
\item {\fontsize{10}{14.8}\selectfont 训练是 4D 并行：Llama 3 405B 从里到外是 TP、CP、PP、DP（FSDP），按通信频率从高到低排；GB200 NVL72 把 72 卡放进一个 NVLink 域，更大的 TP 和 EP 都能走高速链路。\par}
\end{enumerate}\endgroup
\numline{数字}{TP=8 机内$\times$2：17.6 GB、7.0 ms、9083 tok/s；TP=8 组跨机：21.6 ms、2958 tok/s；TP=8$\times$PP=2：7.1 ms、9062 tok/s；TP=4$\times$4：11.6 ms、11,071 tok/s；TP=16 跨机 34.6 ms；NVLink 450 GB/s vs IB 50 GB/s；NVL72}
\fol{rank 编号为什么重要？}{通信组按 rank 编号建立，同一个 TP 组的 rank 没落在同一台机器上，All-Reduce 就被迫走网络；它不会报错，只会让 TPOT 从 7.0 ms 悄悄变成 21.6 ms。}
\fol{卡越多就一定越快吗？}{不一定。TP=16 比 TP=8 每卡权重少一半，但跨机通信 31.95 ms 远超读权重的时间，TPOT 反而从 7.0 变成 34.6 ms；增加并行度之前先算通信。}
\fol{最终方案该怎么定下来？}{先按拓扑和通信模型估一版、排除明显错的配置，再用第 82 讲的压测方法在目标负载下比较几种配置的 goodput，模型估算不能代替实测。}
\numline{别答错}{并行度大小和放在哪些卡上是两件事，后者错了不报错，只会悄悄慢上三倍。}
\chapter{分布式推理与集群}
{\footnotesize\sffamily\color{grayt}第 110--116 讲}\par\vspace{-4pt}
\phantomsection\addcontentsline{toc}{section}{110\quad 推理中的张量并行}
\lec{110}{推理中的张量并行}
\askline{放得下以后还要不要加 TP？}
\askline{TP 翻倍，延迟能降多少？}
\askline{两个阶段的 TP 需求一样吗？}
\vspace{3pt}
\noindent{\qlabel{答}}\begingroup\sloppy\begin{enumerate}[leftmargin=1.6em,itemsep=2.5pt,topsep=0pt,parsep=0pt,label={\sffamily\footnotesize\color{accent}\arabic*.}]
\item {\fontsize{10}{14.8}\selectfont 单卡 decode 一步约 max(计算÷算力, (权重+KV)÷带宽)，TP=t 时每卡只读 1/t、算 1/t，但每层多 2 次 All-Reduce，所以 TPOT$\approx$耗时÷t + 2$\times$层数$\times$All-Reduce(批$\times$d$\times$2 字节)。\par}
\item {\fontsize{10}{14.8}\selectfont Llama-3.1-70B、批 64：TP=4 的 TPOT 15.3 ms、每卡 1043 token/s，TP=8 降到 10.0 ms、每卡 803 token/s；卡数翻倍延迟只降到三分之二，每卡吞吐降 23\%，扩展效率 77\%。\par}
\item {\fontsize{10}{14.8}\selectfont 原因是通信不随 t 缩小：每层 2 次 All-Reduce 的次数不变，每次的数据量（批大小$\times$d）也不变，延迟项 $\alpha$ 一分不少；TP=8 的 10.0 ms 里约三成是通信，批越小这个比例越高。\par}
\item {\fontsize{10}{14.8}\selectfont 8B 批 64：TP=1 是 9.9 ms、6450 token/s，TP=8 是 2.0 ms、4094 token/s；8 份 TP=1 总吞吐约 51,600 token/s，一份 TP=8 仅 32,750 token/s。\par}
\item {\fontsize{10}{14.8}\selectfont 选型顺序是：先保证放得下（70B BF16 至少要 TP=4），再看下限 TP 上 TPOT 是否满足 SLO，满足就别再加，剩下的卡做多份副本提吞吐；TP 大了每卡 KV 空间也更大，一份副本能服务更多请求。\par}
\item {\fontsize{10}{14.8}\selectfont prefill 计算受限，TP 对 TTFT 帮助很大：70B 一个 8192 token 提示，TP=2 是 621+48=669 ms，TP=8 是 155+90=245 ms；但消息是整段提示的隐藏状态，实际要和下一层计算重叠。\par}
\end{enumerate}\endgroup
\numline{公式}{$\text{TPOT}(t)\approx\frac{\text{单卡耗时}}{t}+2\times\text{层数}\times\text{All-Reduce}(\text{批大小}\times d\times2\ \text{字节})$}
\numline{数字}{8B 批 64：TP=1 9.9 ms/6450 tok/s，TP=8 2.0 ms/4094 tok/s；70B TP=8 每卡显存 23.0 GB；TP 4$\rightarrow$8 扩展效率约 77\%；70B 8192 token TTFT：669$\rightarrow$385$\rightarrow$245 ms；8 份 TP=1 共约 51,600 tok/s}
\fol{批大了结论也会变吗？}{会。批大时 KV 读取和计算占主导，通信相对变小，TP 效率好一些；批 1 时延迟项占主导，70B 从 TP=4 到 8 每卡吞吐只从 23 降到 21 token/s。}
\fol{放不下只能加 TP 吗？}{不是。可以先把权重量化到 FP8 或 INT8 让模型变小，也可以用流水并行把层分到别的机器；TP 只是把每一层的计算摊开这一种办法。}
\fol{MLA 模型怎么并行？}{MLA 把所有头的 KV 压成一个共享潜向量，没法按头切，TP 越大每卡都要存完整 KV、越浪费；SGLang 改用 DP 注意力，注意力各卡处理不同请求，MoE 部分才用 EP。}
\numline{别答错}{说「TP 让推理变快」要指明是延迟变快，每卡吞吐和扩展效率其实都在下降。}
\phantomsection\addcontentsline{toc}{section}{111\quad MoE 推理与大规模专家并行}
\lec{111}{MoE 推理与大规模专家并行}
\askline{大规模 MoE 推理时间花在哪？}
\askline{为什么换更快的 GPU 没用？}
\askline{EP 越大负载越不均吗？}
\vspace{3pt}
\noindent{\qlabel{答}}\begingroup\sloppy\begin{enumerate}[leftmargin=1.6em,itemsep=2.5pt,topsep=0pt,parsep=0pt,label={\sffamily\footnotesize\color{accent}\arabic*.}]
\item {\fontsize{10}{14.8}\selectfont 专家算得很快：DeepSeek-V3 每个专家约 4400 万参数，EP=64 时每卡每层放 4 个专家，FP8 权重约 176 MB，读一遍约 0.053 ms，1024 个 token 副本的计算更快，每层专家计算约 0.053 ms。\par}
\item {\fontsize{10}{14.8}\selectfont 慢的是网络：每个 token 要发 8 份 7168 维隐藏状态（FP8）、收 8 份（BF16），每层约 172 KB；一层约 22 MB，50 GB/s 下 All-to-All 0.45 ms，是计算的八倍多。\par}
\item {\fontsize{10}{14.8}\selectfont 58 层累积下来一步约 31 ms，其中通信约 26 ms；带宽从 50 提到 450 GB/s，每卡吞吐从 4560 涨到 22,371 token/s、近 5 倍，而换更快的 GPU 几乎没用，因为它本来就大部分时间在等数据。\par}
\item {\fontsize{10}{14.8}\selectfont 用 roofline 看：每个 token 专家计算约 7 亿次运算、网络流量 172 KB，每传 1 字节有约 4096 次运算；H100 的 FP8 算力除以 50 GB/s 是 39,580 次/字节，差近 10 倍，卡在网络。\par}
\item {\fontsize{10}{14.8}\selectfont 硬上限：每个 token 要走约 10.0 MB（58 层乘 172 KB），和批大小、EP 规模无关，每卡吞吐 $\leq$ 带宽 ÷ 10.0 MB，50 GB/s 时 5011 token/s、450 GB/s 45,100 token/s。\par}
\item {\fontsize{10}{14.8}\selectfont EP 越大越不均：每卡专家越少，专家之间的忙闲越不能互相抵消。一个 4 倍热点专家在 EP8 只让最忙卡多 8\%，EP256 就是 4 倍；完全随机路由 EP8 多 4\%、EP256 多 53\%，所以大 EP 必须配冗余专家和动态负载均衡。\par}
\end{enumerate}\endgroup
\numline{公式}{每卡吞吐 $\le\frac{\text{网络带宽}}{\text{每 token 每步流量}}$（$58$ 层 $\times172\ \text{KB}\approx10.0\ \text{MB}$）}
\numline{数字}{每层专家 0.053 ms vs All-to-All 0.45 ms；一步 31.1 ms、通信约 26 ms；25$\rightarrow$450 GB/s 时 2405$\rightarrow$22,371 tok/s；每传 1 字节 4096 次运算 vs 39,580；上限 5011/45,100 tok/s；冗余专家机内第 3 批回本、跨机第 23 批；CPU/GPU 交点 AVX-512 71、AMX 688 个 token}
\fol{为什么不少用几张卡？}{放不下也算不快：EP=8 时每卡每层 32 个专家，58 层的专家权重加非专家约 99 GB，超过 80 GB；专家集中后每卡读的权重变多，专家计算从 0.053 ms 变成 0.42 ms。}
\fol{冗余专家要多久回本？}{同一台机器走 NVLink 复制 7 份约 0.62 ms，热点每批省 0.24 ms，第 3 批回本；跨机复制要 5.32 ms，第 23 批才回本，热点只持续 16 批时反而亏 1.5 ms。}
\fol{专家该放 CPU 还是 GPU？}{看同一专家这一批收到多少 token：AVX-512 超过约 71 个、AMX 超过约 688 个就该搬到 GPU；更少时 CPU 读本地权重更快，KTransformers 就让 CPU 算。}
\numline{别答错}{MoE 省的是计算不是通信；加批次流量同比例涨、上限不变，只能摊薄读专家权重。}
\phantomsection\addcontentsline{toc}{section}{112\quad PD 分离}
\lec{112}{PD 分离}
\askline{为什么要做 PD 分离？}
\askline{PD 分离的代价是什么？}
\askline{什么场景下不值得分离？}
\vspace{3pt}
\noindent{\qlabel{答}}\begingroup\sloppy\begin{enumerate}[leftmargin=1.6em,itemsep=2.5pt,topsep=0pt,parsep=0pt,label={\sffamily\footnotesize\color{accent}\arabic*.}]
\item {\fontsize{10}{14.8}\selectfont prefill 计算受限、decode 带宽受限，混在一起用分块 prefill 时，一步里带 prefill 块就从约 5 ms 变成约 30 ms，正在 decode 的请求每遇到这样一步就卡一下，TPOT 尾延迟变差。\par}
\item {\fontsize{10}{14.8}\selectfont 混在一起也没法各取所需：prefill 喜欢大算力和更大的张量并行，decode 喜欢大带宽、大显存和更大的批，MoE 的 decode 甚至要几百张卡做专家并行，同卡部署只能选一套折中的配置。\par}
\item {\fontsize{10}{14.8}\selectfont 分开后是三部分：prefill 实例只做 prefill，能用大批；decode 实例只做 decode，每步平稳、TPOT 可预测；中间把请求的 KV Cache 从 prefill 显存传到 decode 显存，走 RDMA。\par}
\item {\fontsize{10}{14.8}\selectfont 收益是尾延迟：4 卡 H100 跑 8B、12 请求/s 时，混合部署 token 间隔 P99 23.7 ms，2P+2D 后 P99 5.8 ms、TTFT P99 77 ms，吞吐约 2360 token/s 不变。\par}
\item {\fontsize{10}{14.8}\selectfont 比例配错反而更差：1P+3D 只有一张卡做 prefill，长提示排队，TTFT P99 从 70 ms 涨到 120 ms；配比取决于请求形状，提示越长输出越短 P 越重，实际系统按流量动态调整两个池的大小。\par}
\item {\fontsize{10}{14.8}\selectfont 异构能提吞吐：4 张 A100 做 P（每请求 0.856 s）、4 张 H20 做 D（0.878 GPU 秒），上限 min(4.67, 4.55, 20.7)=4.55 请求/s，是共置 3.02 的 1.51 倍，对调仅 2.22。\par}
\end{enumerate}\endgroup
\numline{公式}{$\mu=\min(\mu_P,\ \mu_D,\ \text{链路带宽}\div\text{每请求 KV})$}
\numline{数字}{Qwen3-8B 每 token KV 144 KiB、8K 输入 1.125 GiB、25 GB/s 传 48.3 ms；2000 token 的 8B 模型 262 MB、约 5.2 ms；A100 prefill 0.856 s、decode 1.764 GPU 秒，H20 1.805 s、0.878；异构共置 3.02、PD 4.55、对调 2.22；同构 8$\times$A100 共置 3.05、3P5D 2.83、分块上限 4.53}
\fol{同构集群 PD 能提吞吐吗？}{不能。每个请求消耗的 GPU 秒没变，卡数又只能整数分，8 张 A100 共置 3.05、3P5D 反而只有 2.83 请求/s；提吞吐要靠分块 prefill 把空闲算力用起来。}
\fol{PD 会让 TTFT 更好吗？}{不一定。P 实例不够时长提示排队，1P+3D 的 TTFT P99 从 70 ms 涨到 120 ms；KV 传输也会让第二个 token 来得晚一点。}
\fol{和前缀缓存怎么配合？}{多轮对话里上一轮的 KV 在 decode 实例、下一轮的 prefill 在 P 实例，要命中缓存就得把 KV 放到两边都能访问的存储，或者把请求路由到有缓存的那个实例。}
\numline{别答错}{同构集群上 PD 分离不提高吞吐，每个请求 GPU 秒不变、卡只能整数分，它买的是隔离。}
\phantomsection\addcontentsline{toc}{section}{113\quad KV Cache 传输与分层存储}
\lec{113}{KV Cache 传输与分层存储}
\askline{已经算过的 KV，读回来还是重算？}
\askline{KV 能放到内存、SSD 上吗？}
\askline{缓存容量应该加多大？}
\vspace{3pt}
\noindent{\qlabel{答}}\begingroup\sloppy\begin{enumerate}[leftmargin=1.6em,itemsep=2.5pt,topsep=0pt,parsep=0pt,label={\sffamily\footnotesize\color{accent}\arabic*.}]
\item {\fontsize{10}{14.8}\selectfont 一段算过的前缀再用到时只有两条路：把 KV 从某一层存储读回来，或者重新 prefill。两者都和 token 数成正比，所以只需比每个 token 的代价，两者相等时的存储带宽就叫分界带宽，比它快的读回划算，比它慢的不如重算。\par}
\item {\fontsize{10}{14.8}\selectfont 分界带宽 = 每 token KV ÷ 每 token 重算时间。Llama-3-8B 是 8.6 GB/s，70B 是 2.3 GB/s，DeepSeek-V3 用 MLA 只有 1.9 GB/s：模型越大、KV 越小，门槛越低。\par}
\item {\fontsize{10}{14.8}\selectfont 按 8K 前缀算：KV 共 1.05 GB、重算 121 ms；本卡显存 0.3 ms、NVLink 2.3 ms、CPU 内存与 RDMA 21 ms、SSD 105 ms 都更快，只有远端存储 350 ms 不如重算。\par}
\item {\fontsize{10}{14.8}\selectfont 容量下限是 r $\times$ 复用间隔 T，r 正是分界带宽：Qwen3-8B 在 A100 上约 1.41 GB/s，HBM 只留得住 45 秒、内存 195 秒、SSD 2721 秒；10 秒间隔要 14.1 GB，10 分钟要 847 GB。\par}
\item {\fontsize{10}{14.8}\selectfont 读回可以和计算重叠：命中后只算 256 个新 token 约 30.9 ms，整份读回 48.3 ms；先读后算 79.2 ms，逐层预取 49.2 ms，再提前读入多出来的 436 MB（14 层），总时间回到 30.9 ms。\par}
\item {\fontsize{10}{14.8}\selectfont 两个边界：SSD 带宽够但 1 DWPD 每秒只能写 44.4 MB，是全量 KV 的 3.1\%，写入前要准入判断；容量也不是越大越好，Mooncake 线上记录里缓存从 1000 块加到 50000 块，命中率只从 30\% 升到 50\%。\par}
\end{enumerate}\endgroup
\numline{公式}{$\text{分界带宽}=\frac{\text{每 token KV}}{\text{每 token 重算时间}}$；$C\ge r\times T$；$T_{\text{写}}+k\cdot T_{\text{读}}<k\cdot T_{\text{重算}}$；$\text{缓冲}=\text{链路带宽}\times(T_{\text{读}}-T_{\text{算}})$}
\numline{数字}{分界带宽 8.6 / 2.3 / 1.9 GB/s（128 / 320 / 68.6 KiB 每 token）；8K 前缀 1.05 GB、重算 121 ms、CPU 读回 21 ms、SSD 105 ms、远端 350 ms；r = 1.41 GB/s；保留 45 / 195 / 2721 s；10 s 要 14.1 GB、600 s 要 847 GB；25 GB/s 复用 1 次回本、2 GB/s 3 次、1.25 GB/s 不回本；先读后算 79.2 ms、逐层 49.2 ms、提前读入 436 MB 即 14 层；SSD 44.4 MB/s 是 r 的 3.1\%；队列预取 86\% 对 LRU 58\%}
\fol{SSD 上放 KV 现实吗？}{带宽够，10 GB/s 高于 8B 的 8.6 GB/s 门槛，问题在写入寿命：1 DWPD 的盘每秒只能写 44.4 MB，所以只写被复用过的前缀和没结束的多轮会话。}
\fol{远端存储为什么反而不划算？}{因为分界带宽是硬门槛。8B 时远端只有 3 GB/s，读 8000 token 要 350 ms，比重算 121 ms 慢近三倍；换成 70B 重算要 1124 ms，远端 874 ms 就划算了。}
\fol{缓存容量该按什么估？}{先按 C $\geq$ r $\times$ T 定下限，再看命中率曲线是否已经饱和。80\% 的复用发生在 10 秒（消费级对话）或 10 分钟内（企业 API），容量分别只要 14.1 GB 和 847 GB。}
\numline{别答错}{SSD 瓶颈是寿命：1 DWPD 每秒只能写 44.4 MB，只有全量 KV 的 3.1\%。}
\phantomsection\addcontentsline{toc}{section}{114\quad 路由与负载均衡}
\lec{114}{路由与负载均衡}
\askline{请求进来，路由器凭什么选实例？}
\askline{缓存感知路由怎么兼顾负载？}
\askline{有缓存的实例排队很长，还发过去吗？}
\vspace{3pt}
\noindent{\qlabel{答}}\begingroup\sloppy\begin{enumerate}[leftmargin=1.6em,itemsep=2.5pt,topsep=0pt,parsep=0pt,label={\sffamily\footnotesize\color{accent}\arabic*.}]
\item {\fontsize{10}{14.8}\selectfont 大模型推理的路由有两难：把相关请求集中到同一实例能命中前缀缓存，但会造成热点；均匀分散能均衡负载，缓存却命中不了。所以要优化的是整体的吞吐和延迟，命中率只是手段。\par}
\item {\fontsize{10}{14.8}\selectfont 用 12 个应用、300 个会话、1288 个请求、4 个实例的模拟对比：轮询命中率 71.6\%、负载 1.01 倍；按会话固定 84.8\%、1.12 倍；按应用固定 92.1\%，但最忙实例是平均的 1.86 倍。\par}
\item {\fontsize{10}{14.8}\selectfont 缓存感知路由只在负载不超过平均 1.2 倍的实例里挑前缀命中最多的，命中率 96.8\%、负载 1.12 倍：热门应用的请求先集中，一旦过载就溢出到别的实例，在那里也建立起缓存。\par}
\item {\fontsize{10}{14.8}\selectfont 路由器判断谁有缓存有两种做法：自己维护一棵近似前缀树，不精确但不用和实例通信，SGLang 这么做；或者让实例上报块哈希到中心索引，更准但多一层通信，llm-d 这么做。\par}
\item {\fontsize{10}{14.8}\selectfont 不能只看命中率或空闲度，要比首 token 时间 = max(Q, R) + C。8K 前缀的例子里，本地命中排队 250 ms 是 281 ms，空闲实例重算 908 ms，从 200 GbE 远端取回只要 138 ms。\par}
\item {\fontsize{10}{14.8}\selectfont 远端取回要追平本地命中，链路至少 6.30 GB/s。而且缓存可能已经不在：p = 0.9 时平均 196.5 ms 看着很好，但那 10\% 的请求要 968 ms，p99 就由这条重算的慢路径决定。\par}
\end{enumerate}\endgroup
\numline{公式}{$\text{首 token}=\max(Q,\ R)+C$；$R=\text{查找}+\text{KV}\div\text{链路带宽}+\text{PCIe 搬运}$；优化目标是 goodput，命中率只是手段}
\numline{数字}{五种策略命中率/最忙倍数：71.6/1.01、71.4/1.00、84.8/1.12、92.1/1.86、96.8/1.12；四条路径首 token 281 / 908 / 282 / 138 ms；追平需 6.30 GB/s；串行 150 ms 对并行 90 ms；p = 0.9 均值 196.5 ms、慢路径 968 ms；200 GbE 下每秒 16 次取回占链路 77\%}
\fol{缓存感知路由会把热点推得更极端吗？}{会，所以要设负载上限。1.2 倍这条线就是防它：请求先集中到有缓存的实例，超过上限就溢出到别处重算一遍，换来最忙实例只有平均的 1.12 倍。}
\fol{只挑最空闲的实例行不行？}{不行。空闲实例要从头算 8448 个 token 约 888 ms；有缓存的实例即使排队 250 ms，出首 token 也只要 281 ms，差了三倍多。}
\fol{命中率高就说明路由好吗？}{不够。按应用固定命中率 92.1\%，但最忙实例是平均的 1.86 倍，整体延迟被它拖住；而且缓存可能已被淘汰，p99 要看重算的慢路径。}
\numline{别答错}{命中缓存不等于更快：有缓存的实例可能在排长队，要比的是首 token 时间，不是有没有命中。}
\phantomsection\addcontentsline{toc}{section}{115\quad 超节点上的推理}
\lec{115}{超节点上的推理}
\askline{超节点能让推理全面变快吗？}
\askline{什么负载受益最大，为什么？}
\askline{环面和交换网络怎么选？}
\vspace{3pt}
\noindent{\qlabel{答}}\begingroup\sloppy\begin{enumerate}[leftmargin=1.6em,itemsep=2.5pt,topsep=0pt,parsep=0pt,label={\sffamily\footnotesize\color{accent}\arabic*.}]
\item {\fontsize{10}{14.8}\selectfont 超节点只改变卡与卡之间的连法，单卡的显存带宽和算力一点没变。它把 scale-up 域从 8 卡放大到几十上百卡，机内 450 GB/s、延迟约 0.8 $\mu$s，跨机仍要走 50 GB/s 的网卡、延迟约 3.8 $\mu$s。\par}
\item {\fontsize{10}{14.8}\selectfont 估算还是那套：一步时间约等于每卡要读的字节除以显存带宽，再加卡间通信。Ring All-Reduce 走 2(p-1) 轮，T\_AR = 2(p-1) $\times$ ($\alpha$ + m/(p$\cdot$$\beta$))，NVLink 每轮 $\alpha$ 取 0.822 $\mu$s。\par}
\item {\fontsize{10}{14.8}\selectfont 放得进 8 卡的稠密模型几乎不受益。Qwen3-32B 一个 128K 会话：TP8 一步 5.15 ms；跨两台服务器的 TP16 是 6.67 ms；16 卡 NVLink 域里的 TP16 也有 5.64 ms，还是更慢。\par}
\item {\fontsize{10}{14.8}\selectfont 原因是两点：KV 头只有 8 个，权重能切到 1/16，KV 只能切到 1/8；而且一次 All-Reduce 的 11.55 $\mu$s 里带宽项只有 40 ns，耗时几乎全是轮数乘延迟，卡翻倍轮数从 14 变 30。\par}
\item {\fontsize{10}{14.8}\selectfont 放不进 8 卡的模型明显受益：Llama-3.1-405B 从 TP8 $\times$ PP2 换成域内 TP16，一步从 39.13 ms 降到 27.43 ms。流水并行能让模型放下、靠微批次补吞吐，但缩短不了单请求延迟。\par}
\item {\fontsize{10}{14.8}\selectfont 大规模 MoE 收益最大，而且主要来自容量：DeepSeek-V3 的 EP 组从 16 卡扩到 64 卡，每卡专家变少、会话从 34 个增到 87 个，一步仍约 24.8 ms，每卡吞吐从 1436 涨到 3504 tok/s。\par}
\end{enumerate}\endgroup
\numline{公式}{$T_{\text{AR}}=2(p-1)\times(\alpha+m/(p\cdot\beta))$；$\text{可用带宽}=\min(\text{链路带宽},\ u\cdot q/L)$}
\numline{数字}{TP8/跨机TP16/域内TP16：5.15 / 6.67 / 5.64 ms；追平需每次 All-Reduce $\leq$20.9 $\mu$s（域内 24.7、跨机 32.74）；405B 39.13$\rightarrow$27.43 ms；S=16/64/256：34/87/101 会话、1436/3504/3996 tok/s；64 卡分 8 台 2248、重叠 2754，域内重叠反降到 2815；单会话 5.96 ms、168 tok/s；环面 336 对交换 112 $\mu$s；用满 50 GB/s 需 1469 个在途请求}
\fol{超节点上 TP 更大一定更好吗？}{不一定。Qwen3-32B 的 TP16 在域内也比 TP8 慢，因为 KV 头不够分、小消息的轮数翻倍；要追平，每次 All-Reduce 得降到 20.9 $\mu$s 以下。}
\fol{MoE 的收益来自通信还是容量？}{主要是容量。每卡显存被塞满，一步几乎总在 24--25 ms，变的是这 80 GB 里有多少是 KV：S 从 16 到 64 会话翻 2.6 倍，吞吐也翻约 2.4 倍。}
\fol{环面为什么不适合大规模 MoE？}{All-to-All 目的地任意，环面里每份数据平均走 3 跳，链路负载按跳数放大：64 卡时 336 $\mu$s 对交换网络 112 $\mu$s，512 卡时 671 $\mu$s 对 112 $\mu$s。}
\numline{别答错}{超节点不是更快的 GPU：单卡带宽不变，MoE 变快主要靠容量变大，不是通信带宽变高。}
\phantomsection\addcontentsline{toc}{section}{116\quad SLO 与容量规划}
\lec{116}{SLO 与容量规划}
\askline{每秒 50 个请求，要几张卡？}
\askline{SLO 怎么定，又怎么变成卡数？}
\askline{扩容之后积压多久能消掉？}
\vspace{3pt}
\noindent{\qlabel{答}}\begingroup\sloppy\begin{enumerate}[leftmargin=1.6em,itemsep=2.5pt,topsep=0pt,parsep=0pt,label={\sffamily\footnotesize\color{accent}\arabic*.}]
\item {\fontsize{10}{14.8}\selectfont 先写 SLO：TTFT 的 P99 不超过 1 秒，TPOT 的 P99 不超过 30 ms（约每秒 33 个 token）。TPOT 的上限决定一个实例最多能同时 decode 多少个请求，也就决定了它能满载到什么程度。\par}
\item {\fontsize{10}{14.8}\selectfont 三件事联立：prefill 占掉 f 的 GPU 时间，decode 只能用 1 - f；Little 定律给出 b = $\lambda$ $\times$ 输出长度 $\times$ TPOT；TPOT = t(b) ÷ (1 - f)。不断加大 $\lambda$ 直到超限。\par}
\item {\fontsize{10}{14.8}\selectfont 8B、一张 H100、效率 60\%、输入 1000 输出 300：TPOT $\leq$ 30 ms 时每实例每秒 15.3 个请求，约 136 个同时在 decode，prefill 占 39\%，TPOT 29.7 ms。\par}
\item {\fontsize{10}{14.8}\selectfont 实例数 = $\lceil$总请求率 ÷ (单实例请求率 $\times$ 使用率)$\rceil$。50 请求/s、使用率 70\% 要 5 张 H100，每百万请求约 69 美元；TPOT 收紧到 15 ms 要 8 张，输入变 4000 要 18 张。\par}
\item {\fontsize{10}{14.8}\selectfont 非稳态更危险：启动 T\_s 秒留下 $\lambda$$\cdot$T\_s 的积压，排空 = $\lambda$$\cdot$T\_s ÷ ($\mu$ - $\lambda$)。3.5 请求/s、启动 10 s 积压 35 个，PD 的 4.55 请求/s 第 43.3 s 排空；效率掉一成就到 68.8 s。\par}
\item {\fontsize{10}{14.8}\selectfont 成本的分母应该是合格请求：8 卡集群满载时共置每千请求 0.74 元、PD 分离 0.49 元；但 PD 只有四分之一请求满足时限时，每千个合格请求涨到 1.95 元，反而更贵。\par}
\end{enumerate}\endgroup
\numline{公式}{$b=\lambda\times\text{输出长度}\times\text{TPOT}$；$\text{TPOT}=t(b)\div(1-f)$；实例数 $=\lceil\text{QPS}\div(\lambda\times\text{使用率})\rceil$；排空 $=\lambda\cdot T_s\div(\mu-\lambda)$}
\numline{数字}{15.3 请求/s、136 在 decode、prefill 39\%、TPOT 29.7 ms；TPOT 15 ms$\rightarrow$9.7/s、8 张；输入 4000$\rightarrow$4.0/s、18 张；200 请求/s$\rightarrow$19 张；5 张 H100 约 69 美元/百万请求；积压 35 个、排空 43.3 / 62.2 / 68.8 s；$\mu$ 至少 4.20；执行图 50,000 步持平；迁移 1.66 / 6.76 s；重配置 11.6 / 9.7 / 9.0 s；0.74 / 0.49 / 1.95 元}
\fol{为什么不能按单实例上限配卡？}{接近饱和时排队会让延迟急剧上升，P99 守不住；而且排空靠的是余量 $\mu$ - $\lambda$，PD 只比到达率高 8\%，效率差一成就从 43.3 s 拖到 68.8 s。}
\fol{扩容时的瓶颈在哪里？}{搬数据不是大头。TP4 换 TP8 走 50 GbE 要 2.62 s、200 Gbit/s 要 0.66 s，但建通信组、准备执行图要 9 s，总时间 11.6 和 9.7 s。}
\fol{想降成本应该先动哪里？}{先推高单实例请求率：前缀缓存让 prefill 占比变小，量化让一步更快，PD 分离让 TPOT 更稳；最便宜的一招往往是和业务确认 SLO 能不能放宽。}
\numline{别答错}{余量不是浪费：GPU 利用率 100\% 不代表容量够，要看 SLO 下的请求数。}
\chapter{手写迷你推理引擎}
{\footnotesize\sffamily\color{grayt}第 117--125 讲}\par\vspace{-4pt}
\phantomsection\addcontentsline{toc}{section}{117\quad 一个推理引擎的组成}
\lec{117}{一个推理引擎的组成}
\askline{一个推理引擎由哪几部分组成？}
\askline{你会按什么顺序把它写出来？}
\askline{怎么证明自己写的引擎是对的？}
\vspace{3pt}
\noindent{\qlabel{答}}\begingroup\sloppy\begin{enumerate}[leftmargin=1.6em,itemsep=2.5pt,topsep=0pt,parsep=0pt,label={\sffamily\footnotesize\color{accent}\arabic*.}]
\item {\fontsize{10}{14.8}\selectfont 骨架是一个循环：分词、进队列、调度、准备 KV 空间、前向、采样、更新状态，再回到调度，直到所有请求结束。前四章讲的各种优化，都是在这个循环的某一步上做文章。\par}
\item {\fontsize{10}{14.8}\selectfont 按职责分成五部分：模型执行（model.py）是唯一碰矩阵运算的地方，状态（kv\_cache.py）管 KV 存在哪，调度（scheduler.py）决定每步谁进 batch，内存管理管分配回收，加速手段让每步更省。\par}
\item {\fontsize{10}{14.8}\selectfont 优化点可以一一对上：continuous batching 改第 3、7 步，分页改第 4 步，量化改第 5 步用的权重，投机解码改第 5、6 步。所以引擎之间的差别不在数学，而在怎么组织调用。\par}
\item {\fontsize{10}{14.8}\selectfont 实现路线是六个版本，每版只解决上一版最大的一个毛病：v0 让模型跑起来，v1 加 KV Cache，v2 加 continuous batching，v3 换分页 KV，v4 加量化，v5 加投机解码，末版约 400 行。\par}
\item {\fontsize{10}{14.8}\selectfont 选 Qwen3-0.6B：28 层、1024 维、16 查询头配 8 个 KV 头、SwiGLU，权重 1.5 GB，笔记本放得下。用 CPU FP32 不是为快，而是为了能和 transformers 逐 token 比对。\par}
\item {\fontsize{10}{14.8}\selectfont 判断标准从一开始就定死：4 条提示（英文、代码、中文、数字序列）各贪心生成 32 个 token，每个都必须和 HuggingFace 一致，只有 v4 量化允许不同。每讲先过这一关，再测速度。\par}
\end{enumerate}\endgroup
\numline{公式}{循环：$\text{调度}\to\text{准备 KV 位置}\to\text{前向}\to\text{采样}\to\text{更新状态}\to\text{回到调度}$}
\numline{数字}{v0--v5 共 6 版、最后一版约 400 行 Python；28 层、1024 维、16 查询头/8 KV 头、权重 1.5 GB；4 条提示 $\times$ 32 token 逐一对齐；CPU FP32、Apple M2 16 GB}
\fol{为什么用真实模型而不是玩具？}{玩具模型看不出量化和投机解码的效果，也没法和 transformers 对齐。Qwen3-0.6B 结构完整（RMSNorm、RoPE、GQA、SwiGLU），权重 1.5 GB。}
\fol{为什么不直接用现成库？}{这一篇的目的就是把前向、KV 管理、调度自己写一遍，知道每一行在算什么；transformers 只在测试里生成参考答案，引擎本身不依赖它。}
\fol{每一版都要重跑对齐吗？}{要。掩码差一位、位置从 1 开始数、KV 写错位置，输出照样通顺却全错，只有逐 token 比对才暴露；所以每讲先过 4 条提示的检查，再测速度。}
\numline{别答错}{用 FP32 不是为快而是为对齐：掩码或位置差一位，输出照样通顺却全错。}
\phantomsection\addcontentsline{toc}{section}{118\quad 从权重文件到第一个token}
\lec{118}{从权重文件到第一个token}
\askline{1.5 GB 的权重文件里有什么？}
\askline{写 Qwen3 前向最容易错在哪？}
\askline{v0 为什么这么慢？}
\vspace{3pt}
\noindent{\qlabel{答}}\begingroup\sloppy\begin{enumerate}[leftmargin=1.6em,itemsep=2.5pt,topsep=0pt,parsep=0pt,label={\sffamily\footnotesize\color{accent}\arabic*.}]
\item {\fontsize{10}{14.8}\selectfont safetensors 很简单：文件头是一段 JSON，记录每个张量的名字、形状、数据类型和在文件里的位置，后面是连续存放的原始数据。Qwen3-0.6B 一共 311 个张量、752 M 个参数，BF16 存储。\par}
\item {\fontsize{10}{14.8}\selectfont 名字都能对上：嵌入表 [151936, 1024]；每层 q\_proj [2048, 1024]、k\_proj 与 v\_proj [1024, 1024]：q 输出 16$\times$128，k、v 只有 8$\times$128，这就是 GQA。\par}
\item {\fontsize{10}{14.8}\selectfont 要注意头维度不等于 d ÷ 头数：16 $\times$ 128 = 2048 比隐藏维 1024 还大。另外 q\_norm、k\_norm 各只有 128 个数，是 Qwen3 特有的，在 RoPE 之前对每个头做一次 RMSNorm。\par}
\item {\fontsize{10}{14.8}\selectfont lm\_head 和嵌入表是同一份权重，文件里存了两份，所以 752 M 里有 1.56 亿是重复的，实际约 6 亿参数，这就是 0.6B 的来历；PyTorch 把权重存成 [输出, 输入]，所以代码里到处是 x @ W.T。\par}
\item {\fontsize{10}{14.8}\selectfont v0 只有四个文件：config.py 33 行读 10 个尺寸，weights.py 13 行把 311 个张量读进字典并转 FP32，model.py 59 行写出完整前向，engine.py 32 行做贪心生成。\par}
\item {\fontsize{10}{14.8}\selectfont v0 每步都重算整段：32 个 token 用 6.95 s；提示 199 个 token 时第 1 步 419 ms、第 64 步 724 ms，共 34 s，总计算量约 15,000 个 token，真正新算的 263 个。\par}
\end{enumerate}\endgroup
\numline{公式}{$y=x\cdot W^{\top}$（权重存成 $[\text{输出},\text{输入}]$）；$q$ 输出 $16\times128=2048$，$k$、$v$ 输出 $8\times128=1024$}
\numline{数字}{311 张量、752 M 参数；词表 151936、d = 1024、FFN 3072；config 33 行、weights 13 行、model 59 行、engine 32 行；32 token 6.95 s、每 token 217 ms；199 token 提示第 1 步 419 ms、第 64 步 724 ms、共 34.12 s；重算约 15,000 个 token 对真正新的 263 个}
\fol{q、k 上为什么加 RMSNorm？}{Qwen3 在 RoPE 之前对每个头的 q、k 各做一次 RMSNorm，把 128 维的头归一化，让注意力打分更稳定；它只有 128 个可学习参数。}
\fol{为什么头维度不等于 d/头数？}{Qwen3 直接指定 head\_dim = 128，16 个查询头算出来是 2048，比隐藏维 1024 还大，所以不能按 d ÷ 头数去推；k、v 的头数又是 8。}
\fol{v0 到底慢在哪一步？}{慢在每步都把整段序列送进前向：第 n 步要算 199 + n 个 token，计算量随长度线性增长；64 步下来算了约 15,000 个 token。}
\numline{别答错}{RoPE 要前后两半配对、GQA 要用 repeat\_interleave，写错输出仍通顺却全错。}
\phantomsection\addcontentsline{toc}{section}{119\quad KV Cache}
\lec{119}{KV Cache}
\askline{KV Cache 到底缓存了什么？}
\askline{加上 KV Cache 能快多少？}
\askline{decode 每一步为什么这么贵？}
\vspace{3pt}
\noindent{\qlabel{答}}\begingroup\sloppy\begin{enumerate}[leftmargin=1.6em,itemsep=2.5pt,topsep=0pt,parsep=0pt,label={\sffamily\footnotesize\color{accent}\arabic*.}]
\item {\fontsize{10}{14.8}\selectfont 沿一层 Transformer 看：RMSNorm、投影、RoPE、SwiGLU 都是逐位置的，只有注意力跨位置，位置 p 的 q 要和所有 $\leq$ p 的 k 做点积。而 k、v 一旦算出来就不再变，所以每层只要存前面所有位置的 K 和 V。\par}
\item {\fontsize{10}{14.8}\selectfont 数据结构是每个请求一份连续内存，形状 [层数, 最大长度, KV 头数, 头维度]，K、V 各一份。Qwen3-0.6B 在 FP32 下每 token 224 KiB，一个 2048 token 的请求要 448 MiB，BF16 减半。\par}
\item {\fontsize{10}{14.8}\selectfont 注意力的改动只有两处：算出新 token 的 k、v 后先写进缓存，再从缓存读出前 end 个位置的全部；掩码从 T $\times$ T 的下三角改成按位置比较，查询位置 p 只看键位置 $\leq$ p。存的是旋转之后的 k。\par}
\item {\fontsize{10}{14.8}\selectfont 前向多了两个参数：这次新 token 的起始位置 start 和缓存本身；位置从 start 开始编号，输出层只算最后一个位置，因为只有那里要采样。kv\_cache.py 只有 18 行，model.py 加 10 行、删 7 行。\par}
\item {\fontsize{10}{14.8}\selectfont 生成分成一次 prefill 和之后每步一个 token 的 decode。199 token 提示生成 64 个：v0 34.12 s，v1 4.55 s，快 7.5 倍；第 1 步 419 对 357 ms，之后 63--66 ms。\par}
\item {\fontsize{10}{14.8}\selectfont 但这 65 ms 里有大半是在读权重：无论算 1 个还是 199 个 token，每步都要把 FP32 约 3 GB 权重读一遍，token 越少摊到每 token 越贵，这就是 decode 读内存受限。按最大长度一次分配也让空间被白占。\par}
\end{enumerate}\endgroup
\numline{公式}{每 token KV $=2\times\text{层数}\times\text{KV 头数}\times\text{头维度}\times4\ \text{字节}=2\times28\times8\times128\times4=224$ KiB}
\numline{数字}{v0 419 / 417 / 724 ms、共 34.12 s；v1 357 / 63 / 66 ms、共 4.55 s（1/7.5）；每 token KV 224 KiB、2048 token 448 MiB（BF16 减半）；kv\_cache.py 18 行、model.py +10/-7 行；FP32 权重约 3 GB}
\fol{为什么存旋转之后的 k，不存原来的？}{RoPE 只和 token 自己的位置有关，存旋转后的结果，读回来直接和 q 做点积，不用每次再转一遍；历史 token 的位置也不会变。}
\fol{掩码为什么要改成按位置比较？}{v0 每步送整段，用 T $\times$ T 下三角就行；v1 一次只送几个新 token，但键有 end 个，关系只能用 arange(end) $\leq$ pos 表达。}
\fol{decode 一步为什么 65 ms？}{decode 是读内存受限的：每一步都要把全部权重（FP32 约 3 GB）从内存读一遍，只为 1 个 token；prefill 一次算 199 个 token，同样读一遍权重。}
\numline{别答错}{按最大长度一次分配是这版的遗留问题：请求提前结束，没用到的空间照样占着。}
\phantomsection\addcontentsline{toc}{section}{120\quad Continuous Batching}
\lec{120}{Continuous Batching}
\askline{batch 里请求长度不同，怎么一起算？}
\askline{continuous batching 快多少？}
\askline{为什么引擎要拆成几块？}
\vspace{3pt}
\noindent{\qlabel{答}}\begingroup\sloppy\begin{enumerate}[leftmargin=1.6em,itemsep=2.5pt,topsep=0pt,parsep=0pt,label={\sffamily\footnotesize\color{accent}\arabic*.}]
\item {\fontsize{10}{14.8}\selectfont v1 是「一个请求从头跑到尾」，要同时服务多个请求必须拆开：每个请求用自己的 Sequence 记 token、已算位置数和 KV 缓存；调度器维护等待与运行两个队列；引擎只执行一步 step()。\par}
\item {\fontsize{10}{14.8}\selectfont 每一步：运行中的请求各 decode 一个 token；还有空位就从等待队列放新请求进来 prefill；生成完的立刻离开，槽马上给别人，这就是 continuous batching 与静态批的关键区别。\par}
\item {\fontsize{10}{14.8}\selectfont 模型这边把所有请求这一步的新 token 首尾相接拼成一维，记下每个请求占哪一段，这种格式叫 varlen。线性层、RMSNorm、SwiGLU 都是逐 token 的，一次算完；只有注意力必须按段分开。\par}
\item {\fontsize{10}{14.8}\selectfont 拼出来是这样：A 正在 decode、已有 20 个 token；B 刚进来、提示 6 个，拼成 7 个 token，pos 是 [19, 0…5]，segs 是 [(0,1,A), (1,7,B)]，算完各取自己最后一个位置。\par}
\item {\fontsize{10}{14.8}\selectfont 16 个请求（8--96 个 token）792 个：逐个跑 792 步、14.5 token/s；静态批 192 步、22.4 token/s；continuous batching 160 步、27.5 token/s，快 1.9 倍。\par}
\item {\fontsize{10}{14.8}\selectfont 看每步耗时：逐个跑约 69 ms，continuous batching 约 180 ms，请求是 5 倍而耗时只有 2.6 倍，多出来的请求共享了读权重那一遍；上限由 max\_num\_seqs 和 max\_num\_tokens 控制。\par}
\end{enumerate}\endgroup
\numline{公式}{batch 输入：$[\text{tokens}(N),\ \text{pos}(N),\ \text{segs}[(s,e,c)]]$，segs 说明每段属于哪个请求、用哪份 KV}
\numline{数字}{792 个 token（8--96 不等）；三种跑法 792 / 192 / 160 步；14.5 / 22.4 / 27.5 token/s；每步 69 ms $\rightarrow$ 约 180 ms（请求 5 倍、耗时 2.6 倍）；槽空置 48\% 对 38\%；8 槽 $\times$ 160 步 = 1280 个「槽$\cdot$步」，用了 792 个}
\fol{为什么只有注意力不能一起算？}{因为每个请求只能看到自己的上下文。线性层和归一化都是逐 token 的，拼起来互不影响；注意力要和各自那份 KV 做点积，必须按段切开。}
\fol{静态批处理差在哪里？}{整批要等最长的请求结束才能换下一批：两批里都有 96 token 的请求，各约 96 步、共 192 步，短请求的槽空着 48\%；continuous batching 只空 38\%。}
\fol{真实服务里空槽会更多还是更少？}{更少。这次 16 个请求很快全进了 batch，最后只剩几个长请求在跑；线上请求源源不断地来，空槽会被填上，吞吐还会更高。}
\numline{别答错}{v2 按最大长度预留整份 KV：声明 512 个却只生成 50 个，九成空间空着，挡住后面的请求。}
\phantomsection\addcontentsline{toc}{section}{121\quad 分页KV 管理}
\lec{121}{分页KV 管理}
\askline{为什么要给 KV 做分页管理？}
\askline{块表怎么把位置换成块？}
\askline{块用完了怎么办，会算错吗？}
\vspace{3pt}
\noindent{\qlabel{答}}\begingroup\sloppy\begin{enumerate}[leftmargin=1.6em,itemsep=2.5pt,topsep=0pt,parsep=0pt,label={\sffamily\footnotesize\color{accent}\arabic*.}]
\item {\fontsize{10}{14.8}\selectfont v2 给每个请求按「提示 + 最大输出」预留一整块连续的 KV：请求刚进来只用提示那一小段，后面的空间要等生成才用得上，可这块空间一开始就被占住了。\par}
\item {\fontsize{10}{14.8}\selectfont 于是能同时跑的请求数由「最大长度」而不是「实际用了多少」决定。改法是所有请求共用一个块池、每个请求一张块表，记下自己的 KV 在哪些块里。\par}
\item {\fontsize{10}{14.8}\selectfont 块池一次分配好，形状是 [层数, 块数, 块大小, KV 头数, 头维度]，K、V 各一份；块表第 j 项是「我的第 j 段 16 个位置放在池子的第几块」。\par}
\item {\fontsize{10}{14.8}\selectfont 位置换算只有一行：位置 p 在 table[p // 16] 块的第 p \% 16 格。块大小取 16，Qwen3-0.6B 在 FP32 下一个块（28 层合计）是 3.5 MiB。\par}
\item {\fontsize{10}{14.8}\selectfont 模型不用改：每个请求包一个 SeqKV，接口和 v1 的 KVCache 一样（write 和 read），内部按块表换算成块和格子，model.py 一行都不用动。\par}
\item {\fontsize{10}{14.8}\selectfont 调度器改成按空闲块放行：运行中的请求每步先保证下一个位置有块，拿不到就抢占运行队列里最晚进来的那个，还块、回队首重算；最先进来的永不被抢。\par}
\end{enumerate}\endgroup
\numline{公式}{位置 $p\to$ 块号 $\text{table}[\lfloor p/16\rfloor]$，块内偏移 $p\bmod 16$；块池 $\text{shape}=(28,\text{块数},16,8,128)$}
\numline{数字}{块大小 16；单块 28 层 $\times$ 16 $\times$ 224 KiB = 3.5 MiB；压测 20 块（320 位置）共 70 MiB；KV 只有 320 位置时 v2 46.5 s / 17.0 token/s / 240 步 / 每步 3.3 个请求，v3 34.7 s / 22.8 token/s / 167 步 / 每步 4.7 个 / 抢占 9 次（吞吐 +34\%）；块池缩到 6 块时发生 2 次抢占，4/4 条仍与 HuggingFace 逐 token 一致}
\fol{空间足够大时分页还有用吗？}{基本没用。v2 预留的空间本来就够，两者吞吐没有区别；价值只在 KV 成为瓶颈时显现，而真实服务里扣掉权重剩下的显存几乎总是瓶颈。}
\fol{抢占为什么不会算错？}{KV 只是缓存，扔掉了可以从 token 重算。被抢占的请求把提示和已生成的 token 一起重算，K、V 和扔掉的那份完全相同，6 块池抢了 2 次仍逐 token 一致。}
\fol{为什么专挑最晚进来的抢？}{最先进来的请求永远不被抢占，所以总有请求在往前走，不会出现互相抢、谁都完不成；如果最晚的就是自己，就抢自己，这一步不再给它排 decode。}
\numline{别答错}{分页省的是按最大长度预留的浪费，不是 KV 本身；空间足够大时 v2 和 v3 没有区别。}
\phantomsection\addcontentsline{toc}{section}{122\quad 前缀缓存}
\lec{122}{前缀缓存}
\askline{前缀缓存怎么认出同一个块？}
\askline{为什么只登记装满的块？}
\askline{命中了为什么只快一半多？}
\vspace{3pt}
\noindent{\qlabel{答}}\begingroup\sloppy\begin{enumerate}[leftmargin=1.6em,itemsep=2.5pt,topsep=0pt,parsep=0pt,label={\sffamily\footnotesize\color{accent}\arabic*.}]
\item {\fontsize{10}{14.8}\selectfont 真实服务里大量请求开头一样：同一个应用的系统提示、同一段多轮对话的历史、同一份被反复提问的文档。v3 每个请求都把提示从头算一遍，哪怕前面 300 个 token 一字不差。\par}
\item {\fontsize{10}{14.8}\selectfont 要复用一个块，直接比 KV 的数值太慢，我们比它的来源：一个块的 KV 由「它之前的全部 token + 它自己的 16 个 token」唯一决定，所以给每个装满的块算一个哈希。\par}
\item {\fontsize{10}{14.8}\selectfont 每一块的哈希都把前一块的哈希算进去，链式地连到开头，这样两个块哈希相同就说明它们的整个前缀都相同，而不只是这 16 个 token 碰巧一样。\par}
\item {\fontsize{10}{14.8}\selectfont 规则是只登记装满且算完的块，新请求从头逐块查、连续命中就继续、第一次没中就停；含最后一个 token 的块总自己算，不然拿不到最后位置的打分。\par}
\item {\fontsize{10}{14.8}\selectfont 引用计数让一个块被多个请求共用；没人用的块回空闲列表但登记不删，后来的请求还能命中，只有被拿去装别的内容时才作废，空闲列表就成了按释放先后淘汰的缓存。\par}
\item {\fontsize{10}{14.8}\selectfont 实测 8 个请求共用约 330 个 token 的系统提示，后 7 个提示共 2356 个 token，命中 2240 个位置、只 prefill 116 个，用时从 5.5 s 降到 2.4 s，输出完全一致。\par}
\end{enumerate}\endgroup
\numline{公式}{$h_0=\operatorname{hash}(\varnothing,\ \text{第 }0\text{ 块 16 个 token})$；$h_i=\operatorname{hash}(h_{i-1},\ \text{第 }i\text{ 块 16 个 token})$}
\numline{数字}{系统提示约 330 token；后 7 个请求提示共 2356 token；命中 2240 = 7 请求 $\times$ 20 块 $\times$ 16；只 prefill 116 个；5.5 s $\rightarrow$ 2.4 s；输出一致 True；每个请求仍要 decode 8 个 token}
\fol{时间为什么没少 95\%？}{prefill 省掉了，但每个请求还要 decode 8 个 token，这部分时间没变；而且命中的请求在 decode 时，注意力仍要读完整的 330 多个位置的 K、V。}
\fol{多轮对话也能命中吗？}{能。第二轮的输入是第一轮的全部加新问题，第一轮结束时它的块都还在池子里，引用计数归零但登记还在，第二轮进来就能命中。}
\fol{哈希碰撞会读到别人的 KV 吗？}{这是最怕的失效。链式哈希保证相同即整个前缀相同，nano-vllm 更进一步，命中后再比一遍块里存的 token，防止碰撞。}
\numline{别答错}{只登记装满、算完的块，且含最后一个 token 的块不复用，否则拿不到最后位置的打分。}
\phantomsection\addcontentsline{toc}{section}{123\quad 权重量化}
\lec{123}{权重量化}
\askline{W8A16 的 scale 怎么定？}
\askline{为什么按输出通道取 scale？}
\askline{量化了为什么反而更慢？}
\vspace{3pt}
\noindent{\qlabel{答}}\begingroup\sloppy\begin{enumerate}[leftmargin=1.6em,itemsep=2.5pt,topsep=0pt,parsep=0pt,label={\sffamily\footnotesize\color{accent}\arabic*.}]
\item {\fontsize{10}{14.8}\selectfont 只量化每层 7 个线性层（q、k、v、o 投影和 FFN 的三个矩阵），它们占了 28 层权重的绝大部分；嵌入表是查表、输出层决定 15 万词打分最敏感、RMSNorm 只有几千个数，都保持原样。\par}
\item {\fontsize{10}{14.8}\selectfont INT8 的对称量化用 -127 到 127 这 255 个整数，需要一个 scale 把格点映射回实数。这里每个输出通道（权重矩阵的每一行）一个 scale：s 等于这一行最大绝对值 ÷ 127。\par}
\item {\fontsize{10}{14.8}\selectfont scale 管得越细误差越小、存 scale 的开销越大：q\_proj 是 [2048, 1024]，就有 2048 个 scale，平均每 1024 个权重多存 4 字节，开销只有 0.4\%。\par}
\item {\fontsize{10}{14.8}\selectfont 按行取 scale 还有一个好处：缩放可以后置。y 的第 j 个输出 = x$\cdot$(q[j]$\times$s\_j) = (x$\cdot$q[j])$\times$s\_j，所以 y = (x$\cdot$qᵀ) $\odot$ s，不必先把整个矩阵还原成 FP32。\par}
\item {\fontsize{10}{14.8}\selectfont 内存收益很实：线性层 1680 MiB 降到 421 MiB（3.99 倍），整个模型从约 2.8 GiB 降到约 1.6 GiB；128 步里每一步最高分的 token 都相同，平均 KL 散度只有 0.0011。\par}
\item {\fontsize{10}{14.8}\selectfont 但速度反而慢了：CPU 上单请求 64 个 token，FP32 用 4.27 s、INT8 用 5.11 s。因为每次乘之前都要把 int8 矩阵转成临时 FP32，读 1 字节、写 4 字节、再读 4 字节，多一道工序。\par}
\end{enumerate}\endgroup
\numline{公式}{$s=\operatorname{amax}(\text{行})/127$；$q=\operatorname{clamp}(\operatorname{round}(W/s),-127,127)$；$y=(x\cdot q^{\top})\odot s$}
\numline{数字}{线性层 1680 MiB $\rightarrow$ 421 MiB（3.99 倍）；整模型约 2.8 GiB $\rightarrow$ 约 1.6 GiB；scale 开销 0.4\%（2048 个 scale / 1024 个权重）；128 步最高分 token 相同比例 100\%、平均 KL 散度 0.0011；FP32 4.27 s vs INT8 5.11 s（64 token）；4 条检查 3 条逐 token 一致，中文那条第 26 个 token 分叉；省下约 1.2 GiB 可多放约 360 个块}
\fol{逐 token 不一致就是量化坏了？}{不是。这个标准对量化太苛刻，要换成读同一串 token、比每一步打分：128 步最高分 token 全同、平均 KL 散度 0.0011，分布几乎重合，只是少数位置排名翻转。}
\fol{量化什么时候才真的变快？}{要 kernel 直接读 int8。GPU 上融合 kernel 把权重读进片上、在寄存器里换算成 BF16 再乘加，显存只读 FP32 的四分之一，decode 才能跟着快。}
\fol{省下来的显存能干什么？}{给块池。FP32 到 INT8 省下约 1.2 GiB，按一块 3.5 MiB 算能多放约 360 个块，同样显存就能同时服务更多请求。}
\numline{别答错}{没有直接读 int8 的 kernel 时，量化只省内存不省时间，CPU 上甚至慢约 20\%。}
\phantomsection\addcontentsline{toc}{section}{124\quad 投机解码}
\lec{124}{投机解码}
\askline{投机解码为什么结果不变？}
\askline{n-gram 草稿是从哪来的？}
\askline{被拒的草稿 KV 怎么删？}
\vspace{3pt}
\noindent{\qlabel{答}}\begingroup\sloppy\begin{enumerate}[leftmargin=1.6em,itemsep=2.5pt,topsep=0pt,parsep=0pt,label={\sffamily\footnotesize\color{accent}\arabic*.}]
\item {\fontsize{10}{14.8}\selectfont decode 一个 token 要 63 ms，而 prefill 199 个 token 只要 357 ms：一次前向多算几个 token 时间涨得很少，因为大头是把权重读一遍，投机就用这点本来闲着的算力。\par}
\item {\fontsize{10}{14.8}\selectfont 草稿用 n-gram 查找：很多任务会大段重复已有内容，取结尾 n 个 token（3 找不到退到 2、1）从近往远找相同片段，把它后面最多 k 个 token 当草稿；纯列表比较，比一次前向便宜几个数量级。\par}
\item {\fontsize{10}{14.8}\selectfont 验证时草稿接在最后一个真实 token 后面一起送进去：设最后一个真实 token 是 t、草稿是 d$_1$ d$_2$ d₃ d₄，送入的就是 5 个位置连续的 token，模型对每个位置各给一个贪心预测 p$_0$ 到 p₄。\par}
\item {\fontsize{10}{14.8}\selectfont p$_0$ 是看完 t 之后的预测，等于 d$_1$ 就是第一个猜对了；假设 d$_1$、d$_2$ 对了而 d₃ 错了，就接受 d$_1$、d$_2$ 再加上 p$_2$，一步新增 3 个 token；全猜错也还有 p$_0$。\par}
\item {\fontsize{10}{14.8}\selectfont 结果不变是因为每个被接受的 token，都是模型在「同样的前缀」下的贪心选择，和一步一步生成时选的完全相同；起草和验证只改变一步出几个 token，没改变出哪些 token。\par}
\item {\fontsize{10}{14.8}\selectfont 回滚不用专门删：这一步 5 个位置的 K、V 都写进了缓存，被拒位置写的是错的，但只把 num\_computed 推进到被接受的位置，那两格就算没写过、下一步被覆盖。\par}
\end{enumerate}\endgroup
\numline{公式}{送入 $t,d_1,\dots,d_k$ 得 $p_0,\dots,p_k$；$p_{i-1}=d_i$ 则接受 $d_i$，接受到第一个不一致处再加上该位置的 $p_j$}
\numline{数字}{decode 1 token 63 ms vs prefill 199 token 357 ms；改写代码 128 个 token：127 步 $\rightarrow$ 34 步，接受 93 个草稿，平均每次前向 3.66 个 token，9.4 s $\rightarrow$ 6.6 s；自由生成 53 步、2.37 个，9.2 s $\rightarrow$ 8.0 s；每步最多 4 个草稿；两种情况输出都一致}
\fol{草稿全错不是白算了吗？}{不白算。草稿全错时这一步仍然出 1 个 token，只是多算了几个位置；batch 小时多算的位置藏在读权重的时间里，几乎不花时间。}
\fol{步数降了四分之三，时间怎么只降三成？}{CPU 上每一步不再只算 1 个 token 而是 5 个，多算 4 个位置并不免费；GPU 上 decode 时算力大量闲置，多算才接近免费。}
\fol{什么时候该把投机关掉？}{batch 很大时每个请求多出的 k 个位置都要算，硬件已经很忙，多算不再藏在读权重里，工业引擎在高负载时会自动减少甚至关闭投机。}
\numline{别答错}{收益取决于草稿猜得准不准：照抄型任务每次前向 3.66 个 token，自由生成只有 2.37 个。}
\phantomsection\addcontentsline{toc}{section}{125\quad 时间花在哪里}
\lec{125}{时间花在哪里}
\askline{六个版本各换来了什么？}
\askline{现在的时间到底花在哪？}
\askline{离 vLLM 还差什么？}
\vspace{3pt}
\noindent{\qlabel{答}}\begingroup\sloppy\begin{enumerate}[leftmargin=1.6em,itemsep=2.5pt,topsep=0pt,parsep=0pt,label={\sffamily\footnotesize\color{accent}\arabic*.}]
\item {\fontsize{10}{14.8}\selectfont 六个版本约 400 行 Python：v0 手写前向，v1 加 KV Cache，v2 continuous batching，v3 分页与抢占，v3+ 前缀缓存，v4 INT8，v5 投机解码；最大增长是 v3 块管理器 100 行。\par}
\item {\fontsize{10}{14.8}\selectfont 每一版的收益都有前提：KV Cache 要生成长才划算，continuous batching 要请求足够多，分页要 KV 紧张，前缀缓存要请求真共用前缀，投机要草稿猜得准。\par}
\item {\fontsize{10}{14.8}\selectfont 实测的台阶是：v1 把同一任务从 34.1 s 降到 4.55 s（7.5 倍），v2 吞吐 27.5 token/s 是逐个跑的 1.9 倍，v3 在 KV 紧张时再提高 34\%，v3+ 省掉 95\% 的重复 prefill。\par}
\item {\fontsize{10}{14.8}\selectfont 用 v5 关掉投机跑 8 个请求各 64 个 token，共 13.1 s、64 步：线性层等张量运算 10.0 s 占 76\%，注意力 3.1 s 占 24\%，Python 侧调度、拼 batch、采样只有 26 ms、占 0.2\%。\par}
\item {\fontsize{10}{14.8}\selectfont 注意力比它「应该」占的多，是两个为代码清楚付出的代价：按请求用 Python 循环逐个算，8 个请求就是 8 次小矩阵运算；SeqKV.read 每层每步都把块复制成连续内存。\par}
\item {\fontsize{10}{14.8}\selectfont 单请求 64 个 token 时 HuggingFace 4.15 s、miniinfer 4.76 s，几乎一样：transformers 的 generate 也是带 KV Cache 的逐 token 生成，优势要到多请求才显现。\par}
\end{enumerate}\endgroup
\numline{公式}{每步时间 $\approx\frac{\text{权重字节}}{\text{带宽}}$（线性层）$+\frac{\text{该步 KV 字节}}{\text{带宽}}$（注意力），Python 侧约几十微秒}
\numline{数字}{v0 119 行 $\rightarrow$ v5 407 行；v1 34.1 s $\rightarrow$ 4.55 s（7.5 倍）、decode 每步约 65 ms；v2 27.5 token/s（1.9 倍）；v3 +34\%（22.8 vs 17.0 token/s）；v3+ prefill 2356 $\rightarrow$ 116 个（5.5 s $\rightarrow$ 2.4 s）；13.1 s / 64 步：注意力 3.1 s（24\%）、其余 10.0 s（76\%）、Python 26 ms（0.2\%）；HF 4.15 s vs 4.76 s}
\fol{GPU 上 Python 开销重要吗？}{CPU 上一步要 200 ms，几十微秒的调度可以忽略；GPU 上一步 decode 只要几毫秒，同样开销就占到可观比例，vLLM v0.6.0 的剖析里一度占六成。}
\fol{和工业引擎最大的差距在哪？}{kernel、发射开销和服务化：注意力还逐请求循环并复制 KV，写 KV 逐位置下标，量化要转回 FP32；工业引擎用变长 FlashAttention、按块表读 KV 和融合 kernel。}
\fol{怎么验证优化真的有效？}{每一步拆开计时再看 roofline：线性层时间应约等于权重字节除以带宽，注意力约等于 KV 字节除以带宽；偏差超过 30\% 就说明有复制、循环或发射过多。}
\numline{别答错}{量化没变快不是量化的错，是缺一个直接读 int8 的 kernel，这也是和工业引擎差距最大的一处。}
\chapter{解读nano-vllm}
{\footnotesize\sffamily\color{grayt}第 126--128 讲}\par\vspace{-4pt}
\phantomsection\addcontentsline{toc}{section}{126\quad nano-vllm 的整体结构}
\lec{126}{nano-vllm 的整体结构}
\askline{nano-vllm 为什么把调度和执行拆开？}
\askline{它为什么一步只做一种？}
\askline{和 miniinfer 差在哪？}
\vspace{3pt}
\noindent{\qlabel{答}}\begingroup\sloppy\begin{enumerate}[leftmargin=1.6em,itemsep=2.5pt,topsep=0pt,parsep=0pt,label={\sffamily\footnotesize\color{accent}\arabic*.}]
\item {\fontsize{10}{14.8}\selectfont nano-vllm 用约 1200 行 Python 实现了 vLLM 的核心：分页 KV、前缀缓存、continuous batching、分块 prefill、CUDA Graph、张量并行，接口也和 vLLM 一样。\par}
\item {\fontsize{10}{14.8}\selectfont 结构上最大的差别是调度和执行分开了：miniinfer 的 engine.step() 一个函数里又问调度器、又拼 batch、又跑模型，nano-vllm 把它拆成调度器、执行器、引擎三段。\par}
\item {\fontsize{10}{14.8}\selectfont 调度器只决定这一步跑哪些请求、每个算几个 token，顺带管块，只碰 Python 对象不碰张量；执行器把名单变成 GPU 张量、跑模型、采样，每张卡一个。\par}
\item {\fontsize{10}{14.8}\selectfont 拆开的原因是张量并行：有 N 张卡就有 N 份执行器、每张卡一个进程，调度只需要一份放在主进程；rank 0 把名单广播给其余执行器，大家一起算同一步。\par}
\item {\fontsize{10}{14.8}\selectfont 调度规则一句话：能 prefill 就只 prefill，没有才 decode，一步要么全 prefill 要么全 decode。这样 decode 步形状规整、每请求正好一个 token，才能录成 CUDA Graph。\par}
\item {\fontsize{10}{14.8}\selectfont 块管理和 miniinfer v3 几乎一一对应（两个队列、按块放行、抢最晚的、链式哈希、引用计数、倒着还块），不同有三处：一步只做一种、超长提示分块 prefill、块大小 256。\par}
\end{enumerate}\endgroup
\numline{公式}{一步 $=$ 调度器（只碰 Python 对象）$\to$ 执行器（变成 GPU 张量、采样）$\to$ 把新 token 交回调度器}
\numline{数字}{nano-vllm 1195 行 vs miniinfer v5 407 行；调度一侧 247 行 vs 180 行，多出的约 800 行在执行侧（执行器 256、模型 479、注意力 66）；默认每步最多 16384 token、512 个请求、上下文 4096、用 90\% 显存、块 256；README：RTX 4070 + Qwen3-0.6B + 256 请求，vLLM 1361.84 vs nano-vllm 1434.13 token/s}
\fol{块大小 256 有什么代价？}{前缀缓存粒度变粗：两个请求要有 256 个 token 相同的前缀才能共用一块，末块平均浪费约 128 个位置；约 330 token 的系统提示只能命中 1 块，miniinfer 是 20 块。}
\fol{为什么不学 vLLM V1 混着跑？}{vLLM V1 把 prefill 和 decode 统一成「每个请求算几个 token」混在同一步，靠变长查询的注意力 kernel 兼顾；nano-vllm 优先保住 decode 规整、能录图。}
\fol{1200 行真能跑得过 vLLM？}{按作者 README 在 RTX 4070、Qwen3-0.6B、256 个随机长度请求上是 1434.13 对 1361.84 token/s，但这组数字没有第三方复现，规模也不代表生产场景。}
\numline{别答错}{块大小 256 是分页接口要求，代价是前缀缓存变粗、末块浪费约 128 位置。}
\phantomsection\addcontentsline{toc}{section}{127\quad nano-vllm 多做了什么}
\lec{127}{nano-vllm 多做了什么}
\askline{一步的请求怎么变成 GPU 输入？}
\askline{写 KV 为什么要自己写 kernel？}
\askline{decode 怎么录 CUDA Graph？}
\vspace{3pt}
\noindent{\qlabel{答}}\begingroup\sloppy\begin{enumerate}[leftmargin=1.6em,itemsep=2.5pt,topsep=0pt,parsep=0pt,label={\sffamily\footnotesize\color{accent}\arabic*.}]
\item {\fontsize{10}{14.8}\selectfont 启动时先用最坏情况的 prefill 预热一遍：一步的 token 预算（16384）和最长上下文（4096）取小，凑几个这么长的假请求跑一次，让 PyTorch 记下激活峰值，再把剩余显存全切成 KV 块。\par}
\item {\fontsize{10}{14.8}\selectfont 拼输入的产物是一维数组：input\_ids 和 positions 把这一步的 token 拼平，cu\_seqlens\_q/k 是前缀和，slot\_mapping 记每个新 token 写到哪个槽位，block\_tables 补齐成矩形。\par}
\item {\fontsize{10}{14.8}\selectfont 槽位号等于块号 $\times$ 块大小 + 块内偏移，等于把整个块池看成一维数组，所以一个 token 的写入就是一段连续拷贝，访存是合并的。\par}
\item {\fontsize{10}{14.8}\selectfont 写 KV 必须自己写 Triton kernel，关键在 -1：PyTorch 里 -1 表示最后一个，CUDA Graph 补齐的空位会把垃圾写进最后一块、覆盖别人的 KV；Triton 里一句 if slot == -1 就够。\par}
\item {\fontsize{10}{14.8}\selectfont 注意力 prefill 走变长的 flash\_attn\_varlen\_func，decode 走 flash\_attn\_with\_kvcache 直接按块表读 KV；块大小 256 正是这个分页接口的要求。\par}
\item {\fontsize{10}{14.8}\selectfont decode 按 batch 分档录图：1、2、4、8，再每 16 一档到 512，共 36 张；运行时挑不小于实际 batch 的最小一档，多出的行槽位填 -1、上下文长度填 0，等于空操作。\par}
\end{enumerate}\endgroup
\numline{公式}{槽位 $\text{slot}=$ 块号 $\times$ 块大小 $+$ 块内偏移；$\operatorname{argmax}_i\,p_i/E_i$（$E_i\sim\operatorname{Exp}(1)$ 独立）$\iff$ 按概率 $p_i$ 抽出 $i$}
\numline{数字}{36 张图（1、2、4、8，再每 16 一档到 512）；37 个请求用 48 档、补 11 行、多算 23\% 的行；块大小 256（assert \% 256 == 0）；16 查询头 + 2 $\times$ 8 KV 头 $\times$ 128 = 4096 行，双卡各 2048 行（8 查询头、4 K 头、4 V 头）；每层 2 次 All-Reduce，28 层 56 次 + 嵌入 1 次 + 输出 1 次 gather = decode 一步 58 次集合通信}
\fol{补齐出来的空位会不会算错？}{不会。空位写 KV 时槽位是 -1、kernel 直接 return，注意力里 context\_lens 为 0，等于空操作；decode 读权重受限，多算几行几乎不花时间。}
\fol{分档录图有什么坑要注意？}{档位从 16 起每 16 一档到 max\_num\_seqs，512 是 16 的倍数碰不到；设成 24 时 batch 17--24 找不到档位抛 StopIteration，补上 max\_bs 即可。}
\fol{张量并行的通信量大吗？}{每层两次 All-Reduce（o 投影、down 投影），28 层 56 次，加嵌入和输出层，decode 一步 58 次；每次只有 batch $\times$ 1024 个 BF16 数，主要是延迟。}
\numline{别答错}{块大小 256 是分页接口要求；空位靠槽位 -1、上下文长度 0 变成空操作。}
\phantomsection\addcontentsline{toc}{section}{128\quad 动手改 nano-vllm}
\lec{128}{动手改 nano-vllm}
\askline{草稿为什么不放进 token\_ids？}
\askline{验证该走哪条路，为什么？}
\askline{有温度时怎么接受草稿？}
\vspace{3pt}
\noindent{\qlabel{答}}\begingroup\sloppy\begin{enumerate}[leftmargin=1.6em,itemsep=2.5pt,topsep=0pt,parsep=0pt,label={\sffamily\footnotesize\color{accent}\arabic*.}]
\item {\fontsize{10}{14.8}\selectfont 草稿单独放在 Sequence 的 draft 字段里，不放进 token\_ids。token\_ids 牵动太多东西：num\_tokens、已生成数量、前缀缓存的哈希直接切 token\_ids，放进去再删就要求期间没人读它。\par}
\item {\fontsize{10}{14.8}\selectfont 验证走「带前缀缓存的 prefill」那条路：已缓存的是前面全部 token，这一步算最后一个真实 token 加 k 个草稿，拼出 cu\_seqlens 和块表走变长 FlashAttention；代价是不能用 CUDA Graph。\par}
\item {\fontsize{10}{14.8}\selectfont 输出层要在 Context 里加一个 all\_logits 标志：原来 prefill 只算每个请求最后一个位置的打分，验证要每个送入位置的打分，因为第 j 个位置的打分用来验证第 j+1 个草稿。\par}
\item {\fontsize{10}{14.8}\selectfont 有温度时接受规则反而简化了：n-gram 草稿是确定的，规则退化成「每个位置照常采样 x，x 等于草稿就接受，否则输出 x 并丢弃后面的位置」，分布不变。\par}
\item {\fontsize{10}{14.8}\selectfont 块要按 len + k 个位置来要，原来的 can\_append 假定每步只多一个 token；块不够先放弃草稿、再考虑抢占别人，验证后用 trim 把只装了被拒草稿的块还回去。\par}
\item {\fontsize{10}{14.8}\selectfont 回滚只让 num\_cached\_tokens 推进到有效位置，但必须先追加被接受的 token、再登记哈希，因为 hash\_blocks 用 num\_cached\_tokens 加 num\_scheduled\_tokens 决定登记哪些块。\par}
\end{enumerate}\endgroup
\numline{公式}{接受规则：每个位置采样 $x\sim p$，$x=d$ 就接受 $d$，否则输出 $x$ 并停止；缺块数 $=\left\lceil\frac{\text{len}+\text{extra}}{\text{块大小}}\right\rceil-\text{块表已有块数}$}
\numline{数字}{新增 132 行、改 8 个文件；spec.py 新增 27 行，调度器 +42/-11、执行器 +38/-3、块管理器 +17/-4、sequence +3/-2、context +3/-2、embed\_head +1/-1、config +1；miniinfer 只新增 1 个文件、改 4 个；验证时温度取 10$^-$⁵ 接近贪心，3 个请求各 256 token，每步最多 4 个草稿}
\fol{最容易漏掉的两处是什么？}{一是序列化要带上 draft，否则张量并行时其他卡草稿为空、形状对不上，All-Reduce 会卡住，单卡测试发现不了；二是 run\_model 用不用 CUDA Graph 要看 Context。}
\fol{怎么验证移植没有 bug？}{开和关各跑一遍，比前多少个 token 相同。GPU 上验证走 varlen、decode 走 kvcache 加图，BF16 有末位差异，多数请求全同、少数很靠后才不同；前几个就不同就是 bug。}
\fol{移植到 GPU 上一定更快吗？}{不一定。验证步不能用 CUDA Graph，每步多几百次发射，调度器还要为每个请求做一次 n-gram 查找；batch 大时 decode 接近计算受限，草稿位置不再免费。}
\numline{别答错}{改动多来自 nano-vllm 为快做的设计：调度执行分离、张量并行、decode kernel。}
\chapter{主流推理引擎}
{\footnotesize\sffamily\color{grayt}第 129--133 讲}\par\vspace{-4pt}
\phantomsection\addcontentsline{toc}{section}{129\quad vLLM 的架构}
\lec{129}{vLLM 的架构}
\askline{vLLM V1 为什么拆成三类进程？}
\askline{prefill 和 decode 怎么统一调度？}
\askline{异步调度怎么排下一步？}
\vspace{3pt}
\noindent{\qlabel{答}}\begingroup\sloppy\begin{enumerate}[leftmargin=1.6em,itemsep=2.5pt,topsep=0pt,parsep=0pt,label={\sffamily\footnotesize\color{accent}\arabic*.}]
\item {\fontsize{10}{14.8}\selectfont V1 把请求路径拆成三类进程：API 服务器做分词、反分词和流式输出，EngineCore 只做调度和驱动执行，每张卡一个 worker 跑模型，三类活各占一个 CPU 核，GIL 也只锁自己的进程。\par}
\item {\fontsize{10}{14.8}\selectfont 拆的理由是 v0.6.0 时 API 和调度一度占掉一步时间的六成，GPU 大部分时间在等 Python；17 万行里骨架很小，调度器 3569 行只回答跑谁、算几个 token、块不够、何时结束。\par}
\item {\fontsize{10}{14.8}\selectfont 调度器里没有 prefill 和 decode 两个阶段，每个请求只有两个数：已算好的位置数和应有的 token 数，每步把差额补给前者，新请求、分块 prefill、前缀命中、decode、投机全被这一句话统一。\par}
\item {\fontsize{10}{14.8}\selectfont 块哈希是 SHA-256 的链式哈希：父块哈希加本块 token，再加额外键（LoRA 编号、图片哈希、缓存盐），所以 token 相同但 LoRA 不同、图片不同的 KV 不能共用；跨机稳定才能共享 KV。\par}
\item {\fontsize{10}{14.8}\selectfont 空闲块是自写的双向链表，从中间摘一块是 O(1)、不新建 Python 对象；还块时分两类，没有哈希的放前面最先复用，有哈希的按 LRU 留在后面等命中。\par}
\item {\fontsize{10}{14.8}\selectfont 异步调度靠占位：先记下这一步每个 decode 请求会多出 1 个采样 token 加草稿数，把下一步当已有 token 排掉，真实 token 由 worker 从显存取，调度不必等采样回来。\par}
\end{enumerate}\endgroup
\numline{公式}{$\text{num\_tokens\_with\_spec}=$ 提示 $+$ 已生成 $+$ 草稿；每步差额 $=\text{num\_tokens\_with\_spec}-\text{num\_computed\_tokens}$}
\numline{数字}{vllm/v1 共 170876 行；调度器 3569 行；nano-vllm 全仓 1197 行；一步预算 1024 token；串行 17 ms、GPU 忙 71\%，重叠 max(12, 3+2)=12 ms、吞吐 1.42 倍}
\fol{长提示会不会吃光一步的预算？}{会。long\_prefill\_token\_threshold 默认 0 表示不限，一个 3000 token 的提示能把 1024 的预算占满；设成 512 就能给等待队列留出空间。}
\fol{不拆进程行不行？}{小并发下行，高并发不行。分词和反分词是纯 CPU 的 Python 代码，每秒上万 token 时和调度抢同一个 GIL；拆开后各占一个核。}
\fol{怎么验证异步调度生效了？}{对比开关 async\_scheduling 前后的 TPOT：GPU 一步 12 ms、CPU 5 ms 时，开着应是 12 ms，关掉是 17 ms，差的就是没藏住的 CPU 时间。}
\numline{别答错}{没有 prefill 阶段不等于不做 prefill，它只是不按阶段划分，用两个数的差额统一表达。}
\phantomsection\addcontentsline{toc}{section}{130\quad vLLM 的分布式服务}
\lec{130}{vLLM 的分布式服务}
\askline{一条请求在 4 个副本里怎么挑？}
\askline{MoE 的副本为什么要锁步？}
\askline{PD 分离时 KV 是怎么搬的？}
\vspace{3pt}
\noindent{\qlabel{答}}\begingroup\sloppy\begin{enumerate}[leftmargin=1.6em,itemsep=2.5pt,topsep=0pt,parsep=0pt,label={\sffamily\footnotesize\color{accent}\arabic*.}]
\item {\fontsize{10}{14.8}\selectfont 分布式服务是五层套起来：副本内用共享内存广播队列驱动几张卡、副本之间挑一个、MoE 副本之间锁步、PD 分离靠 KV connector 接，每经过一层多一次通信或一次打分。\par}
\item {\fontsize{10}{14.8}\selectfont 副本内是 MultiprocExecutor：调度结果写进共享内存广播队列，所有 worker 执行同一个方法，但只有最后一个流水线阶段里张量并行编号 0 的 worker 回结果，省掉重复序列化。\par}
\item {\fontsize{10}{14.8}\selectfont 副本之间分两种：稠密模型的 DP 副本彼此完全独立，按 DP=1 处理；MoE 的每个副本是 DPEngineCoreProc，因为每步前向都有 All-to-All，必须步调一致，空批次也要陪跑。\par}
\item {\fontsize{10}{14.8}\selectfont 挑副本看四个量：协调器每 100 ms 发布的负载快照、前端自己的本地计数、在途请求数下限和 KV 压力罚分；只看快照时一阵突发平均排队 497 ms，加上本地计数降到 123 ms。\par}
\item {\fontsize{10}{14.8}\selectfont MoE 副本之间用 2 元素张量做一次 All-Reduce：第 0 位报有没有未完成请求，SUM 大于 0 即逻辑或；第 1 位报想不想暂停，SUM 等于副本数才算全体同意。\par}
\item {\fontsize{10}{14.8}\selectfont PD 分离通过 connector 接：调度器一侧回答外部还有多少 token 的 KV、写进哪些块、何时释放，worker 一侧负责加载和保存；NIXL、Mooncake、LMCache 是不同的实现。\par}
\end{enumerate}\endgroup
\numline{公式}{$\text{sync\_dp\_state}$：2 元素张量做一次 SUM All-Reduce，$[0]=1$ 表示有未完成请求（$\text{SUM}>0$ 即或），$[1]=1$ 表示想暂停（$\text{SUM}=\text{dp\_size}$ 全体同意）}
\numline{数字}{两节点 8 卡、4 个副本 $\times$ TP 4；快照间隔最小 100 ms；dp\_sync\_interval 默认 16；平均排队 497 ms $\rightarrow$ 123 ms，95\% 分位 650 ms}
\fol{每 16 步才同步一次会卡住吗？}{不会卡，只会多陪跑。所有副本都空闲后最多还要多跑十几个空批次才停；换来的是正常运行时每 16 步才多一次跨机通信。}
\fol{负载均衡为何不看缓存命中？}{vLLM 这层只看忙不忙。缓存感知路由要放在更外面的路由器（llm-d、SGLang 的路由器），因为它要知道各实例缓存了哪些块；两层可叠加。}
\fol{博客里的 connector 还能用吗？}{不能，这个提交里已经删了，作用相同的是 example\_connector.py；而且现在只有 MoE 副本才锁步，读源码要以手上的提交为准。}
\numline{别答错}{别说 MoE 和稠密副本都锁步，现在只有 MoE 的副本才必须一起前向，稠密副本互不通信。}
\phantomsection\addcontentsline{toc}{section}{131\quad SGLang 的架构}
\lec{131}{SGLang 的架构}
\askline{SGLang 的进程结构与 vLLM 差在哪？}
\askline{基数树怎么做前缀缓存？}
\askline{零开销调度怎么藏住 CPU？}
\vspace{3pt}
\noindent{\qlabel{答}}\begingroup\sloppy\begin{enumerate}[leftmargin=1.6em,itemsep=2.5pt,topsep=0pt,parsep=0pt,label={\sffamily\footnotesize\color{accent}\arabic*.}]
\item {\fontsize{10}{14.8}\selectfont SGLang 也是三个进程，但中间不一样：每张卡一个调度器进程，调度器和执行模型的 worker 在同一个进程里，rank 0 收请求后广播给其余 rank。\par}
\item {\fontsize{10}{14.8}\selectfont 省掉的是每步把调度结果跨进程发给 worker 这一步，代价是每张卡都重复做一遍调度的 CPU 工作，而且各 rank 的调度结果必须完全一致，否则 KV 槽位对不上。\par}
\item {\fontsize{10}{14.8}\selectfont 前缀缓存换成基数树：每个节点存一段 token 和这段 token 的 KV 槽位，孩子按第一个 token 索引，同一段公共前缀只存一次；匹配从根往下走，停在中间就把节点劈成两段。\par}
\item {\fontsize{10}{14.8}\selectfont 淘汰从可淘汰的叶子按 LRU 弹堆：没孩子、没人在用的叶子先出，还槽位、删节点，父节点变成空闲叶子后再入堆，所以总是从末梢往根部推进，公共前缀最后才被淘汰。\par}
\item {\fontsize{10}{14.8}\selectfont 内存是两级池：req\_to\_token 是一个按请求数乘最长上下文的 int32 张量，第 r 行第 i 列是该 token 在 KV 池里的槽位号，粒度是 token 不是块，第 0 行留给补出来的空请求。\par}
\item {\fontsize{10}{14.8}\selectfont 零开销调度先发射这一步、再处理上一步的结果，i+1 步的输入直接引用显存里第 i 步的输出张量，token 不回 CPU；GPU 一步 12 ms、CPU 5 ms 时一步 17 ms 变 12 ms。\par}
\end{enumerate}\endgroup
\numline{公式}{命中 $=$ 沿树逐节点比 $\text{key}$ 的公共前缀；节点 $\text{value}$ 是该段 token 的 KV 槽位，$\text{lock\_ref}>0$ 不可淘汰}
\numline{数字}{GPU 一步 12 ms、CPU 共 5 ms：串行 17 ms（GPU 忙 71\%），重叠 12 ms、吞吐 1.42 倍；块哈希粒度 16--256 token，基数树 1 token；lpm 在等待队列超 128 时退化成 fcfs}
\fol{基数树比块哈希命中率更高吗？}{块足够小时差别不大。差别在工程上：基数树天然适合多轮对话、思维链分叉这类树状前缀；块哈希更简单，哈希跨机稳定，适合全局 KV 池。}
\fol{缓存感知调度是默认开的吗？}{不是。这个提交里默认策略是 fcfs，lpm 这类策略要显式打开；它能提高命中率，但会让前缀罕见的请求等得更久，默认值偏向公平。}
\fol{树上 lock\_ref 是干什么的？}{记有几个正在运行的请求在用这个节点，大于 0 就不能淘汰；加锁要一路加到根，因为祖先节点的槽位也被这个请求引用着。}
\numline{别答错}{别漏了默认策略是 fcfs，缓存感知排序要显式打开，而且它会让前缀罕见的请求等更久。}
\phantomsection\addcontentsline{toc}{section}{132\quad 其他引擎与规模化部署}
\lec{132}{其他引擎与规模化部署}
\askline{各家引擎各适合什么场景？}
\askline{你怎么快速判断引擎定位？}
\askline{大规模专家并行由什么拼成？}
\vspace{3pt}
\noindent{\qlabel{答}}\begingroup\sloppy\begin{enumerate}[leftmargin=1.6em,itemsep=2.5pt,topsep=0pt,parsep=0pt,label={\sffamily\footnotesize\color{accent}\arabic*.}]
\item {\fontsize{10}{14.8}\selectfont 选引擎先看它押注在哪一层：TensorRT-LLM 押 kernel、LMDeploy 押量化 kernel、llama.cpp 押哪里都能跑、MLC-LLM 押编译器，vLLM 和 SGLang 押通用的高并发服务。\par}
\item {\fontsize{10}{14.8}\selectfont TensorRT-LLM 只支持 NVIDIA，强在自家写的 FP8、FP4 矩阵乘和融合的注意力、MoE kernel，新卡上的低精度往往这里最先成熟；调度和 vLLM 大同小异，代价是部署链路更重。\par}
\item {\fontsize{10}{14.8}\selectfont LMDeploy 有两个引擎：C++ 的 TurboMind 押 4 位权重量化和 KV 量化的 kernel，纯 PyTorch 引擎更容易扩到新模型；llama.cpp 用 GGUF 和十几种 k-quant，目标是本地低并发。\par}
\item {\fontsize{10}{14.8}\selectfont MLC-LLM 用 TVM 为每种硬件生成 kernel，所以能到手机和 WebGPU；国产卡厂商一般给 vLLM、SGLang 提供插件或分支，选型先看自己的硬件上哪些引擎有官方支持。\par}
\item {\fontsize{10}{14.8}\selectfont 判断定位看三样：kernel 是自己写的、调的库还是编译器生成，决定在哪块卡上快；调度器是 Python 还是 C++、有没有分页 KV 和前缀缓存，决定并发吞吐；服务层有没有接口和监控，决定上线成本。\par}
\item {\fontsize{10}{14.8}\selectfont 大规模专家并行是四样拼起来：DP 注意力让各卡每步同步 batch、空闲卡跑空批次陪做 All-to-All；DeepEP 这类通信库；EPLB 把热门专家复制到多卡并周期重排；微批重叠藏住通信。\par}
\end{enumerate}\endgroup
\numline{公式}{大规模 $\text{EP}=\text{DP}$ 注意力 $+$ DeepEP 类 All-to-All $+$ EPLB $+$ 微批重叠}
\numline{数字}{TensorRT-LLM 只支持 NVIDIA、押 FP8/FP4；TurboMind 押 W4A16；llama.cpp 十几种 k-quant；MLC-LLM 用 TVM；四个主流引擎都具备 continuous batching、分页 KV、前缀缓存、多卡并行}
\fol{这些结论能当排名用吗？}{不能。它们来自公开文档和源码，是定性判断不是压测结果；引擎更新很快，选型要按官方文档和自己的压测来。}
\fol{为什么要做微批重叠？}{因为 All-to-All 时 GPU 在空等。把一个 batch 切成两个微批，一个通信时另一个计算，vLLM 叫 ubatching，SGLang 叫 two-batch overlap。}
\fol{大规模 MoE 的难点在哪？}{难在组合：DP 注意力要求各卡步调一致，微批重叠要求通信和计算时间配得上，EPLB 重排不能打断服务，PD 分离要求两边块大小和 KV 布局对得上。}
\numline{别答错}{别把这几个引擎的定位说成压测排名，四个主流引擎的调度骨架大同小异，差别只在押注方向。}
\phantomsection\addcontentsline{toc}{section}{133\quad 怎么读一个大型引擎的源码}
\lec{133}{怎么读一个大型引擎的源码}
\askline{17 万行的引擎从哪读起？}
\askline{怎么定位一个功能的实现？}
\askline{怎么验证自己真读懂了？}
\vspace{3pt}
\noindent{\qlabel{答}}\begingroup\sloppy\begin{enumerate}[leftmargin=1.6em,itemsep=2.5pt,topsep=0pt,parsep=0pt,label={\sffamily\footnotesize\color{accent}\arabic*.}]
\item {\fontsize{10}{14.8}\selectfont 先有一张地图再进代码，地图就是自己写过的小引擎：miniinfer 400 行、nano-vllm 1200 行；17 万行里和「一个请求怎么被处理」有关的只占一小部分，其余是模型实现和硬件适配。\par}
\item {\fontsize{10}{14.8}\selectfont 顺着一个请求走而不是顺着目录走：找到入口和主循环，只追它在哪里分词、进哪个队列、哪一行被调度、KV 块在哪里分配、新 token 在哪里接上、在哪里结束，每站只读到知道它做了什么。\par}
\item {\fontsize{10}{14.8}\selectfont 先读设计注释和数据结构，再读函数：调度器开头那段「没有 prefill 阶段」的注释比后面 300 行代码更重要；先看 SchedulerOutput、Request 有哪些字段，函数大多就是在改这些字段。\par}
\item {\fontsize{10}{14.8}\selectfont 定位功能从配置项搜起：long\_prefill\_token\_threshold 只出现在 3 个文件共 21 处，config 里定义它、默认 0 表示不限，engine 把命令行传进来，调度器用它；先读定义处文档串，再跳到使用处。\par}
\item {\fontsize{10}{14.8}\selectfont 验证理解有四级：最小配置跑（一张卡、小模型、关掉 CUDA Graph）、开 DEBUG 日志看每一步调度了多少 token、用 Profiler 抓一段 decode、只改一个参数做对照。\par}
\item {\fontsize{10}{14.8}\selectfont 压测有四个坑：离线吞吐和在线延迟是两回事，随机提示测不出前缀缓存，不控制输出长度就没有可比性、要固定长度，第一次请求要预热，而且要扫一个并发区间而不是只测一个点。\par}
\end{enumerate}\endgroup
\numline{公式}{定位一个功能：配置项定义（文档字符串与默认值）$\to$ 传递 $\to$ 使用处，读两头、略中间}
\numline{数字}{vLLM 17 万行；调度器 3569 行；long\_prefill\_token\_threshold 出现在 3 个文件共 21 处；第五章从 119 行 v0 到 400 行 miniinfer、1200 行 nano-vllm}
\fol{看不懂的代码怎么办？}{用 git log -L 或 git blame 找到引入它的提交，再看 PR 描述：大多数 PR 会写清为什么这样改、解决什么问题、压测数字是多少。}
\fol{怎么确认改动真的变快了？}{压测要能测出你预期的那个变化：只改一个参数，先预热，固定输出长度，扫并发区间画吞吐和延迟曲线，再和理论下限比较。}
\fol{想贡献代码从哪开始？}{从小的开始：文档、测试、good first issue、自己能复现的 bug；改接口或调度策略先开 issue 或 RFC 达成一致，PR 小而聚焦并附上数字。}
\numline{别答错}{不固定输出长度就没有可比性，不同引擎会在不同位置遇到结束符，吞吐数字是假的。}
\chapter{适配原理：把模型跑通}
{\footnotesize\sffamily\color{grayt}第 134--141 讲}\par\vspace{-4pt}
\phantomsection\addcontentsline{toc}{section}{134\quad 换一块卡要重做什么}
\lec{134}{换一块卡要重做什么}
\askline{换一块国产卡要重写整个栈吗？}
\askline{适配工作量怎么估算？}
\askline{兼容路线和重写路线差在哪？}
\vspace{3pt}
\noindent{\qlabel{答}}\begingroup\sloppy\begin{enumerate}[leftmargin=1.6em,itemsep=2.5pt,topsep=0pt,parsep=0pt,label={\sffamily\footnotesize\color{accent}\arabic*.}]
\item {\fontsize{10}{14.8}\selectfont 不用重写整个栈。模型和芯片之间是九层：模型代码、推理引擎、深度学习框架、图编译、算子库、kernel 语言与编译器、集合通信库、运行时与驱动、工具；上面四层站在接口上方，要改的是接口下方。\par}
\item {\fontsize{10}{14.8}\selectfont 拿 nano-vllm 量一下：1195 行里和 NVIDIA 绑定的只有 52 行、占 4.4\%，其中设备与运行时 23 行、自写 kernel 13 行、图捕获 7 行、通信 6 行、第三方 CUDA 算子库 3 行。\par}
\item {\fontsize{10}{14.8}\selectfont 52 行改动小、背后工作量大：23 行只是换个名字，前提是框架已经认识这块卡，背后是几千个算子和显存分配器的移植；最重的是 3 行 flash\_attn，要么等厂商移植，要么自己手写。\par}
\item {\fontsize{10}{14.8}\selectfont 两条路线差别在选哪层对齐 CUDA：兼容路线给一一对应的接口和 hipify 工具，上层 Python 几乎不改，工作量 1 层沿用、7 层翻译、1 层重写，代价是继承 warp 等于 32 这类假设。\par}
\item {\fontsize{10}{14.8}\selectfont 重写路线的硬件不是 SIMT，kernel 要显式安排搬运，CUDA 代码翻不过去；框架用 PrivateUse1 注册新设备名，所有写死 cuda 的地方都要改，工作量 1 层沿用、2 层翻译、6 层重写。\par}
\item {\fontsize{10}{14.8}\selectfont 顺序是跑通、跑对、跑快，不能颠倒：没有基线就分不清错误是优化引入的还是本来就有；适配周期也不由仓库里有多少行 cuda 决定，由这些行背后的库在新硬件上有多成熟决定。\par}
\end{enumerate}\endgroup
\numline{公式}{拐点（$\text{FLOP}/\text{字节}$）$=\frac{\text{峰值算力}}{\text{带宽}}$；decode 一步线性层的算术强度 $\approx M$}
\numline{数字}{1195 行中 52 行、4.4\%；设备运行时 23 行、自写 kernel 13 行、图捕获 7 行、通信 6 行、第三方库 3 行；兼容路线 1/7/1；NPU 路线 1/2/6；MPS 上去掉每层一次同步，decode 每步 55 ms 降到 49 ms}
\fol{哪一层最容易卡住？}{算子库和 kernel 两层。3 行 flash\_attn 是最重的一条；hipify 也翻不动三样：PTX 内联汇编、假设线程组宽度 32 的写法、CUDA 独有的库。}
\fol{为什么跑通跑对跑快不能颠倒？}{因为优化几乎每一步都在改变数值：融合改累加顺序、低精度改舍入。没有跑对的基线，就分不清错误是新引入的还是本来就有的。}
\fol{这分层图什么时候会翻转？}{三种情况：模型用了 MLA、MoE All-to-All 这类专有算子；负载偏向 decode，瓶颈在发射和同步；要多卡，拓扑和 P2P 在单卡阶段看不见。}
\numline{别答错}{别被 4.4\% 骗了，改动量很小，但适配周期由这些行背后的库在新硬件上的成熟度决定。}
\phantomsection\addcontentsline{toc}{section}{135\quad 国产与NVIDIA 硬件差异}
\lec{135}{国产与NVIDIA 硬件差异}
\askline{GPGPU 和 NPU 的结构差在哪？}
\askline{线程组 32 变 64 要改什么？}
\askline{换卡后 Roofline 形状会变吗？}
\vspace{3pt}
\noindent{\qlabel{答}}\begingroup\sloppy\begin{enumerate}[leftmargin=1.6em,itemsep=2.5pt,topsep=0pt,parsep=0pt,label={\sffamily\footnotesize\color{accent}\arabic*.}]
\item {\fontsize{10}{14.8}\selectfont GPGPU 走 SIMT：一个 kernel 启动成千上万个线程，每 w 个编成一个线程组同拍执行同一条指令，搬数据由线程自己写；NPU 放弃 SIMT，标量、向量、矩阵三类单元分工，搬运由专门单元负责。\par}
\item {\fontsize{10}{14.8}\selectfont 两条路线「把卡占满」的办法不同：GPGPU 靠块数，132 个计算单元要 132 个以上线程块，块不够就 split-K；NPU 只有几十个核、每核矩阵单元更大，块太小时吃亏的是搬运和计算无法重叠。\par}
\item {\fontsize{10}{14.8}\selectfont 线程组宽度 32 变 64 要改四处：跨通道归约步数从 5 变 6，表决掩码从 32 位换成 64 位，16$\times$16 的矩阵指令每通道累加器从 8 个变 4 个，块内 256 线程从 8 个组变 4 个组。\par}
\item {\fontsize{10}{14.8}\selectfont 片上存储从 228 KB 降到 64 KB，FlashAttention 一次能装的 K/V 从 304 行掉到 85 行：扫一遍 4096 长上下文的分块次数从 14 次变 49 次，每次都要重做一轮在线 softmax 更新。\par}
\item {\fontsize{10}{14.8}\selectfont Roofline 拐点等于算力除以带宽，三块卡算出来是 295、245、250，差别很小，因为算力和带宽往往一起涨；所以「卡在算力还是带宽」通常不变，decode 线性层都只跑到峰值的 5\% 到 7\%。\par}
\item {\fontsize{10}{14.8}\selectfont 真正变的是三件事：绝对带宽从 3.35 到 5.3 TB/s，同样模型的 TPOT 下限差 1.6 倍；达到上限的难度不同，效率因子要重测；有没有现成实现，注意力在别的卡上可能只有 Triton 版。\par}
\end{enumerate}\endgroup
\numline{公式}{拐点 $=\frac{\text{峰值算力}}{\text{带宽}}$；算术强度 $100$ 时可达 $\min(989,100\times 3.35)=335\ \text{TFLOPS}$（峰值的 $33.9\%$）}
\numline{数字}{H100 989 TFLOPS、3.35 TB/s、132 CU、32 宽、228 KB；GPGPU 示例 1300/5.3/304/64/64；NPU 示例 400/1.6/24/0/512；M=16 只跑到峰值 5.4\%--6.5\%；拐点 295/245/250；K/V 304 行$\rightarrow$85 行}
\fol{矩阵单元越大越好吗？}{不一定。M=16 的 decode 里指令形状要求补齐到 16 或 32 的倍数，库若选 256 行的输出块就有 93.8\% 的行在算补的零；海光为此专门写了小 M 的 kernel。}
\fol{什么时候这些差异会反过来？}{长上下文 prefill 时拼的是峰值算力和片上容量，缓冲大的设计反而占优；通信密集的 MoE 卡在互联和 All-to-All，卡内差异几乎不起作用。}
\fol{参数表上的峰值能直接比吗？}{不能。稀疏算力是稠密的两倍，FP16/BF16/FP8/INT8 的峰值依次翻倍；比两块卡要确认同一种数据类型、同样是稠密值。}
\numline{别答错}{峰值算力别直接比，稀疏值是稠密两倍、不同数据类型依次翻倍，要同类型同稠密。}
\phantomsection\addcontentsline{toc}{section}{136\quad 编程模型与CUDA 迁移}
\lec{136}{编程模型与CUDA 迁移}
\askline{CUDA 生态到底锁在哪？}
\askline{hipify 能翻到什么程度？}
\askline{非 SIMT 的卡怎么写 kernel？}
\vspace{3pt}
\noindent{\qlabel{答}}\begingroup\sloppy\begin{enumerate}[leftmargin=1.6em,itemsep=2.5pt,topsep=0pt,parsep=0pt,label={\sffamily\footnotesize\color{accent}\arabic*.}]
\item {\fontsize{10}{14.8}\selectfont 锁定不在语言，在三样东西：一是调好参数的算子库，接口换名字是一天的事、参数重调是几个月的事；二是写进 kernel 的硬件常数，warp 等于 32、228 KB 共享内存往往是没名字的字面量；三是框架外那圈工具。\par}
\item {\fontsize{10}{14.8}\selectfont 兼容路线给一套一一对应的接口和源到源工具：perl 版按正则替换、不需要能编译就能跑，适合粗扫；clang 版用 C++ 前端解析、更准但要求代码能被完整编译；工具翻完还要人工审查。\par}
\item {\fontsize{10}{14.8}\selectfont 一段 16 行的 kernel 里自动 6 行、要人工确认 7 行、必须重写 3 行：自动的是换名字的分配、发射和通信调用，重写的是 PTX 内联汇编、wmma 矩阵核心指令、cuDNN 调用。\par}
\item {\fontsize{10}{14.8}\selectfont 要人工确认的三类是静默错误：写死 threadIdx.x 取模 32 会把一个 64 通道波前当两个 warp 用，归约循环从 16 开始少一步让 RMSNorm 分母偏小，unsigned 掩码丢掉 32 到 63 号通道。\par}
\item {\fontsize{10}{14.8}\selectfont 真实项目里九成是机械替换，剩下一成集中在最热的 kernel：海光插件 4250 行里架构内建 11 处、按架构分支的宏 4 处、写死 64 共 3 处，全落在一个文件，为 decode 写的小 M 矩阵乘 kernel。\par}
\item {\fontsize{10}{14.8}\selectfont 重写路线没有线程编号：声明片上缓冲、让搬运单元搬进来、矩阵单元算、再搬回显存，然后排成双缓冲流水；每步都要手写搬运时机和粒度，换来的是编译器能在图尺度上做融合。\par}
\end{enumerate}\endgroup
\numline{公式}{归约步数 $=\log_2 w$：$w=32$ 时 5 步、$w=64$ 时 6 步；表决掩码需要 $w$ 位，$\text{uint32\_t}$ 只够 32 条通道}
\numline{数字}{16 行样例：自动 6、人工确认 7、必须重写 3；海光插件 4250 行：架构内建 11 处、架构宏 4 处、写死 64 共 3 处；归约 32 宽度 5 步、64 宽度 6 步}
\fol{迁移该按什么顺序做？}{先不迁 kernel，用框架通用算子跑出基线；再清点绑定写法、跑翻译工具清掉九成机械替换；然后人工核对线程组宽度、掩码、同步和共享内存容量；最后才重写汇编和矩阵指令。}
\fol{第 4 步为什么最容易被跳过？}{因为跑起来看着是对的。实用办法是先把逐层比对脚本准备好，改一个 kernel 就跑一次比对，立刻知道这层输出有没有偏。}
\fol{Triton 算第三条路吗？}{算。它写的是「一个程序处理一整块数据」、不写线程编号，同一份 Python 能编到不同架构，代价是能表达的东西比手写汇编少，块大小换硬件要重调。}
\numline{别答错}{这类错误最坑的是静默：编译通过、不报错、输出有值但是错的，要主动搜取模 32 和掩码。}
\phantomsection\addcontentsline{toc}{section}{137\quad 框架怎么接入新硬件}
\lec{137}{框架怎么接入新硬件}
\askline{PyTorch 怎么才算认识一块新卡？}
\askline{接新卡有哪两条路线？}
\askline{缺一个算子该怎么补？}
\vspace{3pt}
\noindent{\qlabel{答}}\begingroup\sloppy\begin{enumerate}[leftmargin=1.6em,itemsep=2.5pt,topsep=0pt,parsep=0pt,label={\sffamily\footnotesize\color{accent}\arabic*.}]
\item {\fontsize{10}{14.8}\selectfont 技术上就是在那张分发表里填出一列：每个 aten 算子都有一张按分发键排的表，一次调用进来按张量的设备和当前打开的功能算出走哪一行，再调那一行注册的函数。\par}
\item {\fontsize{10}{14.8}\selectfont 键分三类：设备键（CPU、CUDA、MPS、PrivateUse1）、功能键（Autograd、Autocast、量化）、还有兜底的复合键，所以让 PyTorch 认识一块新卡就是在几千张表里各填一列。\par}
\item {\fontsize{10}{14.8}\selectfont 第一种填法沿用 CUDA 键：用 HIP 把 PyTorch 整个重编一遍，仍注册到 CUDA 键，torch.cuda.is\_available() 为真、上层代码一个字不改；代价是报错和环境变量都写着 CUDA，机器里却没有卡。\par}
\item {\fontsize{10}{14.8}\selectfont 第二种注册新设备：厂商把预留的 PrivateUse1 键改个名（npu），自动得到 Tensor.npu() 这些接口，再把算子逐个注册上去；但显存分配器、流、事件、拷贝都要自己在 C++ 里实现，名字干净却要从零铺。\par}
\item {\fontsize{10}{14.8}\selectfont 接入不是补几千个：Qwen3 前向只用到 31 种 aten 算子，要在设备上真正执行的只有 23 种，复合算子只要底下基础算子齐了就自动可用，rms\_norm 三个设备键全空也照样能跑。\par}
\item {\fontsize{10}{14.8}\selectfont 缺算子有四条岔路：报错、回退 CPU、用现有算子组合、自己写 kernel，按先跑通再按占比优化的顺序选；注册自定义算子还要配一份只算形状的 fake 实现给 torch.compile 用，并把 mutates\_args 声明对。\par}
\end{enumerate}\endgroup
\numline{公式}{$x\,\text{@}\,w\to\text{aten::mm}\to$ 分发器查表（设备键 $+$ 功能键）$\to$ 该键上的实现 $\to$ 厂商算子库 $\to$ 运行时发射 kernel}
\numline{数字}{aten::mm 在 CPU、MPS、Meta、SparseCPU 各一行；Qwen3 前向 31 种 aten 算子、设备上真执行 23 种；FP8 矩阵乘回退 CPU 22 ms 对组合 1.5 ms}
\fol{rms\_norm 三个键都空怎么还能跑？}{它注册在复合键上，PyTorch 把它拆成 pow、mean、rsqrt、mul 这些基础算子再各自分发，所以基础算子齐了它自动可用，不需要厂商逐个实现。}
\fol{两条路线的成本差在哪？}{沿用 CUDA 键几乎白拿 CUDA 那一套，写死 torch.cuda 的第三方库直接能用；新设备名字干净，但分配器、流、事件、拷贝和工具链要从零铺，代码里所有 cuda 字样都要改。}
\fol{怎么确认调用真落到了设备上？}{用 torch.\_C.\_dispatch\_has\_kernel\_for\_dispatch\_key 查那个键有没有实现；也可以给 MPS 注册一份 mm，x @ w 就进你的函数，注销后又回自带实现。}
\numline{别答错}{注册自定义算子别漏副作用：原地写 KV Cache 却声明不改输入，图编译会优化掉它，KV 写丢。}
\phantomsection\addcontentsline{toc}{section}{138\quad 图模式与编译}
\lec{138}{图模式与编译}
\askline{一步 decode 要下发多少次算子？}
\askline{torch.compile 厂商接哪一段？}
\askline{图中断为什么是大事？}
\vspace{3pt}
\noindent{\qlabel{答}}\begingroup\sloppy\begin{enumerate}[leftmargin=1.6em,itemsep=2.5pt,topsep=0pt,parsep=0pt,label={\sffamily\footnotesize\color{accent}\arabic*.}]
\item {\fontsize{10}{14.8}\selectfont 一步 decode 是两千多次下发：28 层 $\times$ 79 个算子 = 2212 次，每次都要走分发、算形状、准备参数、交给驱动；CPU 按 8 $\mu$s 一次算要 17.7 ms，而 GPU 的活只有 13.3 ms，GPU 在等 CPU。\par}
\item {\fontsize{10}{14.8}\selectfont 整图下发只付一次入口开销 20 $\mu$s，一步 13.27 ms，比逐算子快 1.33 倍；新硬件上这笔账更糟，因为下发路径更长、为保险加的同步点更多、缺融合算子时算子更碎。\par}
\item {\fontsize{10}{14.8}\selectfont torch.compile 分三段：Dynamo 捕获成 FX 图、AOTAutograd 做函数化和分解、后端生成代码；厂商要接的只是第三段，一个后端就是一个函数，收 (gm, example\_inputs) 返回可调用对象。\par}
\item {\fontsize{10}{14.8}\selectfont 能不能拿到整张图取决于有没有图中断：v1 每层一次 int(pos[-1]) 把张量变成 Python 值，Dynamo 捕不出来，切成 27 张小图、24 处中断；把 end 改由调用方以整数传入，前向是一张 3319 个节点的图。\par}
\item {\fontsize{10}{14.8}\selectfont 形状是第二个坑：默认先按形状特化、发现变了就把那一维变成符号量，连续 prefill 8、16、24、8 加 decode 只编 3 张图；只支持静态形状的后端必须像 vLLM 那样把形状分档，用几张图覆盖所有长度。\par}
\item {\fontsize{10}{14.8}\selectfont 适配顺序是先 eager 跑通跑对，再用 torch.\_dynamo.explain 数中断、把热路径上的中断清到零，然后确认形状策略，最后才开图捕获并和 eager 比对数值。\par}
\end{enumerate}\endgroup
\numline{公式}{后端 $=f(\text{gm},\text{example\_inputs})\to$ 可调用对象；一步下发数 $=28\times 79=2212$}
\numline{数字}{2212 次下发 $\times$ 8 $\mu$s $\approx$ 17.7 ms 对 GPU 13.27 ms；整图快 1.33 倍；v1 27 张图 / 24 处中断，port 版 2 张图 3319 与 3320 个节点；形状实验共 3 张图}
\fol{图中断为什么不是小瑕疵？}{被切成 27 段的前向等于回到逐算子执行：段与段之间不能融合、不能跨段规划片上缓冲、不能一次下发，对依赖图模式的硬件来说等于没有图模式。}
\fol{形状变了会不会反复编译？}{Dynamo 先赌形状固定，一变就退回符号量，这是自动策略；后端支持符号形状才能真正一次编译，否则要像 vLLM 那样 padding 到固定档位，代价是编译时间和缓存内存。}
\fol{开图模式后新错误怎么冒出来？}{典型是原地写 KV Cache 的算子副作用没声明对，图把它优化掉或重排，症状是前几个 token 对、后面开始胡言乱语，排查办法是关掉图模式对比一次再逐段打开。}
\numline{别答错}{别把图中断当性能小瑕疵：它决定有没有图模式，段间不能融合也不能一次下发。}
\phantomsection\addcontentsline{toc}{section}{139\quad 精度对齐}
\lec{139}{精度对齐}
\askline{换卡后怎么算精度对上了？}
\askline{为什么不能要求逐位一致？}
\askline{输出乱码怎么定位到哪一层？}
\vspace{3pt}
\noindent{\qlabel{答}}\begingroup\sloppy\begin{enumerate}[leftmargin=1.6em,itemsep=2.5pt,topsep=0pt,parsep=0pt,label={\sffamily\footnotesize\color{accent}\arabic*.}]
\item {\fontsize{10}{14.8}\selectfont 目标不是逐位一致，而是每一处误差都能解释：要么归因到已知的硬件行为（累加顺序、中间精度、超越函数实现），要么归因到一个具体 bug，解释不了的差异就是还没找到的 bug。\par}
\item {\fontsize{10}{14.8}\selectfont 逐位一致本来就不成立：浮点加法不满足结合律，矩阵乘按 K 怎么切、分几组归约决定了加法顺序；BF16 的 FP32 累加在不同卡上可能是不同实现；exp、rsqrt 也是不同的近似多项式。\par}
\item {\fontsize{10}{14.8}\selectfont 检查分三级：逐层相对误差能定位到层，贪心 token 比对看累积效应，任务级评测最接近用户感受但一次要跑几十分钟；三级盲区不同，必须一起用。\par}
\item {\fontsize{10}{14.8}\selectfont 逐层比对要按位置分别算再取最差：那个 FP16 例子里整张量误差只有 4e-3，可 5 个位置里有 4 个已经坏了，最差位置是 1.0，取平均会把局部失效藏起来。\par}
\item {\fontsize{10}{14.8}\selectfont bug 在第 4 层：激活值 6544 的平方是 4.28e7，超过 FP16 上限 65504，x.pow(2) 变 inf 使整层输出为 0；平方和升到 FP32 后 FP16 落到 3e-3，比 BF16 的 2e-2 还小。\par}
\item {\fontsize{10}{14.8}\selectfont 阈值按精度定：FP32 1e-5、TF32 1e-2、BF16 5e-2、FP16 1e-2，而且要拿同精度的参考比，才分得开「精度的代价」和「硬件的差异」。\par}
\end{enumerate}\endgroup
\numline{公式}{相对误差 $=\frac{\lVert a-b\rVert}{\lVert b\rVert}$；$\text{RMSNorm}(x)=\frac{x}{\sqrt{\operatorname{mean}(x^2)+\varepsilon}}\odot w$}
\numline{数字}{MPS FP32 1e-6 量级且不累积；BF16 2e-2 不随层增；FP16 问题版第 4 层 1.1e0、修好 3e-3；激活 6544、平方 4.28e7、FP16 上限 65504}
\fol{BF16 误差 2e-2 算正常吗？}{不算。BF16 只有 8 位尾数，相对误差天生在 1e-2 到 1e-3 量级，关键看曲线是平的、不随层放大，而且贪心生成和参考逐 token 一致。}
\fol{长提示坏、短提示好是什么原因？}{指向位置相关的实现：RoPE 的角度表用 FP16 存，位置十万以上精度不够；或者因果掩码用 -inf，在某些卡上和乘法结合产生 NaN，症状是某些长度下崩掉。}
\fol{怎么把比对做成例行检查？}{固定几条提示（短、长、中文、代码）和贪心解码，参考答案存成文件，每次改动后跑逐层误差和 token 比对，阈值写进脚本越界就失败，上线前跑一次任务级评测。}
\numline{别答错}{别用整张量平均误差判断：范数大的位置会把局部失效拉回来，必须按位置或按 token 逐点算再取最差。}
\phantomsection\addcontentsline{toc}{section}{140\quad 算子缺口：盘点与自写}
\lec{140}{算子缺口：盘点与自写}
\askline{这块卡支持 PyTorch 吗？}
\askline{一个模型一步用哪些算子？}
\askline{缺的算子用哪种补法？}
\vspace{3pt}
\noindent{\qlabel{答}}\begingroup\sloppy\begin{enumerate}[leftmargin=1.6em,itemsep=2.5pt,topsep=0pt,parsep=0pt,label={\sffamily\footnotesize\color{accent}\arabic*.}]
\item {\fontsize{10}{14.8}\selectfont 这个问题没法用是或否回答：PyTorch 有几千个算子，卡上只实现了一部分，一个具体模型又只用到很小一部分，两个集合的差集才是算子缺口，所以答案要按模型一个个列。\par}
\item {\fontsize{10}{14.8}\selectfont 数算子要在分发器那一层记账：用 TorchDispatchMode 跑一次 prefill 和一步 decode，把每个 aten 调用记成计数器，顺便把「回退到 CPU」的警告捕获下来，那类警告平时被日志淹没，却是性能问题的源头。\par}
\item {\fontsize{10}{14.8}\selectfont 实测 Qwen3-0.6B 一步 decode 调用 3932 次、只用 31 种算子，其中 1711 次是视图类（只改形状和步长、不启 kernel），真正要在设备上算的是 23 种、2221 次。\par}
\item {\fontsize{10}{14.8}\selectfont 盘点时要把视图类、复合算子、原子算子分开，只有原子算子构成缺口，一个模型的最小算子集只有几十种；排序也不能看调用次数，mul 调了 590 次但都是逐元素，mm 只有 197 次却占 62\% 的时间。\par}
\item {\fontsize{10}{14.8}\selectfont 缺口几乎总集中在三处：量化的矩阵乘（\_scaled\_mm）、融合注意力（FlashAttention 内核）、MoE 的分组算子，基础算子和逐元素算子通常齐全。\par}
\item {\fontsize{10}{14.8}\selectfont 三种补法代价差很远：回退 CPU 22273 $\mu$s、用现有算子组合 1459 $\mu$s、不量化直接 BF16 只要 526 $\mu$s；两种补法的误差都是 3.78e-2，因为误差来自 FP8 量化本身。\par}
\end{enumerate}\endgroup
\numline{公式}{整体加速 $=\frac{1}{(1-p)+p/s}$，$p$ 是算子时间占比、$s$ 是它自身的加速倍数}
\numline{数字}{3932 次调用 / 31 种 / 视图 1711 次 / 要算 23 种 2221 次；mm 占 62\%、bmm 9\%、逐元素 8\%；22273 / 1459 / 526 $\mu$s；回退慢 15 倍}
\fol{用现有算子组合能省时间吗？}{不能。组合是把 FP8 反量化成 BF16 再乘，字节读少了却多一次展开、乘法仍按 BF16 算，省下的带宽被反量化吃掉：它保正确性，保不住收益。}
\fol{哪些 kernel 真值得自己写？}{按 Amdahl 算：矩阵乘占 62\%，写到两倍是 1.45 倍，值得；写 KV Cache 只占 5\%，只是因为按块表散写没有现成算子、没有替代品才写。}
\fol{「支持」的两个层次怎么理解？}{一是能跑出正确结果，二是有好实现、接近硬件上限。scaled\_dot\_product\_attention 在 MPS 上能跑但不是 FlashAttention，INT8 能跑也不代表跑在峰值上。}
\numline{别答错}{别拿调用次数排优先级：mul 调 590 次只占 8\%，mm 只 197 次却占 62\%。}
\phantomsection\addcontentsline{toc}{section}{141\quad 实操$\cdot$把 miniinfer 移植到非 CUDA}
\lec{141}{实操$\cdot$把 miniinfer 移植到非 CUDA}
\askline{把引擎迁到新卡分几步？}
\askline{一行 int(pos[-1]) 有多少代价？}
\askline{为什么 CPU 上 BF16 更慢？}
\vspace{3pt}
\noindent{\qlabel{答}}\begingroup\sloppy\begin{enumerate}[leftmargin=1.6em,itemsep=2.5pt,topsep=0pt,parsep=0pt,label={\sffamily\footnotesize\color{accent}\arabic*.}]
\item {\fontsize{10}{14.8}\selectfont 五步：先把设备收口到一个模块，跑通并和参考答案逐 token 比对，盘点算子确认没有东西悄悄回退到 CPU，逐个精度跑逐层比对，最后量性能并且先治同步点。\par}
\item {\fontsize{10}{14.8}\selectfont 收口很便宜：新增 device.py 共 31 行，engine 改 4 处、model 改 3 处；顺序是 CUDA $\rightarrow$ MPS $\rightarrow$ CPU，第一支在海光 DCU 上也真，因为那份 PyTorch 用 HIP 编译、仍自称 cuda。\par}
\item {\fontsize{10}{14.8}\selectfont 正确性用原脚本就能查：四条提示、每条 32 个 token，FP32 下全部与 HuggingFace 逐 token 一致，说明 23 种算子实现都对、1e-6 的数值差异不足以改变贪心选择、KV Cache 读写也没走偏。\par}
\item {\fontsize{10}{14.8}\selectfont FP32 跑通不等于别的精度跑通：换 FP16 直接生成一串句号，逐层比对指到第 4 层，还是 RMSNorm 平方和溢出，把平方和升到 FP32 之后 FP16 和 BF16 都能逐 token 一致。\par}
\item {\fontsize{10}{14.8}\selectfont 性能上最先修的不是 kernel 而是同步：v1 每层一次 int(pos[-1])，MPS 上下发 5.2 $\mu$s、拷回 139.9 $\mu$s，28 层共 3.9 ms；end 改成参数传入，decode 每步 56$\rightarrow$49 ms，值 13\%。\par}
\item {\fontsize{10}{14.8}\selectfont 还有两条反直觉：FP32$\rightarrow$FP16 decode 只快 1.5 倍不是 2 倍，因为固定开销占比上升（带宽 48 对 36 GB/s）；CPU 上 BF16 的 prefill 比 FP32 慢约 10 倍，因为没有原生 BF16 矩阵乘。\par}
\end{enumerate}\endgroup
\numline{公式}{decode 有效带宽 $=\frac{\text{权重字节}}{\text{每步耗时}}$；$28\times 139.9\ \mu\text{s}\approx 3.9\ \text{ms}$}
\numline{数字}{prefill 355$\rightarrow$170$\rightarrow$141 ms；decode 64.9$\rightarrow$55.5$\rightarrow$49.3$\rightarrow$32.9 ms；权重 596 M、FP32 2.38 GB、FP16 1.19 GB；同步 139.9 $\mu$s 对下发 5.2 $\mu$s（27 倍）}
\fol{decode 只快 14\% 为什么？}{decode 受带宽限制，而 M2 的 CPU 和 GPU 共用同一块内存，统一内存架构下换到 GPU 并不能换来更高带宽，所以速度由带宽定、不由谁在算定。}
\fol{移植完离上线还差什么？}{分页 KV、连续批处理、调度这些引擎层逻辑可以直接套 v3；还缺融合算子（现在一步多上千次下发）、FlashAttention、图模式后端、量化和多卡通信。}
\fol{为什么先修同步点再看 kernel？}{同步点是典型的改起来几乎没有成本、收益却百分之十几：把热路径上的 int(tensor)、.item()、tensor.cpu() 去掉，比调一个 kernel 划算得多。}
\numline{别答错}{别只跑通一个精度就宣布完成：FP32 全对而 FP16 溢出，每个精度都要单独跑一遍逐层比对。}
\chapter{调优原理：把模型跑快}
{\footnotesize\sffamily\color{grayt}第 142--147 讲}\par\vspace{-4pt}
\phantomsection\addcontentsline{toc}{section}{142\quad 在新硬件上做性能分析}
\lec{142}{在新硬件上做性能分析}
\askline{没有 Nsight 怎么找瓶颈？}
\askline{时间线上先看哪几样东西？}
\askline{decode 和 prefill 顺序为何反？}
\vspace{3pt}
\noindent{\qlabel{答}}\begingroup\sloppy\begin{enumerate}[leftmargin=1.6em,itemsep=2.5pt,topsep=0pt,parsep=0pt,label={\sffamily\footnotesize\color{accent}\arabic*.}]
\item {\fontsize{10}{14.8}\selectfont 工具名字换了但问题没变，时间线上只看三样：CPU 发射和设备执行两行的关系、最长的几个 kernel 占了多少、空隙的位置和规律。\par}
\item {\fontsize{10}{14.8}\selectfont 空隙的读法：规律的、每层一次的空隙通常是同步点，一次很长的空隙通常是主机侧在做慢操作，比如 Python 循环、重新分配显存、JIT 编译一个 kernel。\par}
\item {\fontsize{10}{14.8}\selectfont 找瓶颈的顺序取决于负载：decode 先看下发和同步，因为一步的 kernel 又短又密、固定开销最先成为瓶颈；prefill 反过来，kernel 本身很长、下发可以忽略，先看算力有没有用上。\par}
\item {\fontsize{10}{14.8}\selectfont 实测按算子汇总第一名不是矩阵乘而是 aten::copy\_：占 63.3\%、844 次，来自 RMSNorm 防溢出的 x.float() 与 .to(dtype) 和 RoPE 转换；按段汇总是注意力 47.5\%、FFN 21.6\%。\par}
\item {\fontsize{10}{14.8}\selectfont 异步后端上分段计时必须显式同步：kernel 排队执行，某次调用会在队列积压时被阻塞，时间就记到恰好等到的那次调用头上；代价是一步从 33 ms 涨到 54 ms，但比例可信。\par}
\item {\fontsize{10}{14.8}\selectfont 没有好 profiler 时有四招：分段计时加同步、删一段再量（差分）、改一个参数看斜率、数 kernel 乘单次下发开销算下限，这些都只依赖计时器。\par}
\end{enumerate}\endgroup
\numline{公式}{实际带宽 $=\frac{\text{读写字节}}{\text{耗时}}$；实际算力 $=\frac{2MNK}{\text{耗时}}$}
\numline{数字}{一步 25.4 ms（FP16、199 token）；每段同步后 61.0 ms；copy\_ 63.3\% / 844 次；注意力 47.5\%、FFN 21.6\%；下发小 add 3.6 $\mu$s 对 int(张量) 158.1 $\mu$s（44 倍）}
\fol{怎么判断一个 kernel 到没到极限？}{只靠三个数：耗时、读写字节、乘加次数。带宽受限的做到标称 80\% 以上接近极限，60\% 以下说明访存不连续、跨步或没对齐；算力受限的矩阵乘 70\% 以上算好。}
\fol{只让归约用 FP32 值不值得改？}{它快 11 $\mu$s，一步调 113 次省约 1.3 ms、大概 4\%，数值差 1\% 还在阈值内；但改的是和溢出安全直接相关的代码，我们选择不改，把它记进待办。}
\fol{一次调优记录该写什么？}{至少写清硬件与软件版本、模型与精度、负载、测的是什么、结果、理论下限与差距、这次改了什么、相比上一次；缺一项几周后就复现不出来。}
\numline{别答错}{别信不显式同步的分段计时：kernel 排队，时间会被记到恰好被阻塞的那次调用上。}
\phantomsection\addcontentsline{toc}{section}{143\quad 矩阵乘法与注意力的适配}
\lec{143}{矩阵乘法与注意力的适配}
\askline{厂商库为什么在 decode 最糟？}
\askline{小 M 矩阵乘怎么选分块？}
\askline{注意力怎么按负载分实现？}
\vspace{3pt}
\noindent{\qlabel{答}}\begingroup\sloppy\begin{enumerate}[leftmargin=1.6em,itemsep=2.5pt,topsep=0pt,parsep=0pt,label={\sffamily\footnotesize\color{accent}\arabic*.}]
\item {\fontsize{10}{14.8}\selectfont 反直觉的结论是厂商的矩阵乘库在 decode 时往往是最糟的选择：不是库写得差，而是它默认的分块为训练和 prefill 那种大矩阵调过，decode 的 M 只有几十，恰好落在它最不擅长的区间。\par}
\item {\fontsize{10}{14.8}\selectfont 一笔账：lm\_head 32224$\times$4096 BF16 权重 264 MB，库选 256$\times$256 块在 M=16 补 94\% 零，285 $\mu$s、0.93 TB/s；专用 kernel 153 $\mu$s、1.73 TB/s，快 1.9 倍。\par}
\item {\fontsize{10}{14.8}\selectfont 补零不只浪费格子（算 8.3 M 个，有用的只有 0.5 M），更关键的是改变访存：256 行的块只为 16 行数据服务，片上存储和寄存器都按 256 行配，能驻留的块和在途读请求变少，而带宽受限的 kernel 恰恰靠它们喂满。\par}
\item {\fontsize{10}{14.8}\selectfont 矩阵单元给 kernel 三个硬约束：指令形状与对齐（海光小 M kernel：M 在 [1,64]、N 只支持 256/512/1024/2048、K 是 128 的倍数）、数据排布（权重预转置、16 字节对齐）、每条通道的累加器数。\par}
\item {\fontsize{10}{14.8}\selectfont 注意力两种实现：prefill 用 FlashAttention 式，矩阵不落地、分块扫 K 和 V、在线更新 m 和 l；decode 用 Flash-Decoding，沿 KV 切分再归约；兜底 bmm+softmax+bmm。\par}
\item {\fontsize{10}{14.8}\selectfont 适配顺序：先通用实现拿基线，再换厂商库并重跑精度比对（累加顺序变了），然后按形状矩阵测（prefill 几个长度、decode M=1 到 64、几种上下文），最后只给占比高且离上限远的格子写专用 kernel。\par}
\end{enumerate}\endgroup
\numline{公式}{实际带宽 $=\frac{\text{权重字节}}{\text{耗时}}$；实际算力 $=\frac{2MNK}{\text{耗时}}$；算术强度 $\approx M$，$M=16$ 时远在拐点左侧}
\numline{数字}{权重 32224 $\times$ 4096 BF16 = 264 MB；M=16 时 285 对 153 $\mu$s、0.93 对 1.73 TB/s、14.8 对 27.6 TFLOPS；M=64 时 190 $\mu$s；容量 101,376 字节最多 82 行 Q，读写 404 MiB、在线更新 12,800 次}
\fol{大矩阵还需要专用 kernel 吗？}{不需要。M $\geq$512 的 prefill 和训练交给厂商库，大块复用充分、库里为这类形状调过；64 到 512 是过渡区要先测，只有 M $\leq$64 的 decode 才往往要专用 kernel。}
\fol{注意力移植最难在哪？}{不在算法而在分块参数：块大小由片上存储容量决定，容量换一块卡就变了；而且块大小和 KV 分页的块大小耦合，不匹配时要么报错要么性能很差。}
\fol{怎么确定该写哪个 kernel？}{做成一张表：行是形状（prefill 长度、decode 的 M=1 到 64、上下文长度），列是实现（库、厂商 kernel、自写），格子里填耗时和实际带宽，表一摆出来就不用争论了。}
\numline{别答错}{别拿大矩阵 TFLOPS 交差：服务里跑的是 M=1 到 64 的 decode 形状，要单独测。}
\phantomsection\addcontentsline{toc}{section}{144\quad 用 Triton 与厂商语言写 kernel}
\lec{144}{用 Triton 与厂商语言写 kernel}
\askline{同一个 RMSNorm 写法差几倍？}
\askline{Triton 真能写一次到处跑？}
\askline{什么时候用库、什么时候手写？}
\vspace{3pt}
\noindent{\qlabel{答}}\begingroup\sloppy\begin{enumerate}[leftmargin=1.6em,itemsep=2.5pt,topsep=0pt,parsep=0pt,label={\sffamily\footnotesize\color{accent}\arabic*.}]
\item {\fontsize{10}{14.8}\selectfont 两条分水岭决定四种写法：一是会不会写线程编号，框架组合连块都不管、Triton 以一个程序处理一整块为单位、厂商语言要写到通道级；二是搬运是否显式，Triton 的 tl.load 由编译器决定，块级 DSL 要你声明缓冲、写出搬运。\par}
\item {\fontsize{10}{14.8}\selectfont 框架组合把 RMSNorm 摊成五个 kernel（pow、mean、rsqrt、两次乘），张量被读写五遍、下发五次；带宽受限的算子这样写，等于白扔几倍的时间。\par}
\item {\fontsize{10}{14.8}\selectfont 实测融合 kernel 在 decode 形状（1 行、4 KB）上快 4.6 倍：时间几乎全是固定开销，下发五次变一次；prefill（199 行、800 KB）快 10.6 倍：张量从读写五遍变成一遍。\par}
\item {\fontsize{10}{14.8}\selectfont 三种写法的误差都在 5e-8 量级，融合 kernel 在片上用 FP32 累加，和组合写法数值行为一致；顺带把上一讲那个微优化量掉了：组合里只让归约升 FP32 只快 14\%，要改就改到位。\par}
\item {\fontsize{10}{14.8}\selectfont 量带宽有个自欺：prefill 那行的 141.3 GB/s 高于机器内存带宽，因为 800 KB 留在缓存里没走内存；换成 4096 行（32 MB）后降到 85.1 GB/s 才是接近内存带宽的真实值。\par}
\item {\fontsize{10}{14.8}\selectfont 选语言按算子分工：适配初期用框架组合跑通；逐元素、归一化、激活、RoPE 这类带宽受限算子用 Triton 或厂商语言；注意力优先厂商移植版；矩阵乘常规形状交给库，decode 的小 M 才手写；跨多家硬件维护才牺牲性能换可移植性。\par}
\end{enumerate}\endgroup
\numline{公式}{有效带宽 $=\frac{\text{读写字节}}{\text{耗时}}$；融合收益 $\approx$ 读写遍数之比（理论 $5\times$，实测 $10\times$）}
\numline{数字}{误差均 5e-8；decode 25.9 / 22.3 / 5.6 $\mu$s；prefill 122.2 / 125.4 / 11.5 $\mu$s；4096 行 3181.6 / 3169.5 / 394.3 $\mu$s；带宽 141.3 对 85.1 GB/s}
\fol{Triton 的折扣在哪两处？}{一是要有目标架构的后端，指令选择、寄存器分配、共享内存布局都要实现一份；二是参数要重调，块大小、num\_warps、流水级数在不同架构上的最优值不同，要在预热阶段 autotune 完。}
\fol{手写 kernel 的维护成本多大？}{要跟着架构代号走：海光那个小 M kernel 有三套分支，gfx928 和 gfx936 用一条内建指令、gfx938 用另一条、其他架构直接陷阱，每出一代新卡都要加一个分支。}
\fol{自写 kernel 最容易出什么错？}{精度问题，它是这类 bug 的主要来源。三条纪律：入口把假设写成检查（形状、对齐、类型不满足就走回退）、归约和同步逐步对照参考实现、每个新 kernel 配一个单算子对比测试。}
\numline{别答错}{别在组合写法里做微调：只值 14\%，写一个融合 kernel 值 4.6 倍，要改就改到位。}
\phantomsection\addcontentsline{toc}{section}{145\quad 低精度与量化的取舍}
\lec{145}{低精度与量化的取舍}
\askline{换一块卡，量化还能不能提速？}
\askline{没有原生 kernel 时量化还剩什么收益？}
\askline{W8A16 为什么反而比 FP32 慢四倍？}
\vspace{3pt}
\noindent{\qlabel{答}}\begingroup\sloppy\begin{enumerate}[leftmargin=1.6em,itemsep=2.5pt,topsep=0pt,parsep=0pt,label={\sffamily\footnotesize\color{accent}\arabic*.}]
\item {\fontsize{10}{14.8}\selectfont 精度支持要分三档：原生支持是矩阵单元直接吃这种类型，少读的字节同时变成少花的时间；能存不能算是能存显存但矩阵单元不认；完全不支持是连类型都没有。少分一档就会把显存减半误判成快一倍。\par}
\item {\fontsize{10}{14.8}\selectfont 第二档最危险，它看起来就是支持的：模型能加载、显存确实降了一半、输出也对，只有压测速度才发现更慢。实测 INT8 两种做法是 4971 和 3922 $\mu$s，FP32 基线只有 1186.6 $\mu$s，权重只剩 17 MB 却慢了四倍。\par}
\item {\fontsize{10}{14.8}\selectfont 原因是反量化写在了矩阵乘外面。W8A16 先把整个权重物化成新的 FP16 张量，这一步要写 34 MB 再读 34 MB，比原来直接读 67 MB 还多，省下的带宽又花回去还多一次写入。\par}
\item {\fontsize{10}{14.8}\selectfont 原生支持该有的样子是 FP16：字节减半、时间也大致减半，1186.6 到 491.4 $\mu$s，2.51 倍。FP16 和 BF16 字节相同但不等价，491.4 对 544.7 $\mu$s 差 11\%，误差 3.6e-4 对 2.9e-3。\par}
\item {\fontsize{10}{14.8}\selectfont 账可以提前算：decode 每步下限等于要读的权重字节除以带宽。3.35 TB/s 下 BF16、FP8、INT4 每步是 4.79、2.40、1.20 ms；只能存不能算时三者挤在 20 到 24 ms，位宽越低越平。\par}
\item {\fontsize{10}{14.8}\selectfont 适配要动四处：权重重排、缩放系数的形状与布局、反量化的位置、回退路径。权重重排错了的症状是输出完全乱掉而不是误差大。KV 量化还要求写 KV 的 kernel 和注意力 kernel 都支持。\par}
\end{enumerate}\endgroup
\numline{公式}{$t=\frac{w}{bw}$（$w=P\times$ 每参数字节）；无原生 kernel：$t_{\text{bad}}=\frac{w+2\cdot P\cdot 2\cdot 2}{bw}$}
\numline{数字}{FP32 1186.6 $\mu$s / FP16 491.4 $\mu$s（2.51 倍）/ BF16 544.7 $\mu$s / W8A16 4971.1 $\mu$s / W8A8 3922.3 $\mu$s / FP8 反量化 2055.4 $\mu$s；3.35 TB/s 下 BF16/FP8/INT4 = 4.79/2.40/1.20 ms、209/417/834 token/s；只能存不能算 23.97/21.57/20.37 ms；FP8 E4M3 最大 448、最小正数 0.001953125；分组缩放相对误差 56.3\%}
\fol{换个带宽只有一半的卡，结论变吗？}{会更糟。1.6 TB/s 下 BF16 每步 10.04 ms、FP8 5.02 ms，只能存不能算时三者是 50.19、45.17、42.66 ms，位宽越低越平，因为主导项变成了反量化的读写。}
\fol{那量化在新卡上到底怎么用？}{按三档给方案。原生支持就直接用，但要核对缩放粒度、累加精度和形状；能存不能算就把反量化融进 kernel，读 INT8、片上转 BF16、立刻喂矩阵单元；完全不支持就换类型，代价是重新校准和验精度。}
\fol{量化模型的精度怎么验收？}{不能和 FP32 原模型比。正确参考是同一套量化方法在参考平台上的输出：同样 checkpoint、同样缩放粒度。两边掉得一样多是量化本身的代价，只有这块卡上掉得更多才是适配问题。}
\numline{别答错}{别把显存减半当成快一倍：权重只剩 17 MB 却比 FP32 慢四倍，量化收益在容量不在时间。}
\phantomsection\addcontentsline{toc}{section}{146\quad 集合通信库与多卡互联}
\lec{146}{集合通信库与多卡互联}
\askline{换一块卡，All-Reduce 结论还成立吗？}
\askline{通信库这一层适配要改多少代码？}
\askline{库有 All-Reduce 了，为什么还要自己写？}
\vspace{3pt}
\noindent{\qlabel{答}}\begingroup\sloppy\begin{enumerate}[leftmargin=1.6em,itemsep=2.5pt,topsep=0pt,parsep=0pt,label={\sffamily\footnotesize\color{accent}\arabic*.}]
\item {\fontsize{10}{14.8}\selectfont 这层适配成本最低：PyTorch 的 dist.all\_reduce、all\_gather 和进程组把库挡在后面，各家又都照 NCCL 的接口做。兼容路线连名字都不改，平台类里 dist\_backend 就写 nccl，底下跑 RCCL。\par}
\item {\fontsize{10}{14.8}\selectfont 重写路线要多做一步：厂商库注册成 PyTorch 的新后端，初始化写 backend 等于 hccl。模型里的 dist.all\_reduce 照旧，但所有写死 nccl 的地方都要改，包括引擎、启动脚本和文档。\par}
\item {\fontsize{10}{14.8}\selectfont 真正要重新量的是四个数：每轮固定开销 $\alpha$ 发几 KB 小消息量、有效带宽 $\beta$ 发几百 MB 大消息量、$\alpha$ 乘 $\beta$ 是小消息分界线、卡间是不是全连接。NVLink 上 $\alpha$ 是 3 $\mu$s、$\beta$ 是 450 GB/s。\par}
\item {\fontsize{10}{14.8}\selectfont 拓扑决定能用哪些算法。全连接时 Ring、Tree、一步到位都能用，小消息一步到位最快；环形直连远端多跳、跳数乘进 $\alpha$，一步到位用不了；跨机用两阶段，先机内 Reduce-Scatter 再机间一轮。\par}
\item {\fontsize{10}{14.8}\selectfont decode 的 All-Reduce 消息只有 256 KB（TP8、隐藏维 4096、BF16、batch 32）。按 $\alpha$ 6 $\mu$s、$\beta$ 200 GB/s，Ring 要 14 轮，开销 84 $\mu$s，搬运只 2.2 $\mu$s，97\% 是开销。\par}
\item {\fontsize{10}{14.8}\selectfont 所以自定义 All-Reduce 只在很窄的区间里更快：卡数在 2、4、6、8 名单内、消息不超过 8 MiB、链路全连接、字节数是 16 的倍数。非全连接时上限压到 512 KB，区间之外退回库。\par}
\end{enumerate}\endgroup
\numline{公式}{$t=\alpha+m/\beta$；小消息分界线 $m=\alpha\beta$；Ring 轮数 $2(p-1)$}
\numline{数字}{NVLink $\alpha$=3 $\mu$s、$\beta$=450 GB/s，$\alpha$$\cdot$$\beta$=1.4 MB；8 卡 512 KB：Ring 44.0 / Tree 20.3 / 一步到位 11.0 $\mu$s；单条 512 KB 消息 4.1 $\mu$s 里延迟占 73\%；16 MB：Ring 227.4 / 一步到位 77.7 $\mu$s；decode 消息 256 KB、Ring 14 轮、84 $\mu$s 中 97\% 是固定开销、28 层五六毫秒；kMaxBlocks=36 而 ROCm 路径用 16；NCCL All-Gather 每次约 55 $\mu$s}
\fol{通信和计算的重叠怎么才算真生效？}{要在时间线上看通信条带和计算条带是否真的并行。三招是独立流加事件同步、按微批次切开、把通信融进算子。常见失败原因是中间有一次隐式同步，把数据拷回了主机。}
\fol{自定义 All-Reduce 什么时候会死锁？}{所有卡对每次调用必须走同一条路径，而这个决定只能依赖全局尺寸信息和公共行布局。否则一部分卡在等自定义 kernel 的标志位、另一部分在库里等，直接死锁。}
\fol{P2P 能不能用为什么必须先查？}{两张卡能不能直接读对方显存取决于硬件、驱动和虚拟化配置。能用时自定义 All-Reduce 才成立，不能用时通信都过主机内存，$\alpha$ 和 $\beta$ 都差一个量级。}
\numline{别答错}{自定义 All-Reduce 不是更快版本：卡数、消息上限、全连接、对齐四条之外要么用不了要么更慢。}
\phantomsection\addcontentsline{toc}{section}{147\quad 推理引擎的硬件插件}
\lec{147}{推理引擎的硬件插件}
\askline{怎么把一块新卡接进 vLLM 又不改源码？}
\askline{一个硬件插件最少要实现哪些接口？}
\askline{什么情况下只能动上游代码？}
\vspace{3pt}
\noindent{\qlabel{答}}\begingroup\sloppy\begin{enumerate}[leftmargin=1.6em,itemsep=2.5pt,topsep=0pt,parsep=0pt,label={\sffamily\footnotesize\color{accent}\arabic*.}]
\item {\fontsize{10}{14.8}\selectfont vLLM 用标准 entry point 留入口，插件在 setup.py 里声明，引擎启动时逐类加载。硬件相关两类：platform\_plugins 返回平台类全限定名，general\_plugins 注册模型和算子。\par}
\item {\fontsize{10}{14.8}\selectfont 入口必须可重入而且足够轻：每个进程都要加载一次，重复调用不能出问题；入口里不能 import torch 或初始化设备，否则发现阶段就把东西拽进来，整模块替换来不及安排。\par}
\item {\fontsize{10}{14.8}\selectfont 平台类最少实现五项：设备类型与名字、改配置、拿注意力后端类、拿通信器类。其中改配置最关键，worker 类就是在这里指定的，其余能力按需覆盖，不实现就退回引擎通用路径。\par}
\item {\fontsize{10}{14.8}\selectfont 清单的性质是引擎核心逻辑一项都不在里面：调度、token 预算、KV 块管理、前缀缓存、请求生命周期全由上游负责，插件只回答这块卡上这件事怎么做。所以工作量只跟硬件差异有关。\par}
\item {\fontsize{10}{14.8}\selectfont 沿用 CUDA 键白拿很多：device\_type 是 cuda、dispatch\_key 是 CUDA、后端叫 nccl，所有 CUDA 分支自动走对，worker 能继承 GPU worker，图捕获复用 CUDA Graph 包装器。\par}
\item {\fontsize{10}{14.8}\selectfont 改上游按代价排：先看能不能变成配置；不行用运行时替换，只拦精确模块名、每个补丁只应用一次、失败锁存并抛出、还能报告；最后才 fork。替换最大的风险是版本错配，必须配版本关卡。\par}
\end{enumerate}\endgroup
\numline{公式}{五个必须实现：$\text{device\_type}$、$\text{device\_name}$、$\text{check\_and\_update\_config}$、$\text{get\_attn\_backend\_cls}$、$\text{get\_device\_communicator\_cls}$}
\numline{数字}{vLLM 约 17 万行、几周一个版本；插件 818 个源文件 255134 行，测试 115629 行，补丁 195 个文件（platform 45 + worker 150），整模块替换 16 对，平台类 1033 行；对照 nano-vllm 只有 1195 行}
\fol{为什么不能直接 fork 上游改？}{维护成本最高。上游几周一个版本，加一个判断自己硬件的分支意味着每次更新都要重新合并，而厂商要支持的模型比上游合并的节奏更快，fork 只能当过渡。}
\fol{运行时替换最怕什么？}{版本错配。上游改了被替换函数的签名或语义，补丁还按老的写，结果可能是跑错而不是报错。所以版本检查要在任何补丁注册之前做完，只读元数据、不导入被检查的包。}
\fol{插件体量为什么这么大？}{25 万行里 11.5 万行是测试。核心算法不大，大的是「在真实硬件、真实模型、真实版本组合上让它一直对」的那部分，这跟 1195 行的 nano-vllm 正好形成对照。}
\numline{别答错}{入口必须可重入，且不能在发现阶段 import torch 或 vLLM，否则整模块替换来不及安排。}
\chapter{案例：海光DCU}
{\footnotesize\sffamily\color{grayt}第 148--153 讲}\par\vspace{-4pt}
\phantomsection\addcontentsline{toc}{section}{148\quad 海光 DCU 的软硬件栈}
\lec{148}{海光 DCU 的软硬件栈}
\askline{DCU 的软件栈怎么和 CUDA 对上？}
\askline{换到 DCU 到底哪几层要重写？}
\askline{波前 64 和 MMAC 这两个常数从哪来？}
\vspace{3pt}
\noindent{\qlabel{答}}\begingroup\sloppy\begin{enumerate}[leftmargin=1.6em,itemsep=2.5pt,topsep=0pt,parsep=0pt,label={\sffamily\footnotesize\color{accent}\arabic*.}]
\item {\fontsize{10}{14.8}\selectfont DCU 走兼容 GPGPU 路线，和 ROCm 一脉相承，编程模型与 CUDA 同构：线程、线程块、线程组、共享内存、矩阵核心都有对应物。所以前两篇的通用原理能逐条套上来，差异也看得清楚。\par}
\item {\fontsize{10}{14.8}\selectfont 九层栈的账是沿用 1 层、翻译 6 层、重写 2 层。沿用模型代码；重写算子库和运行时与驱动；其余六层，引擎、PyTorch、图编译、kernel 语言、通信库、厂商工具，都是翻译。\par}
\item {\fontsize{10}{14.8}\selectfont DTK 对应 CUDA Toolkit，编译器、HIP 运行时、数学库都在里面，默认装 /opt/dtk，版本从 .info/rocm\_version 读，CI 检查 26.04 和 26.10 两条线。驱动侧是 hy-smi。\par}
\item {\fontsize{10}{14.8}\selectfont 第一个常数是波前宽度 64：源码里 kWaveSize 等于 64，而 NVIDIA 的 warp 是 32。它不只是性能参数，通信代码注释说「一个程序就是一个波前」，波前宽度被当成同步语义的依据。\par}
\item {\fontsize{10}{14.8}\selectfont 第二个常数是矩阵核心指令 MMAC：一次吃 16 乘 16 乘 16，BF16 输入、FP32 累加。所以 64 条通道上每条通道持有 256 除以 64 等于 4 个 FP32 累加器，这个数字是硬件定的。\par}
\item {\fontsize{10}{14.8}\selectfont DAS 与 OpenDAS 是框架和组件的适配层，容器脚本一次列出 triton、flash\_attn、lightop、deepgemm、aiter、flash\_mla、lmcache 等十几个组件。适配一块卡的工作量一大半在这些组件上。\par}
\end{enumerate}\endgroup
\numline{公式}{每条通道累加器数 $=$ 矩阵指令输出元素 $\div$ 波前宽度 $=(16\times 16)/64=4$}
\numline{数字}{九层 = 沿用 1 + 翻译 6 + 重写 2；DTK 默认 /opt/dtk、CI 发布线 26.04/26.10；kWaveSize=64；MMAC 16$\times$16$\times$16、输入 BF16 累加 FP32、每通道 4 个累加器；架构代号 gfx928 2 处、gfx936 10 处、gfx938 87 处、gfx950 6 处；插件 25 万行、平台类 1033 行；目标引擎 vLLM 0.28.1}
\fol{gfx938 出现 87 处说明什么？}{说明这份代码主要针对这一代卡调过。同一家的卡不同代之间精度和指令支持都可能不一样，所以代码按架构代号分支，判断也按代号而不是型号系列。}
\fol{为什么官方文档里的版本号不能照抄？}{文档说 DAS 某版本支持 vLLM 0.11 和 SGLang 0.5.7，而手上这份插件源码的目标引擎是 0.28.1，比文档更新。版本矩阵要按仓库里的声明核对。}
\fol{九层里为什么算子库必须重写？}{因为算子库直接贴着硬件：矩阵乘要用本卡的 MMAC 指令和 hipBLASLt，注意力要用移植版 FlashAttention 和 FlashMLA，没有机械替换的余地，只能按本卡指令集重写。}
\numline{别答错}{波前 64 不只是性能参数：一个程序就是一个波前，写入顺序才成立，它是同步语义的依据。}
\phantomsection\addcontentsline{toc}{section}{149\quad 从 CUDA 到 DCU：迁移与版本矩阵}
\lec{149}{从 CUDA 到 DCU：迁移与版本矩阵}
\askline{一份 CUDA 代码搬到 DCU 要改哪些行？}
\askline{迁完之后为什么 torch.cuda 还能用？}
\askline{六层版本怎么才能对上？}
\vspace{3pt}
\noindent{\qlabel{答}}\begingroup\sloppy\begin{enumerate}[leftmargin=1.6em,itemsep=2.5pt,topsep=0pt,parsep=0pt,label={\sffamily\footnotesize\color{accent}\arabic*.}]
\item {\fontsize{10}{14.8}\selectfont 迁移分三类：运行时 API 和库名机械替换，跑一遍 hipify 类工具九成的行就过去了；要人工确认的头号项是波前 64，除 32 取余、0xffffffff 掩码、折半归约都要改；重写的是内联汇编和矩阵指令。\par}
\item {\fontsize{10}{14.8}\selectfont 16 行的样例里自动 6 行、人工确认 7 行、重写 3 行。真实项目自动比例高得多，但剩下的行往往集中在最关键的 kernel 里，所以迁移工作量不能按行数估。\par}
\item {\fontsize{10}{14.8}\selectfont 写死 32 不报错也不崩溃，只是结果偏小：工具不会改这个整数，但波前是 64 条通道，写死 32 等于把一个波前当两个 warp 用，同步粒度和归约范围全错。应写成 warpSize 或具名常量。\par}
\item {\fontsize{10}{14.8}\selectfont 矩阵指令还要按代分支：gfx928 和 gfx936 用 amdgcn 那版内建函数，gfx938 用 hcu 那版。构建侧 sm\_90 变成 gfx938，CUDA\_HOME 的角色由 ROCM\_PATH 承担。\par}
\item {\fontsize{10}{14.8}\selectfont torch.cuda 能用的原因是 PyTorch 的 DCU 版沿用 ROCm 做法，整个后端用 HIP 重编并注册到 CUDA 分发键上。于是分发器查 CUDA 那行拿到的是为 DCU 编译的实现，Python 侧接口照旧。\par}
\item {\fontsize{10}{14.8}\selectfont 六层版本必须成套：插件在打任何补丁前只读元数据检查引擎版本，要求主次补丁号一致、放宽 rc 与 dev 和本地后缀，不通过就抛错并锁存。CI 对 Python 精确、torch 前缀、DTK 发布线。\par}
\end{enumerate}\endgroup
\numline{公式}{$\text{nvcc}\ \text{-arch=sm\_90}\to\text{{-}{-}offload-arch=gfx938}$；$\text{CUDA\_HOME}\to\text{ROCM\_PATH}$；版本关卡：epoch 与 $(\text{major},\text{minor},\text{patch})$ 必须一致}
\numline{数字}{16 行里自动 6 / 人工确认 7 / 重写 3；hipify 类工具过九成行，hip 前缀 36 处而沿用 CUDA 名 47 处；目标版本串 0.28.1rc1.dev491+g462fdb097.das.462fdb0.dtk2604；Python 3.10.12 或 3.12.12、torch 2.11.0+、DTK 26.04/26.10}
\fol{版本检查通过就能跑对吗？}{不能。版本关卡只读元数据，通过检查不代表 PyTorch、DTK 和算子库的二进制接口已经验证兼容。它挡掉的是「肯定不行」的组合，精度和性能是另外两关。}
\fol{device\_count 为什么不能直接用 torch.cuda？}{因为要能在 CUDA\_VISIBLE\_DEVICES 被改写之后仍然给出正确的卡数。插件自己实现了一个无状态的版本，这类看起来一样、细节不一样的地方正是插件存在的理由。}
\fol{从零到跑起来最容易错在哪一步？}{DTK 和 PyTorch 的配对。构建里写着它编译时用的 HIP 版本，必须和 DTK 成套，症状是 import torch 就报符号找不到，或者能 import 但一调 kernel 就崩。}
\numline{别答错}{版本检查只读元数据：通过关卡不代表 PyTorch、DTK 和算子库的二进制接口已验证兼容。}
\phantomsection\addcontentsline{toc}{section}{150\quad 读 vllm-plugin-das：平台与运行时补丁}
\lec{150}{读 vllm-plugin-das：平台与运行时补丁}
\askline{一个真实硬件插件的接入面有多大？}
\askline{不改上游源码，怎么换掉上游模块？}
\askline{补丁机制要防住哪些坑？}
\vspace{3pt}
\noindent{\qlabel{答}}\begingroup\sloppy\begin{enumerate}[leftmargin=1.6em,itemsep=2.5pt,topsep=0pt,parsep=0pt,label={\sffamily\footnotesize\color{accent}\arabic*.}]
\item {\fontsize{10}{14.8}\selectfont 25 万行里和「怎么接进引擎」相关的只有三处：三个 entry point（入口模块 138 行）、平台类 1033 行、补丁机制（195 个补丁文件加几个基础模块）。其余是算子、模型和 11.5 万行测试。\par}
\item {\fontsize{10}{14.8}\selectfont 入口模块顶层只做两件事：配置算子库环境变量、导入兼容性检查，没有 import torch 也没有 import vllm。原因是发现阶段只能布置回调，一旦把实现拽进来，目标模块已先加载，替换来不及安排。\par}
\item {\fontsize{10}{14.8}\selectfont 失败要锁存：平台初始化失败过一次，之后每次调用都直接抛同一个错误而不重试。因为补丁可能只打了一半，重试会在半残的注册表上继续操作，产生更难查的问题，比起不来更糟。\par}
\item {\fontsize{10}{14.8}\selectfont 补丁三种形态：整模块替换 16 对、导入后回调（195 个文件大多是这类）、配置改写。调用方照常写 vllm 的导入路径，由协调器在上游模块加载前解析到插件实现，升级时好比对。\par}
\item {\fontsize{10}{14.8}\selectfont 入口是装在 sys.meta\_path 上的查找器，四道守卫：只拦精确模块名（按前缀匹配会误伤）、已打过补丁的模块禁止重载、正在打补丁时禁止重入导入、失败的补丁拒绝导入而不回退。\par}
\item {\fontsize{10}{14.8}\selectfont 平台类用七行声明回答设备身份：枚举取 ROCM 白拿所有 ROCm 分支，device\_type 是 cuda、dispatch\_key 是 CUDA、后端写 nccl。supports\_fp8 就一行：是不是 gfx938。\par}
\end{enumerate}\endgroup
\numline{公式}{$\text{find\_spec}(\text{fullname})\to$ 精确字典查表 $\to$ 整模块替换 / 挂导入后回调 $\to$ 失败锁存并拒绝导入}
\numline{数字}{818 个源文件 255134 行、测试 115629 行；入口 138 行、平台类 1033 行、patch 下 3.6 万行；补丁 195 个（platform/core\_fix 24、platform/framework\_opt 21、worker/core\_fix 42、worker/framework\_opt 35、worker/op\_opt 73）；整模块替换 16 对}
\fol{supports\_fp8 为什么只有一行？}{因为返回的就是这块卡是不是 gfx938。精度支持按架构代号判断而不是按型号名字，同一家不同代支持可能不一样，这就是那张三档支持表落到代码里的样子。}
\fol{补丁顺序为什么必须显式写出来？}{因为有些目标模块在插件发现之前就被导入了，协调器要按清单顺序立即应用回调；而 KV 适配必须排在 MoE 协调器包装之前，否则基类构造链看到的缓存尺寸不对。}
\fol{这套机制的成本是从哪来的？}{两个约束：不能改上游，和上游一直在动。1195 行的 nano-vllm 没有这类机制，因为它只跑一种硬件；25 万行里 11.5 万行测试就是为这两条付的账。}
\numline{别答错}{补丁失败必须拒绝导入而不是静默回退：半打补丁的状态比起不来更糟，也更难查。}
\phantomsection\addcontentsline{toc}{section}{151\quad DCU 上的矩阵乘：hipBLASLt 与小 M kernel}
\lec{151}{DCU 上的矩阵乘：hipBLASLt 与小 M kernel}
\askline{小 M 矩阵乘为什么要自己写 kernel？}
\askline{这个 kernel 怎么分块、怎么归约？}
\askline{手写 kernel 失败以后怎么办？}
\vspace{3pt}
\noindent{\qlabel{答}}\begingroup\sloppy\begin{enumerate}[leftmargin=1.6em,itemsep=2.5pt,topsep=0pt,parsep=0pt,label={\sffamily\footnotesize\color{accent}\arabic*.}]
\item {\fontsize{10}{14.8}\selectfont 库在小 M 上退化了：decode 时 M 不超过 64，hipBLASLt 选 256 乘 256 宏块，不管 N 多大都要 135 $\mu$s、只有 0.93 TB/s，而只沿 N 分块的 kernel 能到 1.4 到 1.7 TB/s。\par}
\item {\fontsize{10}{14.8}\selectfont 所以 kernel 只覆盖 M 在 1 到 64、N 取 256 到 2048 那四档、K 是 128 的倍数。策略是一个线程块吃 16 行乘 kNPerBlock 列，块里的波前沿 K 分工、最后在共享内存折半归约。\par}
\item {\fontsize{10}{14.8}\selectfont 每个常数都来自硬件：波前 64、MMAC 是 16 乘 16 乘 16、每条通道 4 个 FP32 累加器（256 个输出摊到 64 条通道）、一次读 8 个 BF16 凑 16 字节对齐、内建函数按两代架构分支。\par}
\item {\fontsize{10}{14.8}\selectfont 有个 split-K 变体被宏关掉了：注释说当前 DTK 工具链在这个即时编译单元里遇到那条直接搬进共享内存的路径，会给 gfx938 生成空 code object，所以先排除在设备编译外。\par}
\item {\fontsize{10}{14.8}\selectfont 入口是三层回退：布局不对直接走 torch.mm 并要求 FP32 输出；张量不是 BF16、不在设备、非二维等任一条不满足也退回；小 M 且 supported 通过才走 MMAC kernel，失败就永久关掉这条路。\par}
\item {\fontsize{10}{14.8}\selectfont C++ 侧检查不是防御性样板而是契约：二维、在设备、BF16、连续、K 对得上、M 在 1 到 64、N 是那四个值之一、K 是 128 倍数。任一条不满足，结果是错的而不是慢的。\par}
\end{enumerate}\endgroup
\numline{公式}{tile 为 $16$ 行 $\times\ \text{kNPerBlock}$ 列；$\text{kWarps}$ 个波前沿 $K$ 切，最后在 LDS 里折半归约}
\numline{数字}{hipBLASLt 在 M $\leq$ 64 时约 135 $\mu$s 且与 N 无关；0.93 TB/s 对 1.4--1.7 TB/s；lm\_head 分片 32224$\times$4096：M=16 时 285 $\rightarrow$ 153 $\mu$s、M=64 时 190 $\mu$s；覆盖 M$\in$[1,64]、N$\in$\{256,512,1024,2048\}、K\%128=0；kWaveSize=64、每通道 4 个 FP32 累加器}
\fol{为什么 M 方向不补零？}{因为 16 行正好等于 MMAC 16 乘 16 乘 16 的行数，一块吃 16 行刚好对齐。库那边选 256 乘 256 时 M 等于 16 有 94\% 的行在算补零，时间全浪费在固定开销上。}
\fol{lm\_head 为什么反而用 Triton 写？}{因为它的逻辑简单，是纯流式读权重：M 不超过 64、K 是 128 的倍数、N 任意，只沿 N 分块、M 方向一个块吃完，Triton 够用又好维护。}
\fol{第一次调用为什么要预热？}{kernel 是用 load\_inline 在运行时编译的，外面套一个只缓存一个结果的装饰器，一个进程只编一次。好处是不用预编译所有架构，代价是第一次调用要等几十秒。}
\numline{别答错}{库的 135 $\mu$s 与 N 无关：不是没调好，是为大矩阵调好的库在小 M 上退化，代价在分块开销。}
\phantomsection\addcontentsline{toc}{section}{152\quad DCU 上的注意力、量化与通信}
\lec{152}{DCU 上的注意力、量化与通信}
\askline{DCU 上的注意力后端是怎么挑的？}
\askline{为什么我设的 --block-size 不生效？}
\askline{多卡通信在插件里怎么选路？}
\vspace{3pt}
\noindent{\qlabel{答}}\begingroup\sloppy\begin{enumerate}[leftmargin=1.6em,itemsep=2.5pt,topsep=0pt,parsep=0pt,label={\sffamily\footnotesize\color{accent}\arabic*.}]
\item {\fontsize{10}{14.8}\selectfont 插件注册五个后端：普通注意力的 FlashAttention 版和 Triton 版、MLA 的 FlashMLA 版和 Triton 版、稀疏 MLA 版。注册时只在后端没被覆盖过才注册，用户或第三方的选择优先于插件默认。\par}
\item {\fontsize{10}{14.8}\selectfont 优先级只由三个布尔定：是不是 MLA、是不是稀疏、要不要优先 Triton 版 MLA。稀疏 MLA 给两个候选，MLA 先 FlashMLA 后 Triton，普通 MHA 或 GQA 先 FlashAttention 后 Triton。\par}
\item {\fontsize{10}{14.8}\selectfont 校验逐个调 validate\_configuration，传设备能力、头数、块大小、数据类型，返回不可用的原因列表，ImportError 也算一条。全部不可用就抛错并打印每个候选的原因，不悄悄降级。\par}
\item {\fontsize{10}{14.8}\selectfont 块大小不是自由参数。稀疏 MLA 里有个索引器注意力层，上游可能先给它选 16 token 的块，所以管理器块要保持至少 64 且 64 对齐；DeepSeek-V4 稀疏后端要 256 token，于是取最小公倍数。\par}
\item {\fontsize{10}{14.8}\selectfont 量化启用 23 种、注释掉 4 种，被注释的 mxfp4、bitsandbytes、modelopt\_fp4、gpt\_oss\_mxfp4 都和 4 位浮点有关，注释而不删说明暂时不支持。FP8 按代号判断：是不是 gfx938。\par}
\item {\fontsize{10}{14.8}\selectfont 通信沿用 nccl 和上游的 CudaCommunicator，自定义 All-Reduce 是六项条件的分支：卡数 2 到 8、消息上限 8 MiB、要全连接（否则 512 KB）、16 字节对齐、阶段阈值、块数 36 而默认 16。\par}
\end{enumerate}\endgroup
\numline{公式}{候选列表 $\to\ \text{validate\_configuration}$ 逐个校验 $\to$ 全不可用则抛错并打印原因；$\text{block\_size}=\operatorname{lcm}(\text{后端要求},\text{索引器 }16,\text{压缩比})$}
\numline{数字}{注意力后端 5 个；量化启用 23 种、注释掉 4 种；块大小 64 对齐、DeepSeek-V4 稀疏后端 256 token；自定义 All-Reduce：2/4/6/8 卡、8 MiB、非全连接 512 KB、16 字节对齐、一阶段阈值 512 KB 与 256 KB（PCIe 版 1 MB 与 64 KB）、块数 36 而 ROCm 用 16；NCCL All-Gather 每次约 55 $\mu$s}
\fol{所有卡走了不同路径会怎样？}{直接死锁。一部分卡在等自定义 kernel 的标志位、另一部分在库里等。所以每次调用走自定义还是回退，必须只依赖全局尺寸信息和公共行布局。}
\fol{IPC All-Gather 为什么能被图捕获？}{每个程序自带调用计数器，标志位因此单调递增；两个缓冲槽按调用奇偶交替；一个程序就是一个波前所以写入有序，用硬件常数换掉了一次显式同步。}
\fol{为什么 fp8\_e4m3 的 KV 配置不通用？}{对稀疏 MLA 后端，平台会把它转成内部的 fp8\_ds\_mla，那是该后端专用的打包布局，把缩放系数和数据交错存在一起，不是通用的 cache 格式。}
\numline{别答错}{块大小被注意力后端、索引器层和压缩比共同约束，插件会用最小公倍数改写它并打日志，不是 bug。}
\phantomsection\addcontentsline{toc}{section}{153\quad 模型上线 DCU 的完整清单}
\lec{153}{模型上线 DCU 的完整清单}
\askline{模型从拿到卡到上线，你要验哪些项？}
\askline{怎么保证每一步是真的过关，不是看起来过了？}
\askline{哪几项最容易跳，跳了代价多大？}
\vspace{3pt}
\noindent{\qlabel{答}}\begingroup\sloppy\begin{enumerate}[leftmargin=1.6em,itemsep=2.5pt,topsep=0pt,parsep=0pt,label={\sffamily\footnotesize\color{accent}\arabic*.}]
\item {\fontsize{10}{14.8}\selectfont 清单的价值不在记住别忘了什么，而在每一项都有能机械判定的验收标准：六个阶段共 29 项，环境 5、跑通 4、精度 5、性能 6、多卡 4、服务化 5，前一关没点亮就不进下一关。\par}
\item {\fontsize{10}{14.8}\selectfont 环境关 5 项先排掉：驱动装好、管理工具能看到所有卡，卡数、显存、温度正常且连跑一小时不掉卡；工具包在支持的发布线上，框架能跑通一次 BF16 矩阵乘，算子库版本与目标组合匹配，插件自检全绿。\par}
\item {\fontsize{10}{14.8}\selectfont 跑通关 4 项：短提示生成 32 个 token 不报错；在分发器层记账跑一次前向，CPU 回退计数必须为 0；缺失算子逐个标注组合、自写还是等厂商并记进待办；FP32、BF16、FP16 和量化各生成一次。\par}
\item {\fontsize{10}{14.8}\selectfont 精度关 5 项：以同精度参考平台为基准按位置取最差值，逐层都在阈值内；固定短、长、中文、代码几个提示逐 token 比对；数据集分数记录在案；量化模型和同一套量化方法的参考输出比；每改一个 kernel 就重跑比对。\par}
\item {\fontsize{10}{14.8}\selectfont 性能关 6 项：按每步字节除以带宽、2MNK 除以峰值算出理论下限，数热路径的主机同步次数；decode M=1…64 与 prefill 各测一遍矩阵乘和注意力；图模式收益与图中断都量过；放弃的优化写明占比、收益和理由。\par}
\item {\fontsize{10}{14.8}\selectfont 多卡关 4 项：小消息、大消息各测一次，量出 $\alpha$、$\beta$、是否全连接、逐对 P2P；两卡逐 token 对单卡，允许累加顺序的小差异；扩到目标卡数，TPOT 和吞吐曲线要用 $\alpha$-$\beta$ 模型解释；自定义通信两条路都压测，区间外退回库且不死锁。\par}
\item {\fontsize{10}{14.8}\selectfont 服务化关 5 项：预热覆盖即时编译和图捕获，否则首次编译几十秒落在真实请求上就是超时；按到达率扫区间记 TTFT、TPOT 和 goodput；按 SLO 算实例数并留余量；冻结各层版本和启动参数；回滚与降级路径都验过。\par}
\end{enumerate}\endgroup
\numline{公式}{每步理论下限 $=\frac{\text{每步要读的字节}}{\text{带宽}}$；矩阵乘下限 $=\frac{2MNK}{\text{峰值算力}}$}
\numline{数字}{六阶段 29 项：环境 5、跑通 4、精度 5、性能 6、多卡 4、服务化 5；CPU 回退算子慢 15 倍；主机同步 150 $\mu$s 对 kernel 下发 3.6 $\mu$s（约 44 倍）；即时编译一次几十秒；FP16 溢出时第 4 层最差位置误差 1.0 对整张量 4e-3}
\fol{哪一步最容易看起来过了？}{算子回退和精度分精度验。回退只要在分发器层记一次账、确认计数为 0；精度要 FP32、BF16、FP16 各验一次，FP16 下 RMSNorm 平方和溢出会生成一串句号，只在这个精度暴露。}
\fol{六个阶段的分界线凭什么这么定？}{每条分界都是一个不往回走的承诺：环境过了版本才固定，跑通过后没有东西在悄悄回退，精度过了才有可信基线，单卡瓶颈解释清楚才上多卡，两条通信路径都压过才做服务化。}
\fol{为什么预热、回滚、版本冻结不能当形式？}{它们只在出问题时才体现价值，而新硬件出问题的概率比成熟平台高：即时编译一次几十秒，落在首个真实请求上就是一次超时；版本没冻结，三个月后复现不出当初的结果；回滚路径没验过就切不回去。}
\numline{别答错}{最容易跳过：算子回退计数、按精度分别验、记下取舍理由；跳过它们，性能数字和精度结论都不可信。}
\clearpage
\chapter*{附录\quad 153 讲问题清单}
\addcontentsline{toc}{chapter}{附录\quad 153 讲问题清单}
\markboth{附录：153 讲问题清单}{}
{\footnotesize\sffamily\color{grayt}按讲号排列，用来考前自测：看到问题能说出结论和一个数字就过\par}\vspace{4pt}
\begin{multicols}{2}
\raggedright
\noindent{\sffamily\bfseries 1}\;{\small 大模型怎么把文字变成数字的？}\par\vspace{2.2pt}
\noindent{\sffamily\bfseries 2}\;{\small 判断两个向量像不像，为什么不用距离？}\par\vspace{2.2pt}
\noindent{\sffamily\bfseries 3}\;{\small 矩阵乘法在模型里到底在算什么？}\par\vspace{2.2pt}
\noindent{\sffamily\bfseries 4}\;{\small 矩阵乘法几何上到底在做什么？}\par\vspace{2.2pt}
\noindent{\sffamily\bfseries 5}\;{\small 代码里 [b, s, d] 具体指什么？}\par\vspace{2.2pt}
\noindent{\sffamily\bfseries 6}\;{\small softmax 为什么要先减去最大值？}\par\vspace{2.2pt}
\noindent{\sffamily\bfseries 7}\;{\small 按概率采样在代码里怎么做？}\par\vspace{2.2pt}
\noindent{\sffamily\bfseries 8}\;{\small 梯度下降里的梯度是怎么算出来的？}\par\vspace{2.2pt}
\noindent{\sffamily\bfseries 9}\;{\small RoPE 为什么用旋转来编码位置？}\par\vspace{2.2pt}
\noindent{\sffamily\bfseries 10}\;{\small FP16 和 BF16 到底差在哪，训练用哪个？}\par\vspace{2.2pt}
\noindent{\sffamily\bfseries 11}\;{\small 给你一个模型，怎么估它一次要跑多久？}\par\vspace{2.2pt}
\noindent{\sffamily\bfseries 12}\;{\small 感知机是什么，它为什么学不会异或？}\par\vspace{2.2pt}
\noindent{\sffamily\bfseries 13}\;{\small 神经元和线性层是什么关系？}\par\vspace{2.2pt}
\noindent{\sffamily\bfseries 14}\;{\small 为什么神经网络一定要有激活函数？}\par\vspace{2.2pt}
\noindent{\sffamily\bfseries 15}\;{\small 为什么深网络比宽网络更省参数？}\par\vspace{2.2pt}
\noindent{\sffamily\bfseries 16}\;{\small 分类为什么用交叉熵，不用平方误差？}\par\vspace{2.2pt}
\noindent{\sffamily\bfseries 17}\;{\small 梯度下降到底在做什么？一句话讲清}\par\vspace{2.2pt}
\noindent{\sffamily\bfseries 18}\;{\small 80 亿个参数怎么一次算出梯度？}\par\vspace{2.2pt}
\noindent{\sffamily\bfseries 19}\;{\small SGD、Momentum、Adam 差在哪？}\par\vspace{2.2pt}
\noindent{\sffamily\bfseries 20}\;{\small 一个训练 step 到底做了哪几件事？}\par\vspace{2.2pt}
\noindent{\sffamily\bfseries 21}\;{\small 文字是怎么变成模型能算的数字的？}\par\vspace{2.2pt}
\noindent{\sffamily\bfseries 22}\;{\small 为什么不能把 token 编号直接喂给模型？}\par\vspace{2.2pt}
\noindent{\sffamily\bfseries 23}\;{\small 大模型预训练到底在学什么？}\par\vspace{2.2pt}
\noindent{\sffamily\bfseries 24}\;{\small 一段话是怎么一个词一个词出来的？}\par\vspace{2.2pt}
\noindent{\sffamily\bfseries 25}\;{\small 注意力本身为什么不看顺序？}\par\vspace{2.2pt}
\noindent{\sffamily\bfseries 26}\;{\small 自注意力的公式为什么这么写？}\par\vspace{2.2pt}
\noindent{\sffamily\bfseries 27}\;{\small 语言模型为什么必须要因果掩码？}\par\vspace{2.2pt}
\noindent{\sffamily\bfseries 28}\;{\small 多头注意力多花了几倍参数？}\par\vspace{2.2pt}
\noindent{\sffamily\bfseries 29}\;{\small 深层网络为什么训不动？}\par\vspace{2.2pt}
\noindent{\sffamily\bfseries 30}\;{\small LayerNorm 具体在算什么？}\par\vspace{2.2pt}
\noindent{\sffamily\bfseries 31}\;{\small FFN 到底在算什么？}\par\vspace{2.2pt}
\noindent{\sffamily\bfseries 32}\;{\small 一个 Block 里有哪些部件？}\par\vspace{2.2pt}
\noindent{\sffamily\bfseries 33}\;{\small 输出层到底做了什么？只是打个分吗？}\par\vspace{2.2pt}
\noindent{\sffamily\bfseries 34}\;{\small 一句话从进模型到出概率，经过哪些步骤？}\par\vspace{2.2pt}
\noindent{\sffamily\bfseries 35}\;{\small 一个模型的参数量怎么算出来的？}\par\vspace{2.2pt}
\noindent{\sffamily\bfseries 36}\;{\small 为什么现代大模型都用 Pre-Norm？}\par\vspace{2.2pt}
\noindent{\sffamily\bfseries 37}\;{\small RoPE 凭什么只感知相对位置？}\par\vspace{2.2pt}
\noindent{\sffamily\bfseries 38}\;{\small Llama 的 FFN 为什么有三个矩阵？}\par\vspace{2.2pt}
\noindent{\sffamily\bfseries 39}\;{\small 为什么只有 8 个 KV 头，不是 32 个？}\par\vspace{2.2pt}
\noindent{\sffamily\bfseries 40}\;{\small MLA 怎么做到比 GQA 还省？}\par\vspace{2.2pt}
\noindent{\sffamily\bfseries 41}\;{\small MoE 为什么能参数多、计算少？}\par\vspace{2.2pt}
\noindent{\sffamily\bfseries 42}\;{\small 上下文从 8K 扩到 128K，卡在哪？}\par\vspace{2.2pt}
\noindent{\sffamily\bfseries 43}\;{\small 三代模型的结构差在哪？}\par\vspace{2.2pt}
\noindent{\sffamily\bfseries 44}\;{\small 线性注意力凭什么和 n 无关？}\par\vspace{2.2pt}
\noindent{\sffamily\bfseries 45}\;{\small 大模型预训练数据从哪来？}\par\vspace{2.2pt}
\noindent{\sffamily\bfseries 46}\;{\small 固定算力，模型该做多大？}\par\vspace{2.2pt}
\noindent{\sffamily\bfseries 47}\;{\small 基座模型和指令模型差在哪？}\par\vspace{2.2pt}
\noindent{\sffamily\bfseries 48}\;{\small 为什么用比较而不用打分？}\par\vspace{2.2pt}
\noindent{\sffamily\bfseries 49}\;{\small 推理模型是怎么训出来的？}\par\vspace{2.2pt}
\noindent{\sffamily\bfseries 50}\;{\small LoRA 为什么能省显存？}\par\vspace{2.2pt}
\noindent{\sffamily\bfseries 51}\;{\small 怎么判断一个模型好不好？}\par\vspace{2.2pt}
\noindent{\sffamily\bfseries 52}\;{\small 为什么要自己训分词器？}\par\vspace{2.2pt}
\noindent{\sffamily\bfseries 53}\;{\small 训练样本是怎么切出来的？}\par\vspace{2.2pt}
\noindent{\sffamily\bfseries 54}\;{\small 多头注意力怎么写成矩阵乘？}\par\vspace{2.2pt}
\noindent{\sffamily\bfseries 55}\;{\small Block 由哪些零件组成？}\par\vspace{2.2pt}
\noindent{\sffamily\bfseries 56}\;{\small 预训练到底在优化什么？}\par\vspace{2.2pt}
\noindent{\sffamily\bfseries 57}\;{\small 训练时怎么只留最好的那一份？}\par\vspace{2.2pt}
\noindent{\sffamily\bfseries 58}\;{\small 别人训练好的权重，怎么装进自己的模型？}\par\vspace{2.2pt}
\noindent{\sffamily\bfseries 59}\;{\small SFT 和预训练的训练循环差在哪？}\par\vspace{2.2pt}
\noindent{\sffamily\bfseries 60}\;{\small GPU 为什么比 CPU 适合矩阵乘法？}\par\vspace{2.2pt}
\noindent{\sffamily\bfseries 61}\;{\small H100 的存储分几层，各多快？}\par\vspace{2.2pt}
\noindent{\sffamily\bfseries 62}\;{\small Roofline 怎么判断卡在哪？}\par\vspace{2.2pt}
\noindent{\sffamily\bfseries 63}\;{\small 为什么精度减半算力就翻倍？}\par\vspace{2.2pt}
\noindent{\sffamily\bfseries 64}\;{\small 机内和机间带宽差多少？}\par\vspace{2.2pt}
\noindent{\sffamily\bfseries 65}\;{\small 超节点和普通多机集群差在哪？}\par\vspace{2.2pt}
\noindent{\sffamily\bfseries 66}\;{\small 线上推理慢，你怎么定位瓶颈？}\par\vspace{2.2pt}
\noindent{\sffamily\bfseries 67}\;{\small 规格表的峰值实际能跑到几成？}\par\vspace{2.2pt}
\noindent{\sffamily\bfseries 68}\;{\small prefill 和 decode 瓶颈为什么不同？}\par\vspace{2.2pt}
\noindent{\sffamily\bfseries 69}\;{\small 推理为什么要做 KV Cache？}\par\vspace{2.2pt}
\noindent{\sffamily\bfseries 70}\;{\small TTFT 和 TPOT 分别由什么决定？}\par\vspace{2.2pt}
\noindent{\sffamily\bfseries 71}\;{\small 80 GB 的卡能放多少请求？}\par\vspace{2.2pt}
\noindent{\sffamily\bfseries 72}\;{\small 每百万 token 的成本怎么估？}\par\vspace{2.2pt}
\noindent{\sffamily\bfseries 73}\;{\small 温度、Top-k、Top-p 到底在调什么？}\par\vspace{2.2pt}
\noindent{\sffamily\bfseries 74}\;{\small 线上 batching 和训练时差在哪？}\par\vspace{2.2pt}
\noindent{\sffamily\bfseries 75}\;{\small KV Cache 的显存浪费在哪？}\par\vspace{2.2pt}
\noindent{\sffamily\bfseries 76}\;{\small 多轮对话为什么要重算前缀？}\par\vspace{2.2pt}
\noindent{\sffamily\bfseries 77}\;{\small 长提示为什么把大家卡住？}\par\vspace{2.2pt}
\noindent{\sffamily\bfseries 78}\;{\small KV Cache 池子满了怎么办？}\par\vspace{2.2pt}
\noindent{\sffamily\bfseries 79}\;{\small 100 个客户的模型怎么共用一张卡？}\par\vspace{2.2pt}
\noindent{\sffamily\bfseries 80}\;{\small 怎么保证一定输出合法 JSON？}\par\vspace{2.2pt}
\noindent{\sffamily\bfseries 81}\;{\small 一个请求从发出去到出第一个字，经过哪些模块？}\par\vspace{2.2pt}
\noindent{\sffamily\bfseries 82}\;{\small 一张卡吞吐一万 token/s，这话可信吗？}\par\vspace{2.2pt}
\noindent{\sffamily\bfseries 83}\;{\small 几个逐元素算子，为什么合成一个就快了？}\par\vspace{2.2pt}
\noindent{\sffamily\bfseries 84}\;{\small 标准注意力慢在哪，为什么必须融合？}\par\vspace{2.2pt}
\noindent{\sffamily\bfseries 85}\;{\small decode 的注意力和 prefill 差在哪？}\par\vspace{2.2pt}
\noindent{\sffamily\bfseries 86}\;{\small CUDA 里线程、块、网格是怎么组织的？}\par\vspace{2.2pt}
\noindent{\sffamily\bfseries 87}\;{\small 一个矩阵乘法怎么才能跑到峰值？}\par\vspace{2.2pt}
\noindent{\sffamily\bfseries 88}\;{\small Triton 和写 CUDA 有什么不一样？}\par\vspace{2.2pt}
\noindent{\sffamily\bfseries 89}\;{\small decode 为什么要用 CUDA Graph？}\par\vspace{2.2pt}
\noindent{\sffamily\bfseries 90}\;{\small 量化误差怎么算，上限是多少？}\par\vspace{2.2pt}
\noindent{\sffamily\bfseries 91}\;{\small RTN 掉点，GPTQ 和 AWQ 差在哪？}\par\vspace{2.2pt}
\noindent{\sffamily\bfseries 92}\;{\small 为什么激活量化这么难？}\par\vspace{2.2pt}
\noindent{\sffamily\bfseries 93}\;{\small FP8 两种格式怎么选？}\par\vspace{2.2pt}
\noindent{\sffamily\bfseries 94}\;{\small KV Cache 量化能换来什么？}\par\vspace{2.2pt}
\noindent{\sffamily\bfseries 95}\;{\small 长上下文 KV 太大，除了量化还能怎样？}\par\vspace{2.2pt}
\noindent{\sffamily\bfseries 96}\;{\small 投机解码为什么不掉质量？}\par\vspace{2.2pt}
\noindent{\sffamily\bfseries 97}\;{\small Medusa 和草稿模型比好在哪？}\par\vspace{2.2pt}
\noindent{\sffamily\bfseries 98}\;{\small 蒸馏为什么比只用标准答案好？}\par\vspace{2.2pt}
\noindent{\sffamily\bfseries 99}\;{\small 一张卡为什么不够用？}\par\vspace{2.2pt}
\noindent{\sffamily\bfseries 100}\;{\small 六种集合通信原语怎么区分？}\par\vspace{2.2pt}
\noindent{\sffamily\bfseries 101}\;{\small Ring 怎么把带宽做到下限？}\par\vspace{2.2pt}
\noindent{\sffamily\bfseries 102}\;{\small 一次通信的时间怎么估？}\par\vspace{2.2pt}
\noindent{\sffamily\bfseries 103}\;{\small 数据并行怎么保证各卡不跑偏？}\par\vspace{2.2pt}
\noindent{\sffamily\bfseries 104}\;{\small 16 字节/参数是怎么来的？}\par\vspace{2.2pt}
\noindent{\sffamily\bfseries 105}\;{\small 张量并行是怎么切矩阵的？}\par\vspace{2.2pt}
\noindent{\sffamily\bfseries 106}\;{\small 流水并行怎么切，气泡从哪来？}\par\vspace{2.2pt}
\noindent{\sffamily\bfseries 107}\;{\small 长序列放不下，怎么沿序列切？}\par\vspace{2.2pt}
\noindent{\sffamily\bfseries 108}\;{\small MoE 的专家怎么分到卡上？}\par\vspace{2.2pt}
\noindent{\sffamily\bfseries 109}\;{\small 五种并行怎么映射到硬件拓扑？}\par\vspace{2.2pt}
\noindent{\sffamily\bfseries 110}\;{\small 放得下以后还要不要加 TP？}\par\vspace{2.2pt}
\noindent{\sffamily\bfseries 111}\;{\small 大规模 MoE 推理时间花在哪？}\par\vspace{2.2pt}
\noindent{\sffamily\bfseries 112}\;{\small 为什么要做 PD 分离？}\par\vspace{2.2pt}
\noindent{\sffamily\bfseries 113}\;{\small 已经算过的 KV，读回来还是重算？}\par\vspace{2.2pt}
\noindent{\sffamily\bfseries 114}\;{\small 请求进来，路由器凭什么选实例？}\par\vspace{2.2pt}
\noindent{\sffamily\bfseries 115}\;{\small 超节点能让推理全面变快吗？}\par\vspace{2.2pt}
\noindent{\sffamily\bfseries 116}\;{\small 每秒 50 个请求，要几张卡？}\par\vspace{2.2pt}
\noindent{\sffamily\bfseries 117}\;{\small 一个推理引擎由哪几部分组成？}\par\vspace{2.2pt}
\noindent{\sffamily\bfseries 118}\;{\small 1.5 GB 的权重文件里有什么？}\par\vspace{2.2pt}
\noindent{\sffamily\bfseries 119}\;{\small KV Cache 到底缓存了什么？}\par\vspace{2.2pt}
\noindent{\sffamily\bfseries 120}\;{\small batch 里请求长度不同，怎么一起算？}\par\vspace{2.2pt}
\noindent{\sffamily\bfseries 121}\;{\small 为什么要给 KV 做分页管理？}\par\vspace{2.2pt}
\noindent{\sffamily\bfseries 122}\;{\small 前缀缓存怎么认出同一个块？}\par\vspace{2.2pt}
\noindent{\sffamily\bfseries 123}\;{\small W8A16 的 scale 怎么定？}\par\vspace{2.2pt}
\noindent{\sffamily\bfseries 124}\;{\small 投机解码为什么结果不变？}\par\vspace{2.2pt}
\noindent{\sffamily\bfseries 125}\;{\small 六个版本各换来了什么？}\par\vspace{2.2pt}
\noindent{\sffamily\bfseries 126}\;{\small nano-vllm 为什么把调度和执行拆开？}\par\vspace{2.2pt}
\noindent{\sffamily\bfseries 127}\;{\small 一步的请求怎么变成 GPU 输入？}\par\vspace{2.2pt}
\noindent{\sffamily\bfseries 128}\;{\small 草稿为什么不放进 token\_ids？}\par\vspace{2.2pt}
\noindent{\sffamily\bfseries 129}\;{\small vLLM V1 为什么拆成三类进程？}\par\vspace{2.2pt}
\noindent{\sffamily\bfseries 130}\;{\small 一条请求在 4 个副本里怎么挑？}\par\vspace{2.2pt}
\noindent{\sffamily\bfseries 131}\;{\small SGLang 的进程结构与 vLLM 差在哪？}\par\vspace{2.2pt}
\noindent{\sffamily\bfseries 132}\;{\small 各家引擎各适合什么场景？}\par\vspace{2.2pt}
\noindent{\sffamily\bfseries 133}\;{\small 17 万行的引擎从哪读起？}\par\vspace{2.2pt}
\noindent{\sffamily\bfseries 134}\;{\small 换一块国产卡要重写整个栈吗？}\par\vspace{2.2pt}
\noindent{\sffamily\bfseries 135}\;{\small GPGPU 和 NPU 的结构差在哪？}\par\vspace{2.2pt}
\noindent{\sffamily\bfseries 136}\;{\small CUDA 生态到底锁在哪？}\par\vspace{2.2pt}
\noindent{\sffamily\bfseries 137}\;{\small PyTorch 怎么才算认识一块新卡？}\par\vspace{2.2pt}
\noindent{\sffamily\bfseries 138}\;{\small 一步 decode 要下发多少次算子？}\par\vspace{2.2pt}
\noindent{\sffamily\bfseries 139}\;{\small 换卡后怎么算精度对上了？}\par\vspace{2.2pt}
\noindent{\sffamily\bfseries 140}\;{\small 这块卡支持 PyTorch 吗？}\par\vspace{2.2pt}
\noindent{\sffamily\bfseries 141}\;{\small 把引擎迁到新卡分几步？}\par\vspace{2.2pt}
\noindent{\sffamily\bfseries 142}\;{\small 没有 Nsight 怎么找瓶颈？}\par\vspace{2.2pt}
\noindent{\sffamily\bfseries 143}\;{\small 厂商库为什么在 decode 最糟？}\par\vspace{2.2pt}
\noindent{\sffamily\bfseries 144}\;{\small 同一个 RMSNorm 写法差几倍？}\par\vspace{2.2pt}
\noindent{\sffamily\bfseries 145}\;{\small 换一块卡，量化还能不能提速？}\par\vspace{2.2pt}
\noindent{\sffamily\bfseries 146}\;{\small 换一块卡，All-Reduce 结论还成立吗？}\par\vspace{2.2pt}
\noindent{\sffamily\bfseries 147}\;{\small 怎么把一块新卡接进 vLLM 又不改源码？}\par\vspace{2.2pt}
\noindent{\sffamily\bfseries 148}\;{\small DCU 的软件栈怎么和 CUDA 对上？}\par\vspace{2.2pt}
\noindent{\sffamily\bfseries 149}\;{\small 一份 CUDA 代码搬到 DCU 要改哪些行？}\par\vspace{2.2pt}
\noindent{\sffamily\bfseries 150}\;{\small 一个真实硬件插件的接入面有多大？}\par\vspace{2.2pt}
\noindent{\sffamily\bfseries 151}\;{\small 小 M 矩阵乘为什么要自己写 kernel？}\par\vspace{2.2pt}
\noindent{\sffamily\bfseries 152}\;{\small DCU 上的注意力后端是怎么挑的？}\par\vspace{2.2pt}
\noindent{\sffamily\bfseries 153}\;{\small 模型从拿到卡到上线，你要验哪些项？}\par\vspace{2.2pt}
\end{multicols}
\end{document}